% Generated by the blueprint tool (blueprint-pdf/1) from a blueprint snapshot; edit the
% Markdown sources, not this file.  Compile with `latexmk -pdf` or by running
% `pdflatex` twice (the second run settles the contents and the page references).
% Images, if any, are read from assets/ next to this file.
\documentclass[11pt]{article}
\usepackage[T1]{fontenc}
\usepackage[utf8]{inputenc}
\usepackage{lmodern}
\usepackage{textcomp}
\usepackage[a4paper,margin=2.5cm,headheight=14pt]{geometry}
\usepackage{amsmath,amssymb,amsthm}
\usepackage{amscd}  % KaTeX's CD diagrams; part of amsmath, so always there
\usepackage{mathtools}
\usepackage{bm}
\usepackage{graphicx}
\usepackage{xcolor}
\usepackage{array,booktabs,tabularx}
\usepackage{fancyvrb}
\usepackage[normalem]{ulem}
\usepackage{microtype}
% Optional, and only with their fonts: a small TeX installation has the
% packages without the fonts, which fails at the first formula.
\IfFileExists{mathrsfs.sty}{\IfFileExists{rsfs10.tfm}{\usepackage{mathrsfs}}{}}{}
\IfFileExists{stmaryrd.sty}{\IfFileExists{stmary10.tfm}{\usepackage{stmaryrd}}{}}{}
\IfFileExists{cancel.sty}{\usepackage{cancel}}{}
\usepackage[hyphens]{url}
\usepackage[pdfusetitle,bookmarksnumbered=false]{hyperref}
\usepackage{bookmark}
\definecolor{bplink}{HTML}{1F4E9C}
\definecolor{bpmuted}{HTML}{5F676E}
\hypersetup{colorlinks=true,linkcolor=bplink,urlcolor=bplink,citecolor=bplink}
\setlength{\parindent}{0pt}
\setlength{\parskip}{0.55em plus 0.2em minus 0.1em}
% Lists as tight as the page's.
\makeatletter
\def\@listi{\leftmargin\leftmargini\parsep 0pt\topsep 0.3em plus 0.1em\itemsep 0.1em plus 0.1em}
\let\@listI\@listi
\def\@listii{\leftmargin\leftmarginii\labelwidth\leftmarginii\advance\labelwidth-\labelsep\parsep 0pt\topsep 0.1em\itemsep 0.1em}
% An ordered list that does not start at 1.
\newcommand{\bpstartat}[1]{\setcounter{\@enumctr}{\numexpr#1-1\relax}}
\makeatother
\setcounter{tocdepth}{3}
\emergencystretch=3em
\allowdisplaybreaks

% A numbered statement, set as a paper sets one: "Theorem 2.1 (Title)."
% #1 is the lead ("Theorem~2.1"), #2 what follows it in roman.
\newenvironment{bpstatement}[2]{\par\addvspace{0.6\baselineskip}\noindent\textbf{#1}#2\textbf{.}\enspace\ignorespaces}
  {\par\addvspace{0.4\baselineskip}}
\newenvironment{bpstatementit}[2]{\begin{bpstatement}{#1}{#2}\itshape\ignorespaces}{\end{bpstatement}}
% A step: a relation with prose of its own, read under the object it starts from.
\newenvironment{bpstep}[1]{\par\addvspace{0.3\baselineskip}\begin{list}{}{\leftmargin=1.6em\rightmargin=0pt\topsep=0pt}\item\relax
  {\small\color{bpmuted}#1}\par\nopagebreak}{\end{list}\par}
% Notes about where an object is written out.
\newenvironment{bpnote}{\par\small\color{bpmuted}\itshape}{\par}
\newcommand{\bpomitted}[1]{\par{\small\color{bpmuted}\itshape #1\par}}
\newcommand{\bpbodyheading}[1]{\par\addvspace{0.4\baselineskip}\noindent\textbf{#1}\par\nopagebreak}
\newcommand{\bprule}{\par\medskip\noindent\hfil\rule{0.3\linewidth}{0.4pt}\hfil\par\medskip}
\newcommand{\bpboundary}[1]{\par\nopagebreak{\small\color{bpmuted}#1\par}\nopagebreak}
% A reference to something this export leaves out, and one that resolves nowhere.
\newcommand{\bpout}[1]{#1\textsuperscript{\hyperlink{bp-outside}{\ensuremath{\ast}}}}
\newcommand{\bpbroken}[1]{\textcolor{red!70!black}{[#1?]}}
\newcommand{\bpslug}[1]{\texttt{#1}}
% A bibliography entry: label, anchor, entry; the label hangs in the margin.
\newcommand{\bpbibitem}[3]{\par\setbox0\hbox{[#1]\enspace}\ifdim\wd0<5em\setbox0\hbox to 5em{[#1]\hfil}\fi
  \hangindent=\wd0 \hangafter=1 \noindent\box0\phantomsection\label{#2}#3\par\smallskip}
% A character the export knows no way to set.
\newcommand{\bpglyph}[1]{\fbox{\scriptsize #1}}
% Images, at their own size unless they are wider than the line.
\newsavebox{\bpimagebox}
\newcommand{\bpimage}[1]{\sbox{\bpimagebox}{\includegraphics{#1}}%
  \ifdim\wd\bpimagebox>\linewidth\resizebox{\linewidth}{!}{\usebox{\bpimagebox}}\else\usebox{\bpimagebox}\fi}
\newcommand{\bpmissingimage}[1]{\fbox{\parbox{0.8\linewidth}{\centering\small\color{bpmuted}Image not included: #1}}}
\DefineVerbatimEnvironment{bpverbatim}{Verbatim}{fontsize=\small,xleftmargin=1em}
\renewcommand{\qedsymbol}{\ensuremath{\blacksquare}}
\hypersetup{pdftitle={Weil Conjectures — Part Ⅱ Duality and the trace formula},pdfsubject={Part Ⅱ Duality and the trace formula; all levels},pdfcreator={blueprint}}
% KaTeX commands the website knows
\providecommand{\R}{\mathbb{R}}\providecommand{\N}{\mathbb{N}}\providecommand{\Z}{\mathbb{Z}}
\providecommand{\Q}{\mathbb{Q}}\providecommand{\C}{\mathbb{C}}
\providecommand{\Reals}{\mathbb{R}}\providecommand{\reals}{\mathbb{R}}\providecommand{\Complex}{\mathbb{C}}
\providecommand{\cnums}{\mathbb{C}}\providecommand{\natnums}{\mathbb{N}}\providecommand{\Bbb}[1]{\mathbb{#1}}
\providecommand{\bold}[1]{\mathbf{#1}}\providecommand{\frak}[1]{\mathfrak{#1}}\providecommand{\mathbbm}[1]{\mathbb{#1}}
\providecommand{\lang}{\langle}\providecommand{\rang}{\rangle}\providecommand{\bra}[1]{\langle #1\rvert}
\providecommand{\ket}[1]{\lvert #1\rangle}\providecommand{\braket}[1]{\langle #1\rangle}
\providecommand{\set}[1]{\{#1\}}\providecommand{\Set}[1]{\left\{#1\right\}}
\providecommand{\infin}{\infty}\providecommand{\isin}{\in}\providecommand{\plusmn}{\pm}\providecommand{\sub}{\subset}
\providecommand{\sube}{\subseteq}\providecommand{\supe}{\supseteq}\providecommand{\exist}{\exists}
\providecommand{\bull}{\bullet}\providecommand{\larr}{\leftarrow}\providecommand{\rarr}{\rightarrow}
\providecommand{\lrarr}{\leftrightarrow}\providecommand{\Larr}{\Leftarrow}\providecommand{\Rarr}{\Rightarrow}
\providecommand{\Lrarr}{\Leftrightarrow}\providecommand{\uarr}{\uparrow}\providecommand{\darr}{\downarrow}
\providecommand{\harr}{\leftrightarrow}\providecommand{\lArr}{\Leftarrow}\providecommand{\rArr}{\Rightarrow}
\providecommand{\hArr}{\Leftrightarrow}\providecommand{\alef}{\aleph}\providecommand{\alefsym}{\aleph}
\providecommand{\real}{\Re}\providecommand{\image}{\Im}\providecommand{\weierp}{\wp}
\providecommand{\Alpha}{\mathrm{A}}\providecommand{\Beta}{\mathrm{B}}\providecommand{\Epsilon}{\mathrm{E}}
\providecommand{\Zeta}{\mathrm{Z}}\providecommand{\Eta}{\mathrm{H}}\providecommand{\Iota}{\mathrm{I}}
\providecommand{\Kappa}{\mathrm{K}}\providecommand{\Mu}{\mathrm{M}}\providecommand{\Nu}{\mathrm{N}}
\providecommand{\Omicron}{\mathrm{O}}\providecommand{\Rho}{\mathrm{P}}\providecommand{\Tau}{\mathrm{T}}
\providecommand{\Chi}{\mathrm{X}}\providecommand{\omicron}{o}\providecommand{\gt}{>}\providecommand{\lt}{<}
\providecommand{\And}{\&}\providecommand{\hdashline}{\hline}\providecommand{\operatornamewithlimits}[1]{\operatorname*{#1}}
\providecommand{\llbracket}{[\![}\providecommand{\rrbracket}{]\!]}\providecommand{\cancel}[1]{#1}
\providecommand{\mathscr}{\mathcal}
% the characters beyond ASCII this document uses
\DeclareUnicodeCharacter{2160}{\textup{I}}
\DeclareUnicodeCharacter{2161}{\textup{II}}
\DeclareUnicodeCharacter{2162}{\textup{III}}
% the project's macros ([katex.macros] in blueprint.toml)
\providecommand{\A}{}\renewcommand{\A}{\mathbf A}
\providecommand{\AEn}{}\renewcommand{\AEn}{A_{E,n}}
\providecommand{\Aaf}{}\renewcommand{\Aaf}{\mathbf A}
\providecommand{\Ab}{}\renewcommand{\Ab}{\mathsf{Ab}}
\providecommand{\AffSch}{}\renewcommand{\AffSch}{\operatorname{AffSch}}
\providecommand{\Aone}{}\renewcommand{\Aone}{\mathbf A^1}
\providecommand{\Bl}{}\renewcommand{\Bl}{\operatorname{Bl}}
\providecommand{\Br}{}\renewcommand{\Br}{\mathrm{Br}}
\providecommand{\C}{}\renewcommand{\C}{\mathbb C}
\providecommand{\Cech}{}\renewcommand{\Cech}{\check{H}}
\providecommand{\Coext}{}\renewcommand{\Coext}{\operatorname{Coext}}
\providecommand{\Cone}{}\renewcommand{\Cone}{\operatorname{Cone}}
\providecommand{\Cons}{}\renewcommand{\Cons}{\operatorname{Cons}}
\providecommand{\Cref}{}\renewcommand{\Cref}[1]{\text{#1}}
\providecommand{\D}{}\renewcommand{\D}{D}
\providecommand{\DLAFT}{}\renewcommand{\DLAFT}{\operatorname{FT}}
\providecommand{\DLAFr}{}\renewcommand{\DLAFr}{\operatorname{Fr}}
\providecommand{\DLAQ}{}\renewcommand{\DLAQ}{{\overline{\mathbf Q}_\ell}}
\providecommand{\DLRQ}{}\renewcommand{\DLRQ}{{\overline{\mathbf Q}_\ell}}
\providecommand{\Dbc}{}\renewcommand{\Dbc}{D^b_c}
\providecommand{\Dc}{}\renewcommand{\Dc}{D_{\mathrm{cons}}}
\providecommand{\Dctf}{}\renewcommand{\Dctf}{D^b_{ctf}}
\providecommand{\Dic}{}\renewcommand{\Dic}{D_{\IndCons}}
\providecommand{\Dplus}{}\renewcommand{\Dplus}{D^{+}}
\providecommand{\End}{}\renewcommand{\End}{\operatorname{End}}
\providecommand{\Et}{}\renewcommand{\Et}{\mathrm{\acute{e}t}}
\providecommand{\Ex}{}\renewcommand{\Ex}{\mathrm{Ex}}
\providecommand{\Ext}{}\renewcommand{\Ext}{\operatorname{Ext}}
\providecommand{\F}{}\renewcommand{\F}{\mathbb F}
\providecommand{\FT}{}\renewcommand{\FT}{\operatorname{FT}}
\providecommand{\Fbar}{}\renewcommand{\Fbar}{{\overline{\mathbf F}_q}}
\providecommand{\Fell}{}\renewcommand{\Fell}{\mathbf F_\ell}
\providecommand{\Fib}{}\renewcommand{\Fib}{\operatorname{Fib}}
\providecommand{\Fix}{}\renewcommand{\Fix}{\operatorname{Fix}}
\providecommand{\Fp}{}\renewcommand{\Fp}{{\mathbf F_{p}}}
\providecommand{\Fq}{}\renewcommand{\Fq}{{\mathbf F_q}}
\providecommand{\Fqn}{}\renewcommand{\Fqn}{{\mathbf F_{q^{n}}}}
\providecommand{\Frob}{}\renewcommand{\Frob}{\operatorname{Frob}}
\providecommand{\GL}{}\renewcommand{\GL}{\operatorname{GL}}
\providecommand{\Gal}{}\renewcommand{\Gal}{\operatorname{Gal}}
\providecommand{\Hens}{}\renewcommand{\Hens}{\operatorname{Hens}}
\providecommand{\Hilb}{}\renewcommand{\Hilb}{\operatorname{Hilb}}
\providecommand{\Hom}{}\renewcommand{\Hom}{\operatorname{Hom}}
\providecommand{\IndCons}{}\renewcommand{\IndCons}{\operatorname{IndCons}}
\providecommand{\Kar}{}\renewcommand{\Kar}{\operatorname{Kar}}
\providecommand{\LT}{}\renewcommand{\LT}{\operatorname{LT}}
\providecommand{\Lie}{}\renewcommand{\Lie}{\operatorname{Lie}}
\providecommand{\Loc}{}\renewcommand{\Loc}{\operatorname{Loc}}
\providecommand{\Mod}{}\renewcommand{\Mod}{\mathsf{Mod}}
\providecommand{\N}{}\renewcommand{\N}{\mathbb{N}}
\providecommand{\Needspace}{}\renewcommand{\Needspace}{\par}
\providecommand{\O}{}\renewcommand{\O}{{\mathcal{O}}}
\providecommand{\Ob}{}\renewcommand{\Ob}{\mathrm{Ob}}
\providecommand{\Oe}{}\renewcommand{\Oe}{\mathcal O_E}
\providecommand{\PP}{}\renewcommand{\PP}{\mathbf{P}}
\providecommand{\Perf}{}\renewcommand{\Perf}{\operatorname{Perf}}
\providecommand{\Pic}{}\renewcommand{\Pic}{\mathrm{Pic}}
\providecommand{\Pm}{}\renewcommand{\Pm}{\mathbf P^{m}}
\providecommand{\Pone}{}\renewcommand{\Pone}{\mathbf P^{1}}
\providecommand{\Pqn}{}\renewcommand{\Pqn}{\mathbf P^{n}}
\providecommand{\ProFin}{}\renewcommand{\ProFin}{\operatorname{ProFin}}
\providecommand{\Proet}{}\renewcommand{\Proet}{\mathrm{pro\acute{e}t}}
\providecommand{\Q}{}\renewcommand{\Q}{\mathbb Q}
\providecommand{\Qb}{}\renewcommand{\Qb}{\overline{\mathbf Q}_{\ell}}
\providecommand{\Qbar}{}\renewcommand{\Qbar}{\overline{\mathbf Q}}
\providecommand{\QbarX}{}\renewcommand{\QbarX}{\overline{\mathbf Q}_{\ell,X}}
\providecommand{\Qbarl}{}\renewcommand{\Qbarl}{{\overline{\mathbf Q}_{\ell}}}
\providecommand{\Qell}{}\renewcommand{\Qell}{{\mathbf{Q}_{\ell}}}
\providecommand{\Ql}{}\renewcommand{\Ql}{\mathbf Q_{\ell}}
\providecommand{\Qlb}{}\renewcommand{\Qlb}{{\overline{\mathbf Q}_\ell}}
\providecommand{\Qlbar}{}\renewcommand{\Qlbar}{\overline{\mathbf Q}_\ell}
\providecommand{\R}{}\renewcommand{\R}{\mathbb{R}}
\providecommand{\RG}{}\renewcommand{\RG}{\operatorname{\mathrm{R}\Gamma}}
\providecommand{\RGamma}{}\renewcommand{\RGamma}{R\Gamma}
\providecommand{\RGammac}{}\renewcommand{\RGammac}{R\Gamma_c}
\providecommand{\RHom}{}\renewcommand{\RHom}{R\mathcal{H}om}
\providecommand{\RHomX}{}\renewcommand{\RHomX}{R\mathcal H om_X}
\providecommand{\Rat}{}\renewcommand{\Rat}{\operatorname{IsRat}}
\providecommand{\Res}{}\renewcommand{\Res}{\operatorname{Res}}
\providecommand{\Rlim}{}\renewcommand{\Rlim}{R\!\varprojlim}
\providecommand{\SIm}{}\renewcommand{\SIm}{\operatorname{SIm}}
\providecommand{\Sch}{}\renewcommand{\Sch}{\mathsf{Sch}}
\providecommand{\Sh}{}\renewcommand{\Sh}{\operatorname{Sh}}
\providecommand{\Span}{}\renewcommand{\Span}{\operatorname{span}}
\providecommand{\Spec}{}\renewcommand{\Spec}{\operatorname{Spec}}
\providecommand{\Tr}{}\renewcommand{\Tr}{\operatorname{Tr}}
\providecommand{\Trsm}{}\renewcommand{\Trsm}{\operatorname{Tr}^{\mathrm{sm}}}
\providecommand{\Wt}{}\renewcommand{\Wt}{\operatorname{Wt}}
\providecommand{\Z}{}\renewcommand{\Z}{\mathbb{Z}}
\providecommand{\Zar}{}\renewcommand{\Zar}{\mathrm{zar}}
\providecommand{\Zell}{}\renewcommand{\Zell}{{\mathbf{Z}_{\ell}}}
\providecommand{\Zl}{}\renewcommand{\Zl}{{\mathbf Z_\ell}}
\providecommand{\affet}{}\renewcommand{\affet}[1]{{#1}^{\mathrm{aff}}_{\mathrm{ét}}}
\providecommand{\affproet}{}\renewcommand{\affproet}[1]{{#1}^{\mathrm{aff}}_{\mathrm{proét}}}
\providecommand{\boxt}{}\renewcommand{\boxt}{\boxtimes_A^{\mathrm L}}
\providecommand{\cF}{}\renewcommand{\cF}{\mathcal{F}}
\providecommand{\cHom}{}\renewcommand{\cHom}{R\mathcal{H}\!om}
\providecommand{\cI}{}\renewcommand{\cI}{\mathcal{I}}
\providecommand{\cL}{}\renewcommand{\cL}{\mathcal L}
\providecommand{\cM}{}\renewcommand{\cM}{\mathcal M}
\providecommand{\cN}{}\renewcommand{\cN}{\mathcal N}
\providecommand{\cO}{}\renewcommand{\cO}{\mathcal O}
\providecommand{\cT}{}\renewcommand{\cT}{\mathcal T}
\providecommand{\checked}{}\renewcommand{\checked}{\color{red}\aftergroup\restorecolor}
\providecommand{\coeff}{}\renewcommand{\coeff}{\operatorname{coeff}}
\providecommand{\coker}{}\renewcommand{\coker}{\operatorname{coker}}
\providecommand{\colim}{}\renewcommand{\colim}{\operatorname{colim}}
\providecommand{\constructedby}{}\renewcommand{\constructedby}[1]{\par\smallskip\noindent\textcolor{purple!60!black}{\textsf{Constructed by:} }
  \footnotesize\sloppy\nolinkurl{#1}\par}
\providecommand{\deepresult}{}\renewcommand{\deepresult}[1]{\par\smallskip\noindent\textcolor{orange!45!black}{
  \textsf{Blueprint role: named deep theorem node; see the cluster proof architecture.}}
  \par\noindent\textcolor{orange!45!black}{\textsf{Mathematical source:} }#1\par}
\providecommand{\dt}{}\renewcommand{\dt}{\mathbin{\otimes^{\mathbf L}}}
\providecommand{\dw}{}\renewcommand{\dw}{\operatorname{dw}}
\providecommand{\et}{}\renewcommand{\et}[1]{{#1}_{\mathrm{ét}}}
\providecommand{\ev}{}\renewcommand{\ev}{\operatorname{ev}}
\providecommand{\fib}{}\renewcommand{\fib}{\operatorname{fib}}
\providecommand{\formalizationnote}{}\renewcommand{\formalizationnote}[1]{\begin{quote}\small\raggedright\textbf{Formalization note.} #1\end{quote}}
\providecommand{\fpqc}{}\renewcommand{\fpqc}{X_{\mathrm{fpqc}}}
\providecommand{\gammacont}{}\renewcommand{\gammacont}{\Gamma_{\mathrm{cont}}}
\providecommand{\gm}{}\renewcommand{\gm}{\mathbf G_m}
\providecommand{\hocolim}{}\renewcommand{\hocolim}{\operatorname{hocolim}}
\providecommand{\holim}{}\renewcommand{\holim}{\operatorname{holim}}
\providecommand{\hrB}{}\renewcommand{\hrB}{\@ifnextchar]{\hrEnd}{\hrStep}}
\providecommand{\hrEnd}{}\renewcommand{\hrEnd}{\@firstoftwo{}}
\providecommand{\hrStep}{}\renewcommand{\hrStep}{\@firstoftwo{\hrB}}
\providecommand{\id}{}\renewcommand{\id}{\operatorname{id}}
\providecommand{\kunFl}{}\renewcommand{\kunFl}{\mathbf F_\ell}
\providecommand{\kunZn}{}\renewcommand{\kunZn}{\mathbf Z/n\mathbf Z}
\providecommand{\kuncd}{}\renewcommand{\kuncd}{\operatorname{cd}}
\providecommand{\kundim}{}\renewcommand{\kundim}{\dim}
\providecommand{\kunotimes}{}\renewcommand{\kunotimes}[1]{\otimes^{\mathbf L}_{#1}}
\providecommand{\kunsupp}{}\renewcommand{\kunsupp}{\operatorname{Supp}}
\providecommand{\lotimes}{}\renewcommand{\lotimes}{\otimes^{\mathrm L}}
\providecommand{\m}{}\renewcommand{\m}{{\mathfrak{m}}}
\providecommand{\n}{}\renewcommand{\n}{{\mathfrak{n}}}
\providecommand{\nagataGap}{}\renewcommand{\nagataGap}[1]{\par\noindent\textbf{Formalization gap.} #1\par}
\providecommand{\openobligation}{}\renewcommand{\openobligation}[1]{\par\smallskip\noindent\textcolor{red!65!black}{\textsf{Open obligation:} }
  \footnotesize\sloppy\nolinkurl{#1}\par}
\providecommand{\p}{}\renewcommand{\p}{{\mathfrak{p}}}
\providecommand{\pr}{}\renewcommand{\pr}{\operatorname{pr}}
\providecommand{\preshv}{}\renewcommand{\preshv}[1]{\mathrm{PreShv}(#1)}
\providecommand{\pro}{}\renewcommand{\pro}[1]{\mathrm{Pro}(#1)}
\providecommand{\proet}{}\renewcommand{\proet}{X_{\mathrm{proét}}}
\providecommand{\q}{}\renewcommand{\q}{{\mathfrak{q}}}
\providecommand{\rad}{}\renewcommand{\rad}{\operatorname{rad}}
\providecommand{\rank}{}\renewcommand{\rank}{\operatorname{rank}}
\providecommand{\red}{}\renewcommand{\red}{\color{red}}
\providecommand{\restorecolor}{}\renewcommand{\restorecolor}{\color{black}}
\providecommand{\retiredobligation}{}\renewcommand{\retiredobligation}[1]{\par\smallskip\noindent\textcolor{red!75!black}{\textsf{Retired invalid route:} }
  \footnotesize\sloppy\nolinkurl{#1}\par}
\providecommand{\rk}{}\renewcommand{\rk}{\operatorname{rk}}
\providecommand{\sep}{}\renewcommand{\sep}{\textnormal{sep}}
\providecommand{\shv}{}\renewcommand{\shv}[1]{\mathrm{Shv}(#1)}
\providecommand{\spec}{}\renewcommand{\spec}[1]{\mathrm{Spec}(#1)}
\providecommand{\stacks}{}\renewcommand{\stacks}[1]{\href{https://stacks.math.columbia.edu/tag/#1}{\textnormal{Stacks, Tag~#1}}}
\providecommand{\stacksproject}{}\renewcommand{\stacksproject}[1]{{\cite[\href{https://stacks.math.columbia.edu/tag/#1}{Tag #1}]{stacks-project}}}
\providecommand{\statusnote}{}\renewcommand{\statusnote}[1]{\begin{quote}\small\raggedright\textbf{Status.} #1\end{quote}}
\providecommand{\tildeideal}{}\renewcommand{\tildeideal}[2]{\widetilde{(#1, #2)}}
\providecommand{\tildering}{}\renewcommand{\tildering}[3]{{#1}_{\widetilde{(#2, #3)}}}
\providecommand{\tr}{}\renewcommand{\tr}{\operatorname{tr}}
\providecommand{\wt}{}\renewcommand{\wt}{\operatorname{wt}}

\begin{document}
\pagestyle{plain}
\begin{center}
{\LARGE\bfseries Weil Conjectures\par}
\vspace{0.7em}
{\large Part Ⅱ: Duality and the trace formula\par}
\vspace{0.5em}
{\small\color{bpmuted}Snapshot 2026-09-26, revision 8e4cb88e2a\par}
\end{center}
\vspace{0.5em}
{\small\tableofcontents}
\bigskip\hrule\bigskip

% object: part-Duality-and-the-trace-formula (blueprint/part-Duality-and-the-trace-formula/_section.md)
\part*{\texorpdfstring{Part Ⅱ\quad Duality and the trace formula}{Part Ⅱ Duality and the trace formula}}
\phantomsection\label{bp1}%
\addcontentsline{toc}{part}{\texorpdfstring{Part Ⅱ\quad Duality and the trace formula}{Part Ⅱ Duality and the trace formula}}
% text: prose-P2-01
\phantomsection\label{bp2}%
At the end of Part I we have, for every separated \(X_0\) of finite type over \(\mathbf F_q\), cohomology groups \(H^i_c\) with an action of Frobenius, and they are finite because \(Rf_!\) preserves \(D^b_{\mathrm{ctf}}\). Nothing yet connects them with the numbers \(\#X_0(\mathbf F_{q^m})\) of which the zeta function is made. This part supplies the connection: Weil's axioms (b) and (c), the trace formula and Poincaré duality, and with them rationality and the functional equation (Section \hyperref[bp312]{12.7}).

For the elliptic curve \(y^2=x^3-x\) over \(\mathbf F_3\) these two axioms, together with the single count \(N_1=4\), determine the zeta function and every \(N_m\) (Chapter \hyperref[bp226]{12}, Example \hyperref[bp323]{12.32}).

Both axioms are statements about the diagonal. For example, on a compact oriented manifold \(M\), the Künneth components of the class of the diagonal in \(H^*(M\times M)\) form the copairing inverse to the cup product pairing, and the Lefschetz number of a map \(\varphi\) is the intersection number of its graph with the diagonal. A scheme, however, comes with no orientation and no class of the diagonal, and smooth purity is what provides them. It says that \(f^!\simeq f^*(d)[2d]\) for \(f\) smooth of relative dimension \(d\) (Theorem \hyperref[bp11]{9.3}), and this isomorphism gives the trace map.

Verdier duality (Theorems \hyperref[bp125]{10.8} and \hyperref[bp132]{Theorem 10.10}) says that \(\mathbb D=\cHom(-,\omega)\) is an involution of \(D^b_{\mathrm{ctf}}\) exchanging \(Rf_!\) with \(Rf_*\). For smooth proper \(X\) it turns the trace map into a perfect pairing between \(H^i\) and \(H^{2d-i}(d)\) (Corollary \hyperref[bp152]{10.17}), and this pairing is compatible with Frobenius (Theorem \hyperref[bp165]{10.21}). Verdier duality also gives the vanishing in (a) for such \(X\) (Theorem \hyperref[bp146]{10.15}).

For the trace formula we use the Künneth formula for Verdier duals (Lemma \hyperref[bp191]{11.7}), which turns a correspondence \(u:c_2^*K\to c_1^!K\) into a class on \(C\); for an endomorphism, \(C\) is the graph. Restricting this class to the diagonal gives the Lefschetz–Verdier formula (Theorem \hyperref[bp206]{11.13}), which says that on a proper scheme the trace of \(u\) on cohomology is a sum of local terms over the fixed locus.

This does not yet count points, for two reasons. A local term is computed from \(\cHom(c_2^*K,c_1^!K)\) near a fixed point, and the stalk of \(K\) does not determine it (Remark \hyperref[bp239]{12.4}). Moreover \(X\) need not be proper. Both problems are settled by one computation on a curve: the local term of Frobenius at a rational point is the trace on the stalk (Proposition \hyperref[bp274]{12.16}). The computation works because Frobenius has zero differential, so that its graph meets the diagonal with multiplicity one and is tangent to a coordinate axis. The result is the Grothendieck–Lefschetz trace formula (Theorem \hyperref[bp298]{12.24}): for every Weil complex \(K\), \(\operatorname{tr}(\operatorname{Frob}^m\mid R\Gamma_c(X,K))\) is the sum over \(X_0(\mathbf F_{q^m})\) of the traces of Frobenius on the stalks.

Smooth purity and Verdier duality hold over any finite or algebraically closed field, and their compatibility with Frobenius follows from their construction (Section \hyperref[bp154]{10.5}).

% object: chap-Smooth-purity (blueprint/chap-Smooth-purity/_section.md), content hidden
\section*{\texorpdfstring{9\quad Smooth purity}{9 Smooth purity}}
\phantomsection\label{bp3}%
\addcontentsline{toc}{section}{\texorpdfstring{9\quad Smooth purity}{9 Smooth purity}}
% text: prose-09-01
\phantomsection\label{bp4}%
Poincaré duality for a smooth variety \(X\) of dimension \(d\) pairs \(H^i_c(X,\Lambda)\) with \(H^{2d-i}(X,\Lambda(d))\) through the cup product and a trace map \(H^{2d}_c(X,\Lambda(d))\to\Lambda\). If \(a\) is the structure map of \(X\), such a trace map is the counit of \(Ra_!\dashv a^!\), once we know that \(a^!\Lambda\simeq\Lambda_X(d)[2d]\). But Part I produces \(f^!\) (Theorem \bpout{5.2}) as a right adjoint by Brown representability, with no formula.

A formula exists on charts. Consider \(f:X\to Y\) smooth of pure relative dimension \(d\), and an open subscheme \(U\subset X\) on which \(f\) factors as an étale map \(e:U\to\mathbf A^d_Y\) followed by the projection \(p_d\). With finite coefficients \(f^!K|_U\simeq e^*p_d^*K(d)[2d]\) (Proposition \bpout{6.9}), and under this isomorphism the counit is integration along \(d\) affine lines. This integration comes from the trace map of \(\mathbf P^1\) (Proposition \bpout{6.7}), which sends the class of a point to \(1\). The question is then whether these chart descriptions agree on overlaps.

In topology they need not. For a submersion \(f:X\to Y\) of manifolds with fibres of real dimension \(m\), Verdier's theory gives \(f^!K\simeq f^*K\otimes\mathrm{or}_{X/Y}[m]\). Here the relative orientation sheaf \(\mathrm{or}_{X/Y}\) is locally constant of rank one, and it is trivial exactly when \(f\) can be oriented. A linear change of coordinates \(g\) acts on \(H^m_c(\mathbf R^m)\) by the sign of \(\det g\), and since \(\operatorname{GL}_m(\mathbf R)\) has two components, the local identifications need not glue. For complex manifolds, by contrast, \(\operatorname{GL}_d(\mathbf C)\) is connected and the orientation is canonical.

The same argument works for us on a chart. Indeed, a linear change of coordinates multiplies the trace map of \(\mathbf A^d\) by a locally constant unit on \(\operatorname{GL}_d\), and this unit is \(1\) at the identity, hence everywhere, since \(\operatorname{GL}_d\) has connected geometric fibres. So every vector bundle has a canonical trace map (Proposition \hyperref[bp21]{9.6}). The trace map has source \(H^{2d}_c(-,\Lambda(d))\), so the Tate twist \(\Lambda(d)\) plays the role of the orientation sheaf. Over an algebraically closed field a compatible system of roots of unity trivializes \(\Lambda(d)\), but not canonically, and over \(\mathbf F_q\) the twist is visible: there geometric Frobenius acts on \(\Lambda(1)\) by \(q^{-1}\) (Section \hyperref[bp154]{10.5}), and this is where the factor \(q^d\) in Weil's Poincaré duality comes from.

Connectedness of \(\operatorname{GL}_d\) does not settle the gluing, because two étale charts \(U\to\mathbf A^d_Y\) differ by a map that need not be linear. We therefore construct the isomorphism without charts, from the diagonal, in the way a tubular neighbourhood linearizes a submanifold (Section \hyperref[bp24]{9.4}).

The outcome, Theorem \hyperref[bp11]{9.3}, is a canonical isomorphism

\[
\gamma_f:f^*K(d)[2d]\xrightarrow{\sim}f^!K,
\]

natural in \(K\in D^b_{\mathrm{ctf}}(Y,\Lambda)\) for every admissible \(\Lambda\). It is compatible with base change, composition and change of coefficients, and it is normalized by the trace map of the affine line. This statement is Deligne's [\hyperref[bp330]{SGA4}, XVIII], where it is called Poincaré duality.

% object: sec-Statement (blueprint/chap-Smooth-purity/sec-Statement/_section.md), content hidden
\subsection*{\texorpdfstring{9.1\quad Statement}{9.1 Statement}}
\phantomsection\label{bp5}%
\addcontentsline{toc}{subsection}{\texorpdfstring{9.1\quad Statement}{9.1 Statement}}
% text: prose-09-02
\phantomsection\label{bp6}%
Smooth purity is stated for six kinds of coefficients at once, because the proof moves between them: the geometry is done with finite rings, the case \(\mathcal O_E\) is a limit of finite levels, and the fields are obtained by inverting \(\ell\). Throughout this chapter \(\Lambda\) is an admissible coefficient ring, that is, a finite commutative ring killed by a power of \(\ell\), \(\mathcal O_E\), \(E\) or \(\overline{\mathbf Q}_\ell\) (the rings \(\mathbf Z_\ell\) and \(\mathbf Q_\ell\) being the cases \(E=\mathbf Q_\ell\)), and Tate twists are those of Definition \bpout{2.53}.

The theorem asserts one isomorphism on all of \(X\), whereas on a chart an isomorphism of this shape is already given by Proposition \bpout{6.9} with finite coefficients. A bare isomorphism for each \(f\) separately, however, would not be enough for what follows, and the theorem therefore also lists three properties that the later chapters use. Property (2), compatibility with base change, restricts \(\gamma_f\) to the fibres of \(f\), and it is what makes \(\gamma_f\) commute with Frobenius, because the Frobenius square is a base change (Proposition \hyperref[bp162]{10.20}). Property (3), compatibility with composition, gives the transitivity of the trace map (Proposition \hyperref[bp88]{9.26}), the analogue of integrating over the fibres of a composite in two steps. Property (4), compatibility with change of coefficients, lets us construct \(\gamma_f\) with finite coefficients and pass to \(\mathcal O_E\), \(E\) and \(\overline{\mathbf Q}_\ell\), and it also makes Poincaré duality compatible with change of coefficients (Theorem \hyperref[bp137]{10.12}). Properties (1) and (2) involve maps that we define now, for every admissible \(\Lambda\); the functors are those of Theorem \bpout{5.2}, restricted to \(D^b_{\mathrm{ctf}}\).

The first map compares \(f^!\) before and after a change of base. Let \(f:X\to Y\) and \(g:Y'\to Y\) be morphisms of separated schemes of finite type over \(k\), let \(X'=X\times_YY'\), and let \(f':X'\to Y'\) and \(g':X'\to X\) be the projections. Write \(\varepsilon_f(K):Rf_!f^!K\to K\) for the counit of \(Rf_!\dashv f^!\).

% object: definition-exceptional-base-change-map (blueprint/chap-Smooth-purity/sec-Statement/definition-exceptional-base-change-map/_section.md)
\phantomsection\label{bp7}%
\begin{bpstatement}{Definition~9.1}{ (Exceptional base change map)}
For \(K\in D^b_{\mathrm{ctf}}(Y,\Lambda)\), the \emph{exceptional base change map} \(g'^*f^!K\to f'^!g^*K\) is the adjoint transpose, for \(Rf'_!\dashv f'^!\), of the composite

\[
Rf'_!g'^*f^!K\xrightarrow{\sim}g^*Rf_!f^!K
\xrightarrow{\ g^*\varepsilon_f(K)\ }g^*K,
\]

whose first arrow is the inverse of the base change isomorphism of Proposition proposition-composition-base-change-rf-shriek.
\end{bpstatement}

% text: prose-09-03
\phantomsection\label{bp8}%
The second map is integration along the affine line. Let \(p:\mathbf A^1_Y\to Y\) be the projection, and factor it as \(p=q\circ j\), where \(j:\mathbf A^1_Y\hookrightarrow\mathbf P^1_Y\) is the complement of the section at infinity and \(q:\mathbf P^1_Y\to Y\) is the projection.

% object: definition-trace-map-affine-line (blueprint/chap-Smooth-purity/sec-Statement/definition-trace-map-affine-line/_section.md)
\phantomsection\label{bp9}%
\begin{bpstatement}{Definition~9.2}{ (Trace map of the affine line)}
The \emph{trace map of the affine line} is the composite

\[
\operatorname{Tr}_p:Rp_!\Lambda_{\mathbf A^1_Y}(1)[2]
\simeq Rq_*j_!\Lambda_{\mathbf A^1_Y}(1)[2]
\longrightarrow Rq_*\Lambda_{\mathbf P^1_Y}(1)[2]
\xrightarrow{\ \operatorname{Tr}_q\ }\Lambda_Y,
\tag{9.1}
\]

where the first arrow is induced by the counit \(j_!j^*\to\mathrm{id}\) and \(\operatorname{Tr}_q\) is the trace map of the projective line: for finite \(\Lambda\) it is the map (6.5) of Lemma \bpout{6.4} for the projection \(q\), and otherwise it is the trace of Proposition \bpout{6.22}.
\end{bpstatement}

% text: prose-09-04
\phantomsection\label{bp10}%
Over an algebraically closed field the trace map of the affine line sends the class in \(H^2_c(\mathbf A^1,\Lambda(1))\) of a point to \(1\), as we check with finite coefficients in Section \hyperref[bp17]{9.3}.

% object: theorem-smooth-purity-all-admissible-coefficients (blueprint/chap-Smooth-purity/sec-Statement/theorem-smooth-purity-all-admissible-coefficients/_section.md)
\phantomsection\label{bp11}%
\begin{bpstatementit}{Theorem~9.3}{ (Smooth purity for all admissible coefficients)}
Let \(f:X\to Y\) be smooth of pure relative dimension \(d\). For every admissible coefficient ring \(\Lambda\in\{A,\mathbf Z_\ell,\mathcal O_E,\mathbf Q_\ell,E,\overline{\mathbf Q}_\ell\}\), with \(A\) finite, there are isomorphisms, natural in \(K\in D^b_{\mathrm{ctf}}(Y,\Lambda)\),

\[
\gamma_f(K):
f^{*}K(d)[2d]\xrightarrow{\sim}f^{!}K,
\tag{9.2}
\]

with the following properties.

\begin{enumerate}
\item If \(f\) is the projection \(p:\mathbf A^1_Y\to Y\), then \(\varepsilon_p(\Lambda_Y)\circ Rp_!\bigl(\gamma_p(\Lambda_Y)\bigr)\) is the trace map (9.1).
\item For \(g:Y'\to Y\) as in Definition \hyperref[bp7]{9.1}, with exceptional base change map \(\beta:g'^*f^!K\to f'^!g^*K\), we have \(\beta\circ g'^*\gamma_f(K)=\gamma_{f'}(g^*K)\), the sources being identified by \(g'^*f^*\simeq f'^*g^*\).
\item For a smooth morphism \(h:Y\to Z\) of pure relative dimension \(e\) and \(K\in D^b_{\mathrm{ctf}}(Z,\Lambda)\), the map \(\gamma_{h\circ f}(K)\) is the composite of \(f^*(\gamma_h(K))(d)[2d]\) and \(\gamma_f(h^!K)\), the sources being identified by \((h\circ f)^*\simeq f^*h^*\) and the targets by the composition isomorphism \(f^!h^!\simeq(h\circ f)^!\) of Proposition proposition-composition-open-restriction-compactifications-f.
\item For a coefficient map \(\Lambda\to\Lambda'\) of Proposition \bpout{7.10}, the change-of-coefficients isomorphism \((f^!K)_{\Lambda'}\to f^!(K_{\Lambda'})\) of that proposition carries \((\gamma_{f,\Lambda}(K))_{\Lambda'}\) to \(\gamma_{f,\Lambda'}(K_{\Lambda'})\), the sources being identified by the change-of-coefficients isomorphism for \(f^*\).
\end{enumerate}
\end{bpstatementit}

% text: prose-09-05
\phantomsection\label{bp12}%
We recall what \(D^b_{\mathrm{ctf}}\) means for each coefficient ring. For integral coefficients and for coefficients in a field, \(D^b_{\mathrm{ctf}}=\Dbc\), and we write \(\Dbc\) in calculations that only concern these cases. With integral coefficients finite Tor amplitude is automatic, and membership in \(\Dbc(T,\mathcal O_E)\) can be tested modulo \(\ell\) (Proposition \bpout{3.11}).

The proof passes through finite rings \(A\) with \(\ell^rA=0\), such as \(\mathcal O_E/\ell^n\). For these, \(D^b_{\mathrm{ctf}}\) is the category of Definition \bpout{3.1}, which on the pro-étale site is identified with the pullback of its étale counterpart (Corollary \bpout{2.44}). Here the Tor condition makes derived change of coefficients preserve boundedness, as the example after Proposition \bpout{7.10} shows. When we compare coefficient rings, we write the ring as a subscript of the functor, as in \(Rf_{!,A}\) or \(f^!_{\mathcal O_E}\); otherwise we omit the subscript.

With coefficients \(A\), we have \(Rf_!=R\bar f_*\circ j_!\) for a compactification \(f=\bar f\circ j\), and on \(D^b_{\mathrm{ctf}}\) these functors and their right adjoints are those of Theorem \bpout{5.31}, restricted in Theorem \bpout{6.20}. Proper base change, the projection formula and the localization triangle are the identities of Section \bpout{6.1}, which we may prove on the étale sites by Proposition \bpout{5.15}. For an étale map, extension by zero is exact on sheaves of \(A\)-modules on the pro-étale sites; on the derived categories it is left adjoint to inverse image, and the counit is given on geometric stalks by the sum over the fibre.

Property (2) implies that the exceptional base change map is an isomorphism for smooth \(f\), because \(\gamma_f\) and \(\gamma_{f'}\) are isomorphisms.

For étale \(f\) we have \(d=0\), and the theorem says that \(f^{!}\simeq f^{*}\) for every admissible \(\Lambda\). For a smooth curve \(C\) with structure map \(a:C\to\operatorname{Spec}k\), it gives \(a^!\mathbf Z_{\ell,\operatorname{Spec}k}\simeq\mathbf Z_{\ell,C}(1)[2]\), a single sheaf in degree \(-2\).

For smooth \(f\) the whole functor \(f^!\) is determined by its value on the constant sheaf, through the map \(f^*K\otimes^{\mathbf L}f^!\Lambda_Y\to f^!K\) of Proposition proposition-exceptional-inverse-image-smooth-morphism. Smoothness is needed here, as the following example shows. Let \(i:y\to C\) be the inclusion of a rational point of a smooth curve, with open complement \(j:U\to C\), and take \(K=j_!\Lambda_U\). Then the analogous map \(i^*K\otimes^{\mathbf L}i^!\Lambda_C\to i^!K\) has source zero, since the stalk of \(j_!\Lambda_U\) at \(y\) vanishes, but the localization triangle gives \(i^!j_!\Lambda_U\simeq(i^*Rj_*\Lambda_U)[-1]\neq0\).

We give \(f^!\Lambda_Y\) a name rather than working with \(\Lambda_X(d)[2d]\) directly, because the two are not yet known to be isomorphic. Moreover, the complex \(f^!\Lambda_Y\) is attached to \(f\) without choices, and it satisfies base change (Proposition proposition-exceptional-inverse-image-smooth-morphism). Proving the theorem thus amounts to constructing an isomorphism between it and \(\Lambda_X(d)[2d]\) that does not depend on charts. In this sense the complex is the analogue of the shifted orientation sheaf \(\mathrm{or}_{X/Y}[m]\), and for the structure map of a smooth scheme it is the dualizing complex of Poincaré duality (Corollary \hyperref[bp98]{9.28} and Theorem \hyperref[bp137]{10.12}).

% object: definition-relative-dualizing-complex (blueprint/chap-Smooth-purity/sec-Statement/definition-relative-dualizing-complex/_section.md)
\phantomsection\label{bp13}%
\begin{bpstatement}{Definition~9.4}{ (Relative dualizing complex)}
For smooth \(f:X\to Y\) and any admissible \(\Lambda\), the \emph{relative dualizing complex} is

\[
\omega_f=f^{!}\Lambda_Y
\in D^b_{\mathrm{ctf}}(X,\Lambda).
\tag{9.3}
\]
\end{bpstatement}

% text: prose-09-06
\phantomsection\label{bp14}%
For finite coefficients, \(\omega_f\) can be computed on the pro-étale site or, equivalently, on the étale site (Proposition \bpout{5.15}).


% object: sec-The-tensor-product-map (blueprint/chap-Smooth-purity/sec-The-tensor-product-map/_section.md), content hidden
\subsection*{\texorpdfstring{9.2\quad The tensor product map}{9.2 The tensor product map}}
\phantomsection\label{bp15}%
\addcontentsline{toc}{subsection}{\texorpdfstring{9.2\quad The tensor product map}{9.2 The tensor product map}}
% text: prose-09-07
\phantomsection\label{bp16}%
With finite coefficients, smooth purity says that for \(f\) smooth of pure relative dimension \(d\) the exceptional inverse image \(f^!\) is the ordinary inverse image, twisted by \((d)\) and shifted by \([2d]\) (Proposition \hyperref[bp39]{9.12}). The first step is the map \(f^*K\otimes^{\mathbf L}\omega_f\to f^!K\), an isomorphism by Proposition proposition-exceptional-inverse-image-smooth-morphism, which reduces the problem to the single complex \(\omega_f\).

Fix \(n\geq1\). We carry out the geometric calculation with coefficients \(\mathbf Z/\ell^n\mathbf Z\), and until Section \hyperref[bp41]{9.5} we omit the subscript, so that \(f^*\), \(Rf_!\), \(f^!\), \(\omega_f\) and \(\gamma_f\) refer to this cyclic coefficient ring. Afterwards the calculation extends along \(\mathbf Z/\ell^n\mathbf Z\to A\) when \(\ell^nA=0\), and in this way Proposition \hyperref[bp39]{9.12} gives the result for every finite admissible ring \(A\).

Let \(f:X\to Y\) be smooth of pure relative dimension \(d\) and \(K\in D^b_{\mathrm{ctf}}(Y_{\mathrm{pro}\acute{e}t},\mathbf Z/\ell^n\mathbf Z)\). From the projection formula and the counit \(Rf_!\omega_{f}\to(\mathbf Z/\ell^n\mathbf Z)_{Y}\) we get \(Rf_!\bigl(f^{*}K\otimes^{\mathbf L}\omega_{f}\bigr)
\simeq K\otimes^{\mathbf L} Rf_!\omega_{f}\to K\), whose adjoint is the canonical map

\[
\mu_{f,K}:f^{*}K\otimes^{\mathbf L}_{(\mathbf Z/\ell^n\mathbf Z)_{X}}\omega_{f}
\longrightarrow f^{!}K.
\tag{9.4}
\]

To construct the orientation, we compare \(\omega\) on the fibres of a family of sections, and for this we need base change. Consider a cartesian square

\[
\begin{CD}
X' @>{g'}>> X\\
@V{f'}VV @VV{f}V\\
Y' @>{g}>> Y,
\end{CD}
\tag{9.5}
\]

with \(f\) smooth of pure relative dimension \(d\). For it we have the exceptional base change map of Definition \hyperref[bp7]{9.1}. By Proposition proposition-exceptional-inverse-image-smooth-morphism and (3), with \(\Lambda=\mathbf Z/\ell^n\mathbf Z\), the map (9.4) and the exceptional base change map

\[
g'^{*}f^{!}K\xrightarrow{\sim}(f')^{!}g^{*}K
\tag{9.6}
\]

are isomorphisms for \(K\in D^b_{\mathrm{ctf}}(Y_{\mathrm{pro}\acute{e}t},\mathbf Z/\ell^n\mathbf Z)\), and they are compatible with each other. Moreover, on a chart \(f|_U=p_d\circ e\) we have \(\omega_f|_U\simeq(\mathbf Z/\ell^n\mathbf Z)_U(d)[2d]\). Hence \(\omega_f[-2d]\) is a locally constant sheaf of free \(\mathbf Z/\ell^n\mathbf Z\)-modules of rank one.


% object: sec-The-trace-map-of-affine-space (blueprint/chap-Smooth-purity/sec-The-trace-map-of-affine-space/_section.md), content hidden
\subsection*{\texorpdfstring{9.3\quad The trace map of affine space}{9.3 The trace map of affine space}}
\phantomsection\label{bp17}%
\addcontentsline{toc}{subsection}{\texorpdfstring{9.3\quad The trace map of affine space}{9.3 The trace map of affine space}}
% text: prose-09-08
\phantomsection\label{bp18}%
Every identification of \(f^!\) in this chapter is ultimately an adjunction whose counit is integration along affine lines, so we first make that integration explicit. The affine line is not proper, so we cannot integrate on it directly; instead we integrate on \(\mathbf P^1\). There, on a geometric fibre \(H^2(\mathbf P^1,\mathbf Z/\ell^n(1))\simeq\mathbf Z/\ell^n\) is generated by \(c_1(\mathcal O(1))\), the class of a point, and the trace map of \(\mathbf P^1\) reads off the coefficient of this class. The trace map of \(\mathbf A^1\) is then its composite with the map \(H^2_c(\mathbf A^1,\mathbf Z/\ell^n(1))\to H^2(\mathbf P^1,\mathbf Z/\ell^n(1))\) given by extension by zero. This construction works over an arbitrary base \(S\) and commutes with base change, and iterating it gives the trace map (9.11) of \(\mathbf A^d\), which is the counit of the adjunction on a chart.

For a scheme \(T\), Chern classes are defined by the Kummer sequence on \(T_{\acute{e}t}\),

\[
1\longrightarrow\mu_{\ell^n}\longrightarrow\mathbf G_m
\xrightarrow{\ (\cdot)^{\ell^n}\ }\mathbf G_m
\longrightarrow1,
\tag{9.7}
\]

whose boundary \(H^1(T_{\acute{e}t},\mathbf G_m)\to
H^2(T_{\acute{e}t},(\mathbf Z/\ell^n\mathbf Z)_T(1))\) sends a line bundle \(L\) to \(c_1(L)\). For a scheme \(S\), let \(q:\mathbf P^1_S\to S\) be the projection. By Lemma \bpout{6.4}, the unit and \(c_1(\mathcal O_{\mathbf P^1_S}(1))\) give an isomorphism, compatible with base change,

\[
Rq_*(\mathbf Z/\ell^n\mathbf Z)_{\mathbf P^1_S}
\simeq(\mathbf Z/\ell^n\mathbf Z)_{S}\oplus(\mathbf Z/\ell^n\mathbf Z)_{S}(-1)[-2],
\tag{9.8}
\]

and the \emph{trace map} \(\operatorname{Tr}_q:Rq_*(\mathbf Z/\ell^n\mathbf Z)_{\mathbf P^1_S}(1)[2]
\to(\mathbf Z/\ell^n\mathbf Z)_{S}\) is the projection onto the second summand, after twisting. It sends \(c_1(\mathcal O_{\mathbf P^1_{\bar s}}(1))\) on each geometric fibre to \(1\).

Let \(p_1,p_2\) be the projections of \(\mathbf P^1_S\times_S\mathbf P^1_S\), and let \(x_0,x_1\) and \(y_0,y_1\) be the pullbacks along \(p_1\) and \(p_2\) of the homogeneous coordinates of \(\mathbf P^1_S\). The diagonal \(\Delta\) is the zero locus of the section \(x_0y_1-x_1y_0\) of \(p_1^*\mathcal O(1)\otimes p_2^*\mathcal O(1)\), so that \(\mathcal O(\Delta)\) is this line bundle and \(x_0y_1-x_1y_0\) is its canonical section. Lemma \bpout{6.3} provides the cycle class of \(\Delta\),

\[
\delta:\Delta_*(\mathbf Z/\ell^n\mathbf Z)_{\mathbf P^1_S}
\longrightarrow(\mathbf Z/\ell^n\mathbf Z)_{\mathbf P^1_S\times_S\mathbf P^1_S}(1)[2],
\]

the Kummer boundary of \(\mathcal O(\Delta)\) with its canonical section, compatible with every base change on \(S\). Let \(\operatorname{Tr}_{p_i}\) be the pullback of \(\operatorname{Tr}_q\). Then the diagonal integrates to \(1\) along either factor, that is,

\[
\operatorname{Tr}_{p_i}\circ Rp_{i*}\delta
=\operatorname{id}_{(\mathbf Z/\ell^n\mathbf Z)_{\mathbf P^1_S}}
\qquad(i=1,2).
\tag{9.9}
\]

This is (6.6) of Proposition \bpout{6.7}. To see it directly, note that both sides are endomorphisms of a constant sheaf, so it suffices to check (9.9) on geometric fibres of \(p_i\), where the diagonal is a single rational point with line bundle \(\mathcal O(1)\), and \(\delta\) restricts to \(c_1(\mathcal O(1))\), which \(\operatorname{Tr}_q\) sends to \(1\).

For \(\Lambda=\mathbf Z/\ell^n\mathbf Z\), Proposition \bpout{6.7} gives the adjunction \(Rq_*\dashv q^*(-)(1)[2]\), and its unit and counit, (6.7) and (6.8) there, take the following form. The counit is integration along the fibres,

\[
\varepsilon_K:Rq_*\bigl(q^{*}K(1)[2]\bigr)
\simeq K\otimes^{\mathbf L} Rq_*(\mathbf Z/\ell^n\mathbf Z)_{\mathbf P^1_S}(1)[2]
\xrightarrow{\ \operatorname{id}_K\otimes\operatorname{Tr}_q\ } K,
\tag{9.10}
\]

and the unit is the class of the diagonal,

\[
\eta_M:\ M\simeq Rp_{2*}\Delta_*M
\xrightarrow{\ Rp_{2*}(\delta\otimes p_1^*M)\ } Rp_{2*}p_{1}^{*}M(1)[2]
\simeq q^{*}Rq_*M(1)[2],
\]

using the projection formulas and proper base change. Both commute with base change on \(S\), and the triangle identities reduce to (9.9).

Let \(j:\mathbf A^1_S\hookrightarrow\mathbf P^1_S\) be the complement of the section at infinity, and let \(p=q\circ j\). Composing this adjunction with \(j_!\dashv j^*\) we obtain an adjunction \(Rp_!\dashv p^*(-)(1)[2]\) between \(D^b_{\mathrm{ctf}}(\mathbf A^1_{S,\mathrm{pro}\acute{e}t},\mathbf Z/\ell^n\mathbf Z)\) and \(D^b_{\mathrm{ctf}}(S_{\mathrm{pro}\acute{e}t},\mathbf Z/\ell^n\mathbf Z)\), whose counit is the counit for the open immersion followed by \(\varepsilon_K\). At the constant sheaf this counit is the trace map of the affine line \(\operatorname{Tr}_p:Rp_!(\mathbf Z/\ell^n\mathbf Z)_{\mathbf A^1_S}(1)[2]
\to(\mathbf Z/\ell^n\mathbf Z)_S\) of (9.1), and it is an isomorphism. Indeed, we apply \(Rq_*\) to the localization triangle for the section at infinity. In the resulting triangle, restriction to infinity is the identity on the first summand of (9.8), and it kills the second summand, because the pullback of \(\mathcal O(1)\) along the section at infinity is the trivial line bundle on \(S\), and so the restriction of \(c_1(\mathcal O(1))\) is \(0\). Hence \(Rp_!(\mathbf Z/\ell^n\mathbf Z)_{\mathbf A^1_S}\simeq
(\mathbf Z/\ell^n\mathbf Z)_S(-1)[-2]\) is the second summand. Over a geometric point this is \(H^2_c(\mathbf A^1,(\mathbf Z/\ell^n\mathbf Z)(1))\simeq
\mathbf Z/\ell^n\mathbf Z\), generated by the class \(\kappa\) that maps to the class of a point in \(H^2(\mathbf P^1,(\mathbf Z/\ell^n\mathbf Z)(1))\).

For affine space, let \(p_j:\mathbf A^j_S\to S\) be the projection and let \(r_j:\mathbf A^j_S\to\mathbf A^{j-1}_S\) be the map forgetting the last coordinate. We define \(\operatorname{Tr}_{p_j}\) inductively as the composite

\[
\begin{aligned}
Rp_{j!}(\mathbf Z/\ell^n\mathbf Z)_{\mathbf A^j_S}(j)[2j]
&\simeq Rp_{j-1!}\bigl(Rr_{j!}(\mathbf Z/\ell^n\mathbf Z)_{\mathbf A^j_S}(1)[2]\bigr)(j-1)[2j-2]\\
&\xrightarrow{\ \operatorname{Tr}_{p_{j-1}}\circ Rp_{j-1!}(\operatorname{Tr}_{r_j})\ }
(\mathbf Z/\ell^n\mathbf Z)_S,
\end{aligned}
\]

using \(Rp_{j!}=Rp_{j-1!}Rr_{j!}\) and the projection formula. For \(j=d\) this is the \emph{trace map of affine space}, an isomorphism (the identity when \(d=0\))

\[
\operatorname{Tr}_{p_d}:Rp_{d!}(\mathbf Z/\ell^n\mathbf Z)_{\mathbf A^d_S}(d)[2d]
\xrightarrow{\sim}(\mathbf Z/\ell^n\mathbf Z)_{S}.
\tag{9.11}
\]

By Lemma lemma-affine-space-etale-localization, \(Rp_{d!}\dashv
p_d^*(-)(d)[2d]\), and the counit is the composite of the \(d\) counits for the affine lines, which at the constant sheaf is (9.11). Like those of the affine line, the unit and counit commute with base change. On a geometric fibre, the external product \(\kappa\boxtimes\cdots\boxtimes\kappa\in
H_c^{2d}(\mathbf A^d_{\bar s},(\mathbf Z/\ell^n\mathbf Z)(d))\) integrates to \(1\), since by the projection formula, integrating along the last coordinate replaces the last factor by its integral, which is \(1\).

Objects of a derived category do not glue in general. For example, suppose \(X=U\cup V\) and we know objects on \(U\) and on \(V\) isomorphic to \(f^!K|_U\) and \(f^!K|_V\). Then the two isomorphisms restricted to \(U\cap V\) may differ by an automorphism, and nothing tells us which identification to glue along. On a chart \(U\to\mathbf A^d_Y\) we have an explicit candidate for \(f^!K|_U\), namely \(e^*p_d^*K(d)[2d]\), and we want to compare such candidates with \(f^!K\) and later, for towers, to glue \(f^!\) from them. For this we keep the counit together with the object. The point is that a pair consisting of an object and a counit that satisfies a universal property is unique up to a unique isomorphism compatible with the counit. Hence local pairs can be compared on overlaps without any choice of transition maps. We use this for the comparison following Definition \hyperref[bp19]{9.5} and in the construction of the right adjoint for towers (Proposition \hyperref[bp60]{9.18}).

% object: definition-representing-objects (blueprint/chap-Smooth-purity/sec-The-trace-map-of-affine-space/definition-representing-objects/_section.md)
\phantomsection\label{bp19}%
\begin{bpstatement}{Definition~9.5}{ (Representing objects)}
Let \(f\) be smooth, and let us work in \(D(Y,\mathbf Z/\ell^n\mathbf Z)\) and \(D(X,\mathbf Z/\ell^n\mathbf Z)\) on either the étale or the pro-étale sites. Let \(K\) be an object on \(Y\). We say that a pair \((H,\varepsilon)\), where \(H\) is an object on \(X\) and \(\varepsilon:Rf_!H\to K\) is a morphism, \emph{represents} \(\operatorname{Hom}(Rf_!(-),K)\) if \(b\mapsto\varepsilon\circ Rf_!(b)\) is a bijection \(\operatorname{Hom}(M,H)\to\operatorname{Hom}(Rf_!M,K)\) for all \(M\in D^b_{\mathrm{ctf}}(X,\mathbf Z/\ell^n\mathbf Z)\).
\end{bpstatement}

% text: prose-09-09
\phantomsection\label{bp20}%
The geometric computations with finite coefficients in Section \bpout{6.1} are stated on the pro-étale site and proved by calculations on the étale site. By Corollary \bpout{2.44}, Theorem \bpout{2.38} with (2.10), and Proposition \bpout{5.15}, the pro-étale and étale categories \(D^b_{\mathrm{ctf}}(-,A)\), and the localizing subcategories they generate, are identified compatibly with pullback, direct image and extension by zero. In particular \(Rf_!=R\bar f_*\circ j_!\) is computed on the étale sites. Consequently, an adjunction between \(Rf_!\) and some functor \(G\) on the étale categories \(D^b_{\mathrm{ctf}}(-_{\acute{e}t},A)\) transfers to the pro-étale categories \(D^b_{\mathrm{ctf}}(-_{\mathrm{pro}\acute{e}t},A)\). Indeed, the Hom groups agree by full faithfulness, and the two functors \(Rf_!\) correspond to each other.

By Proposition \bpout{5.15}, the functor \(f^!\) on \(\Loc(D^b_{\mathrm{ctf}}(-,\mathbf Z/\ell^n\mathbf Z))\) of Theorem \bpout{5.31} is the étale functor \(f^!\) on \(D(-_{\acute{e}t},\mathbf Z/\ell^n\mathbf Z)\) of Section \bpout{6.1}, and it preserves \(D^b_{\mathrm{ctf}}\) by Theorem \bpout{6.20}. Hence, for \(f:X\to Y\) smooth of pure relative dimension \(d\), Proposition proposition-exceptional-inverse-image-smooth-morphism gives the local form of \(Rf_!\dashv f^!\): on an open \(U\subset X\) with \(f|_U=p_d\circ e\), where \(e:U\to\mathbf A^d_Y\) is étale,

\[
f^!K|_U\simeq e^*p_{d}^*K(d)[2d],
\]

with counit the sum over the fibres of \(e\) followed by (9.11) tensored with \(K\).

Since a right adjoint is unique up to a unique isomorphism preserving the counit, these counits on charts agree on overlaps with the counit of \(Rf_!\dashv f^!\). With finite coefficients, then, \(f^!\) exists first and the charts describe it. For towers, by contrast, there is no global functor yet, and we glue \(f^!\) from models on charts by a Mayer–Vietoris argument, which uses the triangle for an open cover \(X=U\cup V\) and \(N\in D^b_{\mathrm{ctf}}(X_{\mathrm{pro}\acute{e}t},\mathbf Z/\ell^n\mathbf Z)\):

\[
j_{U\cap V!}(N|_{U\cap V})
\xrightarrow{\ (1,-1)\ }
j_{U!}(N|_U)\oplus j_{V!}(N|_V)
\xrightarrow{\ (1,1)\ }N\longrightarrow{} .
\tag{9.12}
\]

Suppose that \(H\in D^b_{\mathrm{ctf}}(X_{\mathrm{pro}\acute{e}t},\mathbf Z/\ell^n\mathbf Z)\) and \(\varepsilon:Rf_!H\to K\) restrict to pairs that represent \(\operatorname{Hom}(Rf_!(-),K)\) (Definition \hyperref[bp19]{9.5}) on \(U\), \(V\) and \(U\cap V\), where the restriction of \(\varepsilon\) to \(U\) is the composite \(R(f|_U)_!(H|_U)=Rf_!j_{U!}(H|_U)\to Rf_!H\to K\). For \(M\in D^b_{\mathrm{ctf}}(X_{\mathrm{pro}\acute{e}t},\mathbf Z/\ell^n\mathbf Z)\) we apply \(\operatorname{Hom}(-,H)\) and \(\operatorname{Hom}(Rf_!(-),K)\) to (9.12) with \(N=M\). The maps \(b\mapsto\varepsilon\circ Rf_!(b)\) form a morphism between the two long exact sequences, which is bijective on the terms coming from \(U\), \(V\) and \(U\cap V\). By the five lemma it is bijective on \(\operatorname{Hom}(M,H)\), which means that \((H,\varepsilon)\) represents \(\operatorname{Hom}(Rf_!(-),K)\). Since such a pair is unique up to a unique isomorphism compatible with \(\varepsilon\), we get \(H\simeq f^!K\) compatibly with the counits. In Section \hyperref[bp41]{9.5} we use this gluing to realize the compatible counits with finite coefficients as a single map of towers.

On a vector bundle, the choice of a local trivialization transports (9.11) to a trace map. These maps do not depend on the trivialization, so they glue to a global map. This is the model case of the orientation below, because deformation to the normal bundle will reduce the dualizing complex of a smooth morphism to that of a vector bundle.

% object: proposition-purity-vector-bundles (blueprint/chap-Smooth-purity/sec-The-trace-map-of-affine-space/proposition-purity-vector-bundles/_section.md)
\phantomsection\label{bp21}%
\begin{bpstatementit}{Proposition~9.6}{ (Purity for vector bundles)}
Let \(\pi:V\to S\) be a vector bundle of rank \(d\). For \(K\in D^b_{\mathrm{ctf}}(S_{\mathrm{pro}\acute{e}t},\mathbf Z/\ell^n\mathbf Z)\), the trace map induces an isomorphism

\[
\pi^{*}K(d)[2d]\xrightarrow{\sim}\pi^{!}K.
\]
\end{bpstatementit}

% proof: proof/proposition-purity-vector-bundles
\begin{proof}[{Proof}]\phantomsection\label{bp22}%
We transport the trace map of affine space along local trivializations, show that the result does not depend on the trivialization, so that the local maps glue to a trace isomorphism \(\operatorname{Tr}_\pi\) for \(V\), and take the map \(\pi^*K(d)[2d]\to\pi^!K\) to be its adjoint.

\emph{Local trace maps.} Locally on \(S\) the bundle is affine space, where the trace map has already been constructed. Over an open \(U\subset S\) on which \(V\) is trivial, a trivialization \(\tau:V\times_S U\xrightarrow{\sim}\mathbf A^d_U\) transports (9.11) to an isomorphism

\[
\operatorname{Tr}_{\tau}:
R(\pi|_{V\times_S U})_!(\mathbf Z/\ell^n\mathbf Z)_{V\times_S U}(d)[2d]
\xrightarrow{\sim}(\mathbf Z/\ell^n\mathbf Z)_U.
\]

These maps commute with base change in \(U\), because the construction uses only \(c_1(\mathcal O_{\mathbf P^1_U}(1))\), the splitting (9.8) and localization at infinity: pullback preserves the Kummer class, and proper base change and localization identify the pullbacks of the constructions with those over the new base.

\emph{Independence of the trivialization.} A change of trivialization multiplies the trace map by a locally constant unit on \(\GL_d\), and this unit is \(1\). As two trivializations differ by a matrix with values in \(\GL_d\), it suffices to study one universal change of coordinates, parametrized by \(\GL_{d,S}\) itself, and then pull back. Let

\[
\theta:\mathbf A^d_{\GL_{d,S}}\longrightarrow\mathbf A^d_{\GL_{d,S}},
\qquad (g,v)\longmapsto(g,gv),
\]

be the universal linear change of coordinates, an automorphism over \(\GL_{d,S}\). It induces an automorphism of the compactly supported direct image of \((\mathbf Z/\ell^n\mathbf Z)_{\mathbf A^d_{\GL_{d,S}}}(d)[2d]\) to \(\GL_{d,S}\), which the trace map identifies with the constant sheaf \((\mathbf Z/\ell^n\mathbf Z)_{\GL_{d,S}}\) in degree zero, so that this automorphism is multiplication by a section

\[
u\in H^0\bigl((\GL_{d,S})_{\acute{e}t},
(\mathbf Z/\ell^n\mathbf Z)_{\GL_{d,S}}^{\times}\bigr),
\]

a locally constant function with values in \((\mathbf Z/\ell^n\mathbf Z)^{\times}\). It is therefore constant on each geometric fibre \(\GL_{d,\bar s}\), which is connected: for \(d\geq1\) it is a nonempty open subset of the affine space of \(d\times d\) matrices, and for \(d=0\) it is a point. As \(\theta\) is the identity at the identity matrix, which lies on every geometric fibre, \(u=1\).

The local trace maps agree on overlaps of trivializing opens. Indeed, let \(\tau\) and \(\tau'\) be two trivializations over \(U\). Then \(\tau'\circ\tau^{-1}\) is the pullback of \(\theta\) along the transition matrix \(U\to\GL_{d,S}\), so by base change the equality \(u=1\) implies \(\operatorname{Tr}_{\tau'}=\operatorname{Tr}_{\tau}\). Moreover, morphisms between source and target form a sheaf, as both are locally sheaves in degree zero. Hence the local maps glue over a trivializing cover to an isomorphism

\[
\operatorname{Tr}_{\pi}:
R\pi_!(\mathbf Z/\ell^n\mathbf Z)_V(d)[2d]
\xrightarrow{\sim}(\mathbf Z/\ell^n\mathbf Z)_S.
\]

\emph{The map \(\pi^*K(d)[2d]\to\pi^!K\).} For general \(K\), the projection formula and \(\operatorname{Tr}_\pi\) give

\[
R\pi_!\bigl(\pi^{*}K(d)[2d]\bigr)
\simeq K\otimes^{\mathbf L} R\pi_!(\mathbf Z/\ell^n\mathbf Z)_V(d)[2d]
\xrightarrow{\ \mathrm{id}_K\otimes\operatorname{Tr}_{\pi}\ }K.
\]

Its adjoint under \(R\pi_!\dashv\pi^{!}\) (Proposition proposition-exceptional-inverse-image-smooth-morphism) is the asserted map. We show that it is an isomorphism. This may be checked locally on \(V\), because \(\pi^!\) commutes with restriction to opens. Over each trivializing open, \(\operatorname{Tr}_\pi\) is \(\operatorname{Tr}_\tau\) by construction, so there the map is the isomorphism given by the adjunction for affine space following (9.11).
\end{proof}

% text: prose-09-10
\phantomsection\label{bp23}%
The trace maps of vector bundles commute with base change on \(S\), because the local trace maps do.


% object: sec-Finite-coefficients-3 (blueprint/chap-Smooth-purity/sec-Finite-coefficients-3/_section.md), content hidden
\subsection*{\texorpdfstring{9.4\quad Finite coefficients}{9.4 Finite coefficients}}
\phantomsection\label{bp24}%
\addcontentsline{toc}{subsection}{\texorpdfstring{9.4\quad Finite coefficients}{9.4 Finite coefficients}}
% text: prose-09-11
\phantomsection\label{bp25}%
Since \(f^!K\simeq f^*K\otimes^{\mathbf L}\omega_f\), smooth purity with finite coefficients comes down to a single global isomorphism \(\omega_f\simeq(\mathbf Z/\ell^n\mathbf Z)_X(d)[2d]\). We obtain it from the diagonal \(\Delta:X\to X\times_YX\), which is a section of the smooth projection \(p_1\). Deformation to the normal bundle moves \(\Delta\) to the zero section of the relative tangent bundle \(T_{X/Y}\), where Proposition \hyperref[bp21]{9.6} gives the isomorphism.

We use the diagonal because it is attached to \(f\) without choices, and because \(\omega_f\) is recovered from it: \(\omega_f\simeq\Delta^*\omega_{p_1}\) by base change. Near the diagonal, however, \(X\times_YX\) looks like a neighbourhood of the zero section in \(T_{X/Y}\) only to first order, and the deformation makes this precise without choosing a tubular neighbourhood. The comparison it yields, Lemma \hyperref[bp29]{9.9}, is determined by the requirement that it be the identity at \(t=1\).

The deformation argument works for any section of a smooth morphism. Let \(g:V\to S\) be smooth of pure relative dimension \(d\) and let \(s:S\to V\) be a section of \(g\) that is a closed immersion. Then \(s\) is a regular immersion of codimension \(d\), and we let \(N_s\) be its normal bundle, with projection \(\pi:N_s\to S\) and zero section \(z:S\to N_s\).

% object: definition-deformation-normal-bundle (blueprint/chap-Smooth-purity/sec-Finite-coefficients-3/definition-deformation-normal-bundle/_section.md)
\phantomsection\label{bp26}%
\begin{bpstatement}{Definition~9.7}{ (Deformation to the normal bundle)}
The \emph{deformation to the normal bundle} of \(s\) is the open subscheme \(\mathcal D_{s(S)}V\) of the blowup of \(V\times\mathbf A^1\) along \(s(S)\times\{0\}\) obtained by removing the strict transform of \(V\times\{0\}\).
\end{bpstatement}

% text: prose-09-12
\phantomsection\label{bp27}%
It comes with a smooth morphism \(\widetilde g:\mathcal D_{s(S)}V\to\mathbf A^1_S\) and a section \(\widetilde s:\mathbf A^1_S\to\mathcal D_{s(S)}V\) that is a closed immersion. Over \(t=1\) the pair \((\widetilde g,\widetilde s)\) restricts to \((g,s)\), and over \(t=0\) to \((\pi,z)\). To see this in local terms, let \(t\) be the coordinate of \(\mathbf A^1\), and let \(x_1,\ldots,x_d\) be functions on an open neighbourhood \(U\) of \(s(S)\) in \(V\) which are étale coordinates over \(S\) and generate the ideal of \(s(S)\). On the part of \(\mathcal D_{s(S)}V\) lying over \(U\times\mathbf A^1\), the quotients \(u_i=x_i/t\) are regular functions, so that \(x_i=tu_i\) there. At \(t=0\), then, the function \(u_i\) restricts to the linear function on \(N_s\) given by the class of \(x_i\) in the conormal sheaf. Since these classes form a basis of the conormal sheaf, the \(u_i\) restrict to linear coordinates on \(N_s\).

% object: example-deformation-origin-line (blueprint/chap-Smooth-purity/sec-Finite-coefficients-3/example-deformation-origin-line/_section.md)
\phantomsection\label{bp28}%
\begin{bpstatement}{Example~9.8}{ (The origin in the affine line)}
Let \(S=\operatorname{Spec}k\), \(V=\mathbf A^1_k\) with coordinate \(x\), and let \(s\) be the origin. The blowup of \(V\times\mathbf A^1=\mathbf A^2\), with coordinates \((x,t)\), at the origin is covered by two charts. On the chart with coordinates \((u,t)\) and \(x=tu\) the exceptional divisor is \(t=0\); on the chart with coordinates \((x,v)\) and \(t=xv\) it is \(x=0\), and the strict transform of the line \(V\times\{0\}=\{t=0\}\) is \(v=0\). Since a point of the second chart with \(v\neq0\) lies on the first chart, with \(u=1/v\), removing the strict transform leaves exactly the first chart:

\[
\mathcal D_{s(S)}V=\operatorname{Spec}k[u,t]\longrightarrow
V\times\mathbf A^1,\qquad (u,t)\longmapsto(tu,t),
\]

with \(\widetilde g\) the projection to \(t\) and \(\widetilde s\) the line \(u=0\). For \(t=c\neq0\) the fibre is the line \(V\) itself, with coordinate \(u=x/c\) and the origin as section. For \(t=0\) the fibre is the \(u\)-line, mapped to the origin of \(V\); it is the tangent line \(N_s=T_0\mathbf A^1\), and \(u\) is the linear coordinate given by the class of \(x\) in \(\mathfrak m/\mathfrak m^2\). So the family rescales the line by \(1/t\) and in the limit keeps only the first-order information at the section. Here the family is the product \(\mathbf A^1_u\times\mathbf A^1_t\), so \(\widetilde s^*\omega_{\widetilde g}\) is visibly constant along \(t\). In general this holds only étale locally, after a choice of coordinates, and the point of Lemma \hyperref[bp29]{9.9} is that the comparison normalized at \(t=1\) does not depend on that choice.
\end{bpstatement}

% object: lemma-section-normal-bundle (blueprint/chap-Smooth-purity/sec-Finite-coefficients-3/lemma-section-normal-bundle/_section.md)
\phantomsection\label{bp29}%
\begin{bpstatementit}{Lemma~9.9}{ (A section and its normal bundle)}
There is a canonical isomorphism

\[
s^{*}\omega_{g}\xrightarrow{\sim}z^{*}\omega_{\pi}.
\tag{9.13}
\]

It is obtained by restricting to \(t=0\) the unique isomorphism \(\widetilde s^{*}\omega_{\widetilde g}\to a^{*}s^{*}\omega_{g}\) that is the identity at \(t=1\), and inverting. Here \(a:\mathbf A^1_S\to S\) is the projection.
\end{bpstatementit}

% proof: proof/lemma-section-normal-bundle
\begin{proof}[{Proof}]\phantomsection\label{bp30}%
The fibres of \(\widetilde s^{*}\omega_{\widetilde g}\) at \(t=1\) and \(t=0\) are the two sides of (9.13), so it suffices to show that this shifted sheaf of rank one is pulled back from \(S\), étale locally on \(S\). In positive characteristic a locally constant sheaf on \(\mathbf A^1_S\) need not be pulled back from \(S\), so we use coordinates. Étale locally on \(S\), near the section the deformation is étale over \(\mathbf A^d_{\mathbf A^1_S}=\mathbf A^d_S\times\mathbf A^1\), a family pulled back from \(S\), and hence so is \(\widetilde s^{*}\omega_{\widetilde g}[-2d]\). We then normalize the local comparisons at \(t=1\); since the geometric fibres of \(\mathbf A^1_S\to S\) are connected, this makes them agree on overlaps. Finally, they glue because all objects involved are shifted locally constant sheaves.

\emph{Step 1: the objects.} The restrictions of \(\widetilde s^{*}\omega_{\widetilde g}\) to \(t=1\) and \(t=0\) are \(s^{*}\omega_{g}\) and \(z^{*}\omega_{\pi}\), by Proposition proposition-exceptional-inverse-image-smooth-morphism applied to the inclusions of these fibres. Proposition proposition-exceptional-inverse-image-smooth-morphism shows that both \(\widetilde s^{*}\omega_{\widetilde g}[-2d]\) and \(s^{*}\omega_{g}[-2d]\) are locally constant sheaves, free of rank one over \(\mathbf Z/\ell^n\mathbf Z\). Hence morphisms between two such complexes are morphisms of sheaves, and they glue.

\emph{Step 2: local construction.} Étale locally on \(S\) we reduce to the zero section of affine space, whose deformation is pulled back from \(S\). After passing to an étale cover \(\{S_\alpha\to S\}_\alpha\) we may choose, for the base change of the section \(s_\alpha:S_\alpha\to V\times_SS_\alpha\), a neighbourhood \(U_\alpha\) of \(s_\alpha(S_\alpha)\) and an étale map \(e_\alpha:U_\alpha\to\mathbf A^d_{S_\alpha}\) sending \(s_\alpha\) to the zero section. We may assume that the inverse image of the zero section under \(e_\alpha\) is exactly \(s_\alpha(S_\alpha)\). Indeed, this inverse image is separated and étale over \(S_\alpha\), so \(s_\alpha(S_\alpha)\) is open and closed in it, and we remove the complement, which is closed in \(U_\alpha\). Then \(e_\alpha\) is an étale map of pairs whose square of centers is cartesian. Since blowing up and removing the strict transform commute with étale base change, \(e_\alpha\) induces an étale map from a neighbourhood of the deformed section in \(\mathcal D_{s(S)}V\) to the deformation space of the zero section of \(\mathbf A^d_{S_\alpha}\). The latter is \(\mathbf A^d_{\mathbf A^1_{S_\alpha}}\), with coordinates \(u_1,\ldots,u_d\) and \(t\), and its map to \(\mathbf A^d_{S_\alpha}\times\mathbf A^1\) sends \((u,t)\) to \((tu,t)\), that is, \(x_i=tu_i\), and the deformed zero section is \(u=0\).

The adjunction for affine space identifies the relative dualizing complex of the projection \(\mathbf A^d_{\mathbf A^1_{S_\alpha}}\to\mathbf A^1_{S_\alpha}\) with the constant sheaf \((\mathbf Z/\ell^n\mathbf Z)(d)[2d]\). Pulling back along the étale map just constructed, and composing adjunctions for it, gives the same identification for \(\omega_{\widetilde g}\) near the deformed section. Similarly, the chart \(e_\alpha\) gives the analogous identification of \(s^{*}\omega_{g}\) over \(S_\alpha\). Composing the two identifications we obtain an isomorphism over \(\mathbf A^1_{S_\alpha}\)

\[
\widetilde s^{*}\omega_{\widetilde g}\xrightarrow{\sim}
a^{*}s^{*}\omega_{g}.
\]

At \(t=1\) it restricts to an automorphism of \(s^{*}\omega_{g}\) over \(S_\alpha\). Composing with the pullback along \(a\) of the inverse of this automorphism, we may assume that the local isomorphism is the identity at \(t=1\).

\emph{Step 3: gluing.} The local isomorphisms from different charts agree on overlaps, because of the normalization at \(t=1\). Over \(S_\alpha\times_SS_\beta\) the two local isomorphisms that are the identity at \(t=1\) differ by an automorphism of the locally constant sheaf \(a^{*}s^{*}\omega_g[-2d]\) of rank one, that is, by a section of a locally constant sheaf of units of \(\mathbf Z/\ell^n\mathbf Z\). The locus where this section equals \(1\) is open and closed, and it contains the point \(t=1\) of each geometric fibre \(\mathbf A^1_{\bar b}\). As these fibres are connected, the section is \(1\) at every geometric point. So the local isomorphisms glue to a unique isomorphism over \(\mathbf A^1_S\) that is the identity at \(t=1\). Restricting it to \(t=0\) gives an isomorphism \(z^{*}\omega_{\pi}\simeq s^{*}\omega_g\), whose inverse is (9.13).
\end{proof}

% text: prose-09-13
\phantomsection\label{bp31}%
The orientation of \(\omega_f\) will be the comparison (9.13) for the diagonal. Two things have to be checked about it: on a coordinate chart it must be the identification given by coordinates, and it must commute with base change, as property (2) of Theorem \hyperref[bp11]{9.3} requires. Since a chart is an étale map that preserves the section, both checks need to transport (9.13) along base change on \(S\) and along étale maps that preserve the section. In both cases we use the uniqueness in Lemma \hyperref[bp29]{9.9}: any construction that is the identity at \(t=1\) produces the comparison of that lemma.

% object: lemma-naturality-under-base-change-etale (blueprint/chap-Smooth-purity/sec-Finite-coefficients-3/lemma-naturality-under-base-change-etale/_section.md)
\phantomsection\label{bp32}%
\begin{bpstatementit}{Lemma~9.10}{ (Naturality under base change and étale maps)}
Write \(\alpha_{g,s}:s^*\omega_g\xrightarrow{\sim}z^*\omega_\pi\) for the comparison of Lemma \hyperref[bp29]{9.9}.

\begin{enumerate}
\item Let \(h:S'\to S\) be any morphism for which the schemes below satisfy the standing hypotheses. Put \(V'=V\times_SS'\), and denote the base changes of the maps by \(g',s',\pi',z'\). There are canonical identifications \[
   N_{s'}\simeq N_s\times_SS',\qquad
   \mathcal D_{s'(S')}V'
   \simeq\mathcal D_{s(S)}V\times_SS'.
   \] Under these identifications and the exceptional base change isomorphisms of Proposition proposition-exceptional-inverse-image-smooth-morphism, the following square commutes: \[
   \begin{CD}
   h^*s^*\omega_g @>{h^*\alpha_{g,s}}>> h^*z^*\omega_\pi\\
   @V{\sim}VV @VV{\sim}V\\
   (s')^*\omega_{g'} @>{\alpha_{g',s'}}>> (z')^*\omega_{\pi'}.
   \end{CD}
   \]
\item Let \(e:V'\to V\) be an étale \(S\)-morphism and let \(s':S\to V'\) be a closed section such that \(es'=s\) and \(e^{-1}(s(S))=s'(S)\) scheme-theoretically. Put \(g'=ge\). The induced map \(N_{s'}\to N_s\) is a linear isomorphism, and the induced map of deformation spaces is étale. The identifications \[
   (s')^*\omega_{g'}\simeq s^*\omega_g,\qquad
   (z')^*\omega_{\pi'}\simeq z^*\omega_\pi
   \] obtained by composing adjunctions carry \(\alpha_{g',s'}\) to \(\alpha_{g,s}\).
\end{enumerate}

These comparisons respect identities, composition and interchange of base change with an étale map of pairs.
\end{bpstatementit}

% text: prose-09-14
\phantomsection\label{bp33}%
In part (1) the morphism \(h\) is arbitrary. Blowups in general commute only with flat base change, but by its description in local coordinates, the deformation space of a section of a smooth morphism commutes with arbitrary base change.

% proof: proof/lemma-naturality-under-base-change-etale
\begin{proof}[{Proof of Lemma 9.10}]\phantomsection\label{bp34}%
The geometric identifications follow from the coordinate description of the deformation space. Every comparison of dualizing complexes is then the identity at \(t=1\), and so the uniqueness in Lemma \hyperref[bp29]{9.9} identifies it with the comparison of that lemma.

\emph{The geometry.} Since \(g\) is smooth and \(s\) is a section, \(N_s\simeq s^*T_{V/S}\). Since formation of relative differentials for a smooth map commutes with arbitrary base change, and an étale map induces an isomorphism on relative tangent bundles, this proves both assertions about normal bundles.

For deformation spaces, we use the cartesian étale coordinate chart \(U\to\mathbf A^d_S\) of the proof of Lemma \hyperref[bp29]{9.9}. There the deformation is

\[
\mathcal D_{s(S)}U
\simeq U\times_{\mathbf A^d_S}\mathbf A^d_{\mathbf A^1_S},
\qquad x_i=t u_i,
\]

because this is the chart of the blowup along \((t,x_1,\ldots,x_d)\) on which the pullback of the ideal is generated by \(t\), and this description commutes with every change of base. Away from the centre, the deformation is \((V\setminus s(S))\times\mathbf G_m\), which also commutes with base change. After passing to an étale cover of \(S\) if necessary, these opens cover the deformation space, and their identifications glue to the canonical one. For an étale map of pairs as in (2), we obtain the étale map of deformation spaces by flat base change for the blowup and for the removal of the strict transform, and the coordinate description identifies its special fibre with the map of normal bundles.

\emph{The comparison of dualizing complexes.} In both parts the strategy is the same: transport the comparison over \(\mathbf A^1_S\) along the geometric identification, check that the result is still the identity at \(t=1\), and invoke uniqueness. Write

\[
\chi_{g,s}:\widetilde s^*\omega_{\widetilde g}\xrightarrow{\sim}a^*s^*\omega_g
\]

for the isomorphism of Lemma \hyperref[bp29]{9.9} that is the identity at \(t=1\), so that \(\alpha_{g,s}\) is the inverse of its restriction to \(t=0\). For (1), we apply Proposition proposition-exceptional-inverse-image-smooth-morphism to \(g\), \(\pi\) and \(\widetilde g\). These exceptional base change isomorphisms commute with restriction to \(t=1\) and \(t=0\), because they are adjoint transposes of base change for compactly supported direct image applied to the same counit, and this base change respects composition of cartesian squares. Transported through these isomorphisms, the pullback of \(\chi_{g,s}\) is therefore a comparison for \((g',s')\) that is the identity at \(t=1\), and by uniqueness it equals \(\chi_{g',s'}\). Restricting to \(t=0\) and inverting proves (1).

For (2), composing \(Re_!\dashv e^*\) with \(Rg_!\dashv g^!\) gives \(e^*\omega_g\simeq\omega_{g'}\), and likewise for the induced étale map of deformation spaces. On the fibres \(t=1\) and \(t=0\) these identifications are the ones for the original pairs and for the normal bundles, since the counit of a composite adjunction is the composite of the counits. They transport \(\chi_{g,s}\) to a comparison that is the identity at \(t=1\), which is \(\chi_{g',s'}\) by uniqueness.

For the coherence assertions we use the same uniqueness once more. The geometric identifications respect composition, and so do the identifications obtained from adjunctions, as we see by transposing composite counits. Therefore each composite in the coherence assertions is a comparison that is the identity at \(t=1\), and by uniqueness any two of them agree.
\end{proof}

% text: prose-09-15
\phantomsection\label{bp35}%
We apply Lemma \hyperref[bp29]{9.9} to the diagonal. Let \(p_1,p_2:X\times_YX\to X\) be the projections. Since \(f\) is separated, \(\Delta\) is a closed immersion, and it is a section of the smooth morphism \(p_1\). Its conormal sheaf is \(\Omega^1_{X/Y}\), so its normal bundle is the relative tangent bundle \(T_{X/Y}\), and we write \(\pi:T_{X/Y}\to X\) for its projection and \(z:X\to T_{X/Y}\) for its zero section.

% object: proposition-relative-dualizing-complex (blueprint/chap-Smooth-purity/sec-Finite-coefficients-3/proposition-relative-dualizing-complex/_section.md)
\phantomsection\label{bp36}%
\begin{bpstatementit}{Proposition~9.11}{ (The relative dualizing complex)}
Let \(f:X\to Y\) be smooth of pure relative dimension \(d\). There is a canonical isomorphism

\[
\omega_{f}\xrightarrow{\sim}(\mathbf Z/\ell^n\mathbf Z)_{X}(d)[2d].
\tag{9.14}
\]
\end{bpstatementit}

% proof: proof/proposition-relative-dualizing-complex
\begin{proof}[{Proof}]\phantomsection\label{bp37}%
We define the isomorphism as a chain of three steps: restrict \(\omega_f\) to the diagonal of \(X\times_YX\), move the diagonal to the zero section of \(T_{X/Y}\), and use the orientation of a vector bundle.

The projection \(p_1\) is the base change of \(f\) along \(f\), so Proposition proposition-exceptional-inverse-image-smooth-morphism shows that \(\omega_{p_1}\simeq p_2^{*}\omega_f\). Since \(p_2\circ\Delta\) is the identity of \(X\), we get the first isomorphism in

\[
\begin{aligned}
\omega_{f}
&\simeq L\Delta^{*}\omega_{p_1}
\xrightarrow{\sim}z^{*}\omega_{\pi}
\xrightarrow{\sim}z^{*}(\mathbf Z/\ell^n\mathbf Z)_{T_{X/Y}}(d)[2d]
\simeq(\mathbf Z/\ell^n\mathbf Z)_{X}(d)[2d].
\end{aligned}
\]

The second isomorphism is Lemma \hyperref[bp29]{9.9} for the section \(\Delta\) of \(p_1\), and the third is the inverse of the isomorphism of Proposition \hyperref[bp21]{9.6} at the constant sheaf, applied to the vector bundle \(T_{X/Y}\) over \(X\). The last one holds because the pullback of a constant sheaf is constant. Since each step is attached to \(f\) without choices, the composite is canonical.
\end{proof}

% text: prose-09-16
\phantomsection\label{bp38}%
The orientation (9.14) is defined without coordinates, but on a coordinate chart it agrees with the identification given by the trace map of affine space (9.11). This is what makes the map \(f^*K(d)[2d]\to f^!K\) below the expected one, and the trace map in Section \hyperref[bp74]{9.6} relies on it as well.

We prove the agreement for a vector bundle, where both orientations can be computed, and reduce a coordinate chart to the vector bundle \(\mathbf A^d_Y\) by naturality.

\emph{Vector bundles.} Let \(\pi:V\to S\) be a vector bundle of rank \(d\). Here the orientation obtained from the diagonal is the orientation of Proposition \hyperref[bp21]{9.6}, because the construction from the diagonal moves the diagonal to the zero section and then deforms the zero section to itself, and neither step changes integration along the fibres. For the first step, regard \(V\times_SV\) as a vector bundle over the first factor. The automorphism \((v,w)\mapsto(v,w-v)\) of \(V\times_SV\) over the first factor takes the diagonal to the zero section, and it preserves integration in the second factor. To see this, let the translation vector vary over \(V\). Its effect on the top compactly supported direct image, identified with the constant sheaf by integration, is then multiplication by a locally constant unit on \(V\), and this unit is \(1\) at the zero vector, where the translation is the identity. Since each geometric fibre of \(V\to S\) is connected, it is \(1\) everywhere.

Next, the deformation to the normal bundle of the zero section of \(V\) is \(V\times\mathbf A^1\), with the map \(V\times\mathbf A^1\to V\times\mathbf A^1\), \((w,t)\mapsto(tw,t)\), to the original family. This map is defined by scalar multiplication, without a trivialization, and in a local trivialization it is \(x_i=tu_i\), as in the coordinate description above. Integration on this product family is the pullback of the original integration. Therefore the comparison of Lemma \hyperref[bp29]{9.9}, which is the identity at \(t=1\), is the identity under these integrations, in particular at \(t=0\). Combined with the translation, this identifies the orientation obtained from the diagonal for \(\pi\) with the orientation of the vector bundle.

\emph{Coordinate charts.} A chart \(e:U\to\mathbf A^d_Y\) relates the diagonal of \(U\) to the diagonal of \(\mathbf A^d_Y\), where the case of vector bundles applies, so we have to transport the orientation along \(e\). By Lemma \hyperref[bp32]{9.10} it suffices to see that every ingredient of the orientation of a chart obtained from the diagonal is pulled back from affine space. So let \(U\subset X\) be open with an étale map \(e:U\to\mathbf A^d_Y\) over \(Y\). Then the étale map \(e\times e:U\times_YU\to\mathbf A^d_Y\times_Y\mathbf A^d_Y\) carries the diagonal of \(U\) into the diagonal of \(\mathbf A^d_Y\), whose inverse image is \(U\times_{\mathbf A^d_Y}U\). In it the diagonal of \(U\) is open and closed, because \(e\) is separated and étale, and removing the other components, as in the proof of Lemma \hyperref[bp29]{9.9}, we obtain a cartesian square of centers. The induced map of normal bundles is then the differential \(T_{U/Y}\simeq e^*T_{\mathbf A^d_Y/Y}\), a pullback of vector bundles. Likewise, the dualizing complex is pulled back from the affine model, because by composing adjunctions the exceptional inverse image along an étale map is its ordinary inverse image. For the comparison obtained from the deformation, we factor \(e\times e\) as

\[
U\times_YU\xrightarrow{\mathrm{id}_U\times e}
U\times_Y\mathbf A^d_Y\xrightarrow{e\times\mathrm{id}}
\mathbf A^d_Y\times_Y\mathbf A^d_Y.
\]

The second arrow is the base change of the first projection along \(e\), and it carries the diagonal section to the graph of \(e\) over \(U\), so Lemma lemma-naturality-under-base-change-etale applies to it. Part (2), over the common base \(U\), applies to the first arrow after we remove the other components of the inverse image of that graph. In this factorization the map on special fibres is the displayed differential, and the comparison at \(t=1\) is the one obtained from the adjunction in coordinates. Moreover, integration on vector bundles is preserved by the linear isomorphism of normal bundles. By the case of vector bundles, we conclude that the orientation on \(U\) obtained from the diagonal is the orientation of affine space given by its adjunction in coordinates.

Combining the orientation with the isomorphism \(f^!K\simeq f^*K\otimes^{\mathbf L}\omega_f\), we obtain the isomorphism \(\gamma_f\) of Theorem \hyperref[bp11]{9.3} with finite coefficients (Proposition \hyperref[bp39]{9.12}). We state the result for an arbitrary finite admissible ring.

% object: proposition-smooth-purity-finite-coefficients (blueprint/chap-Smooth-purity/sec-Finite-coefficients-3/proposition-smooth-purity-finite-coefficients/_section.md)
\phantomsection\label{bp39}%
\begin{bpstatementit}{Proposition~9.12}{ (Smooth purity with finite coefficients)}
Let \(A\) be any finite admissible coefficient ring, and let \(f:X\to Y\) be smooth of pure relative dimension \(d\). For \(K\in D^b_{\mathrm{ctf}}(Y_{\mathrm{pro}\acute{e}t},A)\) there is a canonical isomorphism, natural in \(K\),

\[
f_A^{*}K(d)[2d]
\xrightarrow{\ \gamma_{f,A}(K)\ }f_A^{!}K.
\tag{9.15}
\]
\end{bpstatementit}

% proof: proof/proposition-smooth-purity-finite-coefficients
\begin{proof}[{Proof}]\phantomsection\label{bp40}%
The calculation so far was made with coefficients \(\mathbf Z/\ell^n\mathbf Z\), and every ingredient in it extends along \(\mathbf Z/\ell^n\mathbf Z\to A\), where we choose \(n\) with \(\ell^nA=0\). Once the Kummer and diagonal classes are extended, the calculation works with \(A\) throughout. In particular, the splitting for the projective line has two free summands, and its integration and triangle identities keep the same normalization, since (9.9) is an identity of endomorphisms of a constant sheaf and is preserved by extension of scalars. So the local adjunctions compute \(f_A^!\), and the arguments for vector bundles and deformations then show that \(\omega_{f,A}\simeq A_X(d)[2d]\). We define \(\gamma_{f,A}(K)\) as the composite

\[
f_A^{*}K\otimes^{\mathbf L}A_X(d)[2d]
\xrightarrow{\sim}f_A^{*}K\otimes^{\mathbf L}\omega_{f,A}
\xrightarrow{\sim}f_A^{!}K.
\]

These are the isomorphisms given by the orientation and by \(\mu_{f,K}\), and both are natural in \(K\). By (9.17) below, \(\gamma_{f,A}\) does not depend on the choice of \(n\) and is compatible with maps of finite rings.
\end{proof}


% object: sec-ell-adic-coefficients-4 (blueprint/chap-Smooth-purity/sec-ell-adic-coefficients-4/_section.md), content hidden
\subsection*{\texorpdfstring{9.5\quad $\ell$-adic coefficients}{9.5 ell-adic coefficients}}
\phantomsection\label{bp41}%
\addcontentsline{toc}{subsection}{\texorpdfstring{9.5\quad $\ell$-adic coefficients}{9.5 ell-adic coefficients}}
% text: prose-09-17
\phantomsection\label{bp42}%
Let \(f:X\to Y\) be smooth of pure relative dimension \(d\). For \(\mathcal O_E\)-coefficients we want an isomorphism of functors \(f^*(-)(d)[2d]\simeq f^!\) on \(\Dbc(-,\mathcal O_E)\). Here \(f^!_{\mathcal O_E}\) is defined by Brown representability, not as a limit of the finite \(f^!\). The reduction principle of Section \bpout{7.3} does not apply, because it checks that a given morphism is an isomorphism, and we have no morphism \(f^*(-)(d)[2d]\to f^!\) with \(\mathcal O_E\)-coefficients yet. Instead we deduce Theorem \hyperref[bp11]{9.3} for \(\mathcal O_E\), \(E\) and \(\overline{\mathbf Q}_\ell\) from Proposition \hyperref[bp39]{9.12} by building the right adjoint ourselves. The key input is that the finite isomorphisms \(\gamma_{f,A}\) commute with change of finite coefficients (9.17), which lets us construct an explicit right adjoint of \(Rf_{!,\mathcal O_E}\) from towers (Proposition \hyperref[bp63]{9.19}). Inverting \(\ell\) then gives the cases \(E\) and \(\overline{\mathbf Q}_\ell\) (Proposition \hyperref[bp68]{9.20} and Theorem \hyperref[bp71]{9.21}). Each right adjoint we build is identified with \(f^!\) by the uniqueness of adjoints, normalized by the counits.

% object: proposition-sp-unique-adjoint (blueprint/chap-Smooth-purity/sec-ell-adic-coefficients-4/proposition-sp-unique-adjoint/_section.md)
\phantomsection\label{bp43}%
\begin{bpstatementit}{Proposition~9.13}{ (Uniqueness of the right adjoint of $Rf_!$)}
Let \(\Lambda\) be an admissible coefficient ring and let \(J\) be a right adjoint of \(Rf_!:D^b_{\mathrm{ctf}}(X,\Lambda)\to D^b_{\mathrm{ctf}}(Y,\Lambda)\) with counit \(\varepsilon^J\). There is a unique natural isomorphism \(\rho:f^!\xrightarrow{\sim}J\) with \(\varepsilon^J(K)\circ Rf_!(\rho(K))=\varepsilon_f(K)\).
\end{bpstatementit}

% proof: proof/proposition-sp-unique-adjoint
\begin{proof}[{Proof}]\phantomsection\label{bp44}%
Let \(\rho(K)\) be the \(J\)-transpose of \(\varepsilon_f(K)\) and \(\rho'(K)\) the \(f^!\)-transpose of \(\varepsilon^J(K)\). Since \(\rho'\rho\) and the identity have the same transpose, they are equal, and by symmetry \(\rho\rho'=\mathrm{id}\). For uniqueness, note that a map compatible with the counits has \(J\)-transpose \(\varepsilon_f(K)\), so it is \(\rho(K)\).
\end{proof}

% object: subsec-Change-of-finite-coefficients (blueprint/chap-Smooth-purity/sec-ell-adic-coefficients-4/subsec-Change-of-finite-coefficients/_section.md), content hidden
\subsubsection*{\texorpdfstring{9.5.1\quad Change of finite coefficients}{9.5.1 Change of finite coefficients}}
\phantomsection\label{bp45}%
\addcontentsline{toc}{subsubsection}{\texorpdfstring{9.5.1\quad Change of finite coefficients}{9.5.1 Change of finite coefficients}}
% text: prose-09-18
\phantomsection\label{bp46}%
The levels \(\mathcal O_E/\ell^n\) of a tower are related by extension of scalars. The first thing to know is therefore how \(Rf_!\) and \(f^!\) behave under a map of finite rings \(A\to B\). The comparison for \(f^!\) is the transpose of a counit, and the counits commute with \(-\otimes^{\mathbf L}B\), because on a chart they are built from the trace of the projective line and sums over fibres.

Let \(A\to B\) be a map of finite commutative rings killed by a power of \(\ell\).

% object: definition-sp-scalar-extension (blueprint/chap-Smooth-purity/sec-ell-adic-coefficients-4/subsec-Change-of-finite-coefficients/definition-sp-scalar-extension/_section.md)
\phantomsection\label{bp47}%
\begin{bpstatement}{Definition~9.14}{ (Extension of scalars)}
For a scheme \(T\) and \(K\in D^b_{\mathrm{ctf}}(T_{\mathrm{pro}\acute{e}t},A)\) put \(K_B=K\otimes^{\mathbf L}_{A_T}B_T\in D^b_{\mathrm{ctf}}(T_{\mathrm{pro}\acute{e}t},B)\), and for a morphism \(u\) put \(u_B=u\otimes^{\mathbf L}\mathrm{id}_{B_T}\). For a further map \(B\to C\) we have \((K_B)_C\simeq K_C\) by associativity of the tensor product.
\end{bpstatement}

% text: prose-09-19
\phantomsection\label{bp48}%
With this notation, the comparison for \(f^!\) is characterized by its counit: it is the unique map that makes the counit over \(B\) agree with the extended counit over \(A\).

% object: lemma-sp-reduction (blueprint/chap-Smooth-purity/sec-ell-adic-coefficients-4/subsec-Change-of-finite-coefficients/lemma-sp-reduction/_section.md)
\phantomsection\label{bp49}%
\begin{bpstatementit}{Lemma~9.15}{ (Extension of scalars commutes with $Rf_!$ and $f^!$)}
For \(M\in D^b_{\mathrm{ctf}}(X_{\mathrm{pro}\acute{e}t},A)\) and \(K\in D^b_{\mathrm{ctf}}(Y_{\mathrm{pro}\acute{e}t},A)\) there are natural isomorphisms \(b(M):(Rf_{!,A}M)_B\xrightarrow{\sim}Rf_{!,B}(M_B)\) and \(c_{A,B}(K):(f_A^!K)_B\xrightarrow{\sim}f_B^!(K_B)\), where \(c_{A,B}(K)\) is determined by

\[
\varepsilon_B(K_B)\circ Rf_{!,B}\bigl(c_{A,B}(K)\bigr)
=(\varepsilon_A(K))_B\circ b(f_A^!K)^{-1}.
\tag{9.16}
\]

For \(A\to B\to C\), the comparisons for \(A\to C\) are the composites of those for \(A\to B\) and \(B\to C\).
\end{bpstatementit}

% proof: proof/lemma-sp-reduction
\begin{proof}[{Proof}]\phantomsection\label{bp50}%
The comparison \(b\) for \(Rf_!\) is a direct computation, and \(c_{A,B}\) is then determined by the adjunction. To see that \(c_{A,B}\) is invertible we use the explicit form of \(f^!\) on a chart, where it is an ordinary inverse image.

The map \(b\) is the change of coefficients for \(Rf_!\) (Proposition \bpout{7.10}). To see that it is an isomorphism, note that extension by zero commutes with \(-\otimes^{\mathbf L}_AB\), and for \(R\bar f_*\) we resolve \(B\) by finite free \(A\)-modules and use the flat models of Lemma \bpout{3.13}. The comparison holds on each free term, which is a finite direct sum of copies of \(A\), and since \(R\bar f_*\) has finite cohomological dimension, only finitely many terms of the resolution contribute in each degree, so the comparison holds for every \(B\). Next, we define \(c_{A,B}(K)\) as the transpose of the right-hand side of (9.16). Whether it is an isomorphism is a local question, so we may restrict to a chart \(f|_U=p_d\circ e\) of Proposition proposition-exceptional-inverse-image-smooth-morphism. There the counit is built from the trace of the projective line and sums over the fibres of \(e\), which commute with extension of coefficients (Proposition \bpout{6.7}). On the chart, therefore, \(c_{A,B}(K)\) is the transpose of the comparison for pullback \((e^*p_d^*K(d)[2d])_B\simeq e^*p_d^*(K_B)(d)[2d]\), which is an isomorphism. Hence \(c_{A,B}(K)\) is an isomorphism on every chart. Transitivity holds because both sides have the same transpose: for \(A\to B\to C\), the counit identity (9.16) for \(A\to C\) is obtained by applying the one for \(B\to C\) after the one for \(A\to B\).
\end{proof}

% text: prose-09-20
\phantomsection\label{bp51}%
It follows that the isomorphisms \(\gamma_{f,A}\) are compatible with change of coefficients: for \(K\in D^b_{\mathrm{ctf}}(Y_{\mathrm{pro}\acute{e}t},A)\),

\[
c_{A,B}(K)\circ(\gamma_{f,A}(K))_B=\gamma_{f,B}(K_B),
\tag{9.17}
\]

the sources being identified by \((f_A^*K(d)[2d])_B\simeq f_B^*(K_B)(d)[2d]\). Indeed, \(\gamma_{f,A}\) is the map \(\mu_{f,K}\) composed with the inverse of the orientation, and each factor commutes with change of coefficients.

The map \(\mu_{f,K}\) is built from the projection formula and the counit. The projection formula commutes with extension of scalars, which is itself a tensor product, and (9.16) relates the counits over \(A\) and \(B\). Hence \(\mu_{f,K}\) commutes with change of coefficients.

The orientation is a map between invertible sheaves placed in degree \(-2d\), so its compatibility is local. For vector bundles it follows from the construction via the projective line (Proposition \hyperref[bp21]{9.6}), since the Kummer class and the splitting (9.8) are compatible with extension of scalars. On the deformation family we use the uniqueness in Lemma \hyperref[bp29]{9.9} over the base of the family. It applies once we know that the extended comparison is a comparison of the kind considered there, and this holds because \(c\) commutes with smooth exceptional base change and with composition with étale maps. Indeed, both sides are transposes of counits, and by the proof of Lemma \hyperref[bp49]{9.15} they agree on charts. Since the original comparison is the identity at \(t=1\), so is the extended one, and by uniqueness the extended comparison is the comparison of Lemma \hyperref[bp29]{9.9} over \(B\).


% object: subsec-Integral-coefficients-2 (blueprint/chap-Smooth-purity/sec-ell-adic-coefficients-4/subsec-Integral-coefficients-2/_section.md), content hidden
\subsubsection*{\texorpdfstring{9.5.2\quad Integral coefficients}{9.5.2 Integral coefficients}}
\phantomsection\label{bp52}%
\addcontentsline{toc}{subsubsection}{\texorpdfstring{9.5.2\quad Integral coefficients}{9.5.2 Integral coefficients}}
% text: prose-09-21
\phantomsection\label{bp53}%
Fix a finite extension \(E/\mathbf Q_\ell\). An object \(M\) of \(\Dbc(T,\mathcal O_E)\) is derived complete, \(M\simeq R\varprojlim_nM/\ell^n\), and the finite theory applies to each \(M/\ell^n\). The units, counits and transition maps at the finite levels, however, live in the different categories \(D(T_{\mathrm{pro}\acute{e}t},\mathcal O_E/\ell^n)\), and a family of them that commutes with the transition maps only up to unspecified isomorphisms does not define a morphism between the derived limits. We therefore work in the derived category \(D(\operatorname{Tow}(T))\) of towers \((C_n)_n\), where \(C_n\) is a sheaf of \(\mathcal O_{E,T}/\ell^n\)-modules on \(T_{\mathrm{pro}\acute{e}t}\) and the transition maps are linear over \(\mathcal O_{E,T}/\ell^{n+1}\to\mathcal O_{E,T}/\ell^n\). These are the modules over \((\mathcal O_{E,T}/\ell^n)_n\) in the topos of inverse systems [Stacks, \href{https://stacks.math.columbia.edu/tag/0940}{0940}]. In this category a compatible family of maps is one morphism. Evaluation \(C_\bullet\mapsto C_n\) is exact, and the evaluations jointly detect isomorphisms, but they do not detect equality of morphisms. This is why the maps below are always compared in \(D(\operatorname{Tow}(T))\). A tower is \emph{constructible} if each \(C_n\) lies in \(D^b_{\mathrm{ctf}}(T_{\mathrm{pro}\acute{e}t},\mathcal O_E/\ell^n)\) with bounds independent of \(n\), and \emph{normalized} if its reduction maps \(C_n\otimes^{\mathbf L}\mathcal O_{E,T}/\ell^m\to C_m\) are isomorphisms for all \(m\le n\).

We need two facts about integral complexes. Part (1) of the next proposition gives the uniform bounds in the definition of a constructible tower, and part (2) shows that the derived limit of such a tower is again an object of \(\Dbc\).

% object: proposition-sp-integral-towers (blueprint/chap-Smooth-purity/sec-ell-adic-coefficients-4/subsec-Integral-coefficients-2/proposition-sp-integral-towers/_section.md)
\phantomsection\label{bp54}%
\begin{bpstatementit}{Proposition~9.16}{ (Integral constructible complexes as towers)}
\begin{enumerate}
\item For \(K\in\Dbc(T,\mathcal O_E)\) there are integers \(a\le b\) such that every reduction \(K/\ell^n\) lies in \(D^b_{\mathrm{ctf}}(T_{\mathrm{pro}\acute{e}t},\mathcal O_E/\ell^n)\), with cohomology and Tor amplitude in \([a,b]\).
\item If \(K\in D(T_{\mathrm{pro}\acute{e}t},\mathcal O_E)\) is derived complete and \(K/\ell\) lies in \(D^b_{\mathrm{ctf}}(T_{\mathrm{pro}\acute{e}t},\mathcal O_E/\ell)\), then \(K\in\Dbc(T,\mathcal O_E)\).
\end{enumerate}
\end{bpstatementit}

% proof: proof/proposition-sp-integral-towers
\begin{proof}[{Proof}]\phantomsection\label{bp55}%
For (1), by Lemma \bpout{5.19}, \(K\) is a finite extension of shifted sheaves \(a_!C\) with \(C\) of Tor amplitude \([-1,0]\) on a stratum. For a single such piece the reductions \((a_!C)/\ell^n=a_!(C/\ell^n)\) are constructible with Tor amplitude in \([-1,0]\) for every \(n\), since reduction commutes with extension by zero. Moreover, Tor amplitude and constructibility of the reductions are preserved by extensions and shifts, and since the number of steps and the shifts do not depend on \(n\), the bounds do not either. For (2), we reduce further from \(\ell\) to the uniformizer. The reduction \(K/\pi_E=(K/\ell)\otimes^{\mathbf L}_{\mathcal O_E/\ell}\mathcal O_E/\pi_E\) is bounded constructible, because \(K/\ell\) has finite Tor amplitude over \(\mathcal O_E/\ell\). Since \(K\) is derived complete, \(K\in\Dbc(T,\mathcal O_E)\) by the criterion (2)\(\Rightarrow\)(1) of Proposition \bpout{3.11}.
\end{proof}

% text: prose-09-22
\phantomsection\label{bp56}%
Nothing is lost by passing from \(\mathcal O_E\)-complexes to towers: a constructible \(\mathcal O_E\)-complex is recovered from its tower of reductions as the derived limit, and the towers that arise in this way are exactly the normalized constructible ones.

% object: lemma-sp-tower-adjunction-foundation (blueprint/chap-Smooth-purity/sec-ell-adic-coefficients-4/subsec-Integral-coefficients-2/lemma-sp-tower-adjunction-foundation/_section.md)
\phantomsection\label{bp57}%
\begin{bpstatementit}{Lemma~9.17}{ (Reduction to towers is left adjoint to the derived limit)}
The functor \(\operatorname{red}_TM=(M/\ell^n)_n\) is left adjoint to \(R\varprojlim:D(\operatorname{Tow}(T))\to D(T_{\mathrm{pro}\acute{e}t},\mathcal O_E)\), with unit the derived completion map \(M\to R\varprojlim_nM/\ell^n\). This adjunction restricts to an equivalence between \(\Dbc(T,\mathcal O_E)\) and the normalized constructible towers.
\end{bpstatementit}

% proof: proof/lemma-sp-tower-adjunction-foundation
\begin{proof}[{Proof}]\phantomsection\label{bp58}%
The adjunction is a composite of two standard adjunctions of ringed topoi. For the equivalence we check that the unit is an isomorphism on \(\Dbc\) and that the counit is an isomorphism on normalized constructible towers. The second is the substantial point, since there we must show that the limit of such a tower is again constructible.

The adjunction is \(Lh^*p^{-1}\dashv Rp_*h_*\), for the change of rings \(h\) from towers to constant inverse systems of \(\mathcal O_{E,T}\)-modules and the morphism \(p\) to \(T_{\mathrm{pro}\acute{e}t}\) with \(p_*=\varprojlim\) [Stacks, \href{https://stacks.math.columbia.edu/tag/06YV}{06YV}, \href{https://stacks.math.columbia.edu/tag/07A6}{07A6}]. By Proposition proposition-sp-integral-towers, \(\operatorname{red}_T\) maps \(\Dbc(T,\mathcal O_E)\) into normalized constructible towers. On \(\Dbc\) the unit is an isomorphism, since objects of \(\Dbc\) are derived complete.

Conversely, let \(C_\bullet\) be normalized and put \(C=R\varprojlim C_n\). The key point is that \(C/\ell^m\simeq C_m\). Indeed, the complex \(C/\ell^m\) differs from \(C_m\) by the derived limit of the tower \((C_m[1])_n\) with transition maps multiplication by \(\ell\), and this limit vanishes, because \(\ell^m\) kills \(C_m\) and so any \(m\) consecutive transition maps compose to zero. Consequently \(C\) is derived complete and \(C/\ell\simeq C_1\) is constructible. Moreover \(C\) is bounded, because the \(C_n\) have cohomology in a fixed interval, and on the replete topos \(T_{\mathrm{pro}\acute{e}t}\) derived limits of such towers increase the range by at most one [Stacks, \href{https://stacks.math.columbia.edu/tag/0940}{0940}]. Therefore \(C\in\Dbc(T,\mathcal O_E)\) by Proposition proposition-sp-integral-towers. Finally, the counit \(\operatorname{red}_TC\to C_\bullet\) is at level \(m\) the reduction map \(C/\ell^m\to C_m\), which is an isomorphism by what we just showed. Since evaluations detect isomorphisms, the counit is an isomorphism.
\end{proof}

% text: prose-09-23
\phantomsection\label{bp59}%
All operations on towers are defined levelwise, and the functor \(Rg_*\) is computed levelwise as well, since an injective tower is levelwise a summand of finite products of injective sheaves. So proper base change and the projection formula hold for towers.

Moreover \(Rf_{!,\bullet}\) preserves constructible towers. Each level lies in \(D^b_{\mathrm{ctf}}\) by Proposition \bpout{7.1}, so we only check that the bounds are independent of \(n\). For cohomology, note that for a compactification \(f=pj\), \(Rf_!\) has amplitude \([0,\operatorname{cd}_\ell(p)]\) for every torsion coefficient ring [Stacks, \href{https://stacks.math.columbia.edu/tag/0F10}{0F10}]. For Tor amplitude we use the projection formula \((Rf_!C_n)\otimes^{\mathbf L}N\simeq Rf_!(C_n\otimes^{\mathbf L}N)\) for every \(\mathcal O_E/\ell^n\)-module \(N\), which follows from the comparison \(b\) and shows that if all \(C_n\) have Tor amplitude in \([a,b]\), then all \(Rf_!C_n\) have Tor amplitude in \([a,b+\operatorname{cd}_\ell(p)]\).

We do not use Brown representability for towers. Instead we construct the right adjoint of \(Rf_{!,\bullet}\) on constructible towers by repeating the geometric argument with finite coefficients in the category of towers.

% object: proposition-sp-tower-geometric (blueprint/chap-Smooth-purity/sec-ell-adic-coefficients-4/subsec-Integral-coefficients-2/proposition-sp-tower-geometric/_section.md)
\phantomsection\label{bp60}%
\begin{bpstatementit}{Proposition~9.18}{ (The adjunction $Rf_!\dashv f^!$ on constructible towers)}
There is an adjunction

\[
Rf_{!,\bullet}\dashv f_\bullet^!
\quad\text{between the constructible towers on }X\text{ and on }Y.
\tag{9.18}
\]

At level \(n\) its unit and counit are those of the finite adjunction, so \((f_\bullet^!C_\bullet)_n\simeq f^!_{\mathcal O_E/\ell^n}C_n\), and the reduction maps of \(f_\bullet^!C_\bullet\) are \(f^!\) of those of \(C_\bullet\) composed with the comparisons \(c\) of Lemma \hyperref[bp49]{9.15}.
\end{bpstatementit}

% proof: proof/proposition-sp-tower-geometric
\begin{proof}[{Proof}]\phantomsection\label{bp61}%
We repeat the construction of Section \hyperref[bp17]{9.3} in the topos of towers. The construction has two parts: on a chart we write down an explicit adjunction, built from the trace of the projective line, and then we glue these local adjunctions over a finite cover. The gluing is done with pairs consisting of an object and a counit, as in Definition \hyperref[bp19]{9.5}, because such pairs are unique up to unique isomorphism and so need no choice of transition maps.

The Kummer sequence \(0\to(\mu_{\ell^n})_n\to(\mathbf G_m,(\cdot)^\ell)\to(\mathbf G_m,\mathrm{id})\to0\) is exact in towers. Hence Chern classes, the diagonal class and the trace of the projective line are single morphisms of towers. Moreover, the identities (9.9) are equalities of maps between sheaves of towers in degree zero, so they can be checked at each level. The construction of (9.10) then gives the adjunction for the projective line in \(D(\operatorname{Tow})\), and extension by zero and the étale adjunction give local adjunctions on charts.

We glue the local adjunctions in \(D(\operatorname{Tow})\), by induction on a finite cover. Let \(X=U\cup V\), and let \((H_U,\varepsilon_U)\) and \((H_V,\varepsilon_V)\) be pairs representing \(\operatorname{Hom}(Rf_{!,\bullet}(-),K_\bullet)\) on \(U\) and on \(V\). Put \(W=U\cap V\). Both pairs restrict to pairs representing this functor on \(W\), so by Yoneda there is a unique counit-preserving isomorphism between the restrictions, and we call the common restriction \(H_W\). Let \(H\) be the cone of \(j_{W!}H_W\to j_{U!}H_U\oplus j_{V!}H_V\), as in the triangle (9.12). Since the counits \(\varepsilon_U,\varepsilon_V\) agree on \(Rf_{!,\bullet}j_{W!}H_W\), the exact Hom sequence shows that they lift to a counit \(\varepsilon:Rf_{!,\bullet}H\to K_\bullet\). The five lemma, applied as in the finite case, then shows that \((H,\varepsilon)\) represents \(\operatorname{Hom}(Rf_{!,\bullet}(-),K_\bullet)\) on \(X\). The tower \(H\) is also constructible. Indeed, \(H|_U\simeq H_U\) and \(H|_V\simeq H_V\), so the union of their intervals bounds its cohomology and Tor amplitude uniformly in \(n\).

Repeating this step over a finite cover of \(X\) by charts gives a pair representing \(\operatorname{Hom}(Rf_{!,\bullet}(-),K_\bullet)\) on all of \(X\), and we let \(f^!_\bullet K_\bullet\) be its object. By uniqueness of representing pairs, this is functorial in \(K_\bullet\) and gives the adjunction (9.18).

The levels are the finite ones, because evaluation at level \(n\) carries the glued pair to a pair representing \(\operatorname{Hom}(Rf_!(-),K_n)\) with coefficients \(\mathcal O_E/\ell^n\). Indeed, on each chart it is the local finite adjunction, and the five lemma argument applies at level \(n\), so by uniqueness this pair is the finite one. As the counit is one morphism of towers, it commutes with the transition maps, and transposing this identity gives the reduction maps.
\end{proof}

% text: prose-09-24
\phantomsection\label{bp62}%
The orientation passes to towers as well. By the finite case each level of \(\omega_{f,\bullet}[-2d]\) is an invertible sheaf, so it is an invertible sheaf of towers. Moreover, the comparisons for vector bundles and for the deformation family are maps of sheaves of towers in degree zero, and equality of such maps can be tested level by level, unlike equality of maps in \(D(\operatorname{Tow})\) in general. Hence the uniqueness in Lemma \hyperref[bp29]{9.9} holds for them, because it holds at each level. Composing the inverse of the orientation with the map \(\mu_{f,K}\) on towers, we obtain an isomorphism, natural in the constructible tower \(K_\bullet\),

\[
\gamma_{f,\bullet}(K_\bullet):f_\bullet^*K_\bullet(d)[2d]\xrightarrow{\sim}f_\bullet^!K_\bullet,
\tag{9.19}
\]

whose level \(n\) is \(\gamma_{f,\mathcal O_E/\ell^n}(K_n)\).

Both functors in (9.18) preserve normalized towers. For \(Rf_{!,\bullet}\) this follows from the comparison \(b\), and for \(f_\bullet^!\) from (9.17) and (9.19).

Since \(R\bar f_*\) and \(j_!\) commute with \(R\varprojlim\) [Stacks, \href{https://stacks.math.columbia.edu/tag/09BP}{09BP}], we have

\[
Rf_{!,\mathcal O_E}M\simeq R\varprojlim_nRf_{!,\mathcal O_E/\ell^n}(M/\ell^n)
\]

for \(M\in\Dbc(X,\mathcal O_E)\). For \(K\in\Dbc(Y,\mathcal O_E)\) we put

\[
H_{f,\mathcal O_E}(K)=R\varprojlim_n f_{\mathcal O_E/\ell^n}^!(K/\ell^n)\in\Dbc(X,\mathcal O_E),
\tag{9.20}
\]

the derived limit of the normalized tower \(f_\bullet^!\operatorname{red}_YK\).

The functor \(H_{f,\mathcal O_E}\) is given by an explicit formula, whereas we only know that \(f^!_{\mathcal O_E}\) exists. The next proposition shows that the formula defines a right adjoint, and so, by uniqueness, computes \(f^!_{\mathcal O_E}\).

% object: proposition-sp-integral-adjunction (blueprint/chap-Smooth-purity/sec-ell-adic-coefficients-4/subsec-Integral-coefficients-2/proposition-sp-integral-adjunction/_section.md)
\phantomsection\label{bp63}%
\begin{bpstatementit}{Proposition~9.19}{ (The integral right adjoint built from towers)}
With the unit and counit obtained from towers, \(H_{f,\mathcal O_E}\) is right adjoint to \(Rf_{!,\mathcal O_E}\) on \(\Dbc(-,\mathcal O_E)\). Hence there is a unique counit-preserving isomorphism \(H_{f,\mathcal O_E}\simeq f^!_{\mathcal O_E}\) with the right adjoint of Theorem \bpout{5.31}.
\end{bpstatementit}

% proof: proof/proposition-sp-integral-adjunction
\begin{proof}[{Proof}]\phantomsection\label{bp64}%
We move the adjunction on towers down to \(\mathcal O_E\)-complexes along the equivalence of Lemma \hyperref[bp57]{9.17}. For this we need both functors to stay inside normalized towers, which is where the equivalence holds.

Since \(Rf_{!,\bullet}\) and \(f^!_\bullet\) preserve normalized towers, the equivalence of Lemma \hyperref[bp57]{9.17} gives \(\operatorname{red}_Y(Rf_{!,\mathcal O_E}M)\simeq Rf_{!,\bullet}\operatorname{red}_XM\) and \(\operatorname{red}_XH_{f,\mathcal O_E}(K)\simeq f^!_\bullet\operatorname{red}_YK\). Using this lemma and (9.18), we get natural bijections

\[
\begin{aligned}
\Hom(Rf_{!,\mathcal O_E}M,K)
&\simeq\Hom(Rf_{!,\bullet}\operatorname{red}_XM,\operatorname{red}_YK)\\
&\simeq\Hom(\operatorname{red}_XM,f_\bullet^!\operatorname{red}_YK)
\simeq\Hom(M,H_{f,\mathcal O_E}(K)).
\end{aligned}
\]

The second statement follows from Proposition \hyperref[bp43]{9.13}.
\end{proof}

% text: prose-09-25
\phantomsection\label{bp65}%
Applying \(R\varprojlim\) to (9.19) at \(\operatorname{red}_YK\), we obtain \(\gamma_{f,\mathcal O_E}:f^*_{\mathcal O_E}K(d)[2d]\xrightarrow{\sim}f^!_{\mathcal O_E}K\), whose reductions are the isomorphisms \(\gamma_{f,\mathcal O_E/\ell^n}\). This gives the isomorphism (9.2) for \(\mathcal O_E\), and for \(E=\mathbf Q_\ell\) the one for \(\mathbf Z_\ell\).


% object: subsec-Field-coefficients (blueprint/chap-Smooth-purity/sec-ell-adic-coefficients-4/subsec-Field-coefficients/_section.md), content hidden
\subsubsection*{\texorpdfstring{9.5.3\quad Field coefficients}{9.5.3 Field coefficients}}
\phantomsection\label{bp66}%
\addcontentsline{toc}{subsubsection}{\texorpdfstring{9.5.3\quad Field coefficients}{9.5.3 Field coefficients}}
% text: prose-09-26
\phantomsection\label{bp67}%
With coefficients in \(E\) there is no tower to pass to the limit over, but none is needed. Indeed, an object of \(\Dbc(T,E)\) comes from an integral complex by inverting \(\ell\), and its Hom groups are the integral Hom groups tensored with \(E\) (9.21). Since an adjunction over \(\mathcal O_E\) is a natural bijection of Hom groups, tensoring it with \(E\) gives an adjunction over \(E\) (Proposition \hyperref[bp68]{9.20}). For \(\overline{\mathbf Q}_\ell\) the Hom groups are colimits over finite extensions of \(\mathbf Q_\ell\), and the adjunctions over these fields pass to the colimit (Theorem \hyperref[bp71]{9.21}). For each coefficient ring we thus construct an explicit right adjoint of \(Rf_!\), and we identify it with \(f^!\) by the uniqueness of adjoints (Proposition \hyperref[bp43]{9.13}).

Every object of \(\Dbc(T,E)\) is \(K_0[1/\ell]\) for some \(K_0\in\Dbc(T,\mathcal O_E)\), called an integral model, and for integral models

\[
\Hom(K_0,L_0[q])\otimes_{\mathcal O_E}E
\xrightarrow{\sim}\Hom_{\Dbc(T,E)}\bigl(K_0[1/\ell],L_0[1/\ell][q]\bigr).
\tag{9.21}
\]

Both statements are Proposition \bpout{3.27}: its essential surjectivity provides the models, and (9.21) is its full faithfulness. To see the latter directly, note that the target is \(\Hom(K_0,L_0[1/\ell][q])\) by adjunction, and \(L_0[1/\ell]\) is the colimit of \(L_0\xrightarrow{\ell}L_0\xrightarrow{\ell}\cdots\). The left side is the colimit of the same system of Hom groups, so the claim is that Hom out of \(K_0\) commutes with this colimit. Since both sides are cohomological functors of \(K_0\), by the five lemma and dévissage we may take \(K_0=a_!C\) with \(C\) dualizable on a stratum (Lemma lemma-structure-integral-constructible-complexes). The stratum formula of Lemma lemma-sections-support-closed-subscheme, whose proof uses only that \(C\) is dualizable, turns Hom out of \(a_!C\) into cohomology on a quasi-compact quasi-separated scheme, and this commutes with the colimit by Corollary \bpout{2.29}. On models \(f^*\) and \(Rf_!\) are the integral functors with \(\ell\) inverted, again because \(R\bar f_*\) commutes with the colimit along \(\ell\).

With (9.21) the adjunction for \(E\) becomes a matter of inverting \(\ell\) in the integral adjunction.

% object: proposition-sp-rational (blueprint/chap-Smooth-purity/sec-ell-adic-coefficients-4/subsec-Field-coefficients/proposition-sp-rational/_section.md)
\phantomsection\label{bp68}%
\begin{bpstatementit}{Proposition~9.20}{ (Smooth purity with coefficients in $E$)}
Let \(E/\mathbf Q_\ell\) be finite. Then \(Rf_{!,E}\dashv\bigl(K\mapsto f_E^*K(d)[2d]\bigr)\) as functors between \(\Dbc(X,E)\) and \(\Dbc(Y,E)\).
\end{bpstatementit}

% proof: proof/proposition-sp-rational
\begin{proof}[{Proof}]\phantomsection\label{bp69}%
An adjunction is a natural bijection of Hom groups, and by (9.21) every Hom group over \(E\) is an integral Hom group tensored with \(E\). So we compute both sides on integral models. Choose integral models \(M_0\) and \(K_0\), so that \(M=M_0[1/\ell]\) and \(K=K_0[1/\ell]\). By (9.21) twice and Proposition \hyperref[bp63]{9.19} with \(\gamma_{f,\mathcal O_E}\),

\[
\begin{aligned}
\Hom(Rf_{!,E}M,K)&\simeq\Hom(Rf_{!,\mathcal O_E}M_0,K_0)\otimes E\\
&\simeq\Hom(M_0,f^*_{\mathcal O_E}K_0(d)[2d])\otimes E
\simeq\Hom(M,f_E^*K(d)[2d]).
\end{aligned}
\]

These bijections are given by the integral unit and counit with \(\ell\) inverted, so they are natural and do not depend on the models.
\end{proof}

% text: prose-09-27
\phantomsection\label{bp70}%
Since \(f^*(-)(d)[2d]\) is now a right adjoint of \(Rf_{!,E}\), as is \(f^!_E\), Proposition \hyperref[bp43]{9.13} gives an isomorphism \(f^!_E\xrightarrow{\sim}f_E^*(-)(d)[2d]\) compatible with the counits, and we let \(\gamma_{f,E}:f_E^*(-)(d)[2d]\xrightarrow{\sim}f^!_E\) be its inverse.

Hom groups over \(\overline{\mathbf Q}_\ell\) are colimits of Hom groups over finite extensions, so the adjunctions over the finite extensions pass to the colimit provided they correspond under the transition maps. To see that they do, let \(E\subset F\) be a finite extension. Then extension of scalars along the finite free map \(\mathcal O_E\to\mathcal O_F\) commutes with \(R\varprojlim\), and at level \(n\) it is extension along \(\mathcal O_E/\ell^n\to\mathcal O_F/\ell^n\), so we may compare the counits in \(D(\operatorname{Tow})\). Extending scalars in the pair of towers \((f_\bullet^!K_\bullet,\varepsilon)\) for \(\mathcal O_E\) gives a pair over \(\mathcal O_F\), whose levels represent \(\operatorname{Hom}(Rf_!(-),K_n)\) by Lemma \hyperref[bp49]{9.15}. By the five lemma, applied levelwise, the pair then represents \(\operatorname{Hom}(Rf_{!,\bullet}(-),K_\bullet)\), and by uniqueness it is isomorphic to the corresponding pair for \(\mathcal O_F\) by a unique counit-preserving isomorphism. Applying \(R\varprojlim\) and inverting \(\ell\), we see that the adjunctions over \(E\) and over \(F\) correspond under extension of scalars, transitively in \(E\subset F\subset F'\).

% object: theorem-sp-qbar (blueprint/chap-Smooth-purity/sec-ell-adic-coefficients-4/subsec-Field-coefficients/theorem-sp-qbar/_section.md)
\phantomsection\label{bp71}%
\begin{bpstatementit}{Theorem~9.21}{ (Smooth purity with coefficients in $\overline{\mathbf Q}_\ell$)}
We have \(Rf_{!,\overline{\mathbf Q}_\ell}\dashv\bigl(K\mapsto f_{\overline{\mathbf Q}_\ell}^*K(d)[2d]\bigr)\) as functors between \(\Dbc(X,\overline{\mathbf Q}_\ell)\) and \(\Dbc(Y,\overline{\mathbf Q}_\ell)\).
\end{bpstatementit}

% proof: proof/theorem-sp-qbar
\begin{proof}[{Proof}]\phantomsection\label{bp72}%
As for \(E\), we compare Hom groups, now as colimits over finite extensions instead of after inverting \(\ell\). By Proposition \bpout{7.11}, \(\Dbc(T,\overline{\mathbf Q}_\ell)\) is the \(2\)-colimit of the categories \(\Dbc(T,F)\) over finite extensions \(F\) of \(\mathbf Q_\ell\), so its Hom groups are the colimits of those over \(F\). The functors \(Rf_!\) and \(f^*\) commute with extension of scalars (Section \bpout{7.2}). Choosing models over some \(E\), the adjunctions of Proposition \hyperref[bp68]{9.20} over the fields \(F\supset E\) are compatible with the transition maps, and their colimit is the asserted bijection.
\end{proof}

% text: prose-09-28
\phantomsection\label{bp73}%
Applying Proposition \hyperref[bp43]{9.13} once more and inverting the resulting isomorphism, we get \(\gamma_{f,\overline{\mathbf Q}_\ell}:f_{\overline{\mathbf Q}_\ell}^*(-)(d)[2d]\xrightarrow{\sim}f^!_{\overline{\mathbf Q}_\ell}\), whose counit is obtained by extension of scalars from a finite extension of \(\mathbf Q_\ell\). What is left of Theorem \hyperref[bp11]{9.3} are the compatibilities with change of coefficients, base change and composition, and the trace normalization.



% object: sec-Compatibilities (blueprint/chap-Smooth-purity/sec-Compatibilities/_section.md), content hidden
\subsection*{\texorpdfstring{9.6\quad Compatibilities}{9.6 Compatibilities}}
\phantomsection\label{bp74}%
\addcontentsline{toc}{subsection}{\texorpdfstring{9.6\quad Compatibilities}{9.6 Compatibilities}}
% text: prose-09-29
\phantomsection\label{bp75}%
The isomorphism \(\gamma_f\) now exists for every admissible coefficient ring, and it remains to prove the properties (2)–(4) of Theorem \hyperref[bp11]{9.3}. In topology, orientations of a submersion restrict to fibres and multiply under composition, and Propositions \hyperref[bp78]{9.22} and \hyperref[bp81]{Proposition 9.23} are the corresponding statements. These propositions also let us define the trace map \(\operatorname{Tr}_f\) (Definition \hyperref[bp85]{9.24}), the analogue of integration over the fibres, on which Poincaré duality rests. With finite coefficients the compatibilities follow from the construction, since every ingredient of the orientation commutes with base change. Passing them to \(\mathbf Z_\ell\), however, is harder, because there the base change map itself has to be built as a limit.

% object: subsec-Base-change-and-composition (blueprint/chap-Smooth-purity/sec-Compatibilities/subsec-Base-change-and-composition/_section.md), content hidden
\subsubsection*{\texorpdfstring{9.6.1\quad Base change and composition}{9.6.1 Base change and composition}}
\phantomsection\label{bp76}%
\addcontentsline{toc}{subsubsection}{\texorpdfstring{9.6.1\quad Base change and composition}{9.6.1 Base change and composition}}
% text: prose-09-30
\phantomsection\label{bp77}%
Fix a cartesian square

\[
\begin{CD}
X' @>{g'}>> X\\
@V{f'}VV @VV{f}V\\
Y' @>{g}>> Y,
\end{CD}
\]

with \(f\) smooth of pure relative dimension \(d\), and a map \(A\to B\) of finite commutative rings killed by a power of \(\ell\). A subscript \(B\) denotes extension of scalars, as in Definition \hyperref[bp47]{9.14}.

\textbf{Base change.} The statement of Proposition \hyperref[bp78]{9.22} needs an exceptional base change isomorphism for \(\mathbf Z_\ell\). It is obtained as a limit of the finite ones, and the limit exists once the finite isomorphisms commute with the reduction maps, that is, once base change commutes with change of coefficients. To state this we name the four comparison maps involved. For \(M\in D^b_{\mathrm{ctf}}(X_{\mathrm{pro}\acute{e}t},A)\) and \(K\in D^b_{\mathrm{ctf}}(Y_{\mathrm{pro}\acute{e}t},A)\), let \(b^f(M):(Rf_{!,A}M)_B\xrightarrow{\sim}Rf_{!,B}M_B\) and \(c^f(K):(f_A^!K)_B\xrightarrow{\sim}f_B^!K_B\) be the change-of-coefficients maps of Lemma \hyperref[bp49]{9.15}. With \(g'^*:D^b_{\mathrm{ctf}}(X_{\mathrm{pro}\acute{e}t},A)\to D^b_{\mathrm{ctf}}(X'_{\mathrm{pro}\acute{e}t},A)\) and \(g^*:D^b_{\mathrm{ctf}}(Y_{\mathrm{pro}\acute{e}t},A)\to D^b_{\mathrm{ctf}}(Y'_{\mathrm{pro}\acute{e}t},A)\), let \(\beta_{f,A}(M):g^*Rf_{!,A}M\xrightarrow{\sim}Rf'_{!,A}g'^*M\) be the base change map for compact support, and let \(e_{f,A}(K):g'^*f_A^!K\xrightarrow{\sim}f'^!_Ag^*K\) be the exceptional base change isomorphism of Proposition proposition-exceptional-inverse-image-smooth-morphism. The claim is

\[
\beta_{f,B}(M_B)\circ g^*(b^f(M))
=b^{f'}(g'^*M)\circ\bigl(\beta_{f,A}(M)\bigr)_B,
\tag{9.22}
\]

\[
c^{f'}(g^*K)\circ(e_{f,A}(K))_B
=e_{f,B}(K_B)\circ g'^*c^f(K).
\tag{9.23}
\]

Both identities say that two ways of passing from \((A,Y)\) to \((B,Y')\) give the same map: changing coefficients first and then base, or the other way round. We prove the first by splitting \(Rf_!\) into an open and a proper part, and deduce the second from the first by transposition.

To prove (9.22), choose a compactification \(f=\bar f\circ j\) and pull it back to \(Y'\). Since \(b\) and \(\beta\) for \(f\) are the composites of those for \(j\) and for \(\bar f\), it suffices to treat these two cases. For the open immersion \(j\), extension by zero commutes with pullback and with extension of scalars, and both sides of (9.22) are this identification. For the proper map \(\bar f\), both \(b^{\bar f}\) and proper base change are transposes of the counit of \(L\bar f^*\dashv R\bar f_*\), the first after extension of scalars and the second after pullback; here we work on the derived categories of sheaves of \(A\)-modules on the pro-étale sites. Therefore both sides of (9.22) transpose to the pulled-back counit with scalars extended.

For the second identity, recall that \(c^f\) and \(e_{f,A}\) are defined by transposing the counit composed with \((b^f)^{-1}\), respectively with \(\beta_{f,A}^{-1}\) (see (9.16)). Transposing both sides of (9.23) under \(Rf'_{!,B}\dashv f'^!_B\) thus gives two maps \(Rf'_{!,B}\bigl(g'^*((f_A^!K)_B)\bigr)\to g^*K_B\). Each is built from the counit \(\varepsilon_{f,A}(K)\), the comparisons \(b\) and the base change maps \(\beta\), and the counit commutes with extension of scalars and with pullback. By naturality of \(b\) and \(\beta\), the two maps therefore agree once (9.22) holds.

These two identities let us assemble the finite base change maps into a single map of towers. Its levels are linked by the change-of-coefficients maps along \(\mathcal O_E/\ell^{n+1}\to\mathcal O_E/\ell^n\).

Let \(E/\mathbf Q_\ell\) be a finite extension and work in the category of towers of (9.18). Extension by zero and proper base change define \(\beta_{f,\bullet}:g^*Rf_{!,\bullet}\to Rf'_{!,\bullet}g'^*\), and transposing \(g^*(\varepsilon_{f,\bullet})\circ\beta_{f,\bullet}^{-1}\) under the adjunction on towers for \(f'\) defines \(e_{f,\bullet}:g'^*f_\bullet^!\to f'^!_\bullet g^*\). At level \(n\) these are \(\beta_{f,\mathcal O_E/\ell^n}\) and the exceptional base change map with coefficients \(\mathcal O_E/\ell^n\), so they are isomorphisms; for the second, this is Proposition proposition-exceptional-inverse-image-smooth-morphism with \(\Lambda=\mathcal O_E/\ell^n\). Moreover, the transition maps of \(f_\bullet^!\) are induced by the change-of-coefficients maps \(c\) (Proposition \hyperref[bp60]{9.18}), and by (9.23) the levels of \(e_{f,\bullet}\) commute with them. So \(e_{f,\bullet}\) is an isomorphism of functors on towers. In particular, for a normalized tower \(K_\bullet\) both \(g'^*f_\bullet^!K_\bullet\) and \(f'^!_\bullet g^*K_\bullet\) are normalized.

The functors \(\operatorname{red}\) of Lemma \hyperref[bp57]{9.17} carry the maps over \(\mathcal O_E\) (base change, composition, counits and transposes) to the maps on towers. We check this in steps. First, the functor \(\operatorname{red}\) is a derived pullback of ringed topoi, so it commutes with \(g^*\) by associativity of derived pullbacks, and it commutes with extension by zero because the right adjoints commute. Next, for proper \(\bar f\), the comparison \(\operatorname{red}_YR\bar f_*M\to R\bar f_{*,\bullet}\operatorname{red}_XM\) is the transpose of \(\operatorname{red}_X\) of the counit of \(\bar f^*\dashv R\bar f_*\), and its level \(n\) is the change-of-coefficients map of Lemma \bpout{7.6} for \(\mathcal O_E\to\mathcal O_E/\ell^n\). As evaluations detect isomorphisms, the comparison is an isomorphism for \(M\in\Dbc(X,\mathcal O_E)\). Composing with the compatibility with \(j_!\) we obtain the comparison \(\operatorname{red}_YRf_{!,\mathcal O_E}M\simeq Rf_{!,\bullet}\operatorname{red}_XM\) for a compactification \(f=\bar f\circ j\). The formal argument in the proof of Proposition \bpout{7.10} shows that, under these comparisons, \(\operatorname{red}\) carries the unit and counit of \(\bar f^*\dashv R\bar f_*\) over \(\mathcal O_E\) to the unit and counit on towers. Since the base change and composition isomorphisms of Proposition \bpout{5.28} are composites of such units and counits and of the identifications of pullbacks and of extensions by zero, \(\operatorname{red}\) carries them, with coefficients \(\mathcal O_E\), to \(\beta_{f,\bullet}\) and to the corresponding composition isomorphisms on towers. For the counit, the adjunction of Proposition \hyperref[bp63]{9.19} is transported from the adjunction on towers, so \(\operatorname{red}_Y\) carries \(\varepsilon_f(K)\) to \(\varepsilon_{f,\bullet}(\operatorname{red}_YK)\), once the sources are identified by this comparison and by the isomorphism \(\operatorname{red}_Xf^!_{\mathcal O_E}K\simeq f^!_\bullet\operatorname{red}_YK\) of its proof. Since a transpose is determined by the counit, \(\operatorname{red}\) also carries transposes under \(Rf_{!,\mathcal O_E}\dashv f^!_{\mathcal O_E}\) to transposes on towers. Finally, by Lemma \hyperref[bp57]{9.17}, \(\operatorname{red}\) is an equivalence on \(\Dbc(-,\mathcal O_E)\) whose inverse is \(R\varprojlim\). A morphism of \(\Dbc(-,\mathcal O_E)\) is therefore \(R\varprojlim\) of its image under \(\operatorname{red}\), up to the unit isomorphisms \(M\simeq R\varprojlim\operatorname{red}M\).

Let \(K\in\Dbc(Y,\mathcal O_E)\), and let \(\operatorname{red}_YK=(K/\ell^n)_n\) be its tower of reductions (Lemma \hyperref[bp57]{9.17}). Since inverse image commutes with reduction modulo \(\ell^n\), and \(f_{\mathcal O_E}^!K\simeq R\varprojlim_nf_{\mathcal O_E/\ell^n}^!(K/\ell^n)\) by Proposition \hyperref[bp63]{9.19} and (9.20), both sides below are \(R\varprojlim_n\) of the corresponding terms at level \(n\) for \(K/\ell^n\). Applying \(R\varprojlim\) to \(e_{f,\bullet}\) at \(\operatorname{red}_YK\) we obtain the integral exceptional base change isomorphism

\[
g'^{*}f_{\mathcal O_E}^{!}K
\xrightarrow{\sim}f'^{!}_{\mathcal O_E}g^{*}K;
\tag{9.24}
\]

for \(E=\mathbf Q_\ell\) the coefficients are \(\mathbf Z_\ell\). With coefficients \(E\) and \(\overline{\mathbf Q}_\ell\) we define exceptional base change by inverting \(\ell\) in (9.24) on integral models and then extending scalars from a finite extension of \(\mathbf Q_\ell\). This is possible because \(f^*\), \(f^!\) and \(Rf_!\) with these coefficients are obtained from the integral functors in the same way (Proposition \hyperref[bp68]{9.20} and Theorem \hyperref[bp71]{9.21}).

In each case the map obtained is the exceptional base change map of Definition \hyperref[bp7]{9.1}. For \(\mathcal O_E\), the map of the definition is the transpose under \(Rf'_{!,\mathcal O_E}\dashv f'^!_{\mathcal O_E}\) of \(g^*(\varepsilon_f(K))\circ\beta^{-1}\). By the preceding paragraph, \(\operatorname{red}\) carries this composite to \(g^*(\varepsilon_{f,\bullet})\circ\beta_{f,\bullet}^{-1}\) at \(\operatorname{red}_YK\), and it carries the transpose to the transpose on towers, which is \(e_{f,\bullet}(\operatorname{red}_YK)\). Hence the map of the definition is \(R\varprojlim\) of \(e_{f,\bullet}(\operatorname{red}_YK)\), that is, (9.24). For \(E\) and \(\overline{\mathbf Q}_\ell\) we argue in the same way, with extension of scalars along \(\mathcal O_E\to E\) and \(E\to\overline{\mathbf Q}_\ell\) in place of \(\operatorname{red}\). Indeed, extension of scalars commutes with \(g^*\), extension by zero and \(R\bar f_*\) through the change-of-coefficients maps of Definition \bpout{7.4}, and by Proposition \bpout{7.10} it carries the units and counits of \(\bar f^*\dashv R\bar f_*\) and of \(Rf_!\dashv f^!\) to the corresponding units and counits. So it carries the base change maps for \(Rf_!\), the counits, and therefore the exceptional base change maps over \(\mathcal O_E\) to those over \(E\). In the same way those over a finite extension of \(\mathbf Q_\ell\) go to those over \(\overline{\mathbf Q}_\ell\).

Since the exceptional base change map is an isomorphism for every admissible coefficient ring, we can state that \(\gamma_f\) commutes with it. For \(g\) the inclusion of a closed point \(y\) of \(Y\), this says that \(\gamma_f\) restricts on the fibre \(X_y\) to the isomorphism \(\gamma\) of the smooth \(\kappa(y)\)-scheme \(X_y\).

% object: proposition-compatibility-purity-base-change (blueprint/chap-Smooth-purity/sec-Compatibilities/subsec-Base-change-and-composition/proposition-compatibility-purity-base-change/_section.md)
\phantomsection\label{bp78}%
\begin{bpstatementit}{Proposition~9.22}{ (Compatibility of purity with base change)}
Let \(\Lambda\) be any admissible coefficient ring and let \(K\in D^b_{\mathrm{ctf}}(Y,\Lambda)\). Identify the sources by \(g'^{*}f^{*}\simeq(f')^{*}g^{*}\) and the targets by the exceptional base change map of Definition \hyperref[bp7]{9.1}. Then

\[
g'^{*}\gamma_f(K)
=\gamma_{f'}(g^{*}K)
\quad\text{in }D^b_{\mathrm{ctf}}(X',\Lambda).
\]
\end{bpstatementit}

% proof: proof/proposition-compatibility-purity-base-change
\begin{proof}[{Proof}]\phantomsection\label{bp79}%
The map \(\gamma_f\) has two factors, the orientation of \(\omega_f\) and the map \(\mu_{f,K}\), and we show that each commutes with base change. For the orientation this is a local statement about rank one sheaves, and it follows from the naturality of the deformation to the normal bundle. For \(\mu_{f,K}\) we compare transposes under the adjunction for \(f'\).

For a finite ring \(A\) we compare the orientations and the maps \(\mu_{f,K}\) directly with coefficients in \(A\), and for \(\mathcal O_E\) we compare the maps on towers \(\gamma_{f,\bullet}\) of (9.19) and apply \(R\varprojlim\). The cases \(E\) and \(\overline{\mathbf Q}_\ell\) follow, since \(\gamma_{f,E}\) and \(\gamma_{f,\overline{\mathbf Q}_\ell}\) are obtained from \(\gamma_{f,\mathcal O_E}\) by inverting \(\ell\) and by extension of scalars from a finite extension of \(\mathbf Q_\ell\). We write out the calculation for towers; dropping the index \(\bullet\) gives the calculation at finite level. There the map \(\gamma_{f,\bullet}\) is the composite of the orientation \((\mathcal O_{E,X}/\ell^\bullet)(d)[2d]\xrightarrow{\sim}\omega_{f,\bullet}\) and the map \(f_\bullet^*K_\bullet\otimes\omega_{f,\bullet}\to f_\bullet^!K_\bullet\) for \(K_\bullet=\operatorname{red}_YK\).

\emph{Orientation.} After removing the shift \([2d]\), the two orientations to be compared are maps between sheaves of towers of rank one in degree zero. Since maps between sheaves form a sheaf, and a map of sheaves of towers is a compatible family of maps of sheaves, equality may be tested locally on \(X'\) and at each level \(n\). At a fixed level, the diagonal, its normal bundle, the deformation to the normal bundle and the trace map of vector bundles all commute with base change (Lemma lemma-naturality-under-base-change-etale), and therefore so does the orientation built from them.

\emph{The map \(f^*K\otimes\omega_f\to f^!K\).} This factor is not a map of sheaves, so we cannot compare the two maps locally. However, a map into \(f'^!_\bullet g^*K_\bullet\) is determined by its transpose under the adjunction on towers for \(f'\). By the definition of \(e_{f,\bullet}\), the transpose of the pulled-back map \(\mu_{f,K}\) for \(f\) is the pulled-back counit for \(f\), preceded by the map of the projection formula. Since the projection formula is natural under pullback, the transpose of the map for \(f'\) is the same composite, as in the counit calculation in the proof of Proposition proposition-exceptional-inverse-image-smooth-morphism, carried out with towers.
\end{proof}

% text: prose-09-31
\phantomsection\label{bp80}%
\textbf{Composition.} For composition we again fix the composition isomorphisms for \(f^!\) with every coefficient ring before comparing the maps \(\gamma\). From now on, \(g\) no longer denotes the base change in a square: let \(X\xrightarrow{f}Y\xrightarrow{g}Z\) be smooth of pure relative dimensions \(d\) and \(e\). Composing the finite adjunctions, we identify \((g\circ f)_A^!\) with \(f_A^!g_A^!\), compatibly with \(R(g\circ f)_{!,A}\simeq Rg_{!,A}Rf_{!,A}\), and the counit of the composite adjunction is the composite of the two counits.

The change-of-coefficients maps are compatible with composition: for \(M\in D^b_{\mathrm{ctf}}(X_{\mathrm{pro}\acute{e}t},A)\) and \(K\in D^b_{\mathrm{ctf}}(Z_{\mathrm{pro}\acute{e}t},A)\),

\[
b^{g\circ f}(M)=Rg_{!,B}(b^f(M))\circ b^g(Rf_{!,A}M),
\qquad
c^{g\circ f}(K)=f_B^!c^g(K)\circ c^f(g_A^!K).
\]

The first identity holds on each free term of a resolution of \(B\) over \(A\), since the composition isomorphism is natural, and the totalization in Lemma \hyperref[bp49]{9.15} extends it to arbitrary \(A\to B\). We get the second by transposing under the composite adjunction.

On towers, the composition isomorphism for \(Rf_{!,\bullet}\) is built from the same constructions. The isomorphism \((g\circ f)^!_\bullet\simeq f^!_\bullet g^!_\bullet\) is the isomorphism of right adjoints given by the composite adjunction on towers. Since evaluation at level \(n\) carries the counits on towers to the finite counits (Proposition \hyperref[bp60]{9.18}), the level \(n\) of these isomorphisms is the finite one, and by the identities above they commute with the transition maps, which are built from \(b\) and \(c\). Applying \(R\varprojlim\) gives the integral composition isomorphisms

\[
R(g\circ f)_{!,\mathcal O_E}\simeq Rg_{!,\mathcal O_E}Rf_{!,\mathcal O_E},
\qquad
(g\circ f)_{\mathcal O_E}^{!}\simeq f_{\mathcal O_E}^{!}g_{\mathcal O_E}^{!},
\]

between \(\Dbc(X,\mathcal O_E)\) and \(\Dbc(Z,\mathcal O_E)\); for \(E=\mathbf Q_\ell\) the coefficients are \(\mathbf Z_\ell\). As for base change, inverting \(\ell\) and extending scalars give the composition isomorphisms with coefficients \(E\) and \(\overline{\mathbf Q}_\ell\). These isomorphisms \((g\circ f)^!\simeq f^!g^!\) are the composition isomorphisms of Proposition proposition-composition-open-restriction-compactifications-f. Indeed, the isomorphism of that proposition is characterized by the following property: it carries the counit of \(R(g\circ f)_!\dashv(g\circ f)^!\), composed with the composition isomorphism for \(Rf_!\), to the composite of the counits for \(f\) and \(g\). By uniqueness of adjoints it therefore suffices to check this property for every coefficient ring. With finite coefficients it holds by definition. On towers the composition isomorphism is the one of the composite adjunction, so in \(D(\operatorname{Tow})\) it carries the counit on towers of \(g\circ f\) to the composite of the counits on towers. By the comparison following the construction of \(e_{f,\bullet}\) in Subsection \hyperref[bp76]{9.6.1}, \(\operatorname{red}\) carries the counits and the composition isomorphism for \(Rf_!\) with coefficients \(\mathcal O_E\) to those on towers, and applying \(R\varprojlim\) gives the property over \(\mathcal O_E\). Similarly, the same comparison shows that extension of scalars along \(\mathcal O_E\to E\) and \(E\to\overline{\mathbf Q}_\ell\) carries the counits and the composition isomorphisms for \(Rf_!\) to the corresponding ones, which gives the property over \(E\) and \(\overline{\mathbf Q}_\ell\).

With these composition isomorphisms in place for every coefficient ring, the statement takes the following form. In topological terms it says that the orientation of a composite of submersions is the product of the orientations.

% object: proposition-compatibility-purity-composition (blueprint/chap-Smooth-purity/sec-Compatibilities/subsec-Base-change-and-composition/proposition-compatibility-purity-composition/_section.md)
\phantomsection\label{bp81}%
\begin{bpstatementit}{Proposition~9.23}{ (Compatibility of purity with composition)}
For every admissible \(\Lambda\) and \(K\in D^b_{\mathrm{ctf}}(Z,\Lambda)\), the map \(\gamma_{g\circ f}(K)\) of (9.2) is the composite

\[
\begin{aligned}
(g\circ f)^{*}K(d+e)[2d+2e]
&\simeq f^{*}\bigl(g^{*}K(e)[2e]\bigr)(d)[2d]\\
&\xrightarrow{\ f^{*}(\gamma_g(K))(d)[2d]\ }
f^{*}\bigl(g^!K\bigr)(d)[2d]\\
&\xrightarrow{\ \gamma_f(g^!K)\ }
f^!g^!K
\simeq(g\circ f)^!K
\end{aligned}
\]

in \(D^b_{\mathrm{ctf}}(X,\Lambda)\), where the coefficient \(\Lambda\) is omitted from the displayed composite.
\end{bpstatementit}

% proof: proof/proposition-compatibility-purity-composition
\begin{proof}[{Proof}]\phantomsection\label{bp82}%
As for base change, we compare orientations and the maps \(\mu_{f,K}\), first over any finite admissible ring and then for towers. The orientations are compared on a chart adapted to both \(f\) and \(g\), where each is iterated integration along affine lines, with the coordinates ordered so that the two iterated integrals are the same. Applying \(R\varprojlim\) then gives the integral identity, and inverting \(\ell\) and descending to a finite extension of \(\mathbf Q_\ell\) give the identities for coefficient fields.

\emph{Orientation, finite coefficients.} Both orientations are maps of rank one locally constant sheaves after removing twist and shift, so it suffices to compare them on a cover of \(X\) by charts, and we choose charts adapted to both morphisms. That is, choose an open \(W\subset Y\) with an étale \(Z\)-map \(W\to\mathbf A^e_Z\), and an open \(U\subset f^{-1}(W)\) with an étale \(W\)-map \(U\to\mathbf A^d_W\). Then the composite \(U\to\mathbf A^d_W\to\mathbf A^d_{\mathbf A^e_Z}=\mathbf A^{d+e}_Z\) is étale, and we order the coordinates of \(\mathbf A^{d+e}_Z\) with the \(e\) coordinates of \(g\) first and the \(d\) coordinates of \(f\) last. Over this chart, both orientations are adjoint to iterated integration along affine lines, in the same order: first the \(d\) coordinates of \(f\), then the \(e\) coordinates of \(g\). Since the recursive construction preceding (9.11) integrates the last coordinate first, the integration on \(\mathbf A^{d+e}_Z\) is exactly this composite. Each diagonal orientation agrees on a chart with the integration in its coordinates, so the two orientations agree on the cover.

\emph{The maps \(\mu_{f,K}\), finite coefficients.} The two maps agree because their transposes under the composite adjunction do. Indeed, the transpose of the iterated map applies the projection formula and the counit for \(f\), then the projection formula and the counit for \(g\), whereas the transpose of the map for \(g\circ f\) applies the projection formula for \(g\circ f\) and its counit. The two isomorphisms given by the projection formula agree by associativity of tensor products, and the counit of the composite adjunction is the composite of the two counits.

\emph{Towers.} The same two arguments apply in the category of towers, with its maps \(\gamma_{f,\bullet}\) and its composite adjunction. Since the orientations are maps of sheaves of towers of rank one, we may test their equality at each level, while the maps (9.4) are compared by the same transposition.
\end{proof}


% object: subsec-The-trace-map (blueprint/chap-Smooth-purity/sec-Compatibilities/subsec-The-trace-map/_section.md), content hidden
\subsubsection*{\texorpdfstring{9.6.2\quad The trace map}{9.6.2 The trace map}}
\phantomsection\label{bp83}%
\addcontentsline{toc}{subsubsection}{\texorpdfstring{9.6.2\quad The trace map}{9.6.2 The trace map}}
% text: prose-09-32
\phantomsection\label{bp84}%
On a compact oriented manifold of dimension \(n\), Poincaré duality pairs \(H^i\) with \(H^{n-i}\) by cup product followed by integration \(H^n(M)\to\mathbf R\), which sends the class of a point to \(1\). To obtain the Poincaré pairing of Weil's axiom (c) we need the analogue of integration. By adjunction we already have a map, the counit \(Rf_!f^!K\to K\), but its source involves \(f^!\), which we only know through its universal property. Composing with \(\gamma_f\) replaces \(f^!K\) by the twisted and shifted inverse image, so the result is a map out of compactly supported cohomology of an explicit complex. For example, for \(X\) smooth of dimension \(d\) over an algebraically closed field \(k\) and \(K=\Lambda\), it is a map \(H_c^{2d}(X_{\mathrm{pro}\acute{e}t},\Lambda_X(d))\to\Lambda\). Here the twist cannot be dropped, for the reason explained before Definition \bpout{2.53}: geometric Frobenius acts on \(H^2(\mathbf P^1_k,\mathbf Z_\ell)\) by \(q\), so only the twisted trace map commutes with Frobenius (Section \hyperref[bp154]{10.5}). The trace map is used in the Poincaré pairing (Theorem \hyperref[bp137]{10.12}) and the trace isomorphism in top degree (Theorem \hyperref[bp146]{10.15}).

Let \(f:X\to Y\) be smooth of pure relative dimension \(d\), and let \(\Lambda\) be any admissible coefficient ring. We omit \(\Lambda\) from the notation and let \(\varepsilon_f(K):Rf_!f^!K\to K\) be the counit of the adjunction \(Rf_!\dashv f^!\) between \(D^b_{\mathrm{ctf}}(X,\Lambda)\) and \(D^b_{\mathrm{ctf}}(Y,\Lambda)\).

% object: definition-trace-map (blueprint/chap-Smooth-purity/sec-Compatibilities/subsec-The-trace-map/definition-trace-map/_section.md)
\phantomsection\label{bp85}%
\begin{bpstatement}{Definition~9.24}{ (Trace map)}
For \(K\in D^b_{\mathrm{ctf}}(Y,\Lambda)\), the \emph{trace map} is the composite

\[
\operatorname{Tr}_{f}(K):\quad
Rf_!\bigl(f^{*}K(d)[2d]\bigr)
\xrightarrow{\ Rf_!\gamma_{f}(K)\ }
Rf_!f^!K
\xrightarrow{\ \varepsilon_{f}(K)\ }K.
\tag{9.25}
\]
\end{bpstatement}

% object: remark-trace-map-special-cases (blueprint/chap-Smooth-purity/sec-Compatibilities/subsec-The-trace-map/remark-trace-map-special-cases/_section.md)
\phantomsection\label{bp86}%
\begin{bpstatement}{Remark~9.25}{ (The trace map in two special cases)}
\begin{enumerate}
\item For the projection \(\mathbf A^d_Y\to Y\) and \(K=\Lambda_Y\), the trace map is the trace map of affine space (9.11) with the corresponding coefficients, which on each geometric fibre sends the class of the origin to \(1\).
\item For a finite étale morphism, where \(d=0\), the trace map is summation over the fibres.
\end{enumerate}
\end{bpstatement}

% text: prose-09-33
\phantomsection\label{bp87}%
By the first part of the remark, for affine space the trace map is the explicit integration along lines built from \(\mathbf P^1\). The trace map of a general smooth morphism is glued from such charts, and the next statement says that the gluing respects composition, as integration over the fibres of a composite of submersions is iterated integration.

% object: proposition-transitivity-trace-map (blueprint/chap-Smooth-purity/sec-Compatibilities/subsec-The-trace-map/proposition-transitivity-trace-map/_section.md)
\phantomsection\label{bp88}%
\begin{bpstatementit}{Proposition~9.26}{ (Transitivity of the trace map)}
Let \(X\xrightarrow{f}Y\xrightarrow{g}Z\) be smooth of pure relative dimensions \(d\) and \(e\), and \(K\in D^b_{\mathrm{ctf}}(Z,\Lambda)\) for any admissible \(\Lambda\). Then \(\operatorname{Tr}_{g\circ f}(K)\) is the composite

\[
\begin{gathered}
R(g\circ f)_!
\bigl((g\circ f)^{*}K(d+e)[2d+2e]\bigr)
\simeq Rg_!Rf_!
\bigl(f^{*}
(g^{*}K(e)[2e])(d)[2d]\bigr)\\
\xrightarrow{\ Rg_!
(\operatorname{Tr}_{f}(g^{*}K(e)[2e]))\ }
Rg_!\bigl(g^{*}K(e)[2e]\bigr)
\xrightarrow{\ \operatorname{Tr}_{g}(K)\ }K.
\end{gathered}
\]
\end{bpstatementit}

% proof: proof/proposition-transitivity-trace-map
\begin{proof}[{Proof}]\phantomsection\label{bp89}%
Each trace map is a map \(\gamma\) followed by a counit. The maps \(\gamma\) compose by Proposition \hyperref[bp81]{9.23}, and the counits compose because the counit of a composite adjunction is the composite of the counits, so the only thing to arrange is the order of the factors. If we expand both traces using (9.25), naturality of \(\gamma_f\) lets us move the map \(\gamma_g(K):g^{*}K(e)[2e]\to g^!K\) past the map \(\gamma_f\) of the first trace. The composite then consists of the two maps \(\gamma_f\) and \(\gamma_g\), which together are \(\gamma_{g\circ f}(K)\) by Proposition \hyperref[bp81]{9.23}, followed by

\[
Rg_!Rf_!f^!g^!K
\xrightarrow{\ Rg_!(\varepsilon_{f}(g^!K))\ }
Rg_!g^!K
\xrightarrow{\ \varepsilon_{g}(K)\ }K.
\]

This is the counit of the composite adjunction \(R(g\circ f)_!\dashv(g\circ f)^!\), so the composite is the map \(\gamma_{g\circ f}(K)\) followed by its counit, that is, \(\operatorname{Tr}_{g\circ f}(K)\).
\end{proof}


% object: subsec-Change-of-coefficients (blueprint/chap-Smooth-purity/sec-Compatibilities/subsec-Change-of-coefficients/_section.md), content hidden
\subsubsection*{\texorpdfstring{9.6.3\quad Change of coefficients}{9.6.3 Change of coefficients}}
\phantomsection\label{bp90}%
\addcontentsline{toc}{subsubsection}{\texorpdfstring{9.6.3\quad Change of coefficients}{9.6.3 Change of coefficients}}
% text: prose-09-34
\phantomsection\label{bp91}%
The last compatibility compares the maps \(\gamma_f\) for different coefficient rings. For a map \(A\to B\) of finite rings the comparison \(c_{A,B}\) of Lemma \hyperref[bp49]{9.15} is the change-of-coefficients map for \(f^!\) of Proposition \bpout{7.10}, since both are defined as the transpose of the extended counit composed with the inverse of the comparison for \(Rf_!\).

% object: proposition-compatibility-purity-change-coefficients (blueprint/chap-Smooth-purity/sec-Compatibilities/subsec-Change-of-coefficients/proposition-compatibility-purity-change-coefficients/_section.md)
\phantomsection\label{bp92}%
\begin{bpstatementit}{Proposition~9.27}{ (Compatibility of purity with change of coefficients)}
Let \(f:X\to Y\) be smooth of pure relative dimension \(d\), let \(\Lambda\to\Lambda'\) be a coefficient map of Proposition \bpout{7.10} and let \(K\in D^b_{\mathrm{ctf}}(Y,\Lambda)\). Then the change-of-coefficients isomorphism \((f^!K)_{\Lambda'}\to f^!(K_{\Lambda'})\) carries \((\gamma_{f,\Lambda}(K))_{\Lambda'}\) to \(\gamma_{f,\Lambda'}(K_{\Lambda'})\), the sources being identified by the change-of-coefficients isomorphism for \(f^*\). The same holds for the trace maps: the change-of-coefficients isomorphism for \(Rf_!\) carries \((\operatorname{Tr}_{f,\Lambda}(K))_{\Lambda'}\) to \(\operatorname{Tr}_{f,\Lambda'}(K_{\Lambda'})\).
\end{bpstatementit}

% proof: proof/proposition-compatibility-purity-change-coefficients
\begin{proof}[{Proof}]\phantomsection\label{bp93}%
Both kinds of comparison maps are compatible with composition of coefficient maps, so it suffices to treat the elementary maps in the list of Proposition \bpout{7.10}. For a map \(A\to B\) of finite rings the assertion for \(\gamma_f\) is (9.17), and the assertion for the traces follows from it and the counit identity (9.16). A continuous map \(\mathcal O_E\to A\) factors through \(\mathcal O_E/\ell^n\) for some \(n\), so it suffices to treat \(\mathcal O_E\to\mathcal O_E/\ell^n\). By construction (Subsection \hyperref[bp52]{9.5.2}), \(\operatorname{red}_X\) carries \(\gamma_{f,\mathcal O_E}(K)\) to \(\gamma_{f,\bullet}(\operatorname{red}_YK)\), whose level \(n\) is \(\gamma_{f,\mathcal O_E/\ell^n}(K/\ell^n)\). Here the targets are identified by the isomorphism \(\operatorname{red}_Xf^!_{\mathcal O_E}K\simeq f^!_\bullet\operatorname{red}_YK\), whose level \(n\) is an isomorphism \((f^!_{\mathcal O_E}K)/\ell^n\simeq f^!_{\mathcal O_E/\ell^n}(K/\ell^n)\), and this identification is the change-of-coefficients map of Proposition \bpout{7.10}. Indeed, the integral counit reduces to the finite one, namely to level \(n\) of the counit on towers (Subsection \hyperref[bp76]{9.6.1} and Proposition \hyperref[bp60]{9.18}), and both maps are transposes of that counit composed with the comparison for \(Rf_!\). So this case follows from the finite one. For \(\mathcal O_E\to E\) and \(E\to\overline{\mathbf Q}_\ell\) the assertion holds by construction, because \(\gamma_{f,E}\) and \(\gamma_{f,\overline{\mathbf Q}_\ell}\) were defined by inverting \(\ell\) in the integral adjunction and by extension of scalars from finite extensions of \(\mathbf Q_\ell\) (Subsection \hyperref[bp66]{9.5.3}). For a finite extension \(F/E\) we showed before Theorem \hyperref[bp71]{9.21} that extension of scalars carries the pair of towers for \(\mathcal O_E\) to the pair for \(\mathcal O_F\) by a counit-preserving isomorphism. By applying \(R\varprojlim\) and inverting \(\ell\) we get the cases \(\mathcal O_E\to\mathcal O_F\) and \(E\to F\). In every case the statement for the traces follows from the one for \(\gamma_f\), because the change-of-coefficients maps for \(Rf_!\) and \(f^!\) are compatible with the counits.
\end{proof}



% object: sec-Proof-of-smooth-purity (blueprint/chap-Smooth-purity/sec-Proof-of-smooth-purity/_section.md), content hidden
\subsection*{\texorpdfstring{9.7\quad Proof of smooth purity}{9.7 Proof of smooth purity}}
\phantomsection\label{bp94}%
\addcontentsline{toc}{subsection}{\texorpdfstring{9.7\quad Proof of smooth purity}{9.7 Proof of smooth purity}}
% text: prose-09-35
\phantomsection\label{bp95}%
The proof of Theorem \hyperref[bp11]{9.3} assembles the constructions above; only the trace normalization (1) with integral coefficients needs a separate argument.

% proof: proof/theorem-smooth-purity-all-admissible-coefficients
\begin{proof}[{Proof of Theorem 9.3}]\phantomsection\label{bp96}%
We recall \(\gamma_f\) for each kind of coefficient ring, prove the trace normalization (1) first with finite and then with integral coefficients, and cite (2)–(4).

\emph{The isomorphisms \(\gamma_f\).} For finite \(A\) it is \(\gamma_{f,A}\) of Proposition \hyperref[bp39]{9.12}. For \(\mathcal O_E\) (and \(\mathbf Z_\ell\)) it is \(R\varprojlim\) of the map on towers (9.19), which is an isomorphism with target \(f^!\) by Proposition \hyperref[bp63]{9.19}. For \(E\) (and \(\mathbf Q_\ell\)) and for \(\overline{\mathbf Q}_\ell\) it is the inverse of the counit-preserving isomorphism of Proposition \hyperref[bp43]{9.13} between \(f^!\) and the right adjoints of Proposition \hyperref[bp68]{9.20} and Theorem \hyperref[bp71]{9.21}. In every case, therefore, \(\gamma_f\) is natural in \(K\).

\emph{Property (1), finite coefficients.} Here (1) holds because the orientation obtained from the diagonal agrees on the chart \(\mathbf A^1_Y\) with the one given by the adjunction in coordinates, whose counit at the constant sheaf is \(\operatorname{Tr}_p\) (the comparison following Proposition \hyperref[bp36]{9.11}).

\emph{Property (1), integral and rational coefficients.} For \(\Lambda=\mathcal O_E\) we write the left side \(\phi\) of (1) as \(u\circ\operatorname{Tr}_p\) and show that \(u=1\) by reducing modulo \(\ell^n\).

First, the trace map \(\operatorname{Tr}_p\) of (9.1) is an isomorphism, and this is what makes \(u\) well defined. Its reduction modulo \(\ell\) is the trace map with coefficients \(\mathcal O_E/\ell\), because the trace of Proposition \bpout{6.22} lifts the finite ones, and this reduction is an isomorphism. Indeed, it is obtained by extension of scalars along the free map \(\mathbf Z/\ell\to\mathcal O_E/\ell\) from the trace map with coefficients \(\mathbf Z/\ell\), which is an isomorphism (Section \hyperref[bp17]{9.3}). So the cone \(C\) of \(\operatorname{Tr}_p\) has \(C/\ell=0\), and hence \(C/\pi_E=(C/\ell)\otimes^{\mathbf L}_{\mathcal O_E/\ell}\mathcal O_E/\pi_E=0\). Since the source and target of \(\operatorname{Tr}_p\) lie in \(\Dbc(Y,\mathcal O_E)\), the complex \(C\) is derived complete, and derived Nakayama (Lemma \bpout{2.62}) gives \(C=0\). Thus \(\phi=u\circ\operatorname{Tr}_p\) for a unique endomorphism \(u\) of \(\mathcal O_{E,Y}\), that is, for an element \(u\in H^0(Y,\mathcal O_E)=\varprojlim_nH^0(Y,\mathcal O_E/\ell^n)\).

Next, we show that \(u\equiv1\) modulo \(\ell^n\) for every \(n\), so that \(u=1\). The trace part of Proposition \hyperref[bp92]{9.27} for \(\mathcal O_E\to\mathcal O_E/\ell^n\) shows that the reduction of \(\phi\) modulo \(\ell^n\) is the left side of (1) with coefficients \(\mathcal O_E/\ell^n\), which by the finite case equals the trace map with these coefficients. This is also the reduction of \(\operatorname{Tr}_p\), and that reduction is an isomorphism. Hence \(u\equiv1\) modulo \(\ell^n\). For \(E\) and \(\overline{\mathbf Q}_\ell\) both sides are obtained from the integral ones by inverting \(\ell\) and extending scalars.

\emph{Properties (2)–(4).} These are Propositions \hyperref[bp78]{9.22}, \hyperref[bp81]{Proposition 9.23} and \hyperref[bp92]{Proposition 9.27}.
\end{proof}

% text: prose-09-36
\phantomsection\label{bp97}%
For the structure morphism of a smooth scheme, Theorem \hyperref[bp11]{9.3} computes the dualizing complex used in Poincaré duality (Chapter \hyperref[bp101]{10}).

% object: corollary-dualizing-complex-smooth-scheme (blueprint/chap-Smooth-purity/sec-Proof-of-smooth-purity/corollary-dualizing-complex-smooth-scheme/_section.md)
\phantomsection\label{bp98}%
\begin{bpstatementit}{Corollary~9.28}{ (The dualizing complex of a smooth scheme)}
Let \(a:X\to\operatorname{Spec}k\) be smooth of pure dimension \(d\), and let \(\Lambda\) be any admissible coefficient ring. There is a natural isomorphism \(a^!\Lambda_{\operatorname{Spec}k}\simeq\Lambda_X(d)[2d]\) in \(D^b_{\mathrm{ctf}}(X,\Lambda)\).
\end{bpstatementit}

% proof: proof/corollary-dualizing-complex-smooth-scheme
\begin{proof}[{Proof}]\phantomsection\label{bp99}%
Apply Theorem \hyperref[bp11]{9.3} with coefficients \(\Lambda\) to \(K=\Lambda_{\operatorname{Spec}k}\), and use \(a^*\Lambda_{\operatorname{Spec}k}=\Lambda_X\).
\end{proof}

% text: prose-09-37
\phantomsection\label{bp100}%
As a check, take \(X=\mathbf A^1_k\) with \(k\) algebraically closed. The adjunction gives \(\Hom(R\Gamma_c(\mathbf A^1,\Lambda),\Lambda[-2])\simeq\Hom(\Lambda_X,a^!\Lambda[-2])
=H^0(\mathbf A^1,\Lambda(1))=\Lambda(1)\), and indeed \(R\Gamma_c(\mathbf A^1,\Lambda)=\Lambda(-1)[-2]\).

If \(k\) is algebraically closed, taking global sections turns the trace \(\operatorname{Tr}_a(\Lambda_{\operatorname{Spec}k})\) into a map \(R\Gamma_c(X_{\mathrm{pro}\acute{e}t},\Lambda_X(d))[2d]\to\Lambda\), and taking \(H^0\) of it gives the trace map \(H_c^{2d}(X_{\mathrm{pro}\acute{e}t},\Lambda_X(d))\to\Lambda\). For example, for \(X=\mathbf P^1_k\) it sends \(c_1(\mathcal O(1))\) to \(1\).



% object: chap-Verdier-and-Poincare-duality (blueprint/chap-Verdier-and-Poincare-duality/_section.md), content hidden
\section*{\texorpdfstring{10\quad Verdier and Poincaré duality}{10 Verdier and Poincaré duality}}
\phantomsection\label{bp101}%
\addcontentsline{toc}{section}{\texorpdfstring{10\quad Verdier and Poincaré duality}{10 Verdier and Poincaré duality}}
% text: prose-10-01
\phantomsection\label{bp102}%
The main estimate in the proof of the Deligne–Laumon theorem, the goal of Part III, is an upper bound: every Frobenius eigenvalue on \(H^i\) of a smooth proper variety has absolute value at most \(q^{i/2}\). The lower bound comes from Weil's axiom (c). Suppose that the cup product pairs \(H^i\) with \(H^{2d-i}\) perfectly and that Frobenius multiplies the pairing by \(q^d\). Then the eigenvalues \(\alpha\) in degree \(i\) correspond to the eigenvalues \(q^d/\alpha\) in degree \(2d-i\), so the upper bound for \(q^d/\alpha\) is a lower bound for \(\alpha\) (Chapter \bpout{19}). The same relation gives the functional equation of the zeta function (Corollary \hyperref[bp325]{12.33}). This chapter proves axiom (c).

For the elliptic curve \(y^2=x^3-x\) over \(\mathbf F_3\) of the introduction, duality says that the two eigenvalues on \(H^1\) have product \(3\) (Example \hyperref[bp169]{10.23}), but it does not bound them separately.

For a compact oriented manifold \(M\) of dimension \(n\) the model is Verdier's. Let \(a:M\to\mathrm{pt}\) be the map to a point. Then \(\omega_M=a^!\mathbf Z\) is the orientation sheaf placed in degree \(-n\), and the functor \(\mathbb D=\cHom(-,\omega_M)\) is an involution on constructible complexes. Duality is the isomorphism \(R\Gamma(M,\mathbb DK)\simeq R\!\operatorname{Hom}(R\Gamma_c(M,K),\mathbf Z)\).

We follow the same route. For a scheme \(X\) with structure map \(a_X:X\to s=\operatorname{Spec}k\), put \(\omega_X=a_X^!\Lambda_s\) and \(\mathbb D_X=\cHom(-,\omega_X)\) (Definition \hyperref[bp108]{10.2}). The adjunction \(Ra_{X!}\dashv a_X^!\) gives for free a pairing of \(R\Gamma_c(X,K)\) with \(R\Gamma(X,\mathbb D_XK)\), and it also gives two exchange formulas

\[
\mathbb D_YRf_!\simeq Rf_*\mathbb D_X,\qquad
f^!\mathbb D_Y\simeq\mathbb D_Xf^*
\]

(Lemma \hyperref[bp110]{10.3}). Over an algebraically closed field the first formula, for \(a_X\), is global duality, \(R\Gamma(X,\mathbb D_XK)\simeq R\!\operatorname{Hom}_\Lambda(R\Gamma_c(X,K),\Lambda)\) (Theorem \hyperref[bp132]{10.10}). For smooth \(X\) of pure dimension \(d\), smooth purity gives \(\omega_X\simeq\Lambda_X(d)[2d]\), and global duality becomes Poincaré duality: the pairing

\[
R\Gamma_c(X,K)\otimes^{\mathbf L}R\Gamma(X,K^\vee(d))[2d]
\longrightarrow\Lambda,\qquad K^\vee=\cHom(K,\Lambda_X),
\]

built from evaluation and the trace map, is perfect (Theorem \hyperref[bp137]{10.12}).

For \(X\) smooth, proper and connected and \(\Lambda\) a field, this will give a perfect cup product pairing \(H^i\times H^{2d-i}\to\Lambda(-d)\) compatible with Frobenius (Corollary \hyperref[bp152]{10.17} and Theorem \hyperref[bp165]{10.21} in Section \hyperref[bp154]{10.5}).

What the adjunction does not give is that the pairing is perfect. For that we need Verdier biduality: the map \(K\to\mathbb D_X\mathbb D_XK\) is an isomorphism for every \(K\in D^b_{\mathrm{ctf}}(X,\Lambda)\) and every admissible \(\Lambda\) (Theorem \hyperref[bp125]{10.8}). Then \(\mathbb D_X\) is an anti-equivalence, and applying the first two exchange formulas to \(\mathbb D_XK\) gives the other two, \(\mathbb D_YRf_*\simeq Rf_!\mathbb D_X\) and \(\mathbb D_Xf^!\simeq f^*\mathbb D_Y\) (Corollary \hyperref[bp128]{10.9}).

Biduality cannot be proved by dévissage, because dualizing a piece \(t_!P\), with \(P\) locally constant on a smooth stratum \(t:T\hookrightarrow X\), twice needs the exchange formula \(\mathbb D_XRt_*\simeq t_!\mathbb D_T\), which is itself a consequence of biduality (Section \hyperref[bp112]{10.2}). We argue instead by induction on dimension, restricting to the geometric generic fibre of a function, where duality commutes with restriction up to the twist \((1)[2]\) (Lemmas \hyperref[bp119]{10.6} and \hyperref[bp122]{Lemma 10.7}).

Global duality was announced by Verdier and proved by Deligne [\hyperref[bp330]{SGA4}, XVIII]. Biduality for constructible complexes over a field is due to Deligne [\hyperref[bp331]{SGA4½}, Th. finitude].

A perfect pairing of complexes induces perfect pairings of cohomology modules over a field (Corollary \hyperref[bp140]{10.13}), but not over \(\mathcal O_E\) (Remark \hyperref[bp143]{10.14}). For this reason we formulate the duality statements in the derived category, as perfect pairings of perfect complexes (Definition \hyperref[bp135]{10.11}).

% object: sec-Exchange-formulas (blueprint/chap-Verdier-and-Poincare-duality/sec-Exchange-formulas/_section.md), content hidden
\subsection*{\texorpdfstring{10.1\quad Exchange formulas}{10.1 Exchange formulas}}
\phantomsection\label{bp103}%
\addcontentsline{toc}{subsection}{\texorpdfstring{10.1\quad Exchange formulas}{10.1 Exchange formulas}}
% text: prose-10-02
\phantomsection\label{bp104}%
After defining the dualizing complex \(\omega_X\), the Verdier dual \(\mathbb D_X\) and the double-dual map \(\delta_K:K\to\mathbb D_X\mathbb D_XK\), we prove the two exchange formulas that need only adjunctions (Lemma \hyperref[bp110]{10.3}). They say that \(\mathbb D\) turns \(Rf_!\) into \(Rf_*\) and \(f^*\) into \(f^!\). For example, for \(f=a_X\) and \(K=\Lambda_X\) the first says that \(R\Gamma(X,\omega_X)\) is the dual of \(R\Gamma_c(X,\Lambda)\), as for the homology of a manifold. The other half, with the roles of \(!\) and \(*\) reversed, needs biduality.

The induction in the proof of Theorem \hyperref[bp125]{10.8} ends on the geometric point, the case \(\dim X=0\). There we want duality to be the duality of finite projective modules, and for this constructible complexes must be perfect. Over \(\mathbf Z/\ell^n\) this is a real condition; for instance, the module \(\mathbf Z/\ell\) has infinite projective dimension over \(\mathbf Z/\ell^2\). This is why we require the coefficients to have finite Tor dimension. With this requirement, Lemma \hyperref[bp105]{10.1} shows that \(D^b_{\mathrm{ctf}}\) of the point is exactly the category of perfect complexes.

% object: lemma-perfect-complexes-geometric-point (blueprint/chap-Verdier-and-Poincare-duality/sec-Exchange-formulas/lemma-perfect-complexes-geometric-point/_section.md)
\phantomsection\label{bp105}%
\begin{bpstatementit}{Lemma~10.1}{ (Perfect complexes at the geometric point)}
Suppose \(k\) is algebraically closed. The category \(D^b_{\mathrm{ctf}}(s,\Lambda)\) identifies with \(\operatorname{Perf}(\Lambda)\), the category of bounded complexes of finite projective \(\Lambda\)-modules up to quasi-isomorphism. For \(K\in D^b_{\mathrm{ctf}}(X,\Lambda)\), both

\[
R\Gamma_c(X,K)=Ra_{X!}K,\qquad R\Gamma(X,K)=Ra_{X*}K
\]

and every geometric stalk \(K_{\bar x}\) are perfect complexes. Their duals \(R\!\operatorname{Hom}_\Lambda(-,\Lambda)\) are perfect, and the evaluation map \(P\to R\!\operatorname{Hom}_\Lambda(R\!\operatorname{Hom}_\Lambda(P,\Lambda),
\Lambda)\) is an isomorphism for each of them.
\end{bpstatementit}

% proof: proof/lemma-perfect-complexes-geometric-point
\begin{proof}[{Proof}]\phantomsection\label{bp106}%
Once \(D^b_{\mathrm{ctf}}(s,\Lambda)\) is identified with \(\operatorname{Perf}(\Lambda)\), the other assertions follow, because the direct images and the stalks of \(K\) are objects of \(D^b_{\mathrm{ctf}}(s,\Lambda)\).

By the comparisons of Section \bpout{2.13}, taking global sections identifies \(D^b_{\mathrm{ctf}}(s,\Lambda)\) with the full subcategory of \(D(\Lambda)\) of complexes with bounded finitely generated cohomology and finite Tor dimension, and we show that these complexes are perfect separately for the three kinds of coefficients. A constructible complex with finite coefficients has bounded finite cohomology and finite Tor amplitude, so it is perfect: take a bounded-above resolution by finite free modules and truncate it below the Tor bound. Then the last syzygy that remains is finitely presented and flat, hence projective, so the truncated resolution is a bounded complex of finite projective modules. With integral coefficients, by contrast, the constructible complexes on the point are those with bounded finitely generated \(\mathcal O_E\)-cohomology, and since a discrete valuation ring has global dimension one, these complexes are perfect. Finally, with field coefficients, bounded complexes with finite-dimensional cohomology are perfect.

By Proposition \bpout{7.1}, the two direct images of \(a_X\) lie in \(D^b_{\mathrm{ctf}}(s,\Lambda)\) and are therefore perfect. If we represent such a complex by a bounded complex of finite projective modules, its termwise dual is again of this form, and the double evaluation map is termwise an isomorphism. The stalk \(K_{\bar x}\) is the complex of global sections of \(\bar x^*K\in D^b_{\mathrm{ctf}}(s,\Lambda)\) (Proposition \bpout{7.1}), so it is perfect as well.
\end{proof}

% text: prose-10-03
\phantomsection\label{bp107}%
On the point, then, duality is the linear duality of perfect complexes, and the double dual is the identity. On a general \(X\), however, we dualize into \(a_X^!\Lambda_s\) instead of \(\Lambda\). By the adjunction \(Ra_{X!}\dashv a_X^!\), maps into \(a_X^!\Lambda_s\) are the same as linear forms on compactly supported cohomology.

% object: definition-verdier-dual-evaluation (blueprint/chap-Verdier-and-Poincare-duality/sec-Exchange-formulas/definition-verdier-dual-evaluation/_section.md)
\phantomsection\label{bp108}%
\begin{bpstatement}{Definition~10.2}{ (Verdier dual and evaluation)}
The \emph{dualizing complex} and the \emph{Verdier dual} are

\[
\omega_X=a_X^!\Lambda_s,\qquad
\mathbb D_XK=\cHom_{\Lambda_X}(K,\omega_X).
\]

Thus \(\mathbb D_X:D(X_{\mathrm{pro}\acute{e}t},\Lambda)^{\mathrm{op}}\to
D(X_{\mathrm{pro}\acute{e}t},\Lambda)\) is an exact functor. The \emph{double-dual morphism}

\[
\delta_K:K\longrightarrow\mathbb D_X\mathbb D_XK
\]

is the tensor–Hom adjoint of the composite \(K\otimes^{\mathbf L}\mathbb D_XK\simeq\mathbb D_XK\otimes^{\mathbf L}K
\to\omega_X\) of the symmetry isomorphism, with its Koszul sign, and evaluation.
\end{bpstatement}

% text: prose-10-04
\phantomsection\label{bp109}%
On the geometric point we have \(\omega_s=\Lambda_s\), so \(\delta\) is the evaluation map of Lemma \hyperref[bp105]{10.1}.

The proof of biduality uses only the exchange formulas that follow from adjunctions. We first write down the comparison maps. Let \(f:X\to Y\) be a morphism, \(K\in D^b_{\mathrm{ctf}}(X,\Lambda)\) and \(L\in D^b_{\mathrm{ctf}}(Y,\Lambda)\). The map

\[
Rf_*\cHom_X(K,f^!L)\longrightarrow\cHom_Y(Rf_!K,L)
\tag{10.1}
\]

is the tensor–Hom adjoint of the composite

\[
\begin{aligned}
Rf_*\cHom_X(K,f^!L)\otimes^{\mathbf L}Rf_!K
&\simeq Rf_!\bigl(f^*Rf_*\cHom_X(K,f^!L)\otimes^{\mathbf L}K\bigr)\\
&\longrightarrow Rf_!\bigl(\cHom_X(K,f^!L)\otimes^{\mathbf L}K\bigr)
\longrightarrow Rf_!f^!L\xrightarrow{\varepsilon_f}L
\end{aligned}
\]

of the projection formula, the counit of \(f^*\dashv Rf_*\), evaluation and the counit \(\varepsilon_f\) of \(Rf_!\dashv f^!\). In particular, for \(L=\omega_Y\), using \(f^!\omega_Y=f^!a_Y^!\Lambda_s=\omega_X\), it is a map \(Rf_*\mathbb D_XK\to\mathbb D_YRf_!K\). The map \(f^!\mathbb D_YL\to\mathbb D_Xf^*L\) is the tensor–Hom adjoint of the map \(f^!\mathbb D_YL\otimes^{\mathbf L}f^*L\to f^!\omega_Y=\omega_X\) that corresponds under \(Rf_!\dashv f^!\) to

\[
Rf_!(f^!\mathbb D_YL\otimes^{\mathbf L}f^*L)
\simeq Rf_!f^!\mathbb D_YL\otimes^{\mathbf L}L
\xrightarrow{\varepsilon_f\otimes1}\mathbb D_YL\otimes^{\mathbf L}L
\longrightarrow\omega_Y,
\]

the last arrow being evaluation.

% object: lemma-internal-adjunction-two-exchange-formulas (blueprint/chap-Verdier-and-Poincare-duality/sec-Exchange-formulas/lemma-internal-adjunction-two-exchange-formulas/_section.md)
\phantomsection\label{bp110}%
\begin{bpstatementit}{Lemma~10.3}{ (Internal adjunction and two exchange formulas)}
Let \(f:X\to Y\) be a morphism, \(K\in D^b_{\mathrm{ctf}}(X,\Lambda)\) and \(L\in D^b_{\mathrm{ctf}}(Y,\Lambda)\). The map (10.1) and the two maps derived from it and defined above are isomorphisms

\begin{align*}
\cHom_Y(Rf_!K,L)&\simeq
Rf_*\cHom_X(K,f^!L),\tag{10.2}\\
e_K:\mathbb D_YRf_!K&\xrightarrow{\ \sim\ } Rf_*\mathbb D_XK,
\tag{10.3}\\
f^!\mathbb D_YL&\simeq\mathbb D_Xf^*L.
\tag{10.4}\end{align*}

If \(f\) is proper, then the composite

\[
Rf_*K\xrightarrow{Rf_*\delta_K}Rf_*\mathbb D_X\mathbb D_XK
\simeq\mathbb D_YRf_*\mathbb D_XK
\]

of \(Rf_*\delta_K\) and \(e_{\mathbb D_XK}\) equals the composite

\[
Rf_*K\xrightarrow{\delta_{Rf_*K}}\mathbb D_Y\mathbb D_YRf_*K
\simeq\mathbb D_YRf_*\mathbb D_XK
\]

of \(\delta_{Rf_*K}\) and \(\mathbb D_Y\) applied to \(e_K\).
\end{bpstatementit}

% proof: proof/lemma-internal-adjunction-two-exchange-formulas
\begin{proof}[{Proof}]\phantomsection\label{bp111}%
We check each isomorphism by Yoneda. After applying \(\Hom(N,-)\), we identify both sides with the same Hom group by the adjunctions and the projection formula. We only have to make sure that the test objects \(N\) lie in the localizing category on which \(Rf_!\dashv f^!\) is defined.

The adjunction \(Rf_!\dashv f^!\) lives on the localizing categories \(\Loc(D^b_{\mathrm{ctf}}(X,\Lambda))\) and \(\Loc(D^b_{\mathrm{ctf}}(Y,\Lambda))\), and the functors in the argument preserve them. Indeed, pullback is exact, preserves direct sums, and carries the constructible generators to constructible objects by Proposition \bpout{7.1}, so it maps \(\Loc(D^b_{\mathrm{ctf}}(Y,\Lambda))\) into \(\Loc(D^b_{\mathrm{ctf}}(X,\Lambda))\). For the same reasons, tensoring with a fixed constructible complex preserves the localizing category. In particular, for \(N\in\Loc(D^b_{\mathrm{ctf}}(Y,\Lambda))\) the objects \(f^*N\otimes^{\mathbf L}K\) and \(N\otimes^{\mathbf L}Rf_!K\) lie in \(\Loc(D^b_{\mathrm{ctf}}(X,\Lambda))\) and \(\Loc(D^b_{\mathrm{ctf}}(Y,\Lambda))\), respectively.

The second step of the computation below is the projection formula

\[
Rf_!(f^*N\otimes^{\mathbf L}K)\simeq N\otimes^{\mathbf L}Rf_!K
\qquad(N\in\Loc(D^b_{\mathrm{ctf}}(Y,\Lambda))).
\]

For finite coefficients it follows from Proposition \bpout{5.15} and [Stacks, \href{https://stacks.math.columbia.edu/tag/0GL5}{0GL5}], and for rational coefficients it is Proposition \bpout{6.23}. For integral coefficients we compactify \(f=pj\), with \(j\) an open immersion and \(p\) proper, and combine Lemma \bpout{6.26} with \(j_!(M\otimes^{\mathbf L}j^*N)
 \simeq j_!M\otimes^{\mathbf L}N\).

By Proposition \bpout{7.1}, both \(\cHom_Y(Rf_!K,L)\) and \(Rf_*\cHom_X(K,f^!L)\) belong to \(D^b_{\mathrm{ctf}}(Y,\Lambda)\), so we may test the comparison map between them by Yoneda inside the full subcategory \(\Loc(D^b_{\mathrm{ctf}}(Y,\Lambda))\). For every object \(N\) of this category we obtain, by the adjunctions and the projection formula,

\[
\begin{aligned}
\Hom_Y(N,\cHom_Y(Rf_!K,L))
&\simeq\Hom_Y(N\otimes^{\mathbf L}Rf_!K,L)\\
&\simeq\Hom_Y(Rf_!(f^*N\otimes^{\mathbf L}K),L)\\
&\simeq\Hom_X(f^*N\otimes^{\mathbf L}K,f^!L)\\
&\simeq\Hom_Y(N,Rf_*\cHom_X(K,f^!L)).
\end{aligned}
\]

The composite bijection is induced by the map (10.1). By Yoneda this map is an isomorphism, which is (10.2). Taking \(L=\omega_Y\) and using \(f^!\omega_Y=\omega_X\), we get \(e_K\). For (10.4), on the other hand, note that both \(f^!\mathbb D_YL\) and \(\mathbb D_Xf^*L\) lie in \(D^b_{\mathrm{ctf}}(X,\Lambda)\). Since tensoring with \(f^*L\) preserves \(\Loc(D^b_{\mathrm{ctf}}(X,\Lambda))\), we may test on objects \(M\) of this category:

\[
\begin{aligned}
\Hom_X(M,f^!\mathbb D_YL)
&\simeq\Hom_Y(Rf_!M,\cHom_Y(L,\omega_Y))\\
&\simeq\Hom_Y(Rf_!M\otimes^{\mathbf L}L,\omega_Y)\\
&\simeq\Hom_X(M\otimes^{\mathbf L}f^*L,f^!\omega_Y)\\
&\simeq\Hom_X(M,\mathbb D_Xf^*L).
\end{aligned}
\]

Again, the bijection is induced by the map defined before the lemma, and by Yoneda we get (10.4).

The comparison of double-dual maps is used in the proof of biduality for \(f=a_X\), where it identifies \(R\Gamma(X,\delta_K)\) with the double-dual map of the perfect complex \(R\Gamma(X,K)\). Both maps are determined by pairings, and we compare the pairings. For proper \(f\) we have \(Rf_!=Rf_*\), and the isomorphism \(e_K\) is adjoint to the pairing

\[
Rf_*K\otimes^{\mathbf L}Rf_*\mathbb D_XK
\longrightarrow Rf_*(K\otimes^{\mathbf L}\mathbb D_XK)
\longrightarrow Rf_*\omega_X\longrightarrow\omega_Y.
\]

Using \(e_K\) and \(e_{\mathbb D_XK}\), we view both \(Rf_*\delta_K\) and \(\delta_{Rf_*K}\) as maps \(Rf_*K\to\mathbb D_YRf_*\mathbb D_XK\). Up to the symmetry of the tensor product, both are adjoint to this pairing. Indeed, for \(\delta_{Rf_*K}\) this holds because under \(e_K\) evaluation on \(Y\) corresponds to the pairing. Similarly, for \(Rf_*\delta_K\) it holds because \(\delta_K\) is adjoint to evaluation on \(X\) and \(e_{\mathbb D_XK}\) is adjoint to the analogous pairing.
\end{proof}


% object: sec-Verdier-biduality (blueprint/chap-Verdier-and-Poincare-duality/sec-Verdier-biduality/_section.md), content hidden
\subsection*{\texorpdfstring{10.2\quad Verdier biduality}{10.2 Verdier biduality}}
\phantomsection\label{bp112}%
\addcontentsline{toc}{subsection}{\texorpdfstring{10.2\quad Verdier biduality}{10.2 Verdier biduality}}
% text: prose-10-05
\phantomsection\label{bp113}%
On a smooth curve \(C\) over an algebraically closed field, with \(\omega_C=\Lambda_C(1)[2]\), biduality is easy to see for the two basic kinds of sheaves. Indeed, for a locally constant sheaf \(L\) of free \(\Lambda\)-modules the dual is \(L^\vee(1)[2]\), and dualizing twice gives back \(L\). For a skyscraper \(i_*M\) at a closed point \(i:x\to C\), the exchange formula \(\mathbb D_Ci_*\simeq i_*\mathbb D_x\) for the proper map \(i\) shows that the dual is \(i_*\mathbb D_xM\), so biduality reduces to biduality on the point. In general, a constructible sheaf is an extension of \(i_*M\) by \(j_!L\), where now \(i\) is the inclusion of a finite set of closed points and \(L\) is locally constant on its complement \(j:U\hookrightarrow C\). The open piece is the difficult one: we have \(\mathbb D_C(j_!L)\simeq Rj_*\mathbb D_UL\), and to dualize \(Rj_*\) we would need the second exchange formula. That formula, however, depends on biduality. In dimension one this can still be checked by hand, using the cohomology of the henselian punctured disc, whereas in higher dimension the direct images \(Rj_*\) along the boundary of a stratum have no such description.

% object: example-biduality-punctured-line (blueprint/chap-Verdier-and-Poincare-duality/sec-Verdier-biduality/example-biduality-punctured-line/_section.md)
\phantomsection\label{bp114}%
\begin{bpstatement}{Example~10.4}{ (The dual of $j_!\Lambda$ on the line)}
Let \(k\) be algebraically closed, \(X=\mathbf A^1_k\), \(i:0\to X\) the origin and \(j:U=\mathbf G_{m,k}\to X\) its complement, and take \(\Lambda\) finite. The sheaf \(j_!\Lambda_U\) has stalk zero at \(0\), but its dual does not. By (10.3) for \(j\) and \(\omega_U=\Lambda_U(1)[2]\),

\[
\mathbb D_X(j_!\Lambda_U)\simeq Rj_*\mathbb D_U\Lambda_U
=Rj_*\Lambda_U(1)[2],
\]

and by Example \bpout{6.25} its stalk at \(0\) is \(\Lambda(1)\) in degree \(-2\) and \(\Lambda\) in degree \(-1\). The term in degree \(-1\) comes from \(R^1j_*\Lambda_U=i_*\Lambda(-1)\). The same answer comes out of the triangle \(j_!\Lambda_U\to\Lambda_X\to i_*\Lambda\to\). Dualizing it, with \(\mathbb D_X\Lambda_X=\Lambda_X(1)[2]\) and \(\mathbb D_X(i_*\Lambda)=i_*\mathbb D_0\Lambda=i_*\Lambda\), gives a triangle

\[
i_*\Lambda\longrightarrow\Lambda_X(1)[2]\longrightarrow
\mathbb D_X(j_!\Lambda_U)\longrightarrow.
\]

On stalks at \(0\) the first map, from degree \(0\) to degree \(-2\), is zero. Hence the cone has \(\Lambda(1)\) in degree \(-2\) and \(\Lambda\) in degree \(-1\).

Dualizing once more, biduality predicts \(\mathbb D_X\mathbb D_X(j_!\Lambda_U)
\simeq j_!\Lambda_U\), with zero stalk at \(0\). In other words, the second dual cancels the stalk that the first created. Here this follows from the two easy cases. The map \(\delta\) is a morphism of triangles, and it is an isomorphism on \(\Lambda_X\) and on \(i_*\Lambda\), so by the five lemma it is an isomorphism on \(j_!\Lambda_U\). The argument used, however, that \(\Lambda_U\) extends to the locally constant sheaf \(\Lambda_X\). A locally constant \(L\) on \(U\) with nontrivial monodromy around \(0\) has no such extension, and that is the case which needs the henselian disc.
\end{bpstatement}

% text: prose-10-06
\phantomsection\label{bp115}%
The proof of Theorem \hyperref[bp125]{10.8} needs two preliminary facts. First, Lemma \hyperref[bp116]{10.5} shows that \(\mathbb D_X\) preserves \(D^b_{\mathrm{ctf}}(X,\Lambda)\). Second, Lemmas \hyperref[bp119]{10.6} and \hyperref[bp122]{Lemma 10.7} show that \(\mathbb D_X\) and \(\delta\) commute with restriction to a geometric generic fibre up to the twist \((1)[2]\), and this twist cancels in the double dual.

% object: lemma-verdier-duality-preserves-constructibility (blueprint/chap-Verdier-and-Poincare-duality/sec-Verdier-biduality/lemma-verdier-duality-preserves-constructibility/_section.md)
\phantomsection\label{bp116}%
\begin{bpstatementit}{Lemma~10.5}{ (Verdier duality preserves constructibility)}
For every admissible \(\Lambda\), \(\mathbb D_X\) restricts to a functor \(D^b_{\mathrm{ctf}}(X,\Lambda)^{\mathrm{op}}\to D^b_{\mathrm{ctf}}(X,\Lambda)\).
\end{bpstatementit}

% proof: proof/lemma-verdier-duality-preserves-constructibility
\begin{proof}[{Proof}]\phantomsection\label{bp117}%
The lemma follows from the stability of \(D^b_{\mathrm{ctf}}\) under the six operations. Proposition \bpout{7.1} shows that \(a_X^!\) and internal Hom preserve \(D^b_{\mathrm{ctf}}\). As \(\Lambda_s\) lies in \(D^b_{\mathrm{ctf}}(s,\Lambda)\), both \(\omega_X=a_X^!\Lambda_s\) and \(\mathbb D_XK\) lie in \(D^b_{\mathrm{ctf}}(X,\Lambda)\) for \(K\in D^b_{\mathrm{ctf}}(X,\Lambda)\). We also give a direct argument with finite coefficients and smooth strata, because the proof of Lemma \hyperref[bp119]{10.6} uses its dévissage to generators \(t_!P\) again.

\emph{Dévissage.} The complex \(K\) is built by finitely many cones from complexes \(t_!P\), for locally closed \(t:T\hookrightarrow X\) with \(T\) smooth and \(P\) a locally constant perfect complex. To see this, choose a finite stratification by smooth strata, disregarding nilpotents, on which the restrictions \(P\) of \(K\) are locally constant perfect complexes, and use the localization triangles. The bounded flat models of Lemma \bpout{3.13} and their stupid filtrations make this precise. Since \(\mathbb D_X\) is exact and \(D^b_{\mathrm{ctf}}(X,\Lambda)\) is closed under cones, it suffices to show that \(\mathbb D_X(t_!P)\) lies in \(D^b_{\mathrm{ctf}}(X,\Lambda)\).

\emph{The generators.} Internal adjunction (10.3) for \(t\) gives

\[
\mathbb D_X(t_!P)\simeq Rt_*\mathbb D_TP.
\]

On a smooth stratum of dimension \(r\), Theorem \hyperref[bp11]{9.3} identifies the dual \(\mathbb D_TP\) with \(P^\vee(r)[2r]\), which is locally a perfect complex. By Proposition \bpout{7.1} its ordinary direct image lies in \(D^b_{\mathrm{ctf}}(X,\Lambda)\).
\end{proof}

% text: prose-10-07
\phantomsection\label{bp118}%
We prove biduality by induction on dimension, restricting the double-dual morphism to the geometric generic fibre of a map to \(\mathbf A^1_k\). This needs comparison maps between the operations on a scheme and on its geometric generic fibre. Let \(B/k\) be a smooth integral curve and \(\bar\eta=\operatorname{Spec}\overline{k(B)}\), and for a \(B\)-scheme \(T\) of finite type let \(g_T:T_{\bar\eta}=T\times_B\bar\eta\to T\) be the projection, with pullback \(g_T^*:D(T_{\mathrm{pro}\acute{e}t},A)\to D(T_{\bar\eta,\mathrm{pro}\acute{e}t},A)\), where \(A\) is a finite coefficient ring. First, for a morphism \(h:T\to X\) of \(B\)-schemes the base change map \(g_X^*Rh_*L\to Rh_{\bar\eta,*}g_T^*L\) is adjoint to \(g_T^*\) of the counit \(h^*Rh_*L\to L\). Second, for a closed immersion \(i:Z\to X\) of \(B\)-schemes, the pullback \(i_{\bar\eta,*}g_Z^*i^!M=g_X^*i_*i^!M\to g_X^*M\) of the counit corresponds under \(i_{\bar\eta,*}\dashv i_{\bar\eta}^!\) to a map \(g_Z^*i^!M\to i_{\bar\eta}^!g_X^*M\). Finally, for internal Hom, \(g_X^*\cHom_X(K,M)\to\cHom_{X_{\bar\eta}}(g_X^*K,g_X^*M)\) is the tensor–Hom adjoint of \(g_X^*\) of evaluation \(\cHom_X(K,M)\otimes^{\mathbf L}K\to M\). The next lemma says that these maps are isomorphisms.

% object: lemma-restriction-geometric-generic-fibre (blueprint/chap-Verdier-and-Poincare-duality/sec-Verdier-biduality/lemma-restriction-geometric-generic-fibre/_section.md)
\phantomsection\label{bp119}%
\begin{bpstatementit}{Lemma~10.6}{ (Restriction to a geometric generic fibre)}
In the setting above, let \(X\) be a \(B\)-scheme of finite type, let \(h:T\to X\) be a separated morphism of \(B\)-schemes of finite type, let \(i:Z\to X\) be a closed immersion, and let \(L\) on \(T\) and \(K,M\) on \(X\) be bounded constructible complexes, with \(K\in D^b_{\mathrm{ctf}}(X,A)\). The three maps

\[
\begin{aligned}
g_X^*Rh_*L&\longrightarrow Rh_{\bar\eta,*}g_T^*L,\\
g_Z^*i^!M&\longrightarrow i_{\bar\eta}^!g_X^*M,\\
g_X^*\cHom_X(K,M)&\longrightarrow
\cHom_{X_{\bar\eta}}(g_X^*K,g_X^*M)
\end{aligned}
\]

defined above are isomorphisms.
\end{bpstatementit}

% proof: proof/lemma-restriction-geometric-generic-fibre
\begin{proof}[{Proof}]\phantomsection\label{bp120}%
We prove the comparisons in the order stated, each using the previous ones: the direct image by continuity along a limit of étale maps, \(i^!\) by the localization triangle, and internal Hom by dévissage to generators \(t_!P\). Throughout, with finite coefficients Proposition \bpout{5.15} lets us compute \(D^b_{\mathrm{ctf}}(X,A)\) and the direct images by their classical étale counterparts.

\emph{Direct image.} The geometric generic point is a limit of étale \(B\)-schemes. Indeed, a separable closure of \(k(B)\) is the inverse limit of nonempty affine opens in the normalizations of \(B\) in its finite separable extensions, with the branch loci removed, and each stage is étale over an open of \(B\). The functor \(Rh_*\) commutes with pullback to each stage, because pullback to an étale \(B\)-scheme is restriction to an object of the small site. It remains to pass to the limit.

Every affine object étale over the limit descends to a stage. Continuity of cohomology (Corollary \bpout{2.29}) therefore shows that, degree by degree, the derived sections with finite coefficients commute with this filtered limit along affine transition maps, provided the complexes vanish below one fixed degree. They do, because \(L\) is bounded. Moreover, all the direct images lie in one fixed range of degrees. The base changes of \(h\) to the stages are of finite type of bounded dimension, and [Stacks, \href{https://stacks.math.columbia.edu/tag/0F10}{0F10}] bounds their \(\ell\)-cohomological dimensions (Definition \bpout{2.45}) independently of the stage. This proves the first comparison over the separable closure. It also holds over the algebraic closure, since a purely inseparable extension induces the same classical étale topos.

\emph{Exceptional inverse image.} Let \(j\) be the open complement of \(i\) and pull back

\[
i_*i^!M\longrightarrow M\longrightarrow Rj_*j^*M
\longrightarrow.
\]

We compare the pulled-back triangle with the triangle with closed support on the generic fibre, whose first term is \(i_{\bar\eta,*}i_{\bar\eta}^!g_X^*M\). Their second terms agree. The third terms agree too, because the first comparison identifies the pullback of \(Rj_*j^*M\) with \(Rj_{\bar\eta,*}j_{\bar\eta}^*g_X^*M\), and the map between them is the same. It follows that the first terms are determined functorially. Indeed, there are no nonzero maps, in any degree, from a complex supported on \(Z_{\bar\eta}\) to one of the form \(Rj_{\bar\eta,*}N\), so the first term of each triangle is the fibre of the second map. Now pullback commutes with direct image along a closed immersion, so the first term of the pulled-back triangle is \(i_{\bar\eta,*}g_Z^*i^!M\). By construction, the induced map \(i_{\bar\eta,*}g_Z^*i^!M\to i_{\bar\eta,*}i_{\bar\eta}^!g_X^*M\) is adjoint to the pulled-back counit, that is, it is \(i_{\bar\eta,*}\) of our comparison map, and it is an isomorphism. Since \(i_{\bar\eta,*}\) is fully faithful, the comparison map itself is an isomorphism.

\emph{Internal Hom.} The complexes \(K\) for which the comparison is an isomorphism are closed under cones and shifts, because both sides are exact functors of \(K\). The dévissage of Lemma \hyperref[bp116]{10.5} therefore reduces us to \(K=t_!P\), where \(t\) is the inclusion of a stratum and \(P\) is locally perfect on it. For these generators, internal adjunction and the dualizability of the locally perfect complex \(P\) express internal Hom through operations we have already compared:

\[
\cHom_X(t_!P,M)
\simeq Rt_*(P^\vee\otimes^{\mathbf L}t^!M).
\]

We check that \(g^*\) commutes with each operation on the right. For \(t^!\) we factor \(t\) as a closed immersion followed by an open immersion: for the open immersion \(t^!\) is restriction, which commutes with \(g^*\), and for the closed immersion we use the second comparison. Similarly, pullback commutes with \(Rt_*\) by the first comparison, and it commutes with tensor products and with \(P^\vee\). The first two comparisons apply here, because by Lemma \hyperref[bp116]{10.5} all complexes used remain bounded constructible. Hence the comparison holds for \(K=t_!P\).
\end{proof}

% text: prose-10-08
\phantomsection\label{bp121}%
For \(B=\mathbf A^1_k\) the dualizing complex changes by a twist under restriction to the geometric generic fibre, and this twist cancels in the double dual.

% object: lemma-double-dual-geometric-generic-fibre (blueprint/chap-Verdier-and-Poincare-duality/sec-Verdier-biduality/lemma-double-dual-geometric-generic-fibre/_section.md)
\phantomsection\label{bp122}%
\begin{bpstatementit}{Lemma~10.7}{ (Double duals on the geometric generic fibre)}
In the setting of Lemma \hyperref[bp119]{10.6}, let \(B=\mathbf A^1_k\) and \(K\in D^b_{\mathrm{ctf}}(X,A)\). For every affine open \(V\) of \(X\) there is an isomorphism \(g_V^*\omega_V\simeq\omega_{V_{\bar\eta}}(1)[2]\), and the isomorphism \(g_V^*\mathbb D_V\mathbb D_V(K|_V)\simeq\mathbb D_{V_{\bar\eta}}
\mathbb D_{V_{\bar\eta}}g_V^*(K|_V)\) obtained from it and from the comparison for internal Hom (applied twice) carries \(g_V^*\delta_{K|_V}\) to \(\delta_{g_V^*(K|_V)}\).
\end{bpstatementit}

% proof: proof/lemma-double-dual-geometric-generic-fibre
\begin{proof}[{Proof}]\phantomsection\label{bp123}%
Replacing \(X\) by an affine open we may assume \(X\) affine. The twist \((1)[2]\) appears because \(X\) has one more dimension than its generic fibre, and we see it by embedding both in affine spaces whose dimensions differ by one. To do this, choose a closed immersion \(i\) over \(\mathbf A^1_k\) of \(X\) into \(\mathbf A^r_{\mathbf A^1_k}\). Since \(\mathbf A^r_{\mathbf A^1_k}\) is smooth of dimension \(r+1\) over \(k\), by composition of exceptional inverse images and Theorem \hyperref[bp11]{9.3} we get \(\omega_X\simeq i^!A(r+1)[2r+2]\). In the same way, the geometric generic fibre is a closed subscheme of \(\mathbf A^r_{\bar\eta}\), so \(\omega_{X_{\bar\eta}}\simeq i_{\bar\eta}^!A(r)[2r]\). By the second comparison, \(g_X^*i^!A(r+1)[2r+2]\simeq i_{\bar\eta}^!A(r+1)[2r+2]\), which is \(\omega_{X_{\bar\eta}}(1)[2]\).

For the double-dual maps, note that the comparison for internal Hom is adjoint to the pullback of the evaluation \(K\otimes^{\mathbf L}\cHom_X(K,M)\to M\), so it carries evaluation to evaluation. Together with the twist it gives \(g_X^*\mathbb D_XK\simeq\mathbb D_{X_{\bar\eta}}(g_X^*K)(1)[2]\). This twist does no harm, because the double-dual map does not change if the dualizing complex is replaced by an isomorphic object \(\omega'\): an isomorphism \(\omega\simeq\omega'\) is used once in each of the two duals, and the two uses cancel in \(\delta\). We apply this to \(\omega'=\omega_{X_{\bar\eta}}(1)[2]\), using that \(\cHom(\cHom(N,\omega(1)[2]),\omega(1)[2])\simeq\cHom(\cHom(N,\omega),\omega)\) compatibly with evaluation. It follows that the isomorphism \(g_X^*\mathbb D_X\mathbb D_XK\simeq\mathbb D_{X_{\bar\eta}}
\mathbb D_{X_{\bar\eta}}g_X^*K\) obtained by applying the comparison twice carries \(g_X^*\delta_K\) to \(\delta_{g_X^*K}\).
\end{proof}

% text: prose-10-09
\phantomsection\label{bp124}%
With these lemmas we can prove biduality.

% object: theorem-verdier-biduality (blueprint/chap-Verdier-and-Poincare-duality/sec-Verdier-biduality/theorem-verdier-biduality/_section.md)
\phantomsection\label{bp125}%
\begin{bpstatementit}{Theorem~10.8}{ (Verdier biduality)}
For every admissible coefficient ring \(\Lambda\) and every scheme \(X\),

\[
\delta_K:K\xrightarrow{\sim}\mathbb D_X\mathbb D_XK
\qquad(K\in D^b_{\mathrm{ctf}}(X,\Lambda)).
\]

Thus \(\mathbb D_X\) is an exact anti-equivalence of \(D^b_{\mathrm{ctf}}(X,\Lambda)\) with itself.
\end{bpstatementit}

% proof: proof/theorem-verdier-biduality
\begin{proof}[{Proof}]\phantomsection\label{bp126}%
We first treat finite coefficients \(A\) over an algebraically closed field \(k\), and deduce the other coefficient rings and a finite base field at the end. In this main case we reduce to proper \(X\), show by induction that the cone of \(\delta_K\) has zero stalks at all nonclosed points, and compute global sections to see that its stalks at the finitely many remaining closed points vanish. This induction is on \(d=\dim X\). Since the induction step passes to a geometric generic fibre over \(\overline{k(t)}\), we run it simultaneously over all algebraically closed fields of characteristic different from \(\ell\). In dimension zero, \(X_{\mathrm{red}}\) is a finite disjoint union of copies of \(\operatorname{Spec}k\), and by nilpotent invariance we may replace \(X\) by it. On each copy, moreover, \(\omega=\Lambda\) and \(\delta_K\) is the evaluation map of a perfect complex, which is an isomorphism by Lemma \hyperref[bp105]{10.1}.

We may assume that \(X\) is proper of dimension \(d\); properness is used only in the last step, to compute global sections. This reduction rests on the fact that the double-dual morphism commutes with restriction to opens. This holds because restriction to an open \(j:U\to X\) commutes with internal Hom, satisfies \(j^*\omega_X=j^!\omega_X=\omega_U\), and identifies all the evaluation maps. Given this, choose a compactification \(j:X\hookrightarrow\overline X\), replace the target by the closure of \(X\), which has dimension \(d\), and extend \(K\) by zero. Then \(j_!K\) is in \(D^b_{\mathrm{ctf}}(\overline X,A)\), and the restriction of its double-dual morphism to \(X\) is \(\delta_K\).

Let \(Q=\Cone(\delta_K)\). We have to show that \(Q=0\). Since \(Q\) is bounded and constructible by Lemma \hyperref[bp116]{10.5}, it is enough to show that all its geometric stalks vanish.

At nonclosed points the induction hypothesis applies. Let \(x\) be a nonclosed point of \(X\). We choose an affine neighbourhood \(U\) of \(x\) and an algebra generator \(t\) whose image in \(\kappa(x)\) is transcendental over \(k\). Such a generator exists because a nonclosed point has positive transcendence degree. The induced map to \(\mathbf A^1_k\) then sends \(x\) to the generic point. Moreover, an irreducible component of \(U\) that meets the generic fibre dominates \(\mathbf A^1_k\), and its generic fibre has dimension one less than the component. Hence the geometric generic fibre has dimension at most \(d-1\), and the induction hypothesis applies to it. Let \(g\) be the projection from the geometric generic fibre of \(U\). By Lemma \hyperref[bp122]{10.7}, locally on \(U\) the complex \(g^*Q\) is isomorphic to the cone of \(\delta_{g^*K}\). This cone vanishes by the induction hypothesis over \(\overline{k(t)}\). As \(x\) lies over the generic point of \(\mathbf A^1_k\), every geometric point above \(x\) factors through the geometric generic fibre. Therefore \(Q\) has zero geometric stalks above every nonclosed point of \(X\).

Only finitely many closed points are left. Indeed, the support of each cohomology sheaf of \(Q\) is constructible, and by the previous paragraph it consists of closed points. A constructible subset of a Noetherian scheme of finite type which contains only closed points is finite, since each positive-dimensional locally closed piece would contain a nonclosed generic point. Thus the finitely many cohomology sheaves of \(Q\) are supported on a single finite set of closed points.

The stalks at these closed points vanish because \(X\) is proper and \(Q\) has no global cohomology. To see that \(R\Gamma(X,Q)=0\), we apply the exchange for proper maps in Lemma \hyperref[bp110]{10.3} to \(a_X\). It identifies \(R\Gamma(X,\delta_K)\) with the double-dual morphism of the perfect complex \(R\Gamma(X,K)\), which is an isomorphism by Lemma \hyperref[bp105]{10.1}. On the other hand, a sheaf \(\mathcal F\) supported on a finite set \(S\) of closed points is \(\bigoplus_{x\in S}x_*\mathcal F_{\bar x}\). Since each residue field is algebraically closed, \(H^0(X,\mathcal F)=\bigoplus_x
\mathcal F_{\bar x}\) and the higher cohomology vanishes. The hypercohomology spectral sequence of the bounded complex \(Q\) therefore degenerates, and it shows that

\[
0=H^mR\Gamma(X,Q)
=\bigoplus_x\mathcal H^m(Q)_{\bar x}.
\]

All these stalks vanish. Since the cohomology sheaves of \(Q\) are supported on this finite set, \(Q=0\).

A finite base field reduces to its algebraic closure. Let \(\bar k\) be an algebraic closure of \(k\) and \(g:X_{\bar k}=X\times_k\bar k\to X\) the projection. The three comparisons of Lemma \hyperref[bp119]{10.6} hold for \(g\), with the same proofs. Indeed, about \(\bar\eta\to B\) these proofs use only that it is a limit of étale \(B\)-schemes, and here, with \(B=\operatorname{Spec}k\) in place of a curve, \(\operatorname{Spec}\bar k\) is the limit of the finite étale \(k\)-schemes \(\operatorname{Spec}k'\), for \(k'\subset\bar k\) finite over \(k\). Unlike in Lemma \hyperref[bp122]{10.7}, the dualizing complex does not change. For an affine open \(V\) of \(X\), choose a closed immersion \(i:V\hookrightarrow\mathbf A^r_k\). By composition of exceptional inverse images and Theorem \hyperref[bp11]{9.3}, \(\omega_V\simeq i^!A(r)[2r]\). The base change \(i_{\bar k}:V_{\bar k}\hookrightarrow\mathbf A^r_{\bar k}\) is a closed immersion into an affine space of the same dimension, so \(\omega_{V_{\bar k}}\simeq i_{\bar k}^!A(r)[2r]\), and the second comparison gives \(g_V^*\omega_V\simeq\omega_{V_{\bar k}}\), with no twist and no shift. The argument of Lemma \hyperref[bp122]{10.7}, with this isomorphism in place of \(g_V^*\omega_V\simeq\omega_{V_{\bar\eta}}(1)[2]\), shows that the isomorphism \(g_V^*\mathbb D_V\mathbb D_V(K|_V)\simeq\mathbb D_{V_{\bar k}}\mathbb D_{V_{\bar k}}g_V^*(K|_V)\) it gives carries \(g_V^*\delta_{K|_V}\) to \(\delta_{g_V^*(K|_V)}\). The latter is an isomorphism by the case already proved, so \(g^*Q\) vanishes on every \(V_{\bar k}\), and \(g^*Q=0\). Finally, every geometric point of \(X\) factors through \(X_{\bar k}\), since its algebraically closed residue field contains an algebraic closure of \(k\). Hence all geometric stalks of \(Q\) vanish, and \(Q=0\) over the finite field as well.

For coefficient rings that are not finite we reduce to finite ones. Indeed, the double-dual map \(\delta:\mathrm{id}\to\mathbb D_X\mathbb D_X\) is a natural transformation of exact functors, built from \(a_X^!\), internal Hom and evaluation. Its extensions of scalars are identified by Proposition \bpout{7.10}. By Proposition \bpout{7.12} the finite case therefore implies the theorem for every admissible \(\Lambda\).
\end{proof}

% text: prose-10-10
\phantomsection\label{bp127}%
Since \(\mathbb D_X\) is now an anti-equivalence, we can dualize the formulas of Lemma \hyperref[bp110]{10.3}. This gives the two exchange formulas with the roles of \(!\) and \(*\) reversed.

% object: corollary-verdier-exchange-formulas (blueprint/chap-Verdier-and-Poincare-duality/sec-Verdier-biduality/corollary-verdier-exchange-formulas/_section.md)
\phantomsection\label{bp128}%
\begin{bpstatementit}{Corollary~10.9}{ (Verdier exchange formulas)}
For a morphism \(f:X\to Y\) there are natural isomorphisms of functors on \(D^b_{\mathrm{ctf}}(X,\Lambda)\), respectively \(D^b_{\mathrm{ctf}}(Y,\Lambda)\),

\[
\begin{aligned}
\mathbb D_YRf_!&\simeq Rf_*\mathbb D_X,&
\mathbb D_YRf_*&\simeq Rf_!\mathbb D_X,\\
\mathbb D_Xf^*&\simeq f^!\mathbb D_Y,&
\mathbb D_Xf^!&\simeq f^*\mathbb D_Y.
\end{aligned}
\]

In particular, for an open \(j:U\hookrightarrow X\) and a closed \(i:Z\hookrightarrow X\),

\[
\mathbb D_Xj_!\simeq Rj_*\mathbb D_U,\qquad
\mathbb D_XRj_*\simeq j_!\mathbb D_U,\qquad
i^*\mathbb D_X\simeq\mathbb D_Zi^!.
\]

The isomorphisms in the left column are (10.3) and (10.4); those in the right column are obtained from them by biduality, as displayed in the proof.
\end{bpstatementit}

% proof: proof/corollary-verdier-exchange-formulas
\begin{proof}[{Proof}]\phantomsection\label{bp129}%
Two of the formulas are Lemma \hyperref[bp110]{10.3}, and we obtain the other two by applying these to the Verdier dual of the argument and using biduality on both schemes. Explicitly, for \(K\in D^b_{\mathrm{ctf}}(X,\Lambda)\) and \(L\in D^b_{\mathrm{ctf}}(Y,\Lambda)\) this reads

\[
\mathbb D_YRf_*K\simeq\mathbb D_YRf_*\mathbb D_X(\mathbb D_XK)
\simeq\mathbb D_Y\mathbb D_YRf_!\mathbb D_XK\simeq Rf_!\mathbb D_XK,
\]

\[
\mathbb D_Xf^!L\simeq\mathbb D_Xf^!\mathbb D_Y(\mathbb D_YL)
\simeq\mathbb D_X\mathbb D_Xf^*\mathbb D_YL\simeq f^*\mathbb D_YL.
\]

The formulas for \(j\) and \(i\) are the special cases \(f=j\) and \(f=i\).
\end{proof}


% object: sec-Poincare-duality (blueprint/chap-Verdier-and-Poincare-duality/sec-Poincare-duality/_section.md), content hidden
\subsection*{\texorpdfstring{10.3\quad Poincaré duality}{10.3 Poincaré duality}}
\phantomsection\label{bp130}%
\addcontentsline{toc}{subsection}{\texorpdfstring{10.3\quad Poincaré duality}{10.3 Poincaré duality}}
% text: prose-10-11
\phantomsection\label{bp131}%
Over an algebraically closed field, \(D^b_{\mathrm{ctf}}(s,\Lambda)\) is the category of perfect \(\Lambda\)-complexes, and on it \(\mathbb D_s\) is \(R\!\operatorname{Hom}_\Lambda(-,\Lambda)\). For \(f=a_X\) the first two Verdier exchange formulas therefore become dualities between compactly supported and ordinary cohomology. This is global Verdier duality (Theorem \hyperref[bp132]{10.10}). It holds for every \(X\) of finite type, smooth or not.

For smooth \(X\) of pure dimension \(d\) we have \(\omega_X\simeq\Lambda_X(d)[2d]\). For locally perfect \(K\) the Verdier dual is therefore the ordinary dual \(K^\vee\) twisted by \((d)[2d]\), and the counit of \(Ra_{X!}\dashv a_X^!\) becomes the trace map of Definition \hyperref[bp85]{9.24}. In this way global duality becomes cup product followed by the trace.

% object: theorem-global-verdier-duality (blueprint/chap-Verdier-and-Poincare-duality/sec-Poincare-duality/theorem-global-verdier-duality/_section.md)
\phantomsection\label{bp132}%
\begin{bpstatementit}{Theorem~10.10}{ (Global Verdier duality)}
Suppose \(k\) is algebraically closed. For every \(K\in D^b_{\mathrm{ctf}}(X,\Lambda)\) there are natural isomorphisms

\begin{align*}
R\Gamma(X,\mathbb D_XK)&\simeq
R\!\operatorname{Hom}_\Lambda(R\Gamma_c(X,K),\Lambda),
\tag{10.5}\\
R\Gamma_c(X,\mathbb D_XK)&\simeq
R\!\operatorname{Hom}_\Lambda(R\Gamma(X,K),\Lambda).
\tag{10.6}\end{align*}

The first is the adjoint of the pairing formed by the cup product with mixed supports, evaluation, and the counit \(\varepsilon_{a_X}\) of \(Ra_{X!}\dashv a_X^!\):

\[
R\Gamma_c(X,K)\otimes^{\mathbf L}R\Gamma(X,\mathbb D_XK)
\longrightarrow R\Gamma_c(X,\omega_X)
\xrightarrow{\varepsilon_{a_X}}\Lambda.
\]
\end{bpstatementit}

% proof: proof/theorem-global-verdier-duality
\begin{proof}[{Proof}]\phantomsection\label{bp133}%
Both isomorphisms are the first two Verdier exchange formulas of Corollary \hyperref[bp128]{10.9} for \(f=a_X\), so we only have to recognize the Verdier dual on the point as the linear dual. On \(s\) the dualizing complex is \(\Lambda_s\), so \(\mathbb D_s=\cHom_s(-,\Lambda_s)\). Since \(k\) is algebraically closed, Lemma \hyperref[bp105]{10.1} identifies \(D^b_{\mathrm{ctf}}(s,\Lambda)\) with \(\operatorname{Perf}(\Lambda)\). Under this identification, then, \(\mathbb D_s\) becomes \(R\!\operatorname{Hom}_\Lambda(-,\Lambda)\), and the two exchange formulas read (10.5) and (10.6). As for the pairing, Lemma \hyperref[bp110]{10.3} constructs the exchange formula (10.3) from evaluation and the counit of \(Rf_!\dashv f^!\). For \(f=a_X\) it is therefore the adjoint of the displayed pairing.
\end{proof}

% text: prose-10-12
\phantomsection\label{bp134}%
Over a finite base field \(k\) the same argument gives us isomorphisms \(Ra_{X*}\mathbb D_XK\simeq\mathbb D_sRa_{X!}K\) and \(Ra_{X!}\mathbb D_XK\simeq\mathbb D_sRa_{X*}K\) in \(D(s_{\mathrm{pro}\acute{e}t},\Lambda)\), with \(\mathbb D_s=\cHom_s(-,\Lambda_s)\). These are the first two formulas of Corollary \hyperref[bp128]{10.9} for \(f=a_X\). After base change to an algebraic closure they become the isomorphisms of the theorem.

The first isomorphism of the theorem is the adjoint of a pairing. We give a name to pairings whose adjoint is an isomorphism.

% object: definition-perfect-pairing (blueprint/chap-Verdier-and-Poincare-duality/sec-Poincare-duality/definition-perfect-pairing/_section.md)
\phantomsection\label{bp135}%
\begin{bpstatement}{Definition~10.11}{ (Perfect pairing)}
A pairing \(P\otimes^{\mathbf L}Q\to\Lambda\) of perfect complexes is \emph{perfect} if its adjoint \(Q\to R\!\operatorname{Hom}_\Lambda(P,\Lambda)\) is an isomorphism.
\end{bpstatement}

% text: prose-10-13
\phantomsection\label{bp136}%
For smooth \(X\), Corollary \hyperref[bp98]{9.28} identifies \(\omega_X\) with \(\Lambda_X(d)[2d]\), and global duality becomes the Poincaré pairing of Weil's axiom (c).

% object: theorem-poincare-duality (blueprint/chap-Verdier-and-Poincare-duality/sec-Poincare-duality/theorem-poincare-duality/_section.md)
\phantomsection\label{bp137}%
\begin{bpstatementit}{Theorem~10.12}{ (Poincaré duality)}
Let \(X/k\) be smooth, separated of pure dimension \(d\), with \(k\) algebraically closed. For every admissible \(\Lambda\) and \(K\in D^b_{\mathrm{ctf}}(X,\Lambda)\), Theorem \hyperref[bp11]{9.3} gives isomorphisms

\[
\omega_X\simeq\Lambda_X(d)[2d],\qquad
\mathbb D_XK\simeq
\cHom_X(K,\Lambda_X)(d)[2d].
\]

Put \(K^\vee=\cHom_X(K,\Lambda_X)\). If \(K\) is locally perfect over \(\Lambda_X\), then so is \(K^\vee\), and it is the ordinary dual. Let \(\operatorname{Tr}_X\) be the map induced on \(R\Gamma_c\) by the trace map of \(a_X\) (Definition \hyperref[bp85]{9.24}). The pairing

\[
\begin{split}
\beta_K:\quad
R\Gamma_c(X,K)\otimes^{\mathbf L}
R\Gamma(X,K^\vee(d))[2d]
&\longrightarrow R\Gamma_c(X,\Lambda_X(d))[2d]\\
&\xrightarrow{\operatorname{Tr}_X}\Lambda
\end{split}
\tag{10.7}
\]

is perfect, and its adjoint is the canonical isomorphism

\[
R\Gamma(X,K^\vee(d))[2d]
\xrightarrow{\sim}
R\!\operatorname{Hom}_\Lambda(R\Gamma_c(X,K),\Lambda).
\tag{10.8}
\]

If \(X\) is proper, \(R\Gamma_c\) may be replaced by \(R\Gamma\). For a map \(\Lambda\to\Lambda'\) of admissible coefficient rings, the change-of-coefficients isomorphisms of Proposition \bpout{7.10} for \(R\Gamma_c\), \(R\Gamma\) and \(\cHom\) carry \(\beta_K\otimes^{\mathbf L}_\Lambda\Lambda'\) to \(\beta_{K_{\Lambda'}}\), where \(K_{\Lambda'}=K\otimes^{\mathbf L}_\Lambda\Lambda'\).
\end{bpstatementit}

% proof: proof/theorem-poincare-duality
\begin{proof}[{Proof}]\phantomsection\label{bp138}%
The theorem is global duality rewritten by means of smooth purity: once \(\omega_X\) and \(\mathbb D_XK\) are identified, the pairing of Theorem \hyperref[bp132]{10.10} becomes \(\beta_K\).

Corollary \hyperref[bp98]{9.28} identifies \(\omega_X\). Since internal Hom commutes with tensoring its second variable by an invertible Tate twist and a shift, we get the displayed description of \(\mathbb D_XK\). If \(K\) is locally perfect, then locally \(K\) is represented by a bounded complex of finite projective modules, whose dual is the termwise dual, with the usual evaluation.

Next, \(\beta_K\) is the pairing of Theorem \hyperref[bp132]{10.10} for \(K\), transported along the isomorphism \(\omega_X\simeq\Lambda_X(d)[2d]\). Indeed, under this isomorphism evaluation \(K\otimes^{\mathbf L}\mathbb D_XK\to\omega_X\) becomes evaluation on coefficients \(K\otimes^{\mathbf L}K^\vee(d)[2d]\to\Lambda_X(d)[2d]\). Moreover, the trace map of Definition \hyperref[bp85]{9.24} is by definition the counit \(\varepsilon_{a_X}\) precomposed with \(\gamma_{a_X}\), so the counit becomes \(\operatorname{Tr}_X\). Hence the adjoint of \(\beta_K\) is the isomorphism (10.5), which is (10.8). By Lemma \hyperref[bp105]{10.1} both factors of \(\beta_K\) are perfect complexes. Since the adjoint of \(\beta_K\) is an isomorphism, \(\beta_K\) is perfect.

For proper \(X\) we have \(Ra_{X!}=Ra_{X*}\), which gives the assertion on proper \(X\). The last statement holds because the pairing is built from tensor products, the two direct images, evaluation, the isomorphism \(\gamma_{a_X}\) and the trace map, all of which Proposition \bpout{7.10} and Theorem \hyperref[bp11]{9.3} show to be compatible with every coefficient map.
\end{proof}

% text: prose-10-14
\phantomsection\label{bp139}%
For \(X=\mathbf A^1\) and \(K=\Lambda\) the theorem pairs \(R\Gamma_c(\mathbf A^1,\Lambda)=\Lambda(-1)[-2]\) with \(R\Gamma(\mathbf A^1,\Lambda(1))[2]=\Lambda(1)[2]\), and the pairing is the trace \(H^2_c(\mathbf A^1,\Lambda(1))\simeq\Lambda\). Over a field we can pass from this statement about complexes to statements about the individual cohomology groups.

% object: corollary-degreewise-duality-over-fields (blueprint/chap-Verdier-and-Poincare-duality/sec-Poincare-duality/corollary-degreewise-duality-over-fields/_section.md)
\phantomsection\label{bp140}%
\begin{bpstatementit}{Corollary~10.13}{ (Degreewise duality over fields)}
Let \(X/k\) be smooth and separated of pure dimension \(d\), with \(k\) algebraically closed, let \(\Lambda\) be a field (a finite coefficient field, \(\mathbf Q_\ell\), \(E\) or \(\overline{\mathbf Q}_\ell\)), let \(K\in D^b_{\mathrm{ctf}}(X,\Lambda)\), and put \(K^\vee=\cHom_X(K,\Lambda_X)\). For every \(i\) the pairings induced by \(\beta_K\),

\[
H_c^i(X,K)\times H^{2d-i}(X,K^\vee(d))\longrightarrow\Lambda,
\qquad
H_c^i(X,K)\times H^{2d-i}(X,K^\vee)\longrightarrow\Lambda(-d),
\]

are perfect pairings of finite-dimensional vector spaces. If \(X\) is proper, the same holds with \(H_c^i\) replaced by \(H^i\).
\end{bpstatementit}

% proof: proof/corollary-degreewise-duality-over-fields
\begin{proof}[{Proof}]\phantomsection\label{bp141}%
We take cohomology of the isomorphism (10.8); over a field this loses nothing, because linear duality is exact. More precisely, a bounded complex with finite-dimensional cohomology is isomorphic in the derived category to the sum of its shifted cohomology spaces, so that \(H^{-i}(R\!\operatorname{Hom}_\Lambda(P,\Lambda))=\Hom_\Lambda(H^i(P),
\Lambda)\). Because of the shift \([2d]\), the source of (10.8) has \(H^{2d-i}(X,K^\vee(d))\) in degree \(-i\), and taking \(H^{-i}\) we get \(H^{2d-i}(X,K^\vee(d))\simeq\Hom_\Lambda(H^i_c(X,K),\Lambda)\). Since this isomorphism is induced by the pairing in degrees \((i,2d-i)\), it says exactly that the first pairing is perfect. Here both spaces are finite-dimensional by Lemma \hyperref[bp105]{10.1}. Finally, the second pairing is the first one with the second factor and the value line twisted by \((-d)\), so it is perfect as well.
\end{proof}

% text: prose-10-15
\phantomsection\label{bp142}%
The only property of a field used in this proof is that \(R\!\operatorname{Hom}_\Lambda(-,\Lambda)\) commutes with taking cohomology.

% object: remark-duality-cohomology-groups-over-oe (blueprint/chap-Verdier-and-Poincare-duality/sec-Poincare-duality/remark-duality-cohomology-groups-over-oe/_section.md)
\phantomsection\label{bp143}%
\begin{bpstatement}{Remark~10.14}{ (Duality on cohomology groups over $\mathcal O_E$ and finite rings)}
Over \(\mathcal O_E\), derived duality gives only the natural universal coefficient exact sequence

\[
0\longrightarrow
\Ext^1_{\mathcal O_E}(H_c^{i+1}(X,K),\mathcal O_E)
\longrightarrow H^{2d-i}(X,K^\vee(d))
\longrightarrow
\Hom_{\mathcal O_E}(H_c^i(X,K),\mathcal O_E)
\longrightarrow0.
\]

It comes from applying derived Hom to a bounded finite free model of \(R\Gamma_c(X,K)\), since \(\mathcal O_E\) has global dimension one. The sequence need not split canonically, and because of its \(\Ext^1\) term the pairings on individual cohomology modules need not be perfect. For example, on a point the perfect complex \([\mathcal O_E\xrightarrow{\pi_E}\mathcal O_E]\) in degrees \(-1,0\) has \(H^0=\mathcal O_E/\pi_E\), whereas its dual has \(H^0=0\) and \(H^1=\mathcal O_E/\pi_E\). Derived evaluation is an isomorphism, but the degree-zero pairing is not perfect.

Over a finite ring \(A\) that is not injective as a module over itself, derived Hom out of a bounded complex of finite projective modules records all the groups \(\Ext^i_A(-,A)\) of the cohomology modules, and this information may be lost in their ordinary duals. Moreover, the cohomology modules of a complex in \(D^b_{\mathrm{ctf}}(X,A)\) need not have finite Tor dimension, so the derived pairing cannot be reduced to pairings of cohomology sheaves.
\end{bpstatement}


% object: sec-Cohomology-in-top-degree (blueprint/chap-Verdier-and-Poincare-duality/sec-Cohomology-in-top-degree/_section.md), content hidden
\subsection*{\texorpdfstring{10.4\quad Cohomology in top degree}{10.4 Cohomology in top degree}}
\phantomsection\label{bp144}%
\addcontentsline{toc}{subsection}{\texorpdfstring{10.4\quad Cohomology in top degree}{10.4 Cohomology in top degree}}
% text: prose-10-16
\phantomsection\label{bp145}%
For constant coefficients on a smooth, proper and connected \(X\), we expect \(H^{2d}(X,\Lambda_X(d))\) to be free of rank one, generated by the class of a point, as for a manifold. We use this normalization in the trace formula and in point counting.

For a compact connected oriented manifold this is the statement that integration \(H^n(M)\to\mathbf R\) is an isomorphism sending a point to \(1\). Here the statement also has an arithmetic meaning, because geometric Frobenius acts on \(\Lambda(-d)\) by \(q^d\). Once we know that the trace map commutes with Frobenius (Theorem \hyperref[bp165]{10.21}), the theorem below says that Frobenius acts on \(H^{2d}\) by \(q^d\) and on \(H^0\) trivially. These give the factors \(1-T\) and \(1-q^dT\) of the zeta function.

Combined with Corollary \hyperref[bp140]{10.13}, the theorem gives the cup product pairing (Corollary \hyperref[bp152]{10.17}).

% object: theorem-trace-isomorphism-top-degree (blueprint/chap-Verdier-and-Poincare-duality/sec-Cohomology-in-top-degree/theorem-trace-isomorphism-top-degree/_section.md)
\phantomsection\label{bp146}%
\begin{bpstatementit}{Theorem~10.15}{ (The trace isomorphism in top degree)}
Let \(X/k\) be smooth, proper and connected of pure dimension \(d\), with \(k\) algebraically closed. For every admissible coefficient ring,

\[
H^i(X,\Lambda_X)=0\quad(i\notin[0,2d]),\qquad
H^0(X,\Lambda_X)=\Lambda,
\]

and the map

\[
\operatorname{tr}_X:H^{2d}(X,\Lambda_X)
\xrightarrow{\sim}\Lambda(-d),
\]

the twist by \((-d)\) of the map \(\operatorname{tr}_X(d):H^{2d}(X,\Lambda_X(d))\to\Lambda\) induced on \(H^{2d}\) by \(\operatorname{Tr}_X\) of Definition \hyperref[bp85]{9.24}, is an isomorphism.
\end{bpstatementit}

% proof: proof/theorem-trace-isomorphism-top-degree
\begin{proof}[{Proof}]\phantomsection\label{bp147}%
We read off the top degree from Poincaré duality in degree zero over \(\mathbf Z/\ell^n\), and then change coefficients: to a finite \(A\) by extension of scalars, to \(\mathcal O_E\) by a limit, and to the fields by inverting \(\ell\).

With finite coefficients, étale cohomology of a \(d\)-dimensional scheme of finite type over an algebraically closed field vanishes above degree \(2d\), that is, \(\operatorname{cd}_\ell(X)\le2d\) in the notation of Definition \bpout{2.45} [\hyperref[bp330]{SGA4}, X, Cor. 4.3]. This gives the vanishing outside \([0,2d]\), and \(H^0=\Lambda\) because \(X\) is connected.

Now let \(\Lambda=\mathbf Z/\ell^n\). Since this ring is injective as a module over itself, \(R\!\operatorname{Hom}_{\mathbf Z/\ell^n}(-,\mathbf Z/\ell^n)\) has cohomology equal to the ordinary dual of the opposite cohomology. Hence Poincaré duality (Theorem \hyperref[bp137]{10.12}) in degree zero identifies \(H^{2d}(X,\mathbf Z/\ell^n(d))\) with \(\Hom(H^0(X,\mathbf Z/\ell^n),\mathbf Z/\ell^n)=\mathbf Z/\ell^n\). By Theorem \hyperref[bp137]{10.12} this identification is given by pairing with \(H^0\) and applying the trace map, so it is \(\operatorname{tr}_X(d)\).

A general finite ring \(A\) need not be injective over itself, so we reach it by extension of scalars instead. Every finite \(A\) is an algebra over some \(\mathbf Z/\ell^n\), and by Proposition \bpout{7.10} the complexes with constant coefficients and their traces over \(A\) are obtained by derived extension of scalars from \(\mathbf Z/\ell^n\). For a complex \(P\) with cohomology in degrees \(\le2d\), right exactness of the tensor product gives \(H^{2d}(P\otimes^{\mathbf L}A)=H^{2d}(P)\otimes A\). The top cohomology is therefore \(\mathbf Z/\ell^n\otimes A=A\), and the trace is an isomorphism. Since derived extension of scalars does not raise the top degree, the same upper bound holds.

For \(\mathcal O_E\), we write \(R\Gamma(X,\mathcal O_{E,X}(d))\) as the derived inverse limit of the complexes \(R\Gamma(X,(\mathcal O_E/\pi_E^n)_X(d))\). Its cohomology fits into the Milnor sequence \(0\to\lim^1_nH^{i-1}\to H^i\to\lim_nH^i\to0\). Since the cohomology groups at each level are finite, they satisfy the Mittag–Leffler condition and there is no \(\lim^1\) contribution to the top group. The traces in top degree are compatible, so we get \(H^{2d}(X,\mathcal O_{E,X}(d))=\mathcal O_E\), with the same vanishing bounds. The lower bound and \(H^0\) also follow directly from ordinary sheaf cohomology and connectedness. Finally, inverting \(\ell\) and then extending scalars along finite extensions of \(E\), we obtain the cases \(\mathbf Q_\ell\), \(E\) and \(\overline{\mathbf Q}_\ell\).
\end{proof}

% text: prose-10-17
\phantomsection\label{bp148}%
By the theorem, \(H^{2d}(X,\Lambda_X(d))\) is free of rank one. It is generated by the class of a point, which we define by the counit of the exceptional adjunction for the inclusion of the point. Let \(x:\operatorname{Spec}k\to X\) be a closed point. Using \(x^!\omega_X=(a_Xx)^!\Lambda_s=\Lambda\) and \(\omega_X\simeq\Lambda_X(d)[2d]\), the counit \(x_*x^!\omega_X\to\omega_X\) induces a map \(\Lambda\to R\Gamma(X,\omega_X)=R\Gamma(X,\Lambda_X(d))[2d]\), and we let \([x]\in H^{2d}(X,\Lambda_X(d))\) be the image of \(1\).

% object: lemma-class-of-point (blueprint/chap-Verdier-and-Poincare-duality/sec-Cohomology-in-top-degree/lemma-class-of-point/_section.md)
\phantomsection\label{bp149}%
\begin{bpstatementit}{Lemma~10.16}{ (The class of a point)}
In the setting of Theorem \hyperref[bp146]{10.15}, we have \(\operatorname{tr}_X(d)([x])=1\) for every closed point \(x\) of \(X\). In particular \([x]\) generates \(H^{2d}(X,\Lambda_X(d))\).
\end{bpstatementit}

% proof: proof/lemma-class-of-point
\begin{proof}[{Proof}]\phantomsection\label{bp150}%
By definition, \([x]\) is the image of \(1\) under the map induced by the counit \(x_*x^!\omega_X\to\omega_X\), and \(\operatorname{tr}_X(d)\) is induced by the counit \(Ra_{X*}\omega_X\to\Lambda\). Because exceptional adjunctions compose, the composite of these two counits is the counit for \(a_Xx=\mathrm{id}_s\), which is the identity of \(\Lambda\). Thus \(\operatorname{tr}_X(d)([x])=1\). Since \(\operatorname{tr}_X(d)\) is an isomorphism onto \(\Lambda\) by Theorem \hyperref[bp146]{10.15}, \([x]\) is a generator.
\end{proof}

% text: prose-10-18
\phantomsection\label{bp151}%
For a field \(\Lambda\) and the constant sheaf, the trace isomorphism in top degree turns Poincaré duality into the classical cup product pairing. On \(\mathbf P^1\), for instance, the only nontrivial instance is \(H^0\times H^2\to\Lambda(-1)\). After twisting by \((1)\) it pairs the generators \(1\) and \([x]\) to \(1\) by Lemma \hyperref[bp149]{10.16}, so it is perfect.

% object: corollary-cup-product-pairing (blueprint/chap-Verdier-and-Poincare-duality/sec-Cohomology-in-top-degree/corollary-cup-product-pairing/_section.md)
\phantomsection\label{bp152}%
\begin{bpstatementit}{Corollary~10.17}{ (The cup product pairing)}
Let \(X/k\) be smooth, proper and connected of pure dimension \(d\), with \(k\) algebraically closed, and let \(\Lambda\) be a field. For every \(i\) the pairing

\[
\langle a,b\rangle=\operatorname{tr}_X(a\cup b),\qquad
H^i(X,\Lambda)\times H^{2d-i}(X,\Lambda)\longrightarrow\Lambda(-d),
\]

is perfect, and \(\langle b,a\rangle=(-1)^i\langle a,b\rangle\). In particular, if \(d\) is odd and \(\Lambda\) has characteristic zero, the pairing on \(H^d(X,\Lambda)\) is alternating, as for a compact complex manifold of odd dimension.
\end{bpstatementit}

% proof: proof/corollary-cup-product-pairing
\begin{proof}[{Proof}]\phantomsection\label{bp153}%
For \(K=\Lambda_X\) we have \(K^\vee=\Lambda_X\), and since \(X\) is proper, Corollary \hyperref[bp140]{10.13} says that the pairing \(H^i\times H^{2d-i}\to\Lambda(-d)\) induced by \(\beta_{\Lambda_X}\) is perfect. This pairing is \(\langle-,-\rangle\), because it is cup product followed by the trace map in top degree, that is, by the isomorphism \(\operatorname{tr}_X\) of Theorem \hyperref[bp146]{10.15}. By graded commutativity of the cup product we have \(\langle b,a\rangle=(-1)^{i(2d-i)}\langle a,b\rangle
=(-1)^i\langle a,b\rangle\).
\end{proof}


% object: sec-Weil-complexes (blueprint/chap-Verdier-and-Poincare-duality/sec-Weil-complexes/_section.md), content hidden
\subsection*{\texorpdfstring{10.5\quad Weil complexes}{10.5 Weil complexes}}
\phantomsection\label{bp154}%
\addcontentsline{toc}{subsection}{\texorpdfstring{10.5\quad Weil complexes}{10.5 Weil complexes}}
% text: prose-10-19
\phantomsection\label{bp155}%
The Weil conjectures concern varieties over \(\mathbf F_q\), while the six operations of Part I live on schemes over the algebraic closure. On \(X=X_0\times_{\mathbf F_q}k\), what remains of a variety \(X_0\) over \(\mathbf F_q\) is an action of Frobenius on \(X\) and on every complex pulled back from \(X_0\).

Here \(k_0=\mathbf F_q\) and \(k\) is an algebraic closure of \(k_0\). In this setting the geometric Frobenius \(\sigma\in\operatorname{Gal}(k/k_0)\) is the inverse of the arithmetic Frobenius \(a\mapsto a^q\). For a scheme \(T\) of the form \(T_0\times_{k_0}k\), with \(T_0\) over \(k_0\), let

\[
\operatorname{Fr}_T:T\longrightarrow T
\]

be the semilinear automorphism induced by \(\sigma\). It is a morphism of schemes covering the automorphism \(\operatorname{Spec}(\sigma)\) of \(\operatorname{Spec}k\).

We use \(\operatorname{Fr}_T\) to define Weil structures, but we prove the trace formula for the \(k\)-linear Frobenius \(\Phi_T=F_{T_0/k_0}\times_{k_0}k\), obtained by base change from the relative \(q\)-power Frobenius of \(T_0\), since its graph is a correspondence over \(k\). The two are compared in Section \hyperref[bp230]{12.1} (Lemma \hyperref[bp234]{12.3}).

Since \(\operatorname{Fr}_T\) is an automorphism of schemes, pullback of sheaves of modules along it is exact and preserves derived completeness. Hence for every admissible coefficient ring \(\Lambda\) it induces an autoequivalence

\[
\operatorname{Fr}_T^*:
D^b_{\mathrm{ctf}}(T,\Lambda)\longrightarrow
D^b_{\mathrm{ctf}}(T,\Lambda),
\]

which agrees with the derived inverse image along \(\operatorname{Fr}_T\).

Instead of complexes on \(T_0\) pulled back to \(T\), we keep only the part of that datum that the trace formula sees: the complex on \(T\), where all cohomology is computed, together with an identification of it with its Frobenius pullback. Every construction on \(T\) (direct images, \(f^!\), truncations, duals) is compatible with such identifications by base change, as we check below, so we never have to descend anything to \(T_0\).

There are also more of these pairs than complexes on \(T_0\), and we use the extra ones in Part III. Indeed, for a \(\mathbf Q_\ell\)-sheaf on \(T_0\), the stalk at a geometric point over a closed point is a continuous representation of a Galois group. It follows that Frobenius preserves a lattice in it, and its eigenvalues are \(\ell\)-adic units. By contrast, on \(\operatorname{Spec}k\) the sheaf \(\mathbf Q_\ell\) with Frobenius acting by \(\ell\) is a pair of this kind that does not come from \(\operatorname{Spec}\mathbf F_q\). Pulled back to a scheme, such pairs are the sheaves \(\mathcal B_b\) that normalize weights in Proposition \bpout{13.7}.

We require an isomorphism so that the Frobenius operator on cohomology is invertible over a field (Theorem \hyperref[bp165]{10.21}). The trace formula itself only needs a morphism (Definition \hyperref[bp232]{12.2}), and it is stated for these pairs (Theorem \hyperref[bp298]{12.24}).

% object: definition-weil-structure (blueprint/chap-Verdier-and-Poincare-duality/sec-Weil-complexes/definition-weil-structure/_section.md)
\phantomsection\label{bp156}%
\begin{bpstatement}{Definition~10.18}{ (Weil structure)}
For \(K\in D^b_{\mathrm{ctf}}(T,\Lambda)\), a \emph{Weil structure} is an isomorphism

\[
\phi_K:\operatorname{Fr}_T^*K\xrightarrow{\sim}K.
\]

A \emph{Weil complex} is a pair \((K,\phi_K)\) with this structure. A morphism \((K,\phi_K)\to(L,\phi_L)\) is a morphism \(h:K\to L\) satisfying

\[
h\circ\phi_K=\phi_L\circ\operatorname{Fr}_T^*(h).
\]
\end{bpstatement}

% text: prose-10-20
\phantomsection\label{bp157}%
A Weil structure is the datum needed to define a Frobenius operator on cohomology. The basic example is a complex \(K\) pulled back from a complex \(K_0\) on \(T_0\). Since the projection \(T\to T_0\) composed with \(\operatorname{Fr}_T\) is again the projection, the pullback of \(K_0\) gives a canonical isomorphism \(\operatorname{Fr}_T^*K\simeq K\).

The triangles, truncations and spectral sequences we form in \(D^b_{\mathrm{ctf}}(T,\Lambda)\) are compatible with Weil structures. Since \(\operatorname{Fr}_T^*\) is exact on sheaves, it commutes with standard truncations and with the triangle of a short exact sequence of sheaves. It also commutes with the localization triangle \(j_!j^*K\to K\to i_*i^*K\) for an open immersion \(j:U\to T\) obtained by base change from an open \(U_0\subset T_0\), with closed complement \(i\). Indeed, for such \(U\) we have \(\operatorname{Fr}_T^{-1}(U)=U\), and \(\operatorname{Fr}_T\) restricts to \(\operatorname{Fr}_U\) and to the Frobenius of the complement. These constructions therefore carry Weil structures to Weil structures, and the induced long exact cohomology sequences are Frobenius-equivariant.

Let \((M,\phi_M)\) be a Weil complex on \(X\) and let \(f:X\to Y\) come from \(f_0\) over \(k_0\). The base change map for the cartesian square with horizontal arrows \(\operatorname{Fr}_X\), \(\operatorname{Fr}_Y\) identifies \(\operatorname{Fr}_Y^*Rf_{!}M\) with \(Rf_{!}\operatorname{Fr}_X^*M\), and the composite

\[
\operatorname{Fr}_Y^*Rf_{!}M
\simeq Rf_{!}\operatorname{Fr}_X^*M
\xrightarrow{\ Rf_{!}(\phi_M)\ }Rf_{!}M
\]

is a Weil structure on \(Rf_!M\). Likewise \(f^*\), \(Rf_*\), \(f^!\), \(\otimes^{\mathbf L}\) and \(\cHom\) carry Weil structures, via the corresponding base change maps. These maps are isomorphisms for every \(f\), because \(\operatorname{Fr}_X\) and \(\operatorname{Fr}_Y\) are isomorphisms of schemes and base change along an isomorphism is transport of structure.

For \(Rf_!\) this needs a word, since \(\operatorname{Fr}_X\) and \(\operatorname{Fr}_Y\) are not \(k\)-morphisms, while \(Rf_!\) is built from compactifications over \(k\). We compute \(Rf_!=R\bar f_*j_!\) with a compactification \(X\overset{j}{\hookrightarrow}\bar X\xrightarrow{\bar f}Y\) obtained by base change from a compactification of \(f_0\) over \(k_0\) (Theorem \bpout{A.2}). Then \(\operatorname{Fr}_{\bar X}\) is an automorphism of the scheme \(\bar X\) with \(\operatorname{Fr}_{\bar X}^{-1}(X)=X\) and \(\bar f\operatorname{Fr}_{\bar X}=\operatorname{Fr}_Y\bar f\). For an automorphism \(h\) of a scheme, \(h^*\simeq(h^{-1})_*\), so transport of structure gives \(\operatorname{Fr}_Y^*R\bar f_*\simeq R\bar f_*\operatorname{Fr}_{\bar X}^*\) and \(\operatorname{Fr}_{\bar X}^*j_!\simeq j_!\operatorname{Fr}_X^*\), and their composite is the base change map for \(Rf_!\). It does not depend on the compactification. Indeed, two compactifications of \(f_0\) are dominated by a third, the closure of \(X_0\) in their fibre product over \(Y_0\), and the refinement isomorphisms of \(Rf_!\) come from proper maps over \(k_0\), which commute with the Frobenii. The map for \(f^!\) is the mate of the map for \(Rf_!\), and it is an isomorphism because both functors \(\operatorname{Fr}^*\) are equivalences.

Pullback along \(\operatorname{Fr}_T\) on compactly supported cohomology is constructed in the same way. For a morphism of schemes \(h:S'\to S\) and \(L\in D(S_{\mathrm{pro}\acute{e}t},\Lambda)\), the \emph{pullback map}

\[
\operatorname{pb}_{h,L}:R\Gamma(S,L)\longrightarrow R\Gamma(S',h^*L)
\]

is \(R\Gamma\) of the unit \(L\to Rh_*h^*L\), followed by the identification \(R\Gamma(S,Rh_*-)=R\Gamma(S',-)\). We write \(\operatorname{pb}_h\) when \(L\) is clear. It is natural in \(L\), and it is transitive: for \(g:S'\to S\) and \(h:S''\to S'\) we have \(\operatorname{pb}_{gh,L}=\operatorname{pb}_{h,g^*L}\circ\operatorname{pb}_{g,L}\) under \((gh)^*\simeq h^*g^*\). Let \(j:T\hookrightarrow\bar T\) be the base change of a compactification of \(T_0\) over \(k_0\). Then \(\operatorname{Fr}_{\bar T}\) is an automorphism of \(\bar T\) with \(\operatorname{Fr}_{\bar T}^{-1}(T)=T\), and \emph{pullback along \(\operatorname{Fr}_T\)} is the composite

\[
\operatorname{pb}_{\operatorname{Fr}_T}:
R\Gamma_c(T,K)=R\Gamma(\bar T,j_!K)
\xrightarrow{\ \operatorname{pb}_{\operatorname{Fr}_{\bar T}}\ }
R\Gamma(\bar T,\operatorname{Fr}_{\bar T}^*j_!K)
\simeq R\Gamma(\bar T,j_!\operatorname{Fr}_T^*K)
=R\Gamma_c(T,\operatorname{Fr}_T^*K).
\]

It does not depend on \(\bar T\). Indeed, a refinement \(\pi:\bar T'\to\bar T\) between two such compactifications comes from \(k_0\), so \(\pi\operatorname{Fr}_{\bar T'}=\operatorname{Fr}_{\bar T}\pi\). Transitivity of pullback maps for this composite, together with the base change \(\operatorname{Fr}_{\bar T}^*R\pi_*\simeq R\pi_*\operatorname{Fr}_{\bar T'}^*\) along the automorphisms \(\operatorname{Fr}\) and with \(R\pi_*j'_!\simeq j_!\), shows that the refinement isomorphism \(R\Gamma(\bar T',j'_!K)\simeq R\Gamma(\bar T,R\pi_*j'_!K)
\simeq R\Gamma(\bar T,j_!K)\) carries one pullback map to the other. In the same way, the pullback map along \(\operatorname{Fr}_T^m\) is defined, and by transitivity it is the \(m\)-fold composite \(\operatorname{pb}_{\operatorname{Fr}_T^m,K}=\operatorname{pb}_{\operatorname{Fr}_T,(\operatorname{Fr}_T^{m-1})^*K}\circ\cdots\circ\operatorname{pb}_{\operatorname{Fr}_T,K}\).

% object: definition-frobenius-operator (blueprint/chap-Verdier-and-Poincare-duality/sec-Weil-complexes/definition-frobenius-operator/_section.md)
\phantomsection\label{bp158}%
\begin{bpstatement}{Definition~10.19}{ (Frobenius operator)}
For a Weil complex \((K,\phi_K)\) on \(T\), the \emph{geometric Frobenius operator} on compactly supported cohomology is the endomorphism

\[
\begin{gathered}
\operatorname{Frob}_{T,K}:R\Gamma_c(T_{\mathrm{pro}\acute{e}t},K)
\xrightarrow{\ \operatorname{pb}_{\operatorname{Fr}_T}\ }
R\Gamma_c(T_{\mathrm{pro}\acute{e}t},\operatorname{Fr}_T^*K)\\
\xrightarrow{\ R\Gamma_c(T_{\mathrm{pro}\acute{e}t},\phi_K)\ }
R\Gamma_c(T_{\mathrm{pro}\acute{e}t},K).
\end{gathered}
\]

Here \(\operatorname{pb}_{\operatorname{Fr}_T}\) is the pullback along \(\operatorname{Fr}_T\) constructed above.
\end{bpstatement}

% text: prose-10-21
\phantomsection\label{bp159}%
When \(T\) and \(K\) are clear, we write \(F\) for \(\operatorname{Frob}_{T,K}\) and for the endomorphisms it induces on the groups \(H^i_c(T,K)\).

For \(r\in\mathbf Z\), write \(\Lambda(r)=\Gamma((\operatorname{Spec}k)_{\mathrm{pro}\acute{e}t},
\Lambda_{\operatorname{Spec}k}(r))\) for the Tate coefficient module with its induced Frobenius operator. Arithmetic Frobenius acts on \(\mu_{\ell^n}(k)\) by \(\zeta\mapsto\zeta^q\). Hence geometric Frobenius acts on \(\Lambda(1)\) by \(q^{-1}\), and on \(\Lambda(-d)\) by \(q^d\). Let \(X_0\) be smooth of pure dimension \(d\). After removing the twist, the trace map of Definition \hyperref[bp85]{9.24} is a Frobenius-equivariant map, by Proposition \hyperref[bp162]{10.20} below:

\[
H_c^{2d}(X_{\mathrm{pro}\acute{e}t},\Lambda_X)\longrightarrow\Lambda(-d).
\]

For \(X=\mathbf P^1_k\) this says that geometric Frobenius acts on \(H^2(\mathbf P^1_k,\Lambda)=\Lambda(-1)\) by \(q\), as checked after Proposition \bpout{6.7}.


% object: sec-Frobenius-and-Poincare-duality (blueprint/chap-Verdier-and-Poincare-duality/sec-Frobenius-and-Poincare-duality/_section.md), content hidden
\subsection*{\texorpdfstring{10.6\quad Frobenius and Poincaré duality}{10.6 Frobenius and Poincaré duality}}
\phantomsection\label{bp160}%
\addcontentsline{toc}{subsection}{\texorpdfstring{10.6\quad Frobenius and Poincaré duality}{10.6 Frobenius and Poincaré duality}}
% text: prose-10-22
\phantomsection\label{bp161}%
Let \(X_0\) be smooth of pure dimension \(d\) and let \(K\) be locally perfect, so that \(\mathbb D_XK\simeq K^\vee(d)[2d]\). In this case cup product pairs \(H^i_c(X,K)\) with \(H^{2d-i}(X,K^\vee)\) and takes values in \(H^{2d}_c(X,\Lambda)\). To use this pairing for the Weil conjectures we need to know how \(F\) acts on both sides. The target is not the trivial Weil complex but \(\Lambda(-d)\), as for \(\mathbf P^1\) above. If we forgot the twist, we would predict the relation \(\alpha\leftrightarrow1/\alpha\) between eigenvalues in complementary degrees instead of \(\alpha\leftrightarrow q^d/\alpha\).

We built Poincaré duality from the purity isomorphism \(\gamma_f\) and the trace map. It therefore suffices to check that these two maps commute with pullback along Frobenius (Proposition \hyperref[bp162]{10.20}). This statement about natural transformations is transport of structure, except for the identification of Tate twists. There we use that pullback along \(\operatorname{Fr}_T\) respects Kummer classes. The compatibility of the duality pairing with Weil structures then follows formally (Theorem \hyperref[bp165]{10.21}).

Let \(f\) be obtained from a smooth morphism \(f_0:X_0\to Y_0\) over \(k_0=\mathbf F_q\) by extension to \(k\), let \(\operatorname{Fr}_X\), \(\operatorname{Fr}_Y\) be the geometric Frobenius maps above, and let \(\gamma_f\) be the isomorphism of (9.2).

% object: proposition-purity-trace-commute-frobenius (blueprint/chap-Verdier-and-Poincare-duality/sec-Frobenius-and-Poincare-duality/proposition-purity-trace-commute-frobenius/_section.md)
\phantomsection\label{bp162}%
\begin{bpstatementit}{Proposition~10.20}{ (Purity and trace commute with Frobenius)}
Fix an admissible \(\Lambda\), and let \(K\in D^b_{\mathrm{ctf}}(Y,\Lambda)\) and \(M\in D^b_{\mathrm{ctf}}(X,\Lambda)\). The square with horizontal arrows \(\operatorname{Fr}_X,\operatorname{Fr}_Y\) and vertical arrows \(f,f\) is cartesian, and its base change comparisons identify

\[
\operatorname{Fr}_X^*f^!K
\simeq f^!\operatorname{Fr}_Y^*K,
\qquad
\operatorname{Fr}_Y^*Rf_!M
\simeq Rf_!\operatorname{Fr}_X^*M.
\tag{10.9}
\]

With these identifications, and the Kummer identification \(\operatorname{Fr}_T^*\Lambda_T(1)\simeq\Lambda_T(1)\) for the twists, we have

\[
\operatorname{Fr}_X^*\gamma_f(K)
=\gamma_f(\operatorname{Fr}_Y^*K),
\qquad
\operatorname{Fr}_Y^*\operatorname{Tr}_f(K)
=\operatorname{Tr}_f(\operatorname{Fr}_Y^*K).
\]
\end{bpstatementit}

% proof: proof/proposition-purity-trace-commute-frobenius
\begin{proof}[{Proof}]\phantomsection\label{bp163}%
Both identities are instances of base change for smooth purity, applied to a square whose horizontal maps are not \(k\)-linear. Since \(\gamma_f\) and \(\operatorname{Tr}_f\) are built from base change maps and from the Kummer class of the twist, we only have to know how \(\operatorname{Fr}_T\) acts on Kummer classes.

The first identity is Proposition \hyperref[bp78]{9.22} for this square. Since \(\operatorname{Fr}_X\) is not a morphism over \(k\), the only point to check is that the identification of twists used there is the Kummer identification, that is, that pullback along \(\operatorname{Fr}_T\) sends the Kummer class of a unit to the Kummer class of its pullback. This holds because the Kummer sequence \(0\to\mu_{\ell^n}\to\mathbf G_m\xrightarrow{\ell^n}\mathbf G_m\to0\) and its boundary map are natural for every morphism of schemes, \(k\)-linear or not.

For the second identity we combine the first with the counit identity defining exceptional base change, which says that after base change for compact support, the pulled-back counit is the counit preceded by the exceptional comparison. As the trace map is the counit composed with \(\gamma_f\), the two identities together give the statement for \(\operatorname{Tr}_f\).
\end{proof}

% text: prose-10-23
\phantomsection\label{bp164}%
Weil structures are transported along \(f^!\), \(f^*\) and \(Rf_!\). On \(f^*K(d)[2d]\) the Weil structure is the pullback of \(\phi_K\) combined with the Kummer identification. By the proposition above and naturality of \(\gamma_f\) and \(\operatorname{Tr}_f\) applied to \(\phi_K\), the maps \(\gamma_f(K)\) and \(\operatorname{Tr}_f(K)\) are morphisms of Weil complexes.

The constant sheaf \(\Lambda_Y\) has the Weil structure given by the identification \(\operatorname{Fr}_Y^*\Lambda_Y\simeq\Lambda_Y\) of pullbacks of constant sheaves. For \(K=\Lambda_Y\) with this structure, the previous paragraph says that the isomorphism \(\omega_f\simeq\Lambda_X(d)[2d]\) and the trace \(\operatorname{Tr}_f(\Lambda_Y)\) are morphisms of Weil complexes.

% object: theorem-poincare-duality-frobenius (blueprint/chap-Verdier-and-Poincare-duality/sec-Frobenius-and-Poincare-duality/theorem-poincare-duality-frobenius/_section.md)
\phantomsection\label{bp165}%
\begin{bpstatementit}{Theorem~10.21}{ (Poincaré duality and Frobenius)}
Let \(X_0/\mathbf F_q\) be smooth and separated of pure dimension \(d\), let \(X=X_0\times\overline{\mathbf F}_q\), and let \((K,\phi_K)\) be a Weil complex on \(X\) (Definition \hyperref[bp156]{10.18}) with \(K\) locally perfect. Equip the dual with the contragredient Weil structure

\[
\operatorname{Fr}_X^*K^\vee
\simeq(\operatorname{Fr}_X^*K)^\vee
\xrightarrow{\ (\phi_K^{-1})^\vee\ }K^\vee.
\]

For every admissible \(\Lambda\) the Poincaré pairing \(\beta_K\) of (10.7) is a morphism of Weil complexes. Thus, with \(F\) the Frobenius operators defined above, the pairing \(H^i_c(X,K)\times H^{2d-i}(X,K^\vee(d))\to\Lambda\) satisfies \(\langle Fx,Fy\rangle=\langle x,y\rangle\), and the pairing \(H^i_c(X,K)\times H^{2d-i}(X,K^\vee)\to\Lambda(-d)\) satisfies

\[
\langle Fx,Fy\rangle=q^d\langle x,y\rangle.
\]

If \(\Lambda\) is a field, then \(F\) is invertible on \(H^i_c(X,K)\), and with respect to the second pairing \(F\) on \(H^{2d-i}(X,K^\vee)\) is \(q^d\) times the inverse transpose of \(F\) on \(H^i_c(X,K)\). Hence the two spaces have the same dimension \(b\), the eigenvalues of \(F\) on \(H^{2d-i}(X,K^\vee)\) are the \(q^d/\alpha\), with multiplicities, where \(\alpha\) ranges over the eigenvalues on \(H^i_c(X,K)\), and \(\det(F\mid H^i_c(X,K))\det(F\mid H^{2d-i}(X,K^\vee))=q^{db}\).
\end{bpstatementit}

% proof: proof/theorem-poincare-duality-frobenius
\begin{proof}[{Proof}]\phantomsection\label{bp166}%
The pairing \(\beta_K\) is a composite of maps each of which commutes with pullback along Frobenius, so it is equivariant. Over a field the remaining statements then follow by linear algebra, because the pairing is perfect.

Cup product with mixed supports and evaluation are natural for pullback by Frobenius, and the isomorphism \(\gamma_{a_X}\) and the trace map are equivariant by Proposition \hyperref[bp162]{10.20}. Moreover, the contragredient Weil structure on \(K^\vee\) is chosen so that evaluation \(K\otimes K^\vee\to\Lambda_X\) is a morphism of Weil complexes, with the trivial structure on \(\Lambda_X\). It follows that the composite \(\beta_K\) is equivariant, which gives \(\langle Fx,Fy\rangle=\langle
x,y\rangle\) for the first pairing. The second pairing is obtained by moving the twist from the second factor to the values. Now geometric Frobenius acts on \(\Lambda(d)\) by \(q^{-d}\) and on \(\Lambda(-d)\) by \(q^d\), so for \(y\) in \(H^{2d-i}(X,K^\vee)\) the operator \(F\) on \(y\otimes\zeta\), with \(\zeta\) a generator of \(\Lambda(d)\), is \(q^{-d}(Fy)\otimes\zeta\). Removing the twist, we get the displayed identity \(\langle Fx,Fy\rangle=q^d\langle x,y\rangle\).

For field coefficients, Corollary \hyperref[bp140]{10.13} identifies \(H^{2d-i}(X,K^\vee)\), via the untwisted pairing, with the dual of \(H_c^i(X,K)\) tensored with \(\Lambda(-d)\). Moreover, Frobenius on \(H_c^i(X,K)\) is invertible, because \(\phi_K\) is an isomorphism and pullback along the automorphism \(\operatorname{Fr}_X\) is invertible on cohomology. In matrix terms, choose dual bases for the perfect pairing, let \(F_i\) be the matrix of \(F\) on \(H^i_c(X,K)\) and \(F_{2d-i}\) that on \(H^{2d-i}(X,K^\vee)\). Then the identity \(\langle Fx,Fy\rangle=q^d\langle x,y\rangle\) reads \({}^tF_iF_{2d-i}=q^d\). In other words, Frobenius on \(H^{2d-i}(X,K^\vee)\) is \(q^d\) times the inverse transpose of Frobenius on \(H_c^i(X,K)\). Since a matrix and its transpose have the same characteristic polynomial, this proves the eigenvalue statement with algebraic multiplicities. Finally, taking determinants of \(F_{2d-i}=q^d\,{}^t(F_i)^{-1}\) on spaces of dimension \(b\), we get \(\det F_{2d-i}=q^{db}(\det F_i)^{-1}\).
\end{proof}

% text: prose-10-24
\phantomsection\label{bp167}%
For \(X_0\) proper and geometrically connected and \(K=\Lambda_X\), Theorem \hyperref[bp146]{10.15} shows that \(F\) is the identity on \(H^0(X,\Lambda)=\Lambda\) and acts by \(q^d\) on \(H^{2d}(X,\Lambda)\simeq\Lambda(-d)\).

% object: example-duality-multiplicative-group (blueprint/chap-Verdier-and-Poincare-duality/sec-Frobenius-and-Poincare-duality/example-duality-multiplicative-group/_section.md)
\phantomsection\label{bp168}%
\begin{bpstatement}{Example~10.22}{ (Duality on the punctured line)}
Let \(X_0=\mathbf G_{m,\mathbf F_q}\), with coordinate \(t\), let \(X\) be its base change to \(k\), and let \(\Lambda=\overline{\mathbf Q}_\ell\), with the constant Weil sheaf. Since \(\operatorname{Pic}(X)=0\) and the units of \(k[t,t^{-1}]\) are the \(ct^m\), the Kummer sequence gives \(H^0(X,\Lambda)=\Lambda\) and \(H^1(X,\Lambda(1))=\Lambda\), generated by the Kummer class of \(t\) (Proposition \bpout{D.4}, with \(g=0\) and two punctures). Thus \(H^1(X,\Lambda)=\Lambda(-1)\), and the other groups vanish. For compact support, on the other hand, we write \(X=\mathbf P^1\setminus\{0,\infty\}\). The long exact sequence of \(j_!\Lambda_X\to\Lambda_{\mathbf P^1}\to\Lambda_{\{0,\infty\}}\to\) reads

\[
0\to H^0_c(X)\to\Lambda\xrightarrow{\ \text{diag}\ }\Lambda^2\to H^1_c(X)
\to H^1(\mathbf P^1)=0\to0\to H^2_c(X)\to H^2(\mathbf P^1)\to0,
\]

so \(H^0_c=0\), \(H^1_c=\Lambda^2/\Lambda\simeq\Lambda\) and \(H^2_c=H^2(\mathbf P^1,\Lambda)=\Lambda(-1)\). Poincaré duality with \(d=1\) pairs \(H^1_c\) with \(H^1\) and \(H^2_c\) with \(H^0\), and the dimensions match.

The Frobenius eigenvalues show the twist at work. On \(H^1_c\), which comes from the two rational points \(0\) and \(\infty\), Frobenius acts by \(1\), whereas on \(H^1=\Lambda(-1)\) it acts by \(q\). Their product is \(1\cdot q=q^d\), as Theorem \hyperref[bp165]{10.21} requires. Without the twist, duality would predict the eigenvalue \(1/1=1\) on \(H^1\) instead of \(q\). Likewise \(H^0\) and \(H^2_c\) have eigenvalues \(1\) and \(q\).
\end{bpstatement}

% object: example-running-example-duality (blueprint/chap-Verdier-and-Poincare-duality/sec-Frobenius-and-Poincare-duality/example-running-example-duality/_section.md)
\phantomsection\label{bp169}%
\begin{bpstatement}{Example~10.23}{ (The running example: duality)}
Let \(C_0\colon y^2=x^3-x\) over \(\mathbf F_3\) be the elliptic curve of the introduction, and \(C=C_0\times_{\mathbf F_3}k\). Since \(C\) has genus one, \(H^1(C,\mathbf Q_\ell)\) is a plane. Indeed, by the Kummer sequence it is the dual of the rational Tate module \(V_\ell J\) of the Jacobian, a \(\mathbf Q_\ell\)-vector space of dimension \(2g=2\) (Corollary \bpout{C.15} and Proposition \bpout{14.2}). For \(d=1\) the middle pairing on \(H^1\) is alternating and perfect, with values in \(\mathbf Q_\ell(-1)\) (Corollary \hyperref[bp152]{10.17}). By Theorem \hyperref[bp165]{10.21} it satisfies \(\langle Fx,Fy\rangle=3\langle x,y\rangle\). On a plane an alternating form is unique up to a scalar, so \(\langle Fx,Fy\rangle=\det(F)\langle
x,y\rangle\). Hence \(\det(F\mid H^1)=3\). If \(\alpha,\beta\) are the eigenvalues, then \(\alpha\beta=3\), and \(\det(1-TF\mid H^1)=1-(\alpha+\beta)T+3T^2\). This is the functional equation of a curve of genus one: replacing \(T\) by \(1/(3T)\) multiplies the polynomial by \((3T^2)^{-1}\).

Duality alone does not give \(|\alpha|=3^{1/2}\). For example, the pair \(\alpha=3\), \(\beta=1\), with \(\det(1-TF)=1-4T+3T^2\), consists of algebraic integers and satisfies \(\alpha\beta=3\). Here the count \(N_1=4\) gives \(\alpha+\beta=0\) and excludes it; in general only purity does.
\end{bpstatement}



% object: chap-The-Lefschetz-Verdier-formula (blueprint/chap-The-Lefschetz-Verdier-formula/_section.md), content hidden
\section*{\texorpdfstring{11\quad The Lefschetz–Verdier formula}{11 The Lefschetz–Verdier formula}}
\phantomsection\label{bp170}%
\addcontentsline{toc}{section}{\texorpdfstring{11\quad The Lefschetz–Verdier formula}{11 The Lefschetz–Verdier formula}}
% text: prose-11-01
\phantomsection\label{bp171}%
The trace formula has to express the trace of Frobenius on \(R\Gamma_c(X,K)\) as a sum over the fixed points of Frobenius, that is, over \(X_0(\mathbf F_{q^m})\). We do not prove it for Frobenius directly. Instead we first prove a fixed point formula for every correspondence on a proper scheme. In this generality the formula is compatible with proper pushforward, and pushing forward to a point turns it into a tautology.

For a general correspondence, however, the fixed locus need not be discrete (Section \hyperref[bp195]{11.4}). The local side is then a sum of contributions, one for each connected component of the fixed locus, each computed in a neighbourhood of that component.

We use Grothendieck's formalism of cohomological correspondences, written up by Illusie in [\hyperref[bp332]{SGA5}, III]. In this formalism, a correspondence on \(X\) is a map \(c=(c_1,c_2):C\to X\times X\), and a cohomological correspondence on \(K\) is a morphism \(u:c_2^*K\to c_1^!K\). For example, for the graph of an endomorphism \(\varphi\) this is a map \(\varphi^*K\to K\). In particular, for Frobenius it is the Weil structure \(\phi_K:\operatorname{Fr}_X^*K\simeq K\) of Section \hyperref[bp154]{10.5}, up to the comparison of \(\operatorname{Fr}_X\) with the \(k\)-linear Frobenius in Chapter \hyperref[bp226]{12}. When \(c_2\) is proper, \(u\) induces an endomorphism \(T_c(u)\) of \(R\Gamma_c(X,K)\). In the same way \(\phi_K\) induces the Frobenius operator \(\operatorname{Frob}_{X,K}\) on \(R\Gamma_c(X,K)\) and on every stalk at a fixed point.

The local contributions come from duality, as \(\operatorname{Hom}(V,V)=V\otimes
V^\vee\) turns an endomorphism into a tensor whose evaluation is the trace. In this way \(u\) gives the trace class \(\operatorname{Tr}_c(u)\in H^0(\Fix(c),\omega_{\Fix(c)})\), supported on the fixed locus. Its integral over a proper open and closed part \(B\) of \(\Fix(c)\) is the local term \(\operatorname{LT}_B(u)\). The Lefschetz–Verdier formula (Theorem \hyperref[bp206]{11.13}) says that if \(X\) and \(C\) are proper over \(k\), then

\[
\operatorname{tr}\bigl(T_c(u)\mid R\Gamma(X,K)\bigr)
=\sum_B\operatorname{LT}_B(u).
\]

The proof involves no computation, because trace classes commute with proper pushforward of correspondences. After pushing forward to the identity correspondence of the point, the trace class becomes the trace of \(T_c(u)\) on the perfect complex \(R\Gamma(X,K)\) (Lemma \hyperref[bp199]{11.10}). The price of this generality, however, is that the formula does not compute the individual local terms; for Frobenius that is done in Chapter \hyperref[bp226]{12}.

With coefficients in a finite ring \(A\) both sides are traces of endomorphisms of perfect complexes (Definition \hyperref[bp174]{11.1}). They are not alternating sums of traces on cohomology modules, since these need not be projective. In fact, the trace is not even determined by the maps induced on cohomology (Example \hyperref[bp181]{11.4}). Moreover, finite rings such as \(\mathcal O_E/\pi_E^n\) cannot be avoided, because the trace formula over \(\mathcal O_E\) is obtained from its reductions.

Throughout, \(k_0=\mathbf F_q\), \(k=\overline{\mathbf F}_q\) and \(s=\operatorname{Spec}k\). For a scheme \(X\) over \(k\), \(a_X:X\to s\) is the structure morphism, and \(X=X_0\times_{k_0}k\) whenever a model \(X_0\) over \(k_0\) is given. From Section \hyperref[bp172]{11.1} until Section \hyperref[bp289]{12.5} the coefficient ring is an arbitrary finite commutative \(\ell\)-primary ring \(A\), with \(\ell\ne\operatorname{char}k\), and we write \(\omega_X=a_X^!A_s\) and \(\mathbb D_XK=\cHom_A(K,\omega_X)\) for the dualizing complex and the Verdier dual of Definition \hyperref[bp108]{10.2} with coefficients \(A\). Some proofs compute on small étale sites and transport the result by means of the comparison theorems for complexes of finite Tor amplitude.

% object: sec-Traces-of-perfect-complexes (blueprint/chap-The-Lefschetz-Verdier-formula/sec-Traces-of-perfect-complexes/_section.md), content hidden
\subsection*{\texorpdfstring{11.1\quad Traces of perfect complexes}{11.1 Traces of perfect complexes}}
\phantomsection\label{bp172}%
\addcontentsline{toc}{subsection}{\texorpdfstring{11.1\quad Traces of perfect complexes}{11.1 Traces of perfect complexes}}
% text: prose-11-02
\phantomsection\label{bp173}%
Both sides of the trace formula are traces of endomorphisms of perfect \(A\)-complexes, namely of \(R\Gamma_c(X,K)\) on the global side and of the stalks of \(K\) on the local side.

Over a field the trace of an endomorphism of a complex is the alternating sum of its traces on cohomology. Over \(A=\mathbf Z/\ell^2\), however, the cohomology modules need not be free, so they have no trace, and Example \hyperref[bp181]{11.4} below shows that the maps on cohomology do not determine the trace. We therefore define the trace on the complex. On a bounded complex of finite projective modules it is the alternating sum of the traces of the components, but the definition below expresses this sum through evaluation and coevaluation, so that it depends only on the morphism in \(D(A)\). In this form it is compatible with every change of rings (Proposition \hyperref[bp176]{11.2}), and this is how we pass from the trace formula over \(\mathcal O_E/\pi_E^n\) to the trace formula over \(\mathcal O_E\) and \(E\) (Section \hyperref[bp289]{12.5}).

% object: definition-categorical-trace (blueprint/chap-The-Lefschetz-Verdier-formula/sec-Traces-of-perfect-complexes/definition-categorical-trace/_section.md)
\phantomsection\label{bp174}%
\begin{bpstatement}{Definition~11.1}{ (Categorical trace)}
For \(P\in\operatorname{Perf}(A)\) put \(P^\vee=R\!\operatorname{Hom}_A(P,A)\). If \(v:P\to P\) is a morphism in \(D(A)\), its trace is the scalar

\[
\operatorname{tr}_A(v\mid P)=
\left[A\xrightarrow{\operatorname{coev}}P\otimes^{\mathbf L} P^\vee
\xrightarrow{v\otimes^{\mathbf L}1}P\otimes^{\mathbf L} P^\vee
\xrightarrow{\sim}P^\vee\otimes^{\mathbf L} P
\xrightarrow{\operatorname{ev}}A\right]\in A.
\]

Here \(\operatorname{coev}\) and \(\operatorname{ev}\) are the coevaluation and evaluation of the dual pair \((P,P^\vee)\) in \(D(A)\), and the unlabelled arrow is the symmetry isomorphism, which includes the Koszul sign.
\end{bpstatement}

% text: prose-11-03
\phantomsection\label{bp175}%
For \(P=A\) in degree zero and \(v\) multiplication by \(a\), coevaluation is \(1\mapsto1\otimes1\) and evaluation is multiplication, so the trace is \(a\). The next proposition lists the properties of the trace that we use. Compatibility with change of rings lets us move the trace formula between coefficient rings, and the chain-level formula is how we compute traces in examples.

% object: proposition-properties-trace (blueprint/chap-The-Lefschetz-Verdier-formula/sec-Traces-of-perfect-complexes/proposition-properties-trace/_section.md)
\phantomsection\label{bp176}%
\begin{bpstatementit}{Proposition~11.2}{ (Properties of the trace)}
The trace is \(A\)-linear in \(v\), invariant under conjugacy, and compatible with derived scalar extension: for every ring map \(A\to B\),

\[
\operatorname{tr}_A(v\mid P)\longmapsto
\operatorname{tr}_B(v\otimes^{\mathbf L} B\mid P\otimes^{\mathbf L} B).
\]

For a bounded complex of finite projective modules and a chain endomorphism \(v^\bullet\) representing \(v\) one has

\[
\operatorname{tr}_A(v\mid P)
=\sum_i(-1)^i\operatorname{Tr}_A(v^i\mid P^i).
\]

If \(P\) is the underlying complex of a two-step filtered object \(F^1P\to P\to P/F^1P\) of the filtered derived category, with \(F^1P\) and \(P/F^1P\) perfect, and \(v\) is the underlying map of an endomorphism \(\tilde v\) of this filtered object, with graded pieces \(v_1\) on \(F^1P\) and \(v_0\) on \(P/F^1P\), then

\[
\operatorname{tr}_A(v\mid P)=\operatorname{tr}_A(v_1\mid F^1P)
+\operatorname{tr}_A(v_0\mid P/F^1P).
\]

Over a field \(L\) it is also

\[
\operatorname{tr}_L(v\mid P)
=\sum_i(-1)^i\operatorname{Tr}_L(H^i(v)\mid H^i(P)).
\]
\end{bpstatementit}

% proof: proof/proposition-properties-trace
\begin{proof}[{Proof}]\phantomsection\label{bp177}%
Everything follows from the chain formula, which we prove first by writing evaluation and coevaluation in a finite dual basis. If \(e_1,\dots,e_n\) and \(e^1,\dots,e^n\) are dual bases of \(A^n\) and its dual, then \(\operatorname{coev}(1)=\sum_je_j\otimes e^j\), and the composite of the definition sends \(1\) to \(\sum_je^j(v(e_j))\), the ordinary trace of a matrix. The same computation applies to a finite projective module, which is a direct summand of such an \(A^n\). On a complex, the symmetry interchanging a term of degree \(i\) and its dual of degree \(-i\) contributes \((-1)^i\), which proves the chain formula. Its right-hand side is invariant under chain homotopy, since the supertrace of \(dh+hd\) is zero by cyclicity and a shift of index.

Tensoring a dual basis with \(B\) preserves evaluation and coevaluation. Hence the trace is compatible with extension of scalars for every ring map \(A\to B\). Linearity and invariance under conjugacy follow from the same description. For additivity, note that a two-step filtered object of the filtered derived category may be represented by a filtered complex of projective modules whose filtration is split in each degree. The endomorphism is then block triangular, so its supertrace is the sum of the supertraces on the two graded pieces. Finally, over a field the boundaries \(B^i\) and cycles \(Z^i\) are \(v\)-invariant subspaces \(B^i\subset Z^i\subset P^i\), with \(Z^i/B^i=H^i(P)\), and the differential induces a \(v\)-equivariant isomorphism \(P^i/Z^i\simeq B^{i+1}\). Hence

\[
\operatorname{Tr}(v^i\mid P^i)=\operatorname{Tr}(v\mid B^i)
+\operatorname{Tr}(H^i(v)\mid H^i(P))+\operatorname{Tr}(v\mid B^{i+1}),
\]

and in the alternating sum the terms for the \(B^i\) cancel in pairs, leaving the cohomological formula.
\end{proof}

% text: prose-11-04
\phantomsection\label{bp178}%
In the additivity statement of the proposition we ask for an endomorphism of the filtered object, and not merely of the three vertices of a triangle. The reason is that the middle component is not determined by the outer ones, and even for a morphism of triangles the trace of the middle component need not be the sum of the outer traces.

% object: remark-coherent-additivity (blueprint/chap-The-Lefschetz-Verdier-formula/sec-Traces-of-perfect-complexes/remark-coherent-additivity/_section.md)
\phantomsection\label{bp179}%
\begin{bpstatement}{Remark~11.3}{ (Additivity needs a filtered endomorphism)}
The coherence condition is needed over a general ring: three endomorphisms of a distinguished triangle whose squares commute only in the triangulated category need not satisfy trace additivity [Stacks, \href{https://stacks.math.columbia.edu/tag/087J}{087J}]. Every localization sequence we use below is obtained functorially from the Frobenius correspondence and the filtration by an open subscheme and its closed complement, and hence it lifts to the filtered derived category.
\end{bpstatement}

% text: prose-11-05
\phantomsection\label{bp180}%
The following example shows why the trace has to be taken on the complex.

% object: example-trace-not-determined-cohomology (blueprint/chap-The-Lefschetz-Verdier-formula/sec-Traces-of-perfect-complexes/example-trace-not-determined-cohomology/_section.md)
\phantomsection\label{bp181}%
\begin{bpstatement}{Example~11.4}{ (A trace not determined by cohomology)}
Let \(A=\mathbf Z/\ell^2\) and let \(P=[A\xrightarrow{\ell}A]\), with the terms in degrees \(0\) and \(1\). This is the perfect complex from Section \bpout{3.2} whose cohomology modules \(H^0(P)=\ell A\) and \(H^1(P)=A/\ell A\) are not perfect. The chain maps \(1_P=(1,1)\) and \(v=(1+\ell,1)\) commute with the differential, since \(\ell(1+\ell)=\ell\) in \(A\), and both induce the identity on \(H^0(P)\) and on \(H^1(P)\). By the chain formula of Proposition \hyperref[bp176]{11.2},

\[
\operatorname{tr}_A(1_P\mid P)=1-1=0,
\qquad
\operatorname{tr}_A(v\mid P)=(1+\ell)-1=\ell\neq0.
\]

Thus, over a finite ring, the trace of an endomorphism of a perfect complex is not determined by the maps it induces on cohomology. The trace formula has to be stated for the operator on the complex.
\end{bpstatement}


% object: sec-Correspondences (blueprint/chap-The-Lefschetz-Verdier-formula/sec-Correspondences/_section.md), content hidden
\subsection*{\texorpdfstring{11.2\quad Correspondences}{11.2 Correspondences}}
\phantomsection\label{bp182}%
\addcontentsline{toc}{subsection}{\texorpdfstring{11.2\quad Correspondences}{11.2 Correspondences}}
% text: prose-11-06
\phantomsection\label{bp183}%
As in topology, a pair of maps \(c_1,c_2:C\to X\) should act on \(R\Gamma_c(X,K)\) by pulling back along \(c_2\) and pushing forward along \(c_1\). The pullback \(R\Gamma_c(X,K)\to R\Gamma_c(C,c_2^*K)\) exists when \(c_2\) is proper. To push forward we then need a map \(Rc_{1!}L\to K\), where \(L\) is the complex on \(C\) we have arrived at. By the adjunction \(Rc_{1!}\dashv c_1^!\) there is a canonical such map exactly when \(L=c_1^!K\), namely the counit. What connects the two steps is therefore a map \(u:c_2^*K\to c_1^!K\). The resulting endomorphism \(T_c(u)\) is defined in Definition \hyperref[bp202]{11.11}.

With \(c_1^*\) in place of \(c_1^!\) the second step fails, since there may be no nonzero map \(Rc_{1!}c_1^*K\to K\). For example, for the inclusion \(i\) of a rational point \(y\) into a smooth curve \(X\),

\[
\operatorname{Hom}(i_*i^*\Lambda_X,\Lambda_X)
=\operatorname{Hom}(\Lambda_y,i^!\Lambda_X)
=\operatorname{Hom}(\Lambda_y,\Lambda_y(-1)[-2])=0,
\]

so the only map \(Ri_!i^*\Lambda_X\to\Lambda_X\) is zero. When \(c_1\) is finite étale, smooth purity gives \(c_1^!=c_1^*\), and the counit is summation over the fibres. Then \(T_c(u)\) is pullback along \(c_2\) followed by a sum over the fibres of \(c_1\), as for a Hecke operator.

The Lefschetz–Verdier formula (Theorem \hyperref[bp206]{11.13}) is stated in this generality. For the trace formula we apply it to the graph of Frobenius (Definition \hyperref[bp232]{12.2}).

% object: definition-correspondence (blueprint/chap-The-Lefschetz-Verdier-formula/sec-Correspondences/definition-correspondence/_section.md)
\phantomsection\label{bp184}%
\begin{bpstatement}{Definition~11.5}{ (Correspondences)}
Let \(X\) be a separated \(k\)-scheme of finite type. A \emph{correspondence} on \(X\) is a morphism of \(k\)-schemes \(c=(c_1,c_2):C\to X\times_sX\) from a separated \(k\)-scheme of finite type (similarly over a base \(S\)). For \(K\in D^b_{\mathrm{ctf}}(X,A)\), a \emph{cohomological correspondence} on \((X,K)\) is a morphism

\[
u:c_2^*K\longrightarrow c_1^!K.
\]

The \emph{fixed locus} of \(c\) is the closed subscheme

\[
\Fix(c)=C\times_{X\times X,\Delta_X}X,
\qquad i_c:\Fix(c)\longrightarrow C,\quad d_c:\Fix(c)\longrightarrow X,
\]

where \(i_c\) and \(d_c\) are the two projections.
\end{bpstatement}

% text: prose-11-07
\phantomsection\label{bp185}%
Here \(c_2\) is the leg along which \(K\) is pulled back and \(c_1\) is the leg along which we take \(c_1^!\); in [\hyperref[bp332]{SGA5}, III, §3] the two names are interchanged. Since \(X\) is separated, \(\Delta_X\) and \(i_c\) are closed immersions.

The example to keep in mind is the graph of an endomorphism \(\varphi:X\to X\), that is, \(C=X\) with \(c_1=\operatorname{id}_X\) and \(c_2=\varphi\). Then \(c_1^!K=K\), and a cohomological correspondence is a morphism \(u:\varphi^*K\to K\). In this case the fixed locus \(\Fix(c)\) is the fixed point scheme of \(\varphi\), which we also denote \(\Fix(\varphi)\). For the Frobenius \(\Phi_X\) and a Weil complex this is the correspondence whose trace formula we want.


% object: sec-The-exceptional-Kunneth-formula (blueprint/chap-The-Lefschetz-Verdier-formula/sec-The-exceptional-Kunneth-formula/_section.md), content hidden
\subsection*{\texorpdfstring{11.3\quad The exceptional Künneth formula}{11.3 The exceptional Künneth formula}}
\phantomsection\label{bp186}%
\addcontentsline{toc}{subsection}{\texorpdfstring{11.3\quad The exceptional Künneth formula}{11.3 The exceptional Künneth formula}}
% text: prose-11-08
\phantomsection\label{bp187}%
A cohomological correspondence \(u:c_2^*K\to c_1^!K\) is a morphism on \(C\), so to obtain from it a class on the fixed locus we must first view it as a section of \(c^!\) of a single complex on \(X\times X\), which can then be restricted to the diagonal. In linear algebra this is the identification \(\operatorname{Hom}(V,W)=W\otimes V^\vee\), under which the trace becomes evaluation. Its sheaf version is the isomorphism \(\cHom(p_2^*K,p_1^!K)\simeq K\boxtimes\mathbb D_XK\) of Lemma \hyperref[bp191]{11.7}. Here \(p_1,p_2:X\times X\to X\) are the projections and \(M\boxtimes N=p_1^*M\otimes^{\mathbf L} p_2^*N\) is the external tensor product in \(D((X\times X)_{\mathrm{pro}\acute{e}t},A)\). Applying \(c^!\) to it, we obtain \(\cHom(c_2^*K,c_1^!K)\). On the point the isomorphism is exactly \(\operatorname{Hom}(V,V)=V\otimes V^\vee\).

The isomorphism comes from Künneth formulas for \(Rf_!\), \(Rf_*\), \(f^!\), \(\cHom\) and Verdier duals. For \(Rf_!\) the formula follows directly from base change for compact supports and the projection formula, and for ordinary direct images from smooth and proper base change (Proposition \bpout{8.18}, Section \bpout{8.9}). The formula for \(f^!\) (Proposition \hyperref[bp188]{11.6}), by contrast, needs a dévissage on the coefficients, starting from constant coefficients, where for smooth \(f\) the functor \(f^!\) is a twist and a shift.

Let \(f_i:X_i\to Y_i\), \(i=1,2\), be morphisms of separated \(k\)-schemes of finite type. For \(M_i\in D^b_{\mathrm{ctf}}(X_i,A)\), using base change for compact supports and the projection formula, we obtain the Künneth isomorphism

\[
R(f_1\times f_2)_!(M_1\boxtimes M_2)
\simeq Rf_{1!}M_1\boxtimes Rf_{2!}M_2.
\]

For the structure maps it reads \(R\Gamma_c(X_1\times X_2,M_1\boxtimes M_2)\simeq R\Gamma_c(X_1,M_1)\otimes^{\mathbf L} R\Gamma_c(X_2,M_2)\). For \(L_i\in D^b_{\mathrm{ctf}}(Y_i,A)\) let

\[
\alpha:f_1^!L_1\boxtimes f_2^!L_2\longrightarrow
(f_1\times f_2)^!(L_1\boxtimes L_2)
\]

be the morphism adjoint, under \(R(f_1\times f_2)_!\dashv(f_1\times f_2)^!\), to the composite

\[
\begin{aligned}
R(f_1\times f_2)_!(f_1^!L_1\boxtimes f_2^!L_2)
&\simeq Rf_{1!}f_1^!L_1\boxtimes Rf_{2!}f_2^!L_2\\
&\xrightarrow{\varepsilon_{f_1}\boxtimes\varepsilon_{f_2}}
L_1\boxtimes L_2,
\end{aligned}
\]

where \(\varepsilon_{f_i}\) is the counit of \(Rf_{i!}\dashv f_i^!\).

% object: proposition-exceptional-external-kunneth (blueprint/chap-The-Lefschetz-Verdier-formula/sec-The-exceptional-Kunneth-formula/proposition-exceptional-external-kunneth/_section.md)
\phantomsection\label{bp188}%
\begin{bpstatementit}{Proposition~11.6}{ (Exceptional external Künneth)}
For all \(L_i\in D^b_{\mathrm{ctf}}(Y_i,A)\), the morphism \(\alpha\) is an isomorphism.
\end{bpstatementit}

% proof: proof/proposition-exceptional-external-kunneth
\begin{proof}[{Proof}]\phantomsection\label{bp189}%
Both sides of \(\alpha\) are exact in each variable, so by dévissage it is enough to check \(\alpha\) on generators. First, for constant coefficients we compute \(f_i^!\) by Theorem \hyperref[bp11]{9.3} for smooth maps and by localization, and check \(\alpha\) directly. Next, general coefficients are built from extensions by zero \(j_!F\) of lisse sheaves. Since \(f^!\) does not commute with \(j_!\), we replace \(j_!\) by \(Rj_*\), which does. The difference has smaller support, and we treat it by induction. For the classical statements we refer to [\hyperref[bp332]{SGA5}, III, (1.6.4), (1.7.3), Proposition 1.7.4 and (3.1.1)].

\emph{Constant and lisse coefficients.} First let the coefficients be constant. For smooth \(f_i\) of relative dimensions \(d_i\), Theorem \hyperref[bp11]{9.3} identifies \(f_i^!L_i\) with \(f_i^*L_i(d_i)[2d_i]\), and \((f_1\times f_2)^!\) with \((f_1\times f_2)^*(d_1+d_2)[2d_1+2d_2]\), so both sides of \(\alpha\) become the same twisted pullback. That \(\alpha\) is this identification follows from the compatibility of the trace map (Definition \hyperref[bp85]{9.24}) with products. Indeed, \(\alpha\) is defined through the counits, and for a smooth map the trace map is the counit composed with \(\gamma_f\).

For closed immersions we use the localization triangle and the case of ordinary external Künneth where one factor is the identity, applied to the complementary open immersions. In general, every map factors locally as a closed immersion followed by a smooth projection of affine spaces. Since \(\alpha\) is compatible with composition of morphisms and with restriction to open subschemes, Zariski localization gives the constant case for general maps. Next, lisse finite flat coefficients become constant with finite projective values on an étale cover, so descent gives the lisse case, and then the case of bounded perfect lisse models. In each of these steps the isomorphism we obtain is \(\alpha\) itself: for smooth maps because \(\gamma_f\) identifies the counit with the trace map, and for localization and descent because \(\alpha\) is natural and compatible with restriction.

\emph{General coefficients.} A constructible complex is lisse on a dense open, and the rest has smaller support. We therefore argue by induction on the sum of the dimensions of the supports of \(L_1,L_2\). Pushforward along closed immersions is compatible with \(\alpha\) by proper base change and the exceptional exchange of Corollary \hyperref[bp128]{10.9}. Hence we may first assume that the support is the entire base. As both sides of \(\alpha\) are exact in each variable, a finite constructible stratification, bounded flat models and two-out-of-three reduce us to generators \(j_{i!}F_i\) for dense open \(j_i:U_i\hookrightarrow Y_i\), with \(F_i\) a lisse sheaf with finite projective stalks (Lemma \bpout{3.13}).

Since \(f_i^!\) does not commute with \(j_{i!}\) in general, we pass through \(Rj_{i*}\). Let \(j'_i\) and \(f_{i,U}\) be the base changes of \(j_i\) and \(f_i\). The exceptional exchange \(f_i^!Rj_{i*}\simeq Rj'_{i*}f_{i,U}^!\) is the mate, under the right adjoints, of base change for \(Rf_{i!}\) along the open immersion. Hence it is compatible with the counits used to define \(\alpha\). By this exchange and the ordinary external Künneth formula, the comparison on \((Rj_{1*}F_1,Rj_{2*}F_2)\) is the pushforward of the lisse comparison. Let \(z_i:Z_i\hookrightarrow Y_i\) be the closed complement of \(U_i\). In the triangles

\[
j_{i!}F_i\longrightarrow Rj_{i*}F_i\longrightarrow
z_{i*}z_i^*Rj_{i*}F_i\xrightarrow{+1},
\]

the third term is still constructible of finite Tor amplitude by Proposition \bpout{7.1}, and its support has smaller dimension. Since \(\alpha\) is natural, it maps the corresponding triangles to each other in each variable. By induction and two-out-of-three, first in one variable and then in the other, we obtain \(\alpha\) for \((j_{1!}F_1,j_{2!}F_2)\).
\end{proof}

% text: prose-11-09
\phantomsection\label{bp190}%
Using the Künneth formulas for both kinds of direct image and for \(f^!\), we now identify cohomological correspondences with sections of \(c^!(K\boxtimes\mathbb D_XK)\).

% object: lemma-external-duality-trace-morphism (blueprint/chap-The-Lefschetz-Verdier-formula/sec-The-exceptional-Kunneth-formula/lemma-external-duality-trace-morphism/_section.md)
\phantomsection\label{bp191}%
\begin{bpstatementit}{Lemma~11.7}{ (External duality)}
Let \(K\in D^b_{\mathrm{ctf}}(X,A)\), let \(p_1,p_2:X\times X\to X\) be the projections, and put \(P_X=K\boxtimes \mathbb D_XK\in D^b_{\mathrm{ctf}}(X\times X,A)\). For a correspondence \(c:C\to X\times X\) there are canonical isomorphisms

\[
P_X\xrightarrow{\sim}\cHom_A(p_2^*K,p_1^!K),
\qquad
\mathcal H_c(K):=\cHom_A(c_2^*K,c_1^!K)
\xrightarrow{\sim}c^!P_X.
\]
\end{bpstatementit}

% proof: proof/lemma-external-duality-trace-morphism
\begin{proof}[{Proof}]\phantomsection\label{bp192}%
The exceptional Künneth formula (Proposition \hyperref[bp188]{11.6}) and the ordinary one give Künneth formulas for internal Hom and for Verdier duals, and these give the two displayed isomorphisms.

\emph{Internal Hom and Verdier duals.} On the point this is the identity \(\Hom(V_1,W_1)\otimes\Hom(V_2,W_2)=\Hom(V_1\otimes V_2,W_1\otimes W_2)\) for finite free modules. For \(E_i,L_i\in D^b_{\mathrm{ctf}}(X_i,A)\) the external comparison \(\cHom_A(E_1,L_1)\boxtimes\cHom_A(E_2,L_2)\to
\cHom_A(E_1\boxtimes E_2,L_1\boxtimes L_2)\) is an isomorphism. By dévissage it suffices to check this on generators \(E_i=Rh_{i!}F_i\), with \(h_i\) the inclusion of a stratum and \(F_i\) a lisse sheaf with finite projective stalks (Lemma \bpout{3.13}), for which

\[
\cHom_A(Rh_{i!}F_i,L_i)
\simeq Rh_{i*}(F_i^\vee\otimes^{\mathbf L} h_i^!L_i).
\]

By the ordinary Künneth formula and \(\alpha\), the exterior product of these expressions is identified with internal Hom on the product. Taking \(L_i=\omega_{X_i}\) and using \(\alpha\) for the structure maps, we get, for \(M_i\in D^b_{\mathrm{ctf}}(X_i,A)\),

\[
\mathbb D_{X_1\times X_2}(M_1\boxtimes M_2)
\simeq \mathbb D_{X_1}M_1\boxtimes \mathbb D_{X_2}M_2.
\]

\emph{The two isomorphisms.} The first isomorphism is the sheaf version of \(\Hom(V,W)\simeq W\otimes V^\vee\), and we obtain it by writing internal Hom as a dual of a tensor product. Indeed, for \(M\in D^b_{\mathrm{ctf}}(U,A)\), \(N\in D^b_{\mathrm{ctf}}(V,A)\) and the projections \(p_U,p_V\) of \(V\times U\), Verdier exchange and biduality let us rewrite \(\cHom_A(p_U^*M,p_V^!N)\) as \(\mathbb D_{V\times U}(\mathbb D_VN\boxtimes M)\). By the formula for duals of exterior products this is \(N\boxtimes \mathbb D_UM\). Taking \(M=N=K\), with the factors of the product ordered as above, we obtain the first isomorphism. The second isomorphism follows from the canonical exchange of exceptional inverse image with Hom, for \(P,Q\in D^b_{\mathrm{ctf}}(X\times X,A)\),

\[
c^!\cHom_A(P,Q)\simeq\cHom_A(c^*P,c^!Q),
\]

applied to \(P=p_2^*K\) and \(Q=p_1^!K\). To prove the exchange, we use Yoneda inside \(\Loc(D^b_{\mathrm{ctf}}(C,A))\), the category on which \(c_!\dashv c^!\) is defined (Theorem \bpout{5.2}). Test objects in all of \(D(C_{\mathrm{pro}\acute{e}t},A)\) would not do, since \(c^!\) is not defined there. Both sides of the exchange lie in \(D^b_{\mathrm{ctf}}(C,A)\) by Proposition \bpout{7.1}, so these test objects suffice. By Proposition \bpout{5.15} the localizing categories are \(D(C_{\acute{e}t},A)\) and \(D((X\times X)_{\acute{e}t},A)\), where the projection formula for \(c_!\) holds for all objects [Stacks, \href{https://stacks.math.columbia.edu/tag/0GL5}{0GL5}], and tensoring with \(c^*P\) preserves \(\Loc(D^b_{\mathrm{ctf}}(C,A))\). For \(M\) in \(\Loc(D^b_{\mathrm{ctf}}(C,A))\) we obtain

\[
\begin{aligned}
\Hom\bigl(M,c^!\cHom_A(P,Q)\bigr)
&\simeq\Hom\bigl(c_!M\otimes^{\mathbf L}P,Q\bigr)
\simeq\Hom\bigl(c_!(M\otimes^{\mathbf L}c^*P),Q\bigr)\\
&\simeq\Hom\bigl(M\otimes^{\mathbf L}c^*P,c^!Q\bigr)
\simeq\Hom\bigl(M,\cHom_A(c^*P,c^!Q)\bigr),
\end{aligned}
\]

by \(c_!\dashv c^!\) and the tensor–Hom adjunction, the projection formula, \(c_!\dashv c^!\) again, and the tensor–Hom adjunction.
\end{proof}

% text: prose-11-10
\phantomsection\label{bp193}%
In linear algebra the trace is evaluation \(W\otimes W^\vee\to A\) on the diagonal. Here evaluation takes values in \(\omega_X\) on the diagonal, and by applying \(c^!\) we move it to the fixed locus.

% object: definition-trace-morphism (blueprint/chap-The-Lefschetz-Verdier-formula/sec-The-exceptional-Kunneth-formula/definition-trace-morphism/_section.md)
\phantomsection\label{bp194}%
\begin{bpstatement}{Definition~11.8}{ (Trace morphism)}
In the setting of Lemma \hyperref[bp191]{11.7}, let \(e_K:P_X\to R\Delta_{X*}\omega_X\) be adjoint to the evaluation \(\Delta_X^*P_X=K\otimes^{\mathbf L} \mathbb D_XK\to\omega_X\). The \emph{trace morphism} is the composite

\[
\mathcal H_c(K)\simeq c^!P_X\xrightarrow{c^!e_K}c^!R\Delta_{X*}\omega_X
\xrightarrow{\sim}Ri_{c*}d_c^!\omega_X
\simeq Ri_{c*}\omega_{\Fix(c)}
\]

of the second isomorphism of the lemma, \(c^!e_K\), exceptional base change for the proper map \(\Delta_X\), and the composition of exceptional inverse images \(d_c^!a_X^!=a_{\Fix(c)}^!\).
\end{bpstatement}


% object: sec-Trace-classes-and-local-terms (blueprint/chap-The-Lefschetz-Verdier-formula/sec-Trace-classes-and-local-terms/_section.md), content hidden
\subsection*{\texorpdfstring{11.4\quad Trace classes and local terms}{11.4 Trace classes and local terms}}
\phantomsection\label{bp195}%
\addcontentsline{toc}{subsection}{\texorpdfstring{11.4\quad Trace classes and local terms}{11.4 Trace classes and local terms}}
% text: prose-11-11
\phantomsection\label{bp196}%
We want to write \(\operatorname{tr}(T_c(u))\) as a sum of contributions from the fixed locus. The first guess, a sum over fixed points of the traces of \(u\) on stalks, does not make sense in general, because the fixed locus need not be finite. For example, for the identity correspondence of \(\mathbf P^1\) and \(K=A\), every point is fixed. The trace on \(R\Gamma(\mathbf P^1,A)=A\oplus A(-1)[-2]\) is \(1+1=2\), the self-intersection of the diagonal.

What does make sense is a class on \(\Fix(c)\). By Lemma \hyperref[bp191]{11.7}, \(u\) is a section of \(c^!(K\boxtimes\mathbb D_XK)\), and composing with the trace morphism we obtain a section of \(\omega_{\Fix(c)}\). Such a section can be integrated over a proper piece \(B\) of the fixed locus. Indeed, \(H^0(B,\omega_B)=\operatorname{Hom}(Ra_{B!}A_B,A)\), and \(Ra_{B*}A_B=Ra_{B!}A_B\) for proper \(B\). Composing with the unit \(A\to Ra_{B*}A_B\) therefore gives an element of \(A\). In the example above, \(H^0(\mathbf P^1,\omega_{\mathbf P^1})=H^2(\mathbf P^1,A(1))\) and the degree map \(H^2(\mathbf P^1,A(1))\to A\) is an isomorphism, so Theorem \hyperref[bp206]{11.13} shows that the trace class is \(2\) times the class of a point.

We will see that the trace class is additive in \(K\) (Proposition \hyperref[bp223]{11.19}) and compatible with proper pushforward (Theorem \hyperref[bp206]{11.13}). However, at an isolated fixed point the local term can still differ from the trace on the stalk (Remark \hyperref[bp239]{12.4}). For Frobenius the two are shown to agree in Proposition \hyperref[bp274]{12.16}.

% object: definition-trace-class-local-terms (blueprint/chap-The-Lefschetz-Verdier-formula/sec-Trace-classes-and-local-terms/definition-trace-class-local-terms/_section.md)
\phantomsection\label{bp197}%
\begin{bpstatement}{Definition~11.9}{ (Trace class and local terms)}
Let \(u:c_2^*K\to c_1^!K\) be a cohomological correspondence on \((X,K)\). By adjunction, \(u\) is a section \(A_C\to\mathcal H_c(K)\). Its composite with the trace morphism of Definition \hyperref[bp194]{11.8}, transposed by \(i_c^*\dashv Ri_{c*}\), is the \emph{trace class}

\[
\operatorname{Tr}_c(u):A_{\Fix(c)}\longrightarrow\omega_{\Fix(c)}.
\]

For a proper \(k\)-scheme \(B\), integration is the counit map

\[
\int_B:H^0(B,\omega_B)
=H^0(Ra_{B!}a_B^!A_s)\longrightarrow A.
\]

If \(B\subset \Fix(c)\) is proper and open and closed, the \emph{local term} of \(u\) on \(B\) is

\[
\operatorname{LT}_B(u)=\int_B\operatorname{Tr}_c(u)|_B.
\]
\end{bpstatement}

% text: prose-11-12
\phantomsection\label{bp198}%
On the point, the local term is the categorical trace of Definition \hyperref[bp174]{11.1}, and we use this later to identify the local terms of Frobenius with traces on stalks.

% object: lemma-local-terms-point (blueprint/chap-The-Lefschetz-Verdier-formula/sec-Trace-classes-and-local-terms/lemma-local-terms-point/_section.md)
\phantomsection\label{bp199}%
\begin{bpstatementit}{Lemma~11.10}{ (Local terms on the point)}
Let \(P\in\operatorname{Perf}(A)\), viewed as an object of \(D^b_{\mathrm{ctf}}(s,A)\), and let \(u:P\to P\) be a cohomological correspondence on \((s,P)\) for the identity correspondence of \(s\), whose fixed locus is \(s\). Then

\[
\operatorname{LT}_s(u)=\operatorname{tr}_A(u\mid P).
\]
\end{bpstatementit}

% proof: proof/lemma-local-terms-point
\begin{proof}[{Proof}]\phantomsection\label{bp200}%
On the point each step of the construction of the trace class is a step of the categorical trace. Here \(\omega_s=A_s\), so \(\mathbb D_sP=P^\vee\), and \(P_s=P\boxtimes
P^\vee=P\otimes^{\mathbf L}P^\vee\). First, under the first isomorphism of Lemma \hyperref[bp191]{11.7}, the map \(u\) becomes the element of \(P\otimes^{\mathbf L}P^\vee\) obtained by applying \(u\otimes1\) to the coevaluation. Next, the diagonal of \(s\) is the identity, so the trace morphism is the evaluation \(P\otimes^{\mathbf L}P^\vee\to A\), taken after the symmetry that puts the factors in the order \(P^\vee\otimes P\). That symmetry is precisely the symmetry in Definition \hyperref[bp174]{11.1}, and integration over \(s\) is the identity of \(A\). The composite is therefore the categorical trace.
\end{proof}

% text: prose-11-13
\phantomsection\label{bp201}%
The global side of the trace formula is the trace of the following endomorphism.

% object: definition-action-correspondence (blueprint/chap-The-Lefschetz-Verdier-formula/sec-Trace-classes-and-local-terms/definition-action-correspondence/_section.md)
\phantomsection\label{bp202}%
\begin{bpstatement}{Definition~11.11}{ (Action of a correspondence on cohomology)}
Let \(u\) be a cohomological correspondence on \((X,K)\) for a correspondence \(c\) with \(c_2\) proper. The endomorphism \(T_c(u)\) of \(R\Gamma_c(X,K)\) is the composite

\[
\begin{aligned}
T_c(u):R\Gamma_c(X,K)&\xrightarrow{\eta_{c_2}}
R\Gamma_c(C,c_2^*K)\xrightarrow{u}R\Gamma_c(C,c_1^!K)\\
&\simeq R\Gamma_c(X,Rc_{1!}c_1^!K)
\xrightarrow{\varepsilon_{c_1}}R\Gamma_c(X,K).
\end{aligned}
\]
\end{bpstatement}

% text: prose-11-14
\phantomsection\label{bp203}%
In the first arrow we use the ordinary unit together with \(Rc_{2*}=Rc_{2!}\), which holds since \(c_2\) is proper, and the displayed isomorphism comes from the composition of direct images with compact support.

The Lefschetz–Verdier formula computes the trace of \(T_c(u)\) in terms of the trace class. We obtain it as the special case, for the map to a point, of the compatibility of trace classes with proper pushforward. This works because \(T_c(u)\) is itself a cohomological correspondence: for proper \(X\) and \(C\) it is a correspondence on the point with coefficients \(R\Gamma(X,K)\), and on the point the trace class is the trace by Lemma \hyperref[bp199]{11.10}. The global trace is therefore the trace class of a pushed-forward correspondence, and it remains to show that trace classes commute with pushforward along proper maps.

To push a cohomological correspondence forward along proper morphisms, we must produce from \(u\) a map \(c_2'^*Rf_*K\to c_1'^!Rf_*K\). The only maps available for this are the units and counits of the adjunctions on the two legs, used as in the construction of \(T_c(u)\). Here properness of \(f\) and \(g\) gives \(Rf_*=Rf_!\) and \(Rg_*=Rg_!\), and this allows us to push forward the counit for \(Rc_{1!}\dashv c_1^!\). We also push forward along closed immersions (Section \hyperref[bp209]{11.5}).

% object: definition-pushforward-correspondence (blueprint/chap-The-Lefschetz-Verdier-formula/sec-Trace-classes-and-local-terms/definition-pushforward-correspondence/_section.md)
\phantomsection\label{bp204}%
\begin{bpstatement}{Definition~11.12}{ (Pushforward of a correspondence)}
Let \(f:X\to Y\) and \(g:C\to C'\) be proper morphisms, and let \(c'=(c'_1,c'_2):C'\to Y\times Y\) be a correspondence with \(c'_ig=fc_i\) for \(i=1,2\). Put \(K'=Rf_*K\in D^b_{\mathrm{ctf}}(Y,A)\). For a cohomological correspondence \(u\) on \((X,K)\), the \emph{pushforward} \(u':c_2'^*K'\to c_1'^!K'\) of \(u\) is the composite \(u'=\beta\,Rg_*u\,\alpha\), where \(\alpha\) is built from the unit and counit of \((\,)^*\dashv R(\,)_*\),

\[
\alpha:c_2'^*Rf_*K\xrightarrow{\eta_g}
Rg_*c_2^*f^*Rf_*K\xrightarrow{\varepsilon_f}Rg_*c_2^*K,
\]

and \(\beta:Rg_*c_1^!K\to c_1'^!Rf_*K\) is built from the counit of \(R(\,)_!\dashv(\,)^!\): it is adjoint to

\[
Rc'_{1!}Rg_*c_1^!K\simeq Rf_*Rc_{1!}c_1^!K
\xrightarrow{Rf_*\varepsilon_{c_1}}Rf_*K.
\]

The map \(g\) induces a proper morphism \(h:\Fix(c)\to \Fix(c')\) of fixed loci, and \(h_*:H^0(\Fix(c),\omega_{\Fix(c)})\to H^0(\Fix(c'),\omega_{\Fix(c')})\) denotes the map given by the ordinary unit, proper pushforward and the exceptional counit.
\end{bpstatement}

% text: prose-11-15
\phantomsection\label{bp205}%
For \(Y=s\), \(C'=s\) and \(f=a_X\), \(g=a_C\) with \(X\) and \(C\) proper, the pushforward \(u'\) is an endomorphism of \(R\Gamma(X,K)\). Unwinding \(\alpha\) and \(\beta\) shows that it is \(T_c(u)\), and the following theorem compares trace classes before and after pushforward. In this special case it is the Lefschetz–Verdier formula.

% object: theorem-proper-pushforward-lefschetz-verdier (blueprint/chap-The-Lefschetz-Verdier-formula/sec-Trace-classes-and-local-terms/theorem-proper-pushforward-lefschetz-verdier/_section.md)
\phantomsection\label{bp206}%
\begin{bpstatementit}{Theorem~11.13}{ (Proper pushforward and Lefschetz–Verdier)}
Let \(f:X\to Y\) and \(g:C\to C'\) be proper, and let \(c'=(c'_1,c'_2):C'\to Y\times Y\) satisfy \(c'_ig=fc_i\). Set \(K'=Rf_*K\), and let \(u'\) be the pushforward of \(u\) to \((Y,K')\) of Definition \hyperref[bp204]{11.12}. With the induced proper map \(h:\Fix(c)\to \Fix(c')\) one has

\[
\operatorname{Tr}_{c'}(u')=h_*\operatorname{Tr}_c(u).
\]

If \(X\) and \(C\) are proper over \(k\), this gives

\[
\operatorname{tr}_A(T_c(u)\mid R\Gamma(X,K))
=\int_{\Fix(c)}\operatorname{Tr}_c(u)
=\sum_B\operatorname{LT}_B(u),
\]

for any finite open-and-closed decomposition \(\Fix(c)=\coprod B\).
\end{bpstatementit}

% proof: proof/theorem-proper-pushforward-lefschetz-verdier
\begin{proof}[{Proof}]\phantomsection\label{bp207}%
The trace class of \(u\) is built in three steps: \(u\) becomes a map into \(P_X=K\boxtimes\mathbb D_XK\), evaluation maps \(P_X\) to the dualizing complex on the diagonal, and \(c^!\) moves the result to the fixed locus. We check that each step commutes with pushforward. For the first two steps this means identifying the map adjoint to \(u'\) with the pushforward of the map adjoint to \(u\), and evaluation for \(K'\) with the pushforward of evaluation for \(K\). The third step needs more care. The squares \(c'_ig=fc_i\) are only required to commute, so in general \(\Fix(c)\ne C\times_{C'}\Fix(c')\), and we compare the two fixed loci through an intermediate scheme.

\emph{Maps into \(P\).} Put \(P_Y=K'\boxtimes \mathbb D_YK'\), and let \(\tau_f:Rf_*\omega_X\to\omega_Y\) be the counit. By the Künneth formula for the proper map \(f\times f\) and the exchange \(\mathbb D_YRf_*\simeq Rf_*\mathbb D_X\) of Verdier duality, we obtain an isomorphism

\[
\rho_f:R(f\times f)_*P_X\xrightarrow{\sim}P_Y.
\]

A cohomological correspondence \(u\) defines a map \(\theta_u:Rc_!A_C\to P_X\), obtained by exceptional adjunction from the section \(A_C\to c^!P_X\). Under \(Rc'_!Rg_*\simeq R(f\times f)_*Rc_!\), the map defined by the pushforward \(u'\) is

\[
\theta_{u'}=
\rho_f\,R(f\times f)_*\theta_u\,
(Rc'_!Rg_*\simeq R(f\times f)_*Rc_!)\,Rc'_!\eta_g.
\]

To verify this identity, we apply tensor–Hom adjunction and then the exceptional adjunction for \(c'_1\). After this both sides become \(Rc'_{1!}\alpha\), followed by \(Rf_*Rc_{1!}u\) and \(Rf_*\varepsilon_{c_1}\). The two sides use the projection formulas for \(c_1\) and for \(c'_1\) in different orders, but they agree because \(c'_1g=fc_1\) and projection formulas are compatible with composition.

\emph{Evaluation.} Evaluation is compatible with \(\rho_f\), in the sense that

\[
e_{K'}\rho_f=
R\Delta_{Y*}\tau_f\,R(f\times f)_*e_K,
\qquad (f\times f)\Delta_X=\Delta_Yf.
\]

After restricting to the diagonal and taking tensor–Hom adjoints, this says that Verdier duality for the proper map \(f\) identifies evaluation against \(Rf_*K\) with the pushforward of evaluation against \(K\), followed by \(\tau_f\). This holds because that duality is defined by means of the counit \(\tau_f\).

\emph{Fixed loci.} The two maps into \(P_Y\) and the two evaluations therefore match. The trace classes, however, live on two different fixed loci, and we must see that passing from \(\Fix(c)\) to \(\Fix(c')\) is the map \(h_*\) of Definition \hyperref[bp204]{11.12}. Recall that the fixed loci are the pullbacks of the diagonals, \(\Fix(c)=C\times_{X\times X}X\) and \(\Fix(c')=C'\times_{Y\times Y}Y\). To compare them, we factor \(h\) through the fibre product \(C\times_{C'}\Fix(c')\):

\[
\Fix(c)\xrightarrow{i'}C\times_{C'}\Fix(c')\xrightarrow{\pi}\Fix(c'),
\]

where the closed immersion \(i'\) is the pullback of \(X\to X\times_YX\) and \(\pi\) is the base change of \(g\). We use proper base change for the square defining \(C\times_{C'}\Fix(c')\), and the exceptional counit for \(i'\). Since counits are compatible with composition, the integration map for \(h\) is \(\tau_h=\tau_\pi R\pi_*\tau_{i'}\). Applying the identity for \(\theta_{u'}\) and the compatibility of evaluation, together with these base change maps, we identify the map \(A_{C'}\to Ri_{c'*}\omega_{\Fix(c')}\) defining the left-hand trace with

\[
A_{C'}\xrightarrow{\eta_g}Rg_*A_C
\longrightarrow Rg_*Ri_{c*}\omega_{\Fix(c)}
=Ri_{c'*}Rh_*\omega_{\Fix(c)}
\xrightarrow{Ri_{c'*}\tau_h}Ri_{c'*}\omega_{\Fix(c')}.
\]

Taking the \(i_{c'}^*\)-adjoint by means of the identity of composite units \(Rg_*\eta_{i_c}\,\eta_g=Ri_{c'*}\eta_h\,\eta_{i_{c'}}\), we obtain the first formula (compare [\hyperref[bp332]{SGA5}, III, Theorem 4.4 and Corollary 4.7]).

\emph{The trace formula.} For proper \(X\) and \(C\), we push forward to the identity correspondence on \(s\), using \(f=a_X\) and \(g=a_C\). Then the pushed cohomological correspondence is \(T_c(u)\). Indeed, for \(g=a_C\) the map \(\alpha\) is the pullback \(R\Gamma(X,K)\to R\Gamma(C,c_2^*K)\) and \(\beta\) is the counit of \(c_1\) after \(R\Gamma\), and these are the outer arrows in the definition of \(T_c(u)\). Moreover, the fixed locus of the identity correspondence of \(s\) is \(s\) itself. Hence the left side of the first formula is the local term of \(T_c(u)\) on \(s\), which by Lemma \hyperref[bp199]{11.10} is the trace of \(T_c(u)\). On the right, \(h_*\) is integration over \(\Fix(c)\), and since integration over a disjoint union is the sum of the integrations, this gives the second assertion.
\end{proof}

% object: example-correspondence-on-point (blueprint/chap-The-Lefschetz-Verdier-formula/sec-Trace-classes-and-local-terms/example-correspondence-on-point/_section.md)
\phantomsection\label{bp208}%
\begin{bpstatement}{Example~11.14}{ (A correspondence on the point)}
Let \(X=s\), let \(C=\coprod_{i=1}^ns\) be \(n\) points, with \(c_1=c_2\) the structure map, and let \(u:c_2^*A=A_C\to c_1^!A=A_C\) be the identity, using \(c_1^!=c_1^*\) for the finite étale map \(c_1\). Then pullback along \(c_2\) is the diagonal \(A\to A^n\), and the counit \(Rc_{1!}c_1^!A=A^n\to A\) is the sum. Hence \(T_c(u)\) is multiplication by \(n\) on \(R\Gamma(s,A)=A\), with trace \(n\). Since \(X\times X=s\), the fixed locus is all of \(C\). On each of these \(n\) points \(u\) restricts to the identity correspondence of \(s\) on \(A\), whose local term is \(\operatorname{tr}_A(1\mid A)=1\) by Lemma \hyperref[bp199]{11.10}. The right side of the theorem is therefore \(n\) as well. Note that the local terms are indexed by the points of \(C\), not of \(X\): here \(X\) has one point and the fixed locus has \(n\).
\end{bpstatement}


% object: sec-Additivity-of-trace-classes (blueprint/chap-The-Lefschetz-Verdier-formula/sec-Additivity-of-trace-classes/_section.md), content hidden
\subsection*{\texorpdfstring{11.5\quad Additivity of trace classes}{11.5 Additivity of trace classes}}
\phantomsection\label{bp209}%
\addcontentsline{toc}{subsection}{\texorpdfstring{11.5\quad Additivity of trace classes}{11.5 Additivity of trace classes}}
% text: prose-11-16
\phantomsection\label{bp210}%
To compute local terms, we cut \(K\) into pieces along a decomposition of \(X\) into an open and a closed part, both invariant under the correspondence. For this the trace class must be additive. For the trace of an endomorphism of a perfect complex, additivity says that a block triangular matrix has trace equal to the sum of the traces of its diagonal blocks. By Remark \hyperref[bp179]{11.3} it can only be expected for filtrations lifted to the filtered derived category. The trace class needs the same care, and we lift each step of the construction of the trace morphism to filtered complexes.

In a two-step filtration, for example, a filtered correspondence is a block triangular matrix \(\left(\begin{smallmatrix}u_0&*\\0&u_1\end{smallmatrix}\right)\) whose off-diagonal entry has weight \(1\). Since \(\omega_X\) sits in filtration degree \(0\), the trace class does not see this entry (Lemma \hyperref[bp214]{11.16} below).

The operations that enter the trace morphism have filtered versions, with explicit formulas for their graded pieces. We use finite decreasing filtrations, the convolution filtration on tensor products, and the filtration on internal Hom given by

\[
F^r\mathcal Hom(\mathcal M,\mathcal N)
=\{v:v(F^a\mathcal M)\subset F^{a+r}\mathcal N
\text{ for every }a\}.
\]

Write \(\cHom_F\), \(\otimes_F\) and \(\mathbb D_F=\cHom_F(-,\omega_X)\) for the derived functors on the filtered derived category, where \(\omega_X\) carries the filtration concentrated in weight zero, and \(M_a\), \(N_a\), \(K_a\) for the graded pieces of \(\mathcal M\), \(\mathcal N\), \(\mathcal K\).

% object: lemma-graded-pieces-filtered-operations (blueprint/chap-The-Lefschetz-Verdier-formula/sec-Additivity-of-trace-classes/lemma-graded-pieces-filtered-operations/_section.md)
\phantomsection\label{bp211}%
\begin{bpstatementit}{Lemma~11.15}{ (Graded pieces of filtered operations)}
For finitely filtered complexes \(\mathcal M,\mathcal N,\mathcal K\) on \(X\) with graded pieces in \(D^b_{\mathrm{ctf}}(X,A)\) there are isomorphisms

\[
\begin{aligned}
\operatorname{gr}^r\cHom_F(\mathcal M,\mathcal N)
&\simeq\bigoplus_a\cHom(M_a,N_{a+r}),\\
\operatorname{gr}^r\mathbb D_F\mathcal K&\simeq \mathbb D_XK_{-r},\\
\operatorname{gr}^r(\mathcal M\otimes_F\mathcal N)
&\simeq\bigoplus_{a+b=r}M_a\otimes^{\mathbf L} N_b.
\end{aligned}
\]

Under these isomorphisms, the restriction of the filtered evaluation \(\mathcal K\otimes_F\mathbb D_F\mathcal K\to\omega_X\) to the summand \(K_a\otimes^{\mathbf L}\mathbb D_XK_a\) of \(\operatorname{gr}^0\) is the evaluation of \(K_a\).
\end{bpstatementit}

% proof: proof/lemma-graded-pieces-filtered-operations
\begin{proof}[{Proof}]\phantomsection\label{bp212}%
The operations are derived using filtered flat and filtered injective resolutions. For finite filtrations these resolutions can be built successively on the graded pieces by mapping cylinders, and on them the three formulas hold termwise. For the evaluation, recall that the filtered evaluation is the tensor–Hom adjoint of \(\operatorname{id}_{\mathcal K}\), whose graded map in weight zero is \(\operatorname{id}_{K_a}\) on each summand. Hence on \(K_a\otimes^{\mathbf L}\mathbb D_XK_a\) it is the ordinary evaluation. Moreover, the Koszul signs are those of the underlying complexes, and the filtration weights contribute no signs.
\end{proof}

% text: prose-11-17
\phantomsection\label{bp213}%
With these formulas we lift the trace morphism of Definition \hyperref[bp194]{11.8} to filtered complexes.

% object: lemma-filtered-trace-morphism (blueprint/chap-The-Lefschetz-Verdier-formula/sec-Additivity-of-trace-classes/lemma-filtered-trace-morphism/_section.md)
\phantomsection\label{bp214}%
\begin{bpstatementit}{Lemma~11.16}{ (Filtered trace morphism)}
Let \(c:C\to X\times X\) be a correspondence and let \(\mathcal K\) be a finitely filtered complex on \(X\) with underlying complex \(K\) and graded pieces \(K_a\) in \(D^b_{\mathrm{ctf}}(X,A)\). There is a morphism of filtered complexes

\[
\operatorname{tr}_F:
\cHom_F(c_2^*\mathcal K,c_1^!\mathcal K)\longrightarrow
Ri_{c*}\omega_{\Fix(c)},
\]

with the target in weight zero, whose underlying morphism is the trace morphism for \(K\), and whose restriction to \(F^0\) factors through

\[
\operatorname{gr}^0\cHom_F(c_2^*\mathcal K,c_1^!\mathcal K)
\simeq\bigoplus_a\cHom(c_2^*K_a,c_1^!K_a)
\]

by a map whose restriction to the summand \(a\) is the trace morphism for \(K_a\).
\end{bpstatementit}

% proof: proof/lemma-filtered-trace-morphism
\begin{proof}[{Proof}]\phantomsection\label{bp215}%
The trace morphism is the composite of four steps: the exceptional tensor–Hom adjunction, external duality, the adjoint of evaluation along the diagonal, and exceptional proper base change for the closed diagonal. Since each is defined by units, counits and tensor evaluation, the composite has a filtered version

\[
\cHom_F(c_2^*\mathcal K,c_1^!\mathcal K)
\simeq c^!(\mathcal K\boxtimes_F\mathbb D_F\mathcal K)
\longrightarrow Ri_{c*}\omega_{\Fix(c)}
\]

whose graded pieces, computed by Lemma \hyperref[bp211]{11.15}, are the unfiltered versions. In particular the filtered comparison for external duality is an isomorphism, because each of its graded maps is a finite sum of the isomorphisms in Lemma \hyperref[bp191]{11.7}, and graded pieces detect isomorphisms of finitely filtered complexes. Forgetting the filtration we recover the trace morphism for \(K\).

The target is pure of weight zero, so \(\operatorname{tr}_F\) vanishes on \(F^1\), and its restriction to \(F^0\) factors through \(\operatorname{gr}^0\). This factorization holds on filtered dg models, hence also in the filtered derived category. By Lemma \hyperref[bp211]{11.15}, \(\operatorname{gr}^0\) of the source is \(\bigoplus_a\cHom(c_2^*K_a,c_1^!K_a)\). Now the steps for internal Hom and external duality preserve the two indices, and on a pair of equal indices evaluation along the diagonal is the evaluation of \(K_a\). Hence the map on the summand \(a\) is the trace morphism for \(K_a\).
\end{proof}

% text: prose-11-18
\phantomsection\label{bp216}%
To apply this to an open subscheme \(U\) and its closed complement \(Z\), we must lift \(u\) to the localization triangle. For this the correspondence must respect the decomposition: we need \(c_1^{-1}(Z)\subset c_2^{-1}(Z)\). For the graph of an endomorphism \(\varphi\), where \(c_1=\operatorname{id}_X\) and \(c_2=\varphi\), this says that \(\varphi(Z)\subset Z\).

It is convenient to work with exceptional adjoints. Put \(\Theta=Rc_{1!}c_2^*:D^b_{\mathrm{ctf}}(X,A)\to D^b_{\mathrm{ctf}}(X,A)\). In these terms a cohomological correspondence \(u\) on \((X,K)\) is the same as a map \(\widehat u:\Theta K\to K\), its adjoint under \(Rc_{1!}\dashv c_1^!\). Under the invariance condition there are no nonzero maps from \(\Theta K_U\) to \(K_Z\), so \(u\) restricts to \(K_U\) in an essentially unique way. Let

\[
C_U=c_2^{-1}(U),\qquad C_Z=(c_1^{-1}(Z))_{\mathrm{red}},\qquad
i':C_Z\longrightarrow C,
\]

with \(i'\) the closed immersion. By the invariance condition both legs of \(C_U\) land in \(U\) and both legs of \(C_Z\) land in \(Z\); we write \(c_U=(c_{1U},c_{2U}):C_U\to U\times U\) and \(c_Z=(c_{1Z},c_{2Z}):C_Z\to Z\times Z\) for the restricted correspondences.

% object: lemma-lift-correspondence-invariant-decomposition (blueprint/chap-The-Lefschetz-Verdier-formula/sec-Additivity-of-trace-classes/lemma-lift-correspondence-invariant-decomposition/_section.md)
\phantomsection\label{bp217}%
\begin{bpstatementit}{Lemma~11.17}{ (Lifting a correspondence to an invariant decomposition)}
Let \(c:C\to X\times X\) be a correspondence, \(i:Z\to X\) a closed immersion with open complement \(j:U\to X\), and suppose

\[
c_1^{-1}(Z)\subset c_2^{-1}(Z),
\qquad\text{equivalently }c_2^{-1}(U)\subset c_1^{-1}(U).
\]

Let \(u\) be a cohomological correspondence on \((X,K)\), and write \(K_U=j_!j^*K\) and \(K_Z=Ri_*i^*K\) for the two pieces of the localization triangle \(K_U\to K\to K_Z\). Then \(u\) has a lift, unique up to contractible choice, to a cohomological correspondence \(\tilde u\) in the filtered derived category on the two-step filtered complex given by this triangle. Its graded pieces are cohomological correspondences \(u_U\) on \((X,K_U)\) and \(u_Z\) on \((X,K_Z)\), and \(\widehat u_Z\) restricted to \(Z\) is

\[
Rc_{1Z!}c_{2Z}^*i^*K\simeq i^*\Theta K
\xrightarrow{i^*\widehat u}i^*K.
\]
\end{bpstatementit}

% proof: proof/lemma-lift-correspondence-invariant-decomposition
\begin{proof}[{Proof}]\phantomsection\label{bp218}%
By the hypothesis \(c_1^{-1}(Z)\subset c_2^{-1}(Z)\), base change of \(Rc_{1!}\) along \(i\), and the adjunction \(i^*\dashv Ri_*\), we have

\[
i^*\Theta K_U=0,\qquad
R\operatorname{Hom}_X(\Theta K_U,K_Z)
\simeq R\operatorname{Hom}_Z(i^*\Theta K_U,i^*K)=0.
\]

The second vanishing is that of the whole mapping complex, including all shifts. Hence the composite \(\Theta K_U\to\Theta K\xrightarrow{\widehat u}K\) factors through \(K_U\to K\), with a contractible space of choices. The coherent cofibre of this factorization is a morphism of two-step filtered objects

\[
\begin{array}{ccccc}
\Theta K_U&\longrightarrow&\Theta K&\longrightarrow&\Theta K_Z\\
\downarrow&&\downarrow \widehat u&&\downarrow\\
K_U&\longrightarrow&K&\longrightarrow&K_Z.
\end{array}
\]

Since this cofibre includes the homotopy, it is a morphism in the filtered derived category, and its adjoint under the filtered adjunction \(Rc_{1!}\dashv c_1^!\) is the lift \(\tilde u\). By base change, composition and nilpotent invariance (which lets us use the reduced scheme \(C_Z\)),

\[
\Theta K_U\simeq j_!Rc_{1U!}c_{2U}^*j^*K,
\qquad i^*\Theta K\simeq Rc_{1Z!}c_{2Z}^*i^*K.
\]

Applying \(i^*\) to the right-hand square and using \(i^*\Theta K_U=0\), so that \(i^*\Theta K\simeq i^*\Theta K_Z\), we obtain the formula for \(\widehat u_Z\) restricted to \(Z\).
\end{proof}

% text: prose-11-19
\phantomsection\label{bp219}%
The graded piece on the closed part is a correspondence on \(Z\) pushed forward to \(X\). This is the form in which Proposition \hyperref[bp274]{12.16} uses it, combined with Theorem \hyperref[bp206]{11.13}.

% object: lemma-closed-piece-pushforward (blueprint/chap-The-Lefschetz-Verdier-formula/sec-Additivity-of-trace-classes/lemma-closed-piece-pushforward/_section.md)
\phantomsection\label{bp220}%
\begin{bpstatementit}{Lemma~11.18}{ (The closed piece as a pushforward)}
In the situation of Lemma \hyperref[bp217]{11.17}, put \(M=i^*K\in D^b_{\mathrm{ctf}}(Z,A)\), and let \(u_{|Z}\) be the cohomological correspondence on \((Z,M)\) for \(c_Z\) whose exceptional adjoint is \(Rc_{1Z!}c_{2Z}^*M\simeq i^*\Theta K\xrightarrow{i^*\widehat u}M\). Then \(u_Z\) is the pushforward of \(u_{|Z}\) along the proper maps \(i\) and \(i'\) (Definition \hyperref[bp204]{11.12}).
\end{bpstatementit}

% proof: proof/lemma-closed-piece-pushforward
\begin{proof}[{Proof}]\phantomsection\label{bp221}%
The maps \(\alpha\) and \(\beta\) of Definition \hyperref[bp204]{11.12} are here

\[
c_2^*Ri_*M\xrightarrow{\eta_{i'}}
Ri'_*i'^*c_2^*Ri_*M
=Ri'_*c_{2Z}^*i^*Ri_*M
\xrightarrow{Ri'_*c_{2Z}^*\varepsilon_i}Ri'_*c_{2Z}^*M
\]

and the map \(Ri'_*c_{1Z}^!M\to c_1^!Ri_*M\) adjoint to

\[
Rc_{1!}Ri'_*c_{1Z}^!M
=Ri_*Rc_{1Z!}c_{1Z}^!M\xrightarrow{Ri_*\varepsilon_{c_{1Z}}}Ri_*M.
\]

The exceptional adjoint of the pushforward is thus a map \(\Theta Ri_*M\to Ri_*M\). After applying \(i^*\) and the base change identity \(i^*\Theta Ri_*M\simeq Rc_{1Z!}c_{2Z}^*M\), it becomes the exceptional adjoint of \(u_{|Z}\). By Lemma \hyperref[bp217]{11.17} the same holds for \(\widehat u_Z\). Since \(K_Z=Ri_*M\) and a map into \(Ri_*M\) is determined by its adjoint under \(i^*\dashv Ri_*\), the two correspondences agree.
\end{proof}

% text: prose-11-20
\phantomsection\label{bp222}%
Additivity takes the following form; in Proposition \hyperref[bp274]{12.16} we apply it with \(Z\) a single fixed point of Frobenius.

% object: proposition-additivity-trace-classes (blueprint/chap-The-Lefschetz-Verdier-formula/sec-Additivity-of-trace-classes/proposition-additivity-trace-classes/_section.md)
\phantomsection\label{bp223}%
\begin{bpstatementit}{Proposition~11.19}{ (Additivity of trace classes)}
Let \(c:C\to X\times X\) be a correspondence.

\begin{enumerate}
\item Let \(\mathcal K\) be a finitely filtered complex on \(X\) with underlying complex \(K\) and graded pieces \(K_a\) in \(D^b_{\mathrm{ctf}}(X,A)\), and let \(\tilde u:c_2^*\mathcal K\to c_1^!\mathcal K\) be a cohomological correspondence in the filtered derived category, with underlying correspondence \(u\) on \(K\) and graded pieces \(u_a\) on \(K_a\). Then \[
   \operatorname{Tr}_c(u)=\sum_a\operatorname{Tr}_c(u_a)
   \quad\text{in }H^0(\Fix(c),\omega_{\Fix(c)}).
   \]
\item Let \(i:Z\to X\) be a closed immersion with open complement \(U\) such that \(c_1^{-1}(Z)\subset c_2^{-1}(Z)\), let \(K\in D^b_{\mathrm{ctf}}(X,A)\), and let \(u\) be a cohomological correspondence on \((X,K)\). Then \[
   \operatorname{Tr}_c(u)=\operatorname{Tr}_c(u_U)+\operatorname{Tr}_c(u_Z),
   \] where \(u_U\) and \(u_Z\) are the graded pieces of the lift of Lemma \hyperref[bp217]{11.17}, and \(u_Z\) is the pushforward along \(i\) and \(C_Z\to C\) of the correspondence \(u_{|Z}\) on \((Z,i^*K)\) for \(c_Z:C_Z=(c_1^{-1}(Z))_{\mathrm{red}}\to
Z\times Z\) (Lemma \hyperref[bp220]{11.18}).
\end{enumerate}
\end{bpstatementit}

% proof: proof/proposition-additivity-trace-classes
\begin{proof}[{Proof}]\phantomsection\label{bp224}%
For (1), the section \(A_C\to\cHom_F(c_2^*\mathcal K,c_1^!\mathcal K)\) defined by \(\tilde u\) comes from the unit in weight zero, so it lands in \(F^0\), and its image in \(\operatorname{gr}^0\) is the tuple \((u_a)_a\). Composing with the filtered trace morphism of Lemma \hyperref[bp214]{11.16} and forgetting the filtration gives on the one hand \(\operatorname{Tr}_c(u)\), and on the other hand, by the factorization through \(\operatorname{gr}^0\), the sum of the trace classes \(\operatorname{Tr}_c(u_a)\). For (2) we apply (1) to the lift \(\tilde u\) of Lemma \hyperref[bp217]{11.17}.
\end{proof}

% text: prose-11-21
\phantomsection\label{bp225}%
Restricting the formula to a proper open and closed part \(B\) of the fixed locus and integrating, we get \(\operatorname{LT}_B(u)=\operatorname{LT}_B(u_U)+\operatorname{LT}_B(u_Z)\).



% object: chap-The-trace-formula (blueprint/chap-The-trace-formula/_section.md), content hidden
\section*{\texorpdfstring{12\quad The trace formula}{12 The trace formula}}
\phantomsection\label{bp226}%
\addcontentsline{toc}{section}{\texorpdfstring{12\quad The trace formula}{12 The trace formula}}
% text: prose-12-01
\phantomsection\label{bp227}%
Let \(X_0\) be a separated scheme of finite type over \(\mathbf F_q\), let \(X=X_0\times_{\mathbf F_q}\overline{\mathbf F}_q\), and let \(\ell\nmid q\). By Lemma \hyperref[bp240]{12.5}, the points of \(X_0\) over \(\mathbf F_{q^m}\) are the fixed points of the \(m\)-th power of the Frobenius \(\Phi_X\) on \(X\). The Lefschetz–Verdier formula (Theorem \hyperref[bp206]{11.13}) writes a trace on cohomology as a sum of local terms, but it says nothing about what a local term is. To count points we need the local term of Frobenius at a fixed point to be the trace on the stalk.

For a self-map \(\varphi\) of a compact manifold with transverse fixed points, the Lefschetz formula counts each fixed point \(x\) with the sign of \(\det(1-d\varphi_x)\). The differential of Frobenius is zero, so we expect every sign to be \(+1\). On a curve this can be seen directly. If \(y\) is a fixed point of \(\Phi_X^m\), \(u\) is a uniformizer at \(y\) defined over \(\mathbf F_{q^m}\), and \(v\) is the same function on the second factor, then near \((y,y)\) the graph of \(\Phi_X^m\) is \(v=u^{q^m}\), and it meets the diagonal \(v=u\) with multiplicity one, since \(u^{q^m}-u\) has a simple zero at \(u=0\).

% object: question-lefschetz-formula-frobenius (blueprint/chap-The-trace-formula/question-lefschetz-formula-frobenius/_section.md)
\phantomsection\label{bp228}%
\begin{bpstatement}{Question~12.1}{ (The Lefschetz formula for Frobenius)}
Does the Lefschetz formula hold for Frobenius?
\end{bpstatement}

% text: prose-12-02
\phantomsection\label{bp229}%
For constant coefficients the local term at a transverse fixed point is \(1\), so on a proper \(X\) transversality would be enough. Two things stand in the way in general. First, \(X\) need not be proper. If we compactify, the new fixed points at the boundary have zero stalk for \(j_!K\), and we need their local terms to vanish. Second, a local term is computed from \(\cHom(c_2^*K,c_1^!K)\) near the fixed point, where \(c_1,c_2\) are the two projections of the graph of Frobenius to \(X\) (Definition \hyperref[bp184]{11.5}). The stalk \(K_{\bar y}\) does not determine this complex (Remark \hyperref[bp239]{12.4}), so the vanishing of the local term at a point where the stalk vanishes is a real statement, which transversality alone does not give.

One computation on a curve settles both problems: Proposition \hyperref[bp274]{12.16} says that for any morphism \(\psi:(\Phi_X^m)^*K\to K\) on a smooth curve, the local term at a rational point \(y\) is \(\operatorname{tr}(\psi_{\bar y}\mid K_{\bar y})\). The proof, which is Grothendieck's [\hyperref[bp332]{SGA5}, XV, §3], uses \(d\Phi_X=0\) a second time: the graph \(v=u^{q^m}\) is tangent to the axis \(v=0\) and the diagonal is not, so blowing up \((y,y)\) separates them, and near the exceptional divisor the local term becomes a restriction map that vanishes when \(K_{\bar y}=0\) (Remark \hyperref[bp243]{12.6}). Since both sides of the formula are additive along an open subscheme and its complement (Lemma \hyperref[bp283]{12.19}) and compatible with \(Rf_!\) (Lemma \hyperref[bp280]{12.18}), we never need to compactify anything but \(\mathbf A^1\subset\mathbf P^1\).

The answer is Theorem \hyperref[bp298]{12.24}: for every Weil complex \(K\in D^b_{\mathrm{ctf}}(X,\Lambda)\) (Section \hyperref[bp154]{10.5}), where \(\Lambda\in\{A,\mathbf Z_\ell,\mathcal O_E,\mathbf Q_\ell,E,\overline{\mathbf Q}_\ell\}\) and \(A\) is a finite commutative \(\ell\)-primary ring, and for every \(m\geq1\), we have in \(\Lambda\)

\[
\operatorname{tr}_\Lambda\bigl(\operatorname{Frob}_{X,K}^m\mid
R\Gamma_c(X_{\mathrm{pro}\acute{e}t},K)\bigr)
=\sum_{y\in X_0(\mathbf F_{q^m})}
\operatorname{tr}_\Lambda\bigl(\operatorname{Frob}_{y,m}\mid
K_{\bar y}\bigr).
\]

Here \(\operatorname{tr}_\Lambda\) is the trace of an endomorphism of a perfect complex (Definition \hyperref[bp174]{11.1}), and \(\operatorname{Frob}_{y,m}\) is the \(m\)-th iterate of the Weil structure on the stalk at the fixed point \(\bar y\). For \(K=\Lambda_X\) and \(\Lambda\) a field the formula is Weil's axiom (b), and it gives the rationality of the zeta function (Corollary \hyperref[bp320]{12.31}).

Purity, however, needs the formula for general \(K\). For \(\Lambda\in\{\mathbf Q_\ell,E,\overline{\mathbf Q}_\ell\}\), in the form of an Euler product (Theorem \hyperref[bp317]{12.30}), it reads

\[
\prod_{x\in|X_0|}\det\bigl(1-T^{\deg x}\operatorname{Frob}_x\mid K_{\bar x}\bigr)^{-1}
=\prod_i\det\bigl(1-T\operatorname{Frob}_{X,K}\mid H^i_c(X,K)\bigr)^{(-1)^{i+1}},
\]

which relates the local eigenvalues of \(K\) to the eigenvalues on cohomology.

The conventions of Chapter \hyperref[bp170]{11} remain in force.

% object: sec-Frobenius-correspondences (blueprint/chap-The-trace-formula/sec-Frobenius-correspondences/_section.md), content hidden
\subsection*{\texorpdfstring{12.1\quad Frobenius correspondences}{12.1 Frobenius correspondences}}
\phantomsection\label{bp230}%
\addcontentsline{toc}{subsection}{\texorpdfstring{12.1\quad Frobenius correspondences}{12.1 Frobenius correspondences}}
% text: prose-12-03
\phantomsection\label{bp231}%
Two Frobenius maps act on \(X=X_0\times_{k_0}k\): the \(k\)-linear scheme Frobenius \(\Phi_X\) and the semilinear geometric Frobenius \(\operatorname{Fr}_X\) that defines Weil structures. We prove the trace formula for \(\Phi_X\), whose graph is a correspondence over \(k\), and Lemma \hyperref[bp234]{12.3} translates it to \(\operatorname{Fr}_X\).

The operator \(T_{c_m}\) on cohomology with compact support defined below makes sense even when \(X\) is not proper, because the second leg \(\Phi_X^m\) of the graph \(c_m\) is proper. Indeed, \(\Phi_X\) and all its iterates are finite, since a finitely generated algebra over a perfect field is finite over its subring of \(q\)-th powers.

We cannot work with \(\operatorname{Fr}_X\) directly. Indeed, the Weil structure \(\phi_K\) is a map \(\operatorname{Fr}_X^*K\to K\), but \(\operatorname{Fr}_X\) lies over the automorphism \(\operatorname{Spec}(\sigma)\) of \(s\). Hence its graph is not a morphism \(X\to X\times_sX\) of \(k\)-schemes. Since the formalism of Chapter \hyperref[bp170]{11} uses \(c_1^!\) only for \(k\)-morphisms, it does not apply to this graph. The \(k\)-linear \(\Phi_X=F_X\circ\operatorname{Fr}_X\), on the other hand, has a graph of the required kind. Pullback along the absolute Frobenius \(F_X\) is the identity on étale sheaves, and this is how Lemma \hyperref[bp234]{12.3} turns \(\phi_K\) into a map \(\Phi_X^*K\to K\). For example, on \(X=\mathbf A^1_k\) the map \(\Phi_X\) is \(t\mapsto t^q\), and the fixed points of \(\Phi_X^m\) are the \(q^m\) elements of \(\mathbf F_{q^m}\). Lemma \hyperref[bp240]{12.5} describes them in general.

We allow maps \(\Phi_X^*K\to K\) that are not isomorphisms, because the proof over \(E\) clears denominators: it multiplies a Weil structure by a power of \(\pi_E\) (Lemma \hyperref[bp294]{12.22}). The operators below are the two sides of the trace formula in Theorem \hyperref[bp287]{12.20}.

% object: definition-frobenius-correspondences-their-operators (blueprint/chap-The-trace-formula/sec-Frobenius-correspondences/definition-frobenius-correspondences-their-operators/_section.md)
\phantomsection\label{bp232}%
\begin{bpstatement}{Definition~12.2}{ (Frobenius correspondences and their operators)}
Let \(X_0/k_0\) be separated and of finite type, and put \(X=X_0\times_{k_0}k\). Denote by \(\Phi_X:X\to X\) the base change of the \(q\)-power scheme Frobenius of \(X_0\). It is a finite \(k\)-morphism and a universal homeomorphism. A \emph{Frobenius correspondence} on \(K\) is a morphism, not necessarily invertible,

\[
\psi:\Phi_X^*K\longrightarrow K,
\qquad K\in D^b_{\mathrm{ctf}}(X,A).
\]

Its iterates are the morphisms

\[
\psi_m=\psi\circ\Phi_X^*\psi\circ\cdots\circ
(\Phi_X^{m-1})^*\psi:(\Phi_X^m)^*K\longrightarrow K.
\]

Let \(c_m=(\operatorname{id}_X,\Phi_X^m):X\to X\times_sX\) be the graph of \(\Phi_X^m\). Since its first leg is the identity, \(\psi_m\) is a cohomological correspondence on \((X,K)\) for \(c_m\). Its operator on cohomology with compact support (Definition \hyperref[bp202]{11.11}) is the endomorphism

\[
T_{c_m}(\psi_m)
=R\Gamma_c(X,\psi_m)\circ(\Phi_X^m)^*,
\]

where \((\Phi_X^m)^*:R\Gamma_c(X,K)\to R\Gamma_c(X,(\Phi_X^m)^*K)\) is the pullback map induced by the unit of \((\Phi_X^m)^*\dashv R(\Phi_X^m)_*\). For \(x\in X_0(\mathbf F_{q^m})\) let \(\bar x:s\to X\) be the associated fixed point. The stalk of \(\psi_m\) at \(\bar x\) is an endomorphism \(\psi_{m,\bar x}:K_{\bar x}\to K_{\bar x}\).
\end{bpstatement}

% text: prose-12-04
\phantomsection\label{bp233}%
The operator \(T_{c_m}(\psi_m)\) is defined with the \(k\)-linear \(\Phi_X\), while the Frobenius operator of Section \hyperref[bp154]{10.5}, whose trace the theorem is about, is defined with \(\operatorname{Fr}_X\). The next lemma says that the two agree, so that the trace formula proved for \(\Phi_X\) is a statement about \(\operatorname{Frob}_{X,K}\).

% object: lemma-scheme-frobenius-geometric-frobenius (blueprint/chap-The-trace-formula/sec-Frobenius-correspondences/lemma-scheme-frobenius-geometric-frobenius/_section.md)
\phantomsection\label{bp234}%
\begin{bpstatementit}{Lemma~12.3}{ (Scheme Frobenius and geometric Frobenius)}
Let \(\operatorname{Fr}_X\) be the semilinear geometric Frobenius of Definition \hyperref[bp156]{10.18}, and for a morphism \(\phi:\operatorname{Fr}_X^*K\to K\) put \(\phi_m=\phi\circ \operatorname{Fr}_X^*\phi\circ\cdots\circ(\operatorname{Fr}_X^{m-1})^*\phi\). There are natural isomorphisms

\[
\kappa_m:(\Phi_X^m)^*K\xrightarrow{\sim}(\operatorname{Fr}_X^m)^*K
\]

compatible with iteration. For a Weil structure \(\phi:\operatorname{Fr}_X^*K\xrightarrow{\sim}K\), put \(\psi=\phi\kappa_1\). Then

\[
\psi_m=\phi_m\kappa_m,
\qquad T_{c_m}(\psi_m)=\operatorname{Frob}_{X,K}^{m},
\]

where \(\operatorname{Frob}_{X,K}\) is the geometric Frobenius operator on compact cohomology of Definition \hyperref[bp158]{10.19}. The same identity of operators holds for an arbitrary morphism \((\operatorname{Fr}_X^m)^*K\to K\) in place of \(\phi_m\).
\end{bpstatementit}

% proof: proof/lemma-scheme-frobenius-geometric-frobenius
\begin{proof}[{Proof}]\phantomsection\label{bp235}%
The maps \(\Phi_X\) and \(\operatorname{Fr}_X\) differ by the absolute Frobenius, and pullback along the absolute Frobenius is naturally the identity on étale sheaves. This gives \(\kappa_m\), and we then compare the two pullback maps on one compactification. On coordinate rings the two maps are

\[
\Phi_X^\#(a\otimes b)=a^q\otimes b,
\qquad \operatorname{Fr}_X^\#(a\otimes b)=a\otimes b^{1/q}
\]

so that \(\Phi_X=F_X\operatorname{Fr}_X\), where \(F_X\) is the absolute \(q\)-power Frobenius of \(X\). The absolute Frobenius commutes with every morphism of schemes over \(\mathbf F_p\), so also \(\Phi_X^m=\operatorname{Fr}_X^mF_X^m\). On the classical constructible category, pullback along absolute Frobenius is naturally isomorphic to the identity, because for an étale object its relative Frobenius is an isomorphism, functorial in that object. Transporting this natural isomorphism to \(D^b_{\mathrm{ctf}}(X,A)\) by the comparison with the pro-étale site, we obtain \(\theta:F_X^*\simeq\mathrm{id}\). It is compatible with pullback along every morphism \(h\), that is, \(h^*\theta=\theta h^*\) under \(h^*F^*=F^*h^*\), since relative Frobenius of étale objects is compatible with base change. Let \(\theta^m:(F_X^m)^*\simeq\mathrm{id}\) be its \(m\)-fold composite and put

\[
\kappa_m=\theta^m_{(\operatorname{Fr}_X^m)^*K}:
(\Phi_X^m)^*K=(F_X^m)^*(\operatorname{Fr}_X^m)^*K
\xrightarrow{\sim}(\operatorname{Fr}_X^m)^*K .
\]

Compatibility with iteration means that \(\kappa_{m+1}=\kappa_{m,\operatorname{Fr}_X^*K}\circ(\Phi_X^m)^*\kappa_1\), and it follows from the compatibility of \(\theta\) with pullback along \(\Phi_X^m\). Now \(\psi_m=\phi_m\kappa_m\) follows by induction on \(m\). If it holds for \(m\), then

\[
\begin{aligned}
\psi_{m+1}&=\psi_m\circ(\Phi_X^m)^*\psi
=\phi_m\circ\kappa_m\circ(\Phi_X^m)^*\phi\circ(\Phi_X^m)^*\kappa_1\\
&=\phi_m\circ(\operatorname{Fr}_X^m)^*\phi\circ\kappa_{m,\operatorname{Fr}_X^*K}\circ(\Phi_X^m)^*\kappa_1
=\phi_{m+1}\kappa_{m+1},
\end{aligned}
\]

where the third equality is the naturality of \(\kappa_m\) in \(K\).

For the pullback maps on compact cohomology, let \(j:X\hookrightarrow\bar X\) be the base change of a compactification of \(X_0\) over \(k_0\), and use the pullback maps \(\operatorname{pb}_h\) of Section \hyperref[bp154]{10.5}. The endomorphisms \(\Phi_{\bar X}\), \(\operatorname{Fr}_{\bar X}\) and \(F_{\bar X}\) of the scheme \(\bar X\) restrict to \(\Phi_X\), \(\operatorname{Fr}_X\) and \(F_X\). Each of them has inverse image \(X\) of \(X\), so each commutes with \(j_!\), and \(\Phi_{\bar X}^m=\operatorname{Fr}_{\bar X}^mF_{\bar X}^m\). Put \(L=j_!K\), so that \(R\Gamma_c(X,K)=R\Gamma(\bar X,L)\).

First, the map \((\Phi_X^m)^*\) of Definition \hyperref[bp232]{12.2} is \(\operatorname{pb}_{\Phi_{\bar X}^m,L}\). Indeed, \(\Phi_X^m\) is finite, so \(R(\Phi_X^m)_!=R(\Phi_X^m)_*\), and the square \(j\Phi_X^m=\Phi_{\bar X}^mj\) is cartesian, so \(j_!R(\Phi_X^m)_*\simeq R(\Phi_{\bar X}^m)_*j_!\). Under this isomorphism the identification \(R\Gamma_c(X,R(\Phi_X^m)_*-)=R\Gamma_c(X,-)\) used there is the composition of direct images \(R\Gamma(\bar X,R(\Phi_{\bar X}^m)_*-)=R\Gamma(\bar X,-)\), and the unit for \(\Phi_X^m\) corresponds to the unit for \(\Phi_{\bar X}^m\).

Second, by transitivity, \(\operatorname{pb}_{\Phi_{\bar X}^m,L}
 =\operatorname{pb}_{F_{\bar X}^m,(\operatorname{Fr}_{\bar X}^m)^*L}\circ
 \operatorname{pb}_{\operatorname{Fr}_{\bar X}^m,L}\). The absolute Frobenius \(F_{\bar X}^m\) satisfies the hypothesis of [Stacks, \href{https://stacks.math.columbia.edu/tag/03SP}{03SP}], since its relative Frobenius over each étale object is an isomorphism, functorial in that object. Since \(A\) is finite, we may compute on the étale site (Proposition \bpout{5.15}). There the formal argument of [Stacks, \href{https://stacks.math.columbia.edu/tag/03SP}{03SP}] identifies \((F_{\bar X}^m)_*\) and \((F_{\bar X}^m)^{-1}\) with the identity functors of the topos, compatibly with the unit, and this identification is \(\theta^m\), since both come from the relative Frobenius isomorphisms of étale objects. Hence for every \(N\in D^b_{\mathrm{ctf}}(\bar X,A)\) the pullback map \(\operatorname{pb}_{F_{\bar X}^m,N}\) followed by \(R\Gamma(\bar X,\theta^m_N)\) is the identity of \(R\Gamma(\bar X,N)\). With \(N=(\operatorname{Fr}_{\bar X}^m)^*L\) this gives

\[
R\Gamma_c(X,\kappa_m)\circ(\Phi_X^m)^*=\operatorname{pb}_{\operatorname{Fr}_X^m} .
\]

Finally, the naturality of \(\operatorname{pb}_{\operatorname{Fr}_X}\) moves the maps \(R\Gamma_c(X,(\operatorname{Fr}_X^i)^*\phi)\) past the pullback maps, and transitivity gives \(\operatorname{Frob}_{X,K}^m=R\Gamma_c(X,\phi_m)\circ\operatorname{pb}_{\operatorname{Fr}_X^m}\). Hence

\[
T_{c_m}(\psi_m)=R\Gamma_c(X,\phi_m)\circ R\Gamma_c(X,\kappa_m)\circ(\Phi_X^m)^*
=R\Gamma_c(X,\phi_m)\circ\operatorname{pb}_{\operatorname{Fr}_X^m}
=\operatorname{Frob}_{X,K}^m .
\]

The same computation with an arbitrary morphism \(\chi:(\operatorname{Fr}_X^m)^*K\to K\) in place of \(\phi_m\) gives \(T_{c_m}(\chi\kappa_m)=R\Gamma_c(X,\chi)\circ\operatorname{pb}_{\operatorname{Fr}_X^m}\).
\end{proof}

% text: prose-12-05
\phantomsection\label{bp236}%
For example, if \(X_0\) is smooth and geometrically connected of dimension \(d\) and \(K=A_X\) with its constant Weil structure, then \(T_{c_1}(\psi)\) acts on \(H^{2d}_c(X,A)=A(-d)\) by \(q^d\), which agrees with the action of geometric Frobenius on \(A(-d)\) recorded after Definition \hyperref[bp158]{10.19}.


% object: sec-Local-terms-on-a-curve (blueprint/chap-The-trace-formula/sec-Local-terms-on-a-curve/_section.md), content hidden
\subsection*{\texorpdfstring{12.2\quad Local terms on a curve}{12.2 Local terms on a curve}}
\phantomsection\label{bp237}%
\addcontentsline{toc}{subsection}{\texorpdfstring{12.2\quad Local terms on a curve}{12.2 Local terms on a curve}}
% text: prose-12-06
\phantomsection\label{bp238}%
The local terms of the trace formula live on the fixed point scheme \(\Fix(\Phi_X^m)=\Fix(c_m)\) of the graph \(c_m\). More precisely, a local term is defined on each open and closed part of this scheme that is proper over \(k\). The next lemma shows that the fixed scheme is a finite reduced set of points, one for each \(\mathbf F_{q^m}\)-rational point of \(X_0\), so every such point has its own local term. For Frobenius on a smooth curve we then show that this local term is the trace on the stalk.

On \(\mathbf A^1\) the fixed scheme of \(\Phi^m\) is \(\operatorname{Spec}k[t]/(t^{q^m}-t)\). Its \(q^m\) points are reduced because \(t^{q^m}-t\) has derivative \(-1\). This is the algebraic form of the transversality seen at the start of Chapter \hyperref[bp226]{12}. Reducedness alone does not determine the local term, because the stalk does not determine \(\cHom(c_2^*K,c_1^!K)\) near the fixed point (Remark \hyperref[bp239]{12.4} below).

% object: remark-locality-warning (blueprint/chap-The-trace-formula/sec-Local-terms-on-a-curve/remark-locality-warning/_section.md)
\phantomsection\label{bp239}%
\begin{bpstatement}{Remark~12.4}{ (Local terms are not computed stalk by stalk)}
The vanishing of the local term at a rational point \(y\) with \(K_{\bar y}=0\), proved below, is a geometric fact about the restriction from cohomology with supports in the graph of \(\Phi_X^m\) to cohomology with supports in \(y\), and it does not follow from a locality property of internal Hom. For example, if \(j:\mathbf G_m\hookrightarrow\mathbf A^1\) and \(K=j_!A\), then the stalk of \(K\) at the origin is zero, but \(\cHom_A(j_!A,j_!A)\simeq Rj_*A\) has stalk \(A\) in degree zero there. Hence a local term cannot be computed by commuting a geometric stalk functor past this internal Hom.
\end{bpstatement}

% object: lemma-frobenius-fixed-scheme (blueprint/chap-The-trace-formula/sec-Local-terms-on-a-curve/lemma-frobenius-fixed-scheme/_section.md)
\phantomsection\label{bp240}%
\begin{bpstatementit}{Lemma~12.5}{ (The Frobenius fixed scheme)}
For \(m\geq1\), the fixed point scheme of \(\Phi_X^m\) satisfies

\[
\Fix(\Phi_X^m)
\simeq\coprod_{x\in X_0(\mathbf F_{q^m})}s.
\]

In particular, the fixed scheme is finite, reduced and proper over \(k\).
\end{bpstatementit}

% proof: proof/lemma-frobenius-fixed-scheme
\begin{proof}[{Proof}]\phantomsection\label{bp241}%
We reduce to a closed subscheme of affine space, where the equations \(t^{q^m}=t\) have simple roots.

For every open \(U_0\subset X_0\) the restriction to \(U=U_0\times_{k_0}k\) of the fixed scheme of \(\Phi_X^m\) is the fixed scheme of \(\Phi_U^m\), because \(\Phi_X\) is the base change of a morphism of \(X_0\) and so maps \(U\) into itself. As \(X_0\) is quasi-compact, we may therefore assume that \(X_0\) is a closed subscheme of \(\mathbf A^r_{\mathbf F_q}\), with coordinates \(t_1,\ldots,t_r\).

In this case the fixed scheme is the pullback of the diagonal along the graph. Since the equations of \(X_0\) have coefficients in \(\mathbf F_q\), we have \((\Phi_X^m)^\#t_i=t_i^{q^m}\), and the fixed scheme is cut out by the defining equations of \(X_0\) together with \(t_i^{q^m}-t_i=0\). Each polynomial \(t^{q^m}-t\) is separable, so the algebra \(k[t_1,\ldots,t_r]/(t_i^{q^m}-t_i)_i\) is a product of copies of \(k\) indexed by \(\mathbf F_{q^m}^r\). Its quotient by the defining equations of \(X_0\) retains exactly the factors at the rational solutions of those equations, as asserted. In particular the fixed scheme is a closed subscheme of a finite étale \(k\)-scheme, hence étale. The argument uses only the equations of \(X_0\), so \(X_0\) need not be reduced.
\end{proof}

% text: prose-12-07
\phantomsection\label{bp242}%
Thus every fixed point \(y\) is an open and closed point of the fixed scheme, and its local term \(\operatorname{LT}_y\) is defined.

% object: remark-outline-local-computation (blueprint/chap-The-trace-formula/sec-Local-terms-on-a-curve/remark-outline-local-computation/_section.md)
\phantomsection\label{bp243}%
\begin{bpstatement}{Remark~12.6}{ (Outline of the local computation)}
Let \(X_0\) be a smooth curve and \(y\in X_0(\mathbf F_{q^m})\). The trace formula on \(\mathbf P^1\) needs Proposition \hyperref[bp274]{12.16} below, which says that the local term at \(y\) of any morphism \(\psi:(\Phi_X^m)^*K\to K\) is the naive trace \(\operatorname{tr}_A(\psi_{\bar y}\mid K_{\bar y})\), also when \(K\) is not lisse at \(y\). To prove it, we filter \(K\) by the open complement of \(y\) and the point \(y\). By additivity of traces this reduces the proposition to two cases: the point itself, where Lemma \hyperref[bp199]{11.10} applies, and the case \(K_{\bar y}=0\), where we must show that the local term vanishes.

The second case needs a geometric argument. Indeed, the local term is defined on a neighbourhood of \((y,y)\) in \(X\times X\), where the Frobenius graph meets the diagonal, and the vanishing of the stalk at \(y\) does not by itself control it. For this we follow [\hyperref[bp332]{SGA5}, III B, §§2–4]. In this approach the local term is a restriction map from the graph to the diagonal, and blowing up \((y,y)\) separates the two curves. After the blowup the local term becomes a restriction to punctured curves near the point where two transverse divisors meet. There it vanishes when the monodromy is an \(\ell\)-group. In general we reduce to that case by passing to finite covers of degree prime to \(\ell\) (Appendix \bpout{D}); for such a cover the trace map splits the pullback.

What vanishes is a morphism in \(D(A)\), before any trace is taken. For this reason the vanishing holds for every \(\psi\), for every finite \(\ell\)-primary ring \(A\), and for every \(K\) with \(K_{\bar y}=0\). In particular it holds for \(j_!K\) on \(\mathbf P^1\), which is not lisse at infinity. We forget the \(A\)-action only in Lemma \hyperref[bp253]{12.9}, to detect a zero object.
\end{bpstatement}

% object: subsec-The-local-term-as-a-restriction-map (blueprint/chap-The-trace-formula/sec-Local-terms-on-a-curve/subsec-The-local-term-as-a-restriction-map/_section.md), content hidden
\subsubsection*{\texorpdfstring{12.2.1\quad The local term as a restriction map}{12.2.1 The local term as a restriction map}}
\phantomsection\label{bp244}%
\addcontentsline{toc}{subsubsection}{\texorpdfstring{12.2.1\quad The local term as a restriction map}{12.2.1 The local term as a restriction map}}
% text: prose-12-08
\phantomsection\label{bp245}%
The vanishing of a local term when \(K_{\bar y}=0\) is hard to see from the definition through the trace morphism (Definition \hyperref[bp194]{11.8}), so we first replace that definition by a description familiar from topology. There the intersection number of two submanifolds \(\Gamma\) and \(\Delta\) is computed by restricting the Thom class of \(\Gamma\) to \(\Delta\). Since the Thom class is supported on \(\Gamma\), the result is a class supported on \(\Gamma\cap\Delta\). Similarly, in our setting \(\psi\) defines a class supported on the graph. A restriction map \(\rho\) takes it to cohomology of the diagonal supported at \(y\), and the local term is the evaluation of the result. Note that the map \(\rho\) does not involve \(\psi\). It therefore suffices to show that \(\rho\) itself is zero when \(K_{\bar y}=0\) (Lemma \hyperref[bp270]{12.15}), and then the vanishing holds for every \(\psi\), invertible or not.

% object: lemma-local-term-restriction-map (blueprint/chap-The-trace-formula/sec-Local-terms-on-a-curve/subsec-The-local-term-as-a-restriction-map/lemma-local-term-restriction-map/_section.md)
\phantomsection\label{bp246}%
\begin{bpstatementit}{Lemma~12.7}{ (The local term as a restriction map)}
Let \(X_0\) be a smooth curve over \(k_0\), let \(X=X_0\times_{k_0}k\), let \(K\in D^b_{\mathrm{ctf}}(X,A)\), and let \(y\) be a fixed point of \(\Phi_X^m\). Let \(\Gamma\subset X\times X\) be the transposed graph of \(\Phi_X^m\), the image of the closed immersion \(\gamma=(\Phi_X^m,\operatorname{id}_X)\), and let \(\Delta\subset X\times X\) be the diagonal, the image of \(\delta=(\operatorname{id}_X,\operatorname{id}_X)\). On the strict localization \(T\) of \(X\times X\) at \((y,y)\) put

\[
P=(\mathbb D_XK\boxtimes K)|_T.
\]

A morphism \(\psi:(\Phi_X^m)^*K\to K\) determines a class \([\psi]\in H^0_\Gamma(T,P)\), and the local term \(\operatorname{LT}_y(\psi)\) of Definition \hyperref[bp197]{11.9} is the image of this class under the composite

\[
R\Gamma_\Gamma(T,P)
\xrightarrow{\rho}
R\Gamma_y(\Delta,\delta^*P)
\xrightarrow{\mathrm{ev}_K}
R\Gamma_y(X,\omega_X)
\simeq R\Gamma(y,\omega_y)
\xrightarrow{\int_y} A,
\]

where \(\rho\) is restriction from \(\Gamma\)-supported to \(y\)-supported cohomology.
\end{bpstatementit}

% text: prose-12-09
\phantomsection\label{bp247}%
In this lemma the two factors of the product are interchanged compared with the convention \(c_m=(\operatorname{id}_X,\Phi_X^m)\) of Definition \hyperref[bp232]{12.2}.

% proof: proof/lemma-local-term-restriction-map
\begin{proof}[{Proof of Lemma 12.7}]\phantomsection\label{bp248}%
We show that the local term and the displayed composite are both adjoint to one map \(R\gamma_*\gamma^!P\to R\delta_*\omega_X\). Once the class \([\psi]\) and the map \(\rho\) are built, the comparison is formal, and computing on étale models changes nothing.

By external duality, \(\gamma^!P\) is identified with \(\cHom((\Phi_X^m)^*K,K)\), so exceptional adjunction identifies \(\psi\) with \([\psi]\). The map \(\rho\) is constructed from the two canonical arrows

\[
R\gamma_*\gamma^!P\longrightarrow P\longrightarrow R\delta_*\delta^*P.
\]

After restriction to \(\Delta\), the source is supported on \(\Gamma\cap\Delta\). By Lemma \hyperref[bp240]{12.5} the scheme \(\Gamma\cap\Delta\) is reduced, and on \(T\) it is the single point \((y,y)\). Hence the resulting map factors through \(R\Gamma_y(\Delta,\delta^*P)\). Equivalently, it is the ordinary pullback of the object with supports, followed by the base change map to local cohomology with support in the intersection.

The adjoint of \(\mathrm{ev}_K:\delta^*P\to\omega_X\) is the map \(P\to R\delta_*\omega_X\) used to define the trace morphism in Definition \hyperref[bp194]{11.8}. Hence the displayed composite, before \(\int_y\), is adjoint to the map

\[
R\gamma_*\gamma^!P\longrightarrow P\longrightarrow R\delta_*\omega_X.
\]

The local term is built from the same map. Indeed, by exceptional proper base change for \(\delta\), the adjoint of this map, which is supported on \(y\), is the trace map \(\gamma^!P\to Ri_*\omega_y\) of Definition \hyperref[bp197]{11.9}. The triangle identities of the ordinary and exceptional adjunctions then show that the two constructions give the same element of \(A\).

The constructions above also make sense on the constructible étale models, where restriction to the strict localization is the filtered limit over étale neighbourhoods. The comparison theorems identify inverse images, exceptional adjunctions, duality and the support maps on the two sites, and hence they identify the class obtained on the étale models with the pro-étale local term of Definition \hyperref[bp197]{11.9}.
\end{proof}


% object: subsec-Vanishing-on-a-strictly-local-surface (blueprint/chap-The-trace-formula/sec-Local-terms-on-a-curve/subsec-Vanishing-on-a-strictly-local-surface/_section.md), content hidden
\subsubsection*{\texorpdfstring{12.2.2\quad Vanishing on a strictly local surface}{12.2.2 Vanishing on a strictly local surface}}
\phantomsection\label{bp249}%
\addcontentsline{toc}{subsubsection}{\texorpdfstring{12.2.2\quad Vanishing on a strictly local surface}{12.2.2 Vanishing on a strictly local surface}}
% text: prose-12-10
\phantomsection\label{bp250}%
The vanishing of the local term is explained by the position of the curves at \((y,y)\). Since \(d\Phi_X=0\), the transposed graph \(\Gamma\) is tangent to the vertical axis \(V=\{y\}\times X\). The diagonal \(\Delta\) and the horizontal axis \(H\) have other tangent directions. This persists on suitable covers of degree prime to \(\ell\) (Lemma \hyperref[bp265]{12.13} below). Blowing up the closed point separates the curves: the strict transforms of \(\Gamma\cup V\), of \(\Delta\) and of \(H\) meet the exceptional divisor \(E\simeq\mathbf P^1\) in different points.

Let \(M\) be the direct image of the coefficients from the complement of \(E\). We consider \(M\) on \(E\) minus the points of \(\Gamma\cup V\), and let \(z\) be the point where \(H\) meets \(E\). There \(M\) is locally constant away from \(z\), and its stalk at \(z\) is zero. Both facts use that the monodromy is an \(\ell\)-group (Lemma \hyperref[bp253]{12.9}). Now a branch of \(\Delta\) meets \(E\) in one point \(w\), and the restriction to this branch factors through the stalk \(M_w\). Composing the restriction from \(E\) to \(M_w\) with the specialization isomorphism from \(M_w\) to the generic stalk of \(E\), we get the restriction to the zero stalk \(M_z\) followed by specialization. Hence the restriction to \(M_w\) is zero.

The transverse case is Lemma \hyperref[bp253]{12.9}, and Lemma \hyperref[bp259]{12.11} reduces to it by blowing up the closed point. We compute on classical étale models of strictly local schemes, which are filtered limits of schemes of finite type, and in the proof of Lemma \hyperref[bp270]{12.15} we transport the result to the pro-étale site.

% object: lemma-punctured-surface-punctured-curve (blueprint/chap-The-trace-formula/sec-Local-terms-on-a-curve/subsec-Vanishing-on-a-strictly-local-surface/lemma-punctured-surface-punctured-curve/_section.md)
\phantomsection\label{bp251}%
\begin{bpstatementit}{Lemma~12.8}{ (A punctured surface and a punctured curve)}
Let \(S\) be the strict localization of a smooth surface over an algebraically closed field \(k\) of characteristic \(p\) at a closed point \(z\), let \(\ell\neq p\), and let \(D_1,D_2\subset S\) be regular divisors meeting transversely at \(z\). Then the restriction map

\[
R\Gamma(S-D_2,\mathbf F_\ell)\longrightarrow R\Gamma(D_1-z,\mathbf F_\ell)
\]

is an isomorphism.
\end{bpstatementit}

% proof: proof/lemma-punctured-surface-punctured-curve
\begin{proof}[{Proof}]\phantomsection\label{bp252}%
Both complexes have cohomology \(\mathbf F_\ell\) in degree zero, \(\mathbf F_\ell(-1)\) in degree one, and zero in all other degrees. For the surface this follows from localization, the acyclicity of the strictly local schemes \(S,D_2\), and the purity identity \(i_{D_2}^!\mathbf F_\ell=\mathbf F_\ell(-1)[-2]\). This identity, in turn, comes by transitivity from the dualizing complexes of the smooth schemes \(S\) and \(D_2\).

For the punctured trait \(D_1-z\), vanishing in degrees \(\geq2\) follows from Proposition \bpout{D.4}, by passing to the filtered limit of punctured neighbourhoods that are curves. In degree one the Kummer sequence gives \(H^1(-,\mu_\ell)=k'^*/k'^{*\ell}\), where \(k'\) is the fraction field. Moreover, the units of the strictly henselian trait are \(\ell\)-divisible. Hence the valuation induces \(k'^*/k'^{*\ell}\simeq\mathbf Z/\ell\), and the class of a uniformizer is a generator.

To see that the restriction map itself, and not only the two sides, is an isomorphism, we find a generator of \(H^1\) on each side. The restriction is clearly the identity in degree zero. In degree one, let \(g\) be a local equation of \(D_2\). Since the regular local ring of \(S\) is factorial, \(\operatorname{Pic}(S-D_2)=0\) and the units on \(S-D_2\) are the units of \(S\) times the powers of \(g\). As the units of \(S\) are \(\ell\)-divisible, the Kummer sequence shows that \(H^1(S-D_2,\mu_\ell)\simeq\mathbf Z/\ell\), generated by the class of \(g\). Since \(D_1\) and \(D_2\) meet transversely, \(g\) restricts to a uniformizer of \(D_1\), whose class generates \(H^1(D_1-z,\mu_\ell)\). So the restriction is an isomorphism in degree one as well, and since both complexes vanish in the other degrees, it is an isomorphism.
\end{proof}

% object: lemma-cohomology-near-two-transverse-divisors (blueprint/chap-The-trace-formula/sec-Local-terms-on-a-curve/subsec-Vanishing-on-a-strictly-local-surface/lemma-cohomology-near-two-transverse-divisors/_section.md)
\phantomsection\label{bp253}%
\begin{bpstatementit}{Lemma~12.9}{ (Cohomology near two transverse divisors)}
Let \(X\) be a smooth surface over an algebraically closed field of characteristic \(p\), or an étale limit with affine transition maps of such surfaces. Let \(D_1,D_2\) be regular divisors meeting transversely at exactly one point \(z\), and let

\[
j:X-D_2\hookrightarrow X,\qquad
j':X-(D_1\cup D_2)\hookrightarrow X-D_2.
\]

Let \(\ell\ne p\), let \(A\) be a finite commutative \(\ell\)-primary ring, and let \(N\) be a bounded complex of \(A\)-modules on \(X-(D_1\cup D_2)\) with finite locally constant cohomology whose monodromy factors through a finite \(\ell\)-group. Put \(M=Rj_*j'_!N\). Then \(M_z=0\), and \(M|_{D_2-z}\) has finite locally constant cohomology. If \(D_2\) is connected, then for every closed geometric point \(w\in D_2-z\) the restriction map

\[
R\Gamma(D_2,M|_{D_2})\longrightarrow M_w
\]

is zero in \(D(A)\).
\end{bpstatementit}

% proof: proof/lemma-cohomology-near-two-transverse-divisors
\begin{proof}[{Proof}]\phantomsection\label{bp254}%
The idea is that, since the monodromy is an \(\ell\)-group, we may filter the coefficients until they become the constant sheaf \(\mathbf F_\ell\). For constant coefficients, \(M_z=0\) is a statement on the strict localization of \(X\) at \(z\): restriction from the complement of \(D_2\) to the punctured curve \(D_1-z\) is an isomorphism on cohomology. This holds because a local equation of \(D_2\) gives the generator of \(H^1\) on both sides. Once we know that \(M\) is locally constant along \(D_2-z\), the vanishing of the restriction map follows from \(M_z=0\) by specialization along \(D_2\).

For the assertion at \(z\), pass to the strict localization \(S\) of \(X\) at \(z\) and write \(D_{1,S},D_{2,S}\) for the two transverse divisors. By the formula for stalks of a direct image along an open immersion [Stacks, \href{https://stacks.math.columbia.edu/tag/03Q7}{03Q7}] and localization we have

\[
M_z=R\Gamma(S-D_{2,S},j'_!N).
\]

To prove that this object is zero we may forget the \(A\)-action, since the exact forgetful functor is conservative and the cohomology of sheaves of modules agrees with the cohomology of the underlying abelian sheaves [Stacks, \href{https://stacks.math.columbia.edu/tag/03FD}{03FD}]. A finite abelian \(\ell\)-group with an action of a finite \(\ell\)-group has a finite filtration whose quotients are \(\mathbf F_\ell\) with trivial action. To see this, filter by powers of \(\ell\) and use the nilpotence of the augmentation ideal of the group algebra over \(\mathbf F_\ell\). Since \(j'_!\) and \(R\Gamma(S-D_{2,S},-)\) are exact, the class of \(N\) for which \(R\Gamma(S-D_{2,S},j'_!N)=0\) is closed under extensions. Using truncations and these composition series we reduce the assertion to \(N=\mathbf F_\ell\) on \(S-(D_{1,S}\cup D_{2,S})\).

By Lemma \hyperref[bp251]{12.8}, applied to \(S\) with the divisors \(D_{1,S},D_{2,S}\), the restriction map \(R\Gamma(S-D_{2,S},\mathbf F_\ell)\to R\Gamma(D_{1,S}-z,\mathbf F_\ell)\) is an isomorphism. The localization sequence for \(D_{1,S}-z\) and its complement in \(S-D_{2,S}\) identifies \(R\Gamma(S-D_{2,S},j'_!\mathbf F_\ell)\) with the cone of this map, shifted by \(-1\). Hence this complex is zero, and \(M_z=0\) in \(D(A)\) by conservativity.

We turn to the local constancy of \(M\) along \(D_2-z\), where only the divisor \(D_2\) remains. Since \(D_1\) meets \(D_2\) only at \(z\), we have \(j'_!N=N\) near \(D_2-z\), so \(M=Rj_*N\) there. To compute \(Rj_*N\) we use Abhyankar's lemma along the regular divisor \(D_2\) [Stacks, \href{https://stacks.math.columbia.edu/tag/0EYG}{0EYG}]: étale locally on \(D_2\), a finite tame local system is trivialized by a Kummer cover of a local parameter of \(D_2\) and an unramified cover. The lemma applies here because our finite \(\ell\)-group monodromy is tame, wild inertia being pro-\(p\). Consequently, the stalks of \(Rj_*N\) are the continuous cohomology of the same representation of tame inertia, and their specialization maps along \(D_2-z\) are isomorphisms. Explicitly, the pro-\(\ell\) part of tame inertia is \(\mathbf Z_\ell(1)\). Its cohomology with coefficients in a finite \(\ell\)-primary module is computed by the two-term complex given by a topological generator minus one. On the other hand, the finite quotients of the remaining factors of tame inertia have order prime to \(\ell\), so taking invariants under them is exact on these coefficients. Finally, we may choose the parameter and the unramified cover on an étale neighbourhood. Hence this description, including the specialization maps, is locally constant. Applying the argument to the cohomology sheaves of \(N\) and using the bounded spectral sequence proves the assertion for \(N\).

The last assertion combines the two facts just proved. A global section of \(M|_{D_2}\) specializes to every stalk. Its image at \(z\) is zero, and along \(D_2-z\) specialization is an isomorphism on stalks. To make this precise, choose a geometric generic point \(\bar\eta\) of the connected regular curve \(D_2\), together with specialization maps from \(\bar\eta\) to \(z\) and to \(w\). By naturality of restriction we get a commutative square

\[
\begin{CD}
R\Gamma(D_2,M|_{D_2}) @>>> M_w\\
@VVV @VV{\sim}V\\
M_z=0 @>>> M_{\bar\eta}.
\end{CD}
\]

The right arrow is an isomorphism by the local constancy just proved, so the upper arrow is zero.
\end{proof}

% text: prose-12-11
\phantomsection\label{bp255}%
Since the restriction to a curve meeting \(D_2\) in one point factors through the stalk of \(M\) at that point, the lemma implies that restrictions to punctured curves vanish.

% object: lemma-restriction-punctured-curve (blueprint/chap-The-trace-formula/sec-Local-terms-on-a-curve/subsec-Vanishing-on-a-strictly-local-surface/lemma-restriction-punctured-curve/_section.md)
\phantomsection\label{bp256}%
\begin{bpstatementit}{Lemma~12.10}{ (Restriction to a punctured curve)}
In the situation of Lemma \hyperref[bp253]{12.9} with \(D_2\) connected, let \(C\subset X\) be a closed curve with \(C\cap D_2=\{w\}\), and let \(C_{(w)}\) be its strict localization at \(w\). Then the restriction map

\[
R\Gamma(X-D_2,j'_!N)\longrightarrow
R\Gamma(C_{(w)}-w,(j'_!N)|_{C_{(w)}-w})
\]

is zero in \(D(A)\).
\end{bpstatementit}

% proof: proof/lemma-restriction-punctured-curve
\begin{proof}[{Proof}]\phantomsection\label{bp257}%
The source is \(R\Gamma(X-D_2,j'_!N)=R\Gamma(X,M)\), by the definition \(M=Rj_*j'_!N\). Since \(C\) meets \(D_2\) only at \(w\), the punctured curve \(C_{(w)}-w\) lies over the strict localization of \(X\) at \(w\) and avoids \(D_2\). The map therefore factors as the restriction to \(R\Gamma(D_2,M|_{D_2})\), followed by the restriction to the stalk \(M_w=R\Gamma(X_{(w)}-D_2,j'_!N)\) (by the stalk formula), followed by the restriction to the displayed target. Since all these maps are restrictions of the same complex, the composite is the given map. If \(w\ne z\), the middle map is zero by Lemma \hyperref[bp253]{12.9}. If \(w=z\), the intermediate object \(M_z\) is zero.
\end{proof}

% text: prose-12-12
\phantomsection\label{bp258}%
Near \((y,y)\) the curves we care about do not meet transversely, since the graph, the diagonal and the axis \(H\) all pass through the same point. The next lemma reduces this configuration to the transverse one by a single blowup.

% object: lemma-vanishing-restrictions-after-blowing-up (blueprint/chap-The-trace-formula/sec-Local-terms-on-a-curve/subsec-Vanishing-on-a-strictly-local-surface/lemma-vanishing-restrictions-after-blowing-up/_section.md)
\phantomsection\label{bp259}%
\begin{bpstatementit}{Lemma~12.11}{ (Vanishing of restrictions after blowing up)}
Let \(T\) be the strict localization of a smooth surface at a closed point \(t\). Let \(H\) be a regular curve through \(t\), and let \(D\) and \(C\) be reduced closed curves through \(t\) such that the tangent cone of \(D\) at \(t\) meets the tangent cones of \(C\) and of \(H\) only in \(0\). Let \(j':T-(D\cup H)\hookrightarrow T-D\) be the open immersion and let \(N\) be a bounded complex on \(T-(D\cup H)\) with finite locally constant cohomology and finite \(\ell\)-group monodromy. Then the restriction map

\[
R\Gamma(T-D,j'_!N)\longrightarrow
R\Gamma(C-t,(j'_!N)|_{C-t})
\]

is zero in \(D(A)\).
\end{bpstatementit}

% proof: proof/lemma-vanishing-restrictions-after-blowing-up
\begin{proof}[{Proof}]\phantomsection\label{bp260}%
After blowing up \(t\), the three curves are separated along the exceptional divisor, which then plays the role of \(D_2\) in the previous lemma.

Let \(\widetilde T\to T\) be the blowup in \(t\), with exceptional divisor \(E\simeq\mathbf P^1_k\), and let \(D',H',C'\) be the strict transforms. In these terms, the hypotheses on tangent cones say that \(D'\cap E\) is disjoint from both \(C'\cap E\) and \(H'\cap E\). We therefore apply Lemmas \hyperref[bp253]{12.9} and \hyperref[bp256]{Lemma 12.10} to the surface \(X=\widetilde T-D'\) with the divisors \(D_1=H'\cap X\) and \(D_2=E\cap X=E-D'\). Here the divisor \(E-D'\) is connected, since it is a projective line with finitely many points removed, and \(H'\cap X\) meets it transversely in the single point \(z\) corresponding to the tangent direction of \(H\). Since the blowup is an isomorphism off \(E\), the source of the map is \(R\Gamma(X-E,j'_!N')\) for the pulled-back complex \(N'\).

The strict transform \(C'\) is proper and quasi-finite over the strictly local curve \(C\), hence finite. Its closed fibre is the finite set of tangent directions of its branches, so \(C'\) is a finite disjoint union of strictly local curves \(C'_w\), each meeting \(E\) at a single point \(w\) of \(E-D'\). Each restriction to \(C'_w-w\) is zero by Lemma \hyperref[bp256]{12.10}. Since \(C'-E\simeq C-t\), the given restriction map is the direct sum of these zero maps.
\end{proof}


% object: subsec-The-local-term-at-a-zero-stalk (blueprint/chap-The-trace-formula/sec-Local-terms-on-a-curve/subsec-The-local-term-at-a-zero-stalk/_section.md), content hidden
\subsubsection*{\texorpdfstring{12.2.3\quad The local term at a zero stalk}{12.2.3 The local term at a zero stalk}}
\phantomsection\label{bp261}%
\addcontentsline{toc}{subsubsection}{\texorpdfstring{12.2.3\quad The local term at a zero stalk}{12.2.3 The local term at a zero stalk}}
% text: prose-12-13
\phantomsection\label{bp262}%
In the setting of Lemma \hyperref[bp246]{12.7} we compute on the classical étale models. This is legitimate because the classical comparison identifies tensor products, duality, exceptional adjunctions and the pairing between graph and diagonal with their versions there. Let \(X_1\) and \(X_2\) be two copies of the strict localization of \(X\) at \(y\). The strict localization of \(X_1\times X_2\) at \((y,y)\) is the strictly local surface \(T\) of Lemma \hyperref[bp246]{12.7}. Write

\[
H=X_1\times\{y\},\qquad V=\{y\}\times X_2,
\]

and keep \(\Gamma\) and \(\Delta\) for the transposed Frobenius graph and the diagonal in \(T\). Let \(u\) and \(v\) be the pullbacks to the two factors \(X_1\) and \(X_2\) of the same uniformizer of \(X\) at \(y\), so that \(\Delta\), \(V\) and \(H\) have equations \(u-v=0\), \(u=0\) and \(v=0\). Write \(\varphi(v)=(\Phi_X^m)^\#u\), a power series in \(v\); the graph \(\Gamma\) has equation \(u-\varphi(v)=0\). Since \(d\Phi_X=0\), the graph \(\Gamma\) is tangent to \(V\), while \(\Delta\) and \(H\) have other tangent directions.

The next three lemmas prepare the proof of Lemma \hyperref[bp270]{12.15}.

% object: lemma-covers-degree-prime-to-ell (blueprint/chap-The-trace-formula/sec-Local-terms-on-a-curve/subsec-The-local-term-at-a-zero-stalk/lemma-covers-degree-prime-to-ell/_section.md)
\phantomsection\label{bp263}%
\begin{bpstatementit}{Lemma~12.12}{ (Covers of degree prime to $\ell$)}
Let \(S\) be the strict localization of a smooth curve over the algebraically closed field \(k\) at a closed point \(s\), and let \(\mathcal F_1,\dots,\mathcal F_r\) be finite locally constant sheaves on \((S-s)_{\acute{e}t}\). There is an integer \(d_1\) prime to \(\ell\) with the following property. For every multiple \(d\) of \(d_1\) that is prime to \(\ell\) there is a finite flat morphism \(S'\to S\) of degree \(d\) such that

\begin{enumerate}
\item \(S'\) is the strict localization of a smooth curve over \(k\) at a closed point \(s'\), and \(s'\) is the only point of \(S'\) over \(s\); more precisely, \(S'\to S\) is the base change of a finite flat morphism \(C'\to C\) of smooth curves over \(k\), where \(C\) is an étale neighbourhood of \(s\), which is étale away from \(s\) and has a single point over \(s\),
\item \(S'-s'\to S-s\) is finite étale, and
\item the monodromy of the pullback of each \(\mathcal F_i\) to \(S'-s'\) factors through a finite \(\ell\)-group.
\end{enumerate}
\end{bpstatementit}

% proof: proof/lemma-covers-degree-prime-to-ell
\begin{proof}[{Proof}]\phantomsection\label{bp264}%
We take the fixed field of an \(\ell\)-Sylow subgroup of the monodromy, adjust the degree by an Eisenstein extension, and descend to curves.

Let \(F\) be the fraction field of the discrete valuation ring of \(S\). Then the punctured trait \(S-s\) is \(\operatorname{Spec}(F)\), so the \(\mathcal F_i\) are finite sets with a continuous action of \(\operatorname{Gal}(F^{\mathrm{sep}}/F)\). We choose a finite Galois extension \(L/F\) over which all of them become constant, let \(P\) be an \(\ell\)-Sylow subgroup of \(G=\operatorname{Gal}(L/F)\), and put \(F_1=L^P\) and \(d_1=[F_1:F]=[G:P]\), which is prime to \(\ell\).

Let \(d=d_1e\) be a multiple of \(d_1\) prime to \(\ell\). The normalization \(S_1\) of \(S\) in the separable extension \(F_1\) is finite over \(S\), and it is local because it is a domain finite over a henselian local ring. Hence \(S_1\) is a strictly henselian discrete valuation ring with residue field \(k\), and it has degree \(d_1\) over \(S\). If \(e>1\), let \(t\) be a uniformizer of \(S_1\). The Eisenstein polynomial \(z^e+tz-t\) is irreducible, and its derivative \(ez^{e-1}+t\) is nonzero, also when \(p\mid e\), so it is separable. Hence \(S'=S_1[z]/(z^e+tz-t)\) is a discrete valuation ring, totally ramified of degree \(e\) over \(S_1\), with separable fraction field \(F'\) over \(F\). If \(e=1\) we put \(S'=S_1\). Then \(S'\to S\) is finite flat of degree \(d\), with a single point \(s'\) over \(s\), and \(S'-s'=\operatorname{Spec}(F')\to
\operatorname{Spec}(F)\) is finite étale. The group \(\operatorname{Gal}(F^{\mathrm{sep}}/F')\) acts on the stalks of the \(\mathcal F_i\) through its image in \(G\), which lies in \(\operatorname{Gal}(L/F_1)=P\) since \(F_1\subset F'\).

Finally, by a limit argument \(S'\to S\) descends to a finite flat morphism \(C'\to C\) onto an étale neighbourhood \(C\) of \(s\), which we may take with normal source. Since \(k\) is perfect, the normal curve \(C'\) is smooth. Moreover, the morphism \(C'\to C\) is étale over a dense open of \(C\), and after shrinking \(C\) around \(s\) it is étale away from \(s\). Since \(k\) is algebraically closed, \(C'\times_CS\) is the product of the strict localizations of \(C'\) at the points over \(s\). As it equals the local scheme \(S'\), there is a single point over \(s\), and \(S'\) is the strict localization of \(C'\) there.
\end{proof}

% object: lemma-tangent-cones-on-cover (blueprint/chap-The-trace-formula/sec-Local-terms-on-a-curve/subsec-The-local-term-at-a-zero-stalk/lemma-tangent-cones-on-cover/_section.md)
\phantomsection\label{bp265}%
\begin{bpstatementit}{Lemma~12.13}{ (Tangent cones on the cover)}
For \(i=1,2\) let \(X_i'\to X_i\) be a finite flat morphism of degree \(d\) as in Lemma \hyperref[bp263]{12.12}, with \(y_i'\) the point over \(y\). Let \(T'\) be the strict localization of \(X_1'\times X_2'\) at \(y'=(y_1',y_2')\), and write \(\Gamma',V',\Delta',H'\) for the inverse images of \(\Gamma,V,\Delta,H\) in \(T'\). Then \(T'\) is the strict localization of a smooth surface over \(k\) at a closed point, the morphism \(T'\to T\) is finite flat of degree \(d^2\) and étale over \(T-(H\cup V)\), the reduction \(H'_{\mathrm{red}}\) is regular, and the tangent cone of \((\Gamma'\cup V')_{\mathrm{red}}\) at \(y'\) meets the tangent cones of \(\Delta'_{\mathrm{red}}\) and of \(H'_{\mathrm{red}}\) only in \(0\).
\end{bpstatementit}

% proof: proof/lemma-tangent-cones-on-cover
\begin{proof}[{Proof}]\phantomsection\label{bp266}%
The two covers descend to finite flat morphisms \(C_i'\to C_i\) of degree \(d\) between smooth curves, where \(C_i\) is an étale neighbourhood of \(y\) in \(X\), étale away from \(y\) and with a single point \(y_i'\) over \(y\) (Lemma lemma-covers-degree-prime-to-ell). Their product \(C_1'\times C_2'\to C_1\times C_2\) is a finite flat morphism of degree \(d^2\) of smooth surfaces, étale away from the two fibres through \(y\), with \(y'\) the only point over \((y,y)\). Base change to the strict localization \(T\) of \(C_1\times C_2\) at \((y,y)\) gives a scheme finite over \(T\) with one point over the closed point, hence local and strictly henselian, and this is \(T'\). This proves the assertions about \(T'\) and \(T'\to T\).

Choose uniformizers \(\xi\) of \(X_1'\) and \(\eta\) of \(X_2'\), so that the completed local ring of \(T'\) at \(y'\) is \(k[[\xi,\eta]]\). Since the residue field \(k\) is algebraically closed, both covers are totally ramified, so \(u\) is a unit of \(k[[\xi]]\) times \(\xi^d\), and similarly for \(v\). Multiplying \(\xi\) and \(\eta\) by suitable \(d\)th roots of nonzero constants, we may assume that

\[
u\equiv\xi^d,\qquad v\equiv\eta^d\pmod{(\xi,\eta)^{d+1}}.
\]

Since \(\Phi_X^m\) is purely inseparable of degree \(q^m\) and fixes \(y\), we have \(\operatorname{ord}_v\varphi=q^m>1\), and hence \(\varphi(v)\in(\eta)^{dq^m}\subseteq(\xi,\eta)^{d+1}\). It follows that the pulled-back equations of \(\Gamma'\), \(V'\), \(\Delta'\) and \(H'\) have initial forms

\[
\xi^d,\qquad\xi^d,\qquad\xi^d-\eta^d,\qquad\eta^d,
\]

respectively. The support of the tangent cone of a curve is the zero set of the initial form of its equation, and it does not change when we pass to the reduction, since if \(J\) is the radical of \(I\) and \(J^n\subset I\) then the ideals of initial forms satisfy \(\mathrm{in}(J)^n\subset\mathrm{in}(I)
\subset\mathrm{in}(J)\). Thus the tangent cone of \((\Gamma'\cup V')_{\mathrm{red}}\) is supported on the line \(\xi=0\), that of \(H'_{\mathrm{red}}\) on \(\eta=0\), and that of \(\Delta'_{\mathrm{red}}\) on the lines \(\xi=\zeta\eta\) with \(\zeta^d=1\). This holds also when \(p\mid d\) and \(\xi^d-\eta^d\) is a \(p\)th power. Away from the origin the first support meets neither of the others. Finally \(H'_{\mathrm{red}}\) is the curve \(\eta=0\), which is regular.
\end{proof}

% object: lemma-descent-vanishing-finite-etale (blueprint/chap-The-trace-formula/sec-Local-terms-on-a-curve/subsec-The-local-term-at-a-zero-stalk/lemma-descent-vanishing-finite-etale/_section.md)
\phantomsection\label{bp267}%
\begin{bpstatementit}{Lemma~12.14}{ (Descent of vanishing along a finite étale cover)}
Let \(f:W'\to W\) be a finite étale morphism of schemes of constant degree \(n\), with \(n\) invertible in \(A\), and let \(L\in D(W_{\acute{e}t},A)\). Then the pullback map \(R\Gamma(W,L)\to R\Gamma(W',f^{-1}L)\) has a left inverse in \(D(A)\). In particular, a morphism \(\alpha:M\to R\Gamma(W,L)\) in \(D(A)\) is zero if its composite with this pullback map is zero.
\end{bpstatementit}

% proof: proof/lemma-descent-vanishing-finite-etale
\begin{proof}[{Proof}]\phantomsection\label{bp268}%
Since \(f\) is finite étale, \(f_*=f_!\) is exact, and the trace map \(f_*f^{-1}\mathcal F\to\mathcal F\) composed with the adjunction unit \(\mathcal F\to f_*f^{-1}\mathcal F\) is multiplication by \(n\) for every sheaf of \(A\)-modules \(\mathcal F\) [Stacks, \href{https://stacks.math.columbia.edu/tag/03SH}{03SH}]. Since both maps are natural and the functors are exact, applying them termwise to a complex representing \(L\) gives maps \(L\to f_*f^{-1}L\to L\) in \(D(W_{\acute{e}t},A)\) whose composite is \(n\). Applying \(R\Gamma(W,-)\) and using \(R\Gamma(W,f_*f^{-1}L)=R\Gamma(W',f^{-1}L)\), we see that \(n^{-1}\) times the map induced by the trace is a left inverse of the pullback.
\end{proof}

% text: prose-12-14
\phantomsection\label{bp269}%
The map \(\rho\) of Lemma \hyperref[bp246]{12.7} depends on \(K\) and on the geometry of \(\Gamma\) and \(\Delta\) near \((y,y)\), but not on \(\psi\). We therefore state the vanishing of local terms at a zero stalk as a statement about \(\rho\) alone.

% object: lemma-vanishing-local-term-zero-stalk (blueprint/chap-The-trace-formula/sec-Local-terms-on-a-curve/subsec-The-local-term-at-a-zero-stalk/lemma-vanishing-local-term-zero-stalk/_section.md)
\phantomsection\label{bp270}%
\begin{bpstatementit}{Lemma~12.15}{ (Vanishing of the local term at a zero stalk)}
In the setting of Lemma \hyperref[bp246]{12.7}, let \(y\) be a fixed point of \(\Phi_X^m\), and suppose \(K_{\bar y}=0\). Then

\[
\rho:R\Gamma_\Gamma(T,P)\longrightarrow R\Gamma_y(\Delta,\delta^*P)
\]

is zero in \(D(A)\). Consequently \(\operatorname{LT}_y(\psi)=0\) for every morphism \(\psi:(\Phi_X^m)^*K\to K\).
\end{bpstatementit}

% proof: proof/lemma-vanishing-local-term-zero-stalk
\begin{proof}[{Proof}]\phantomsection\label{bp271}%
We use the notation introduced before Lemma \hyperref[bp263]{12.12} and compute on the étale model: we identify \(\rho\) with a restriction map \(r\) away from \(y\), kill \(r\) on a cover with \(\ell\)-group monodromy, and descend.

\emph{Step 1: reduction to \(r\).} Let \(N=(\mathbb D_XK)|_{X_1-y}\), the dual of \(K|_{X_1-y}\) with respect to the restriction of \(\omega_X\), and let \(Q=K|_{X_2-y}\). Write \(j_1:X_1-y\hookrightarrow X_1\) and \(j_2:X_2-y\hookrightarrow X_2\) for the inclusions. As \(K_{\bar y}=0\), localization gives \(K|_{X_2}=j_{2!}Q\), and \((\mathbb D_XK)|_{X_1}=Rj_{1*}N\) because the dual of an extension by zero is the direct image of the dual. The ordinary external Künneth comparison therefore identifies \(P=(\mathbb D_XK\boxtimes K)|_T\) with \(Rj_*j'_!(N\boxtimes Q)\), where

\[
j:T-V\hookrightarrow T,\qquad
j':T-(H\cup V)\hookrightarrow T-V.
\]

Both \(R\Gamma(T,P)\) and \(R\Gamma(\Delta,\delta^*P)\) are zero, because the schemes are strictly local and the stalk at the closed point is \((\mathbb D_XK)_{\bar y}\otimes^{\mathbf L} K_{\bar y}=0\). So the localization triangles identify \(\rho\), shifted by one, with the ordinary restriction \(R\Gamma(T-\Gamma,P)\to R\Gamma(\Delta-y,\delta^*P)\), which under the Künneth identification is

\[
r:R\Gamma(T-(V\cup\Gamma),j'_!(N\boxtimes Q))
\longrightarrow R\Gamma(\Delta-y,(N\boxtimes Q)|_{\Delta-y}).
\]

\emph{Step 2: the restriction vanishes on a cover.} Choose bounded finite projective lisse models of \(N\) and \(Q\) by Lemma \bpout{3.13}; their terms are finite locally constant, as a punctured strictly local trait is its generic point. Then Lemma \hyperref[bp263]{12.12}, applied to the terms of the model of \(N\) on \(X_1-y\) and to those of \(Q\) on \(X_2-y\), gives integers \(d_1,d_2\) prime to \(\ell\). With \(d=\operatorname{lcm}(d_1,d_2)\) we then obtain covers \(X_1'\to X_1\) and \(X_2'\to X_2\) of the same degree \(d\). We want equal degrees so that \(u\) and \(v\) become \(d\)th powers for the same \(d\) in Lemma \hyperref[bp265]{12.13}. Let \(T'\), \(\Gamma'\), \(V'\), \(\Delta'\), \(H'\) be as in that lemma, and let \(N'\) and \(Q'\) be the pullbacks of \(N\) and \(Q\) to \(X_1'-y_1'\) and \(X_2'-y_2'\). By construction, the complex \(N'\boxtimes Q'\) has finite \(\ell\)-group monodromy on \(T'-(H'\cup V')\), hence also on the smaller open set obtained by removing \(\Gamma'\). Then the restriction \(r'\) on \(T'\), defined like \(r\), is zero by Lemma \hyperref[bp259]{12.11}, applied to \(T'\) with \(D=(\Gamma'\cup V')_{\mathrm{red}}\), the regular curve \(H'_{\mathrm{red}}\) and \(C=\Delta'_{\mathrm{red}}\); its hypotheses hold by Lemma \hyperref[bp265]{12.13}, and passing to reductions changes neither the open complements nor the étale cohomology.

\emph{Step 3: descent.} The map \(r\) followed by pullback to \(\Delta'-y'\) equals pullback to \(T'\) followed by \(r'=0\), because pullback along \(T'\to T\) commutes with restriction and with extension by zero. As \(\Delta\) meets \(H\cup V\) only at \(y\), Lemma \hyperref[bp265]{12.13} shows that \(\Delta'-y'\to\Delta-y\) is finite étale of degree \(d^2\). This degree is invertible in \(A\), so Lemma \hyperref[bp267]{12.14} gives \(r=0\), and hence \(\rho=0\). By the comparison theorems, this identity on the classical étale models of these strictly local schemes also holds in the pro-étale formulation. Finally, by Lemma \hyperref[bp246]{12.7}, \(\operatorname{LT}_y(\psi)\) is the image of \([\psi]\) under a composite that begins with \(\rho\), so it is zero.
\end{proof}


% object: subsec-The-local-term-at-a-fixed-point (blueprint/chap-The-trace-formula/sec-Local-terms-on-a-curve/subsec-The-local-term-at-a-fixed-point/_section.md), content hidden
\subsubsection*{\texorpdfstring{12.2.4\quad The local term at a fixed point}{12.2.4 The local term at a fixed point}}
\phantomsection\label{bp272}%
\addcontentsline{toc}{subsubsection}{\texorpdfstring{12.2.4\quad The local term at a fixed point}{12.2.4 The local term at a fixed point}}
% text: prose-12-15
\phantomsection\label{bp273}%
The point \(y\) is stable under \(\Phi_X^m\), so additivity of trace classes applies to the localization triangle for \(X\setminus\{y\}\) and \(\{y\}\). Here the open piece has zero stalk at \(y\), so its local term vanishes by Lemma \hyperref[bp270]{12.15}. The closed piece, in turn, is pushed forward from the point, where the local term is a trace by Lemma \hyperref[bp199]{11.10}. Note that the argument never uses that \(K\) is lisse at \(y\). This is why, in the proof of Theorem \hyperref[bp287]{12.20}, we can apply the result to \(j_!K\) at the point at infinity of \(\mathbf P^1\).

% object: proposition-frobenius-local-term-curve (blueprint/chap-The-trace-formula/sec-Local-terms-on-a-curve/subsec-The-local-term-at-a-fixed-point/proposition-frobenius-local-term-curve/_section.md)
\phantomsection\label{bp274}%
\begin{bpstatementit}{Proposition~12.16}{ (The Frobenius local term on a curve)}
Let \(X_0\) be a smooth curve over \(k_0\), let \(X=X_0\times_{k_0}k\), let \(K\in D^b_{\mathrm{ctf}}(X,A)\) and let \(\psi:(\Phi_X^m)^*K\to K\) be any morphism. For \(y\in X_0(\mathbf F_{q^m})\) one has

\[
\operatorname{LT}_y(\psi)
=\operatorname{tr}_A(\psi_{\bar y}\mid K_{\bar y}).
\]
\end{bpstatementit}

% proof: proof/proposition-frobenius-local-term-curve
\begin{proof}[{Proof}]\phantomsection\label{bp275}%
Let \(i:\{y\}\hookrightarrow X\) and \(j:U=X\setminus\{y\}\hookrightarrow X\). We split \(\psi\) along the localization triangle and compute the two local terms separately.

\emph{Splitting \(\psi\).} Both pieces are stable under \(\Phi_X^m\), and in fact \(((\Phi_X^m)^{-1}\{y\})_{\mathrm{red}}=\{y\}\). In the enhanced derived category, the functorial localization triangle is a two-step filtered object,

\[
K_U=j_!j^*K\longrightarrow K\longrightarrow i_*i^*K
\xrightarrow{+1},
\]

and \(\psi\) induces a coherent map from the \((\Phi_X^m)^*\)-pullback of this filtration to the filtration itself, with induced correspondences \(\psi_U\) and \(i_*\psi_{\bar y}\) on the pieces. This rests on two base change identities. For the open part, \((\Phi_X^m)^*j_!\simeq j_!(\Phi_X^m|_U)^*\) by base change along inverse images in the cartesian square. The square for the point, by contrast, is cartesian only after passing to reductions, so there proper base change followed by nilpotent invariance gives \((\Phi_X^m)^*i_*\simeq i_*(\Phi_X^m|_y)^*\). Both comparisons are mates of the adjunction units for \(j\) and \(i\) used in the triangle, which is why the induced maps are compatible with the filtration.

\emph{Additivity.} The second half of Proposition \hyperref[bp223]{11.19} applies to the graph \(c_m\), because its hypothesis \(\Phi_X^m(\{y\})\subset\{y\}\) holds for the fixed point \(y\). It gives

\[
\operatorname{LT}_y(\psi)
=\operatorname{LT}_y(\psi_U)
+\operatorname{LT}_y(i_*\psi_{\bar y}).
\]

\emph{The two local terms.} The first term is zero by Lemma \hyperref[bp270]{12.15}, since \((K_U)_{\bar y}=0\). For the second, the description of the closed piece \(u_Z\) in Proposition \hyperref[bp223]{11.19} shows that \(i_*\psi_{\bar y}\) is the pushforward along \(i\) of the cohomological correspondence \(\psi_{\bar y}\) on \((y,K_{\bar y})\) for the identity correspondence of \(y\). Its local term is therefore \(\operatorname{tr}_A(\psi_{\bar y}\mid K_{\bar y})\), the trace on the perfect complex \(K_{\bar y}\), by the compatibility with proper pushforward in Theorem \hyperref[bp206]{11.13} and Lemma \hyperref[bp199]{11.10}.
\end{proof}

% text: prose-12-16
\phantomsection\label{bp276}%
For \(K=A_X\) on \(X=\mathbf A^1_k\) the proposition gives \(\operatorname{LT}_y(\psi_m)=1\) at each of the \(q^m\) fixed points of \(\Phi^m_X\). This matches the intersection count at the start of the chapter. For general \(K\) the proposition is not a formal consequence of the definitions, as Remark \hyperref[bp239]{12.4} explains.

% object: example-frobenius-projective-line (blueprint/chap-The-trace-formula/sec-Local-terms-on-a-curve/subsec-The-local-term-at-a-fixed-point/example-frobenius-projective-line/_section.md)
\phantomsection\label{bp277}%
\begin{bpstatement}{Example~12.17}{ (Frobenius on the projective line)}
Let \(X=\mathbf P^1_k\), \(K=A_X\), and let \(\psi\) be the constant Frobenius correspondence, so that \(\psi_m\) is the identity on every stalk. Both sides of the Lefschetz–Verdier formula for the graph \(c_m\) can be computed explicitly. The fixed points of \(\Phi_X^m\) are the \(q^m+1\) points of \(\mathbf P^1(\mathbf F_{q^m})\). By the proposition each has local term \(\operatorname{tr}_A(1\mid A)=1\), so the local side is \(q^m+1\). On the global side, \(R\Gamma(\mathbf P^1,A)=A\oplus A(-1)[-2]\). By Lemma \hyperref[bp234]{12.3} the operator \(T_{c_m}(\psi_m)\) is the \(m\)-th power of geometric Frobenius. It acts by \(1\) on \(H^0\) and by \(q^m\) on \(H^2=A(-1)\), and both degrees are even, so its trace is \(1+q^m\) with no signs. Compare the identity correspondence at the start of Section \hyperref[bp195]{11.4}. There the fixed locus is all of \(\mathbf P^1\) and the answer \(2\) is a self-intersection. Here the fixed locus is finite and each point counts once.
\end{bpstatement}



% object: sec-Frobenius-and-Rf (blueprint/chap-The-trace-formula/sec-Frobenius-and-Rf/_section.md), content hidden
\subsection*{\texorpdfstring{12.3\quad Frobenius and $Rf_!$}{12.3 Frobenius and Rf!}}
\phantomsection\label{bp278}%
\addcontentsline{toc}{subsection}{\texorpdfstring{12.3\quad Frobenius and $Rf_!$}{12.3 Frobenius and Rf!}}
% text: prose-12-17
\phantomsection\label{bp279}%
The trace formula with finite coefficients is proved by induction on the number of affine coordinates. We project a closed subscheme of \(\mathbf A^d\) to \(\mathbf A^{d-1}\), apply the induction hypothesis to \(Rf_!K\), and identify each stalk of \(Rf_!K\) at a rational point with the cohomology of the fibre, a closed subscheme of \(\mathbf A^1\) over a larger finite field. For this, \(Rf_!K\) must carry a Frobenius correspondence with two properties: its operator on \(R\Gamma_c\) is that of \(\psi\), and its stalk maps are the Frobenius operators on the cohomology of the fibres.

The square formed by \(f\) and the two Frobenii commutes, but it is usually not cartesian. Indeed, already for \(f:\mathbf A^1\to s\) the Frobenius of \(s\) is the identity, so the fibre product is \(\mathbf A^1\) itself, while \(\Phi_{\mathbf A^1}\) has degree \(q\). The difference is measured by the relative Frobenius. It is a universal homeomorphism, so pullback along it is an equivalence on étale sheaves, and the comparison map built from it is an isomorphism.

% object: lemma-comparison-via-relative-frobenius (blueprint/chap-The-trace-formula/sec-Frobenius-and-Rf/lemma-comparison-via-relative-frobenius/_section.md)
\phantomsection\label{bp280}%
\begin{bpstatementit}{Lemma~12.18}{ (Comparison via relative Frobenius)}
Let \(f_0:X_0\to Y_0\) be a morphism of separated schemes of finite type over \(k_0\), write \(f:X\to Y\) for its geometric base change, and let \(m\geq1\). Put

\[
X'_m=Y\times_{\Phi_Y^m,Y,f}X,
\quad p_m:X'_m\to X,\quad f'_m:X'_m\to Y,
\quad r_m=(f,\Phi_X^m):X\to X'_m,
\]

where \(p_m\) and \(f'_m\) are the two projections.

\begin{enumerate}
\item The relative Frobenius \(r_m\) is a finite universal homeomorphism, \(p_mr_m=\Phi_X^m\) and \(f'_mr_m=f\), and there are natural isomorphisms \[
   \begin{aligned}
   \theta_{f,m}:(\Phi_Y^m)^*Rf_!K
   &\simeq Rf'_{m!}p_m^*K
   \xrightarrow{Rf'_{m!}\eta_{r_m}}
   Rf'_{m!}Rr_{m*}r_m^*p_m^*K\\
   &\simeq Rf_!(\Phi_X^m)^*K,
   \end{aligned}
   \] compatible with iteration in \(m\).
\item For a Frobenius correspondence \(\psi:\Phi_X^*K\to K\) on \(K\in D^b_{\mathrm{ctf}}(X,A)\), the object \(L=Rf_!K\) of \(D^b_{\mathrm{ctf}}(Y,A)\) carries the Frobenius correspondence \(\psi'=Rf_!\psi\circ\theta_{f,1}:\Phi_Y^*L\to L\), and its iterates are \[
   \psi'_m=Rf_!\psi_m\circ\theta_{f,m}.
   \]
\item The composition isomorphism \(R\Gamma_c(Y,L)\simeq R\Gamma_c(X,K)\) intertwines the endomorphisms \(T_{c_m}(\psi'_m)\) on \(Y\) and \(T_{c_m}(\psi_m)\) on \(X\) of Definition \hyperref[bp232]{12.2}.
\item At \(y\in Y_0(\mathbf F_{q^m})\), base change for compact supports identifies \((L_{\bar y},\psi'_{m,\bar y})\) with \[
   \left(R\Gamma_c(X_{\bar y},K|_{X_{\bar y}}),
   T_{c_1}(\psi_m|_{X_{\bar y}})\right),
   \] where the fibre is regarded as a scheme over \(\mathbf F_{q^m}\), with the \(q^m\)-power Frobenius as its scheme Frobenius.
\end{enumerate}
\end{bpstatementit}

% proof: proof/lemma-comparison-via-relative-frobenius
\begin{proof}[{Proof}]\phantomsection\label{bp281}%
The only geometric fact we use is that, since \(r_m\) is a universal homeomorphism, pullback along it is an equivalence and its unit is an isomorphism. The rest is base change and composition of units, and each part uses the previous ones.

\emph{Proof of (1).} The identities \(p_mr_m=\Phi_X^m\) and \(f'_mr_m=f\) hold by the definition of \(r_m=(f,\Phi_X^m)\), since \(f\Phi_X^m=\Phi_Y^mf\). The map \(r_m\) is the geometric base change of the relative \(q^m\)-power Frobenius of \(X_0/Y_0\). It is integral, radicial and surjective [Stacks, \href{https://stacks.math.columbia.edu/tag/0CCB}{0CCB}], and finite because it is integral and of finite type. For affine \(Y_0=\operatorname{Spec}R\) and \(X_0=\operatorname{Spec}S\), with \(F\) the \(q\)-power Frobenius of \(R\), it is given on coordinate rings by

\[
R\otimes_{F^m,R}S\longrightarrow S,
\qquad b\otimes a\longmapsto f^\#(b)a^{q^m}.
\]

By topological invariance for classical objects of finite Tor amplitude, \(r_m^*\) is an equivalence, so the unit \(\eta_{r_m}\) in the construction is an isomorphism. In the construction, the first arrow is base change for compact supports, and the last comes from the composition of direct images with compact support, since \(r_m\) is proper. Moreover, pasting the cartesian squares defining \(X'_m\) and \(X'_n\), we get the comparison for \(m+n\), because the relative Frobenii compose and the unit of their composite is the composite of the units. Thus

\[
\theta_{f,m+n,K}=
\theta_{f,m,(\Phi_X^n)^*K}\circ
(\Phi_Y^m)^*\theta_{f,n,K},
\]

which is the compatibility with iteration.

\emph{Proof of (2).} The formula for \(\psi'_m\) follows by induction on \(m\) from the displayed iteration identity and the naturality of \(\theta\).

\emph{Proof of (3).} Both endomorphisms are a pullback map followed by \(R\Gamma_c\) of the correspondence, so we must see that the pullback along \(\Phi_Y^m\) on \(R\Gamma_c(Y,L)\) becomes the pullback along \(\Phi_X^m\) on \(R\Gamma_c(X,K)\). For a finite map \(h\) we also write \(h^*\) for the pullback map on cohomology with compact support, defined by the unit of \(h^*\dashv Rh_*\). Since base change is compatible with these units, the pullback along \(\Phi_Y^m\), followed by the base change isomorphism, is the pullback along \(p_m\), and the next arrow in \(\theta_{f,m}\) becomes \(r_m^*\). Their composite is \(r_m^*p_m^*=(p_mr_m)^*=(\Phi_X^m)^*\), and composing with \(R\Gamma_c(\psi_m)\) we see that the operators are intertwined.

\emph{Proof of (4).} We pull the construction back to the fixed point \(y\). Since \(\Phi_Y^m\) fixes \(\bar y\), the fibre of \(X'_m\to Y\) over \(\bar y\) is again the fibre \(X_{\bar y}\) of \(f\), and \(r_m\) restricts to a self-map of \(X_{\bar y}\). For \(f:\mathbf A^1\to s\) and \(m=1\), for instance, \(X'_1=\mathbf A^1\) and \(r_1\) is \(t\mapsto t^q\). In general, the base change of \(r_m\) to this fibre is its \(q^m\)-power scheme Frobenius, and by proper base change its unit restricts to the unit on the fibre. Hence the same calculation identifies \(\psi'_{m,\bar y}\) with the Frobenius operator on the cohomology with compact support of the fibre. Moreover, at a rational point \(x\) of the fibre, the restricted correspondence \(\psi_m|_{X_{\bar y}}\) has stalk map \(\psi_{m,\bar x}\), by composition of pullbacks.
\end{proof}

% text: prose-12-18
\phantomsection\label{bp282}%
To reduce the trace formula to pieces of \(X_0\), we cut \(X_0\) into an open subscheme and its closed complement.

% object: lemma-additivity-open-subscheme-complement (blueprint/chap-The-trace-formula/sec-Frobenius-and-Rf/lemma-additivity-open-subscheme-complement/_section.md)
\phantomsection\label{bp283}%
\begin{bpstatementit}{Lemma~12.19}{ (Additivity for an open subscheme and its complement)}
Let \(j_0:U_0\hookrightarrow X_0\) be open with reduced closed complement \(i_0:Z_0\to X_0\), and let \(K\in D^b_{\mathrm{ctf}}(X,A)\). A Frobenius correspondence \(\psi\) on \(K\) induces Frobenius correspondences on \(j^*K\), \(i^*K\), \(j_!j^*K\) and \(Ri_*i^*K\), whose stalk maps are the restrictions of \(\psi_m\). Writing \(T_{c_m}\) also for the endomorphisms of \(R\Gamma_c\) induced by the Frobenius correspondences on \(U\) and \(Z\), for every \(m\geq1\)

\[
\operatorname{tr}_A(T_{c_m}(\psi_m)\mid R\Gamma_c(X,K))
=\operatorname{tr}_A(T_{c_m}\mid R\Gamma_c(U,j^*K))
+\operatorname{tr}_A(T_{c_m}\mid R\Gamma_c(Z,i^*K)).
\]
\end{bpstatementit}

% proof: proof/lemma-additivity-open-subscheme-complement
\begin{proof}[{Proof}]\phantomsection\label{bp284}%
The four induced correspondences come from base change, and \(R\Gamma_c\) of the localization triangle is a two-step filtered object with an endomorphism, on which traces add.

Since Frobenius preserves opens defined over the finite field, its square with \(j\) is cartesian, and by base change we get the correspondences on \(j^*K\) and \(j_!j^*K\). The Frobenius square of the closed complement, however, is cartesian only after passing to reductions, and we obtain the correspondences on \(i^*K\) and \(Ri_*i^*K\) from Lemma lemma-comparison-via-relative-frobenius–(2) and nilpotent invariance. We then apply the functorial localization triangle \(j_!j^*K\to K\to Ri_*i^*K\) to the Frobenius correspondence and take direct images with compact support. Since the triangle is functorial, the Frobenius correspondence becomes an endomorphism of a two-step filtered object, and hence the traces add by Proposition \hyperref[bp176]{11.2}. Finally, Lemma lemma-comparison-via-relative-frobenius for \(j\) and \(i\) identifies the endomorphisms \(T_{c_m}\) of the cohomology of \(j_!j^*K\) and \(Ri_*i^*K\) with those on \(U\) and \(Z\).
\end{proof}


% object: sec-Finite-coefficients-4 (blueprint/chap-The-trace-formula/sec-Finite-coefficients-4/_section.md), content hidden
\subsection*{\texorpdfstring{12.4\quad Finite coefficients}{12.4 Finite coefficients}}
\phantomsection\label{bp285}%
\addcontentsline{toc}{subsection}{\texorpdfstring{12.4\quad Finite coefficients}{12.4 Finite coefficients}}
% text: prose-12-19
\phantomsection\label{bp286}%
The proof uses a single compactification, \(\mathbf A^1\subset\mathbf P^1\), so the Lefschetz–Verdier formula and the curve computation of Proposition \hyperref[bp274]{12.16} are applied only on \(\mathbf P^1\). This suffices because every affine scheme of finite type over \(\mathbf F_q\) is a closed subscheme of some \(\mathbf A^d\), and the fibres of a coordinate projection lie in \(\mathbf A^1\). On \(\mathbf P^1\) the only new fixed point is the rational point \(\infty\), where \(j_!K\) has zero stalk. The fibres over points of \(Y_0(\mathbf F_{q^m})\) are schemes over \(\mathbf F_{q^m}\), so we prove the theorem for all finite fields at once. The other coefficient rings are treated in Section \hyperref[bp289]{12.5}.

The base of the induction is the case of finite schemes. There we can see directly what the formula counts. For example, take \(X_0=\operatorname{Spec}\mathbf F_{q^2}\) and \(K=A\). Then the scheme \(X\) consists of two points interchanged by \(\Phi_X\). Hence \(T_{c_1}\) is the swap of \(A\oplus A\), with trace \(0=\#X_0(\mathbf F_q)\), and \(T_{c_2}\) is the identity, with trace \(2=\#X_0(\mathbf F_{q^2})\). In general \(T_{c_m}\) is a block permutation matrix, and only the points fixed by \(\Phi_X^m\) contribute to its trace.

The iterates \(\psi_m\), the endomorphisms \(T_{c_m}(\psi_m)\) and the stalk maps \(\psi_{m,\bar x}\) are those of Definition \hyperref[bp232]{12.2}.

% object: theorem-trace-formula-frobenius-correspondence (blueprint/chap-The-trace-formula/sec-Finite-coefficients-4/theorem-trace-formula-frobenius-correspondence/_section.md)
\phantomsection\label{bp287}%
\begin{bpstatementit}{Theorem~12.20}{ (Trace formula for a Frobenius correspondence)}
Let \(X_0/\mathbf F_q\) be separated and of finite type, put \(X=X_0\times_{\mathbf F_q}k\), let \(A\) be any finite commutative \(\ell\)-primary ring and \(K\in D^b_{\mathrm{ctf}}(X,A)\), and let \(\psi:\Phi_X^*K\to K\) be any Frobenius correspondence on \(K\). For every \(m\geq1\), in \(A\),

\[
\operatorname{tr}_A(T_{c_m}(\psi_m)\mid R\Gamma_c(X_{\mathrm{pro}\acute{e}t},K))
=\sum_{x\in X_0(\mathbf F_{q^m})}
\operatorname{tr}_A(\psi_{m,\bar x}\mid K_{\bar x}).
\]
\end{bpstatementit}

% proof: proof/theorem-trace-formula-frobenius-correspondence
\begin{proof}[{Proof}]\phantomsection\label{bp288}%
Both sides are defined, because the complexes are perfect by Lemma \hyperref[bp105]{10.1}. We prove the assertion simultaneously for all finite fields, all Frobenius correspondences and all \(m\), so that \(\mathbf F_q\) may be replaced by \(\mathbf F_{q^m}\) when we pass to a rational fibre. We go in four steps: finite schemes; the affine line, where we use the Lefschetz–Verdier formula on \(\mathbf P^1\); closed subschemes of \(\mathbf A^d\), by projecting away one coordinate; and arbitrary \(X_0\), by cutting it into affine pieces.

\emph{Finite schemes.} If \(X_0\) is finite, both sides of the formula can be read off from a block matrix. Indeed, up to nilpotents, \(X\) is the finite disjoint union of its geometric points. By nilpotent invariance, \(R\Gamma_c(X,K)=\bigoplus_{z\in X(k)}K_z\). In this decomposition, the pullback along \(\Phi_X^m\) given by the unit permutes these summands, and the block of \(T_{c_m}(\psi_m)\) from \(K_{\Phi_X^m(z)}\) to \(K_z\) is the stalk of \(\psi_m\) at \(z\), a map \(K_{\Phi_X^m(z)}=((\Phi_X^m)^*K)_z\to K_z\). On a finite projective model only the diagonal blocks contribute to the trace. These occur exactly at the points \(z\in X_0(\mathbf F_{q^m})\), with the required stalk maps.

\emph{The affine line.} This is the only step where we compute a local term. We pass to \(\mathbf P^1\), where the Lefschetz–Verdier formula applies, and extend \(K\) by zero, so that the new fixed point \(\infty\) contributes nothing. Let \(j:\mathbf A^1_k\hookrightarrow\mathbf P^1_k\). The correspondence \(\psi_m\) extends to a correspondence \(\bar\psi_m\) on \(j_!K\), compatible with iteration. Indeed, \(j\) is defined over \(\mathbf F_q\), so its square with the Frobenius maps is cartesian, and base change for this square gives \((\Phi_{\mathbf P^1}^m)^*j_!K\simeq j_!(\Phi_{\mathbf A^1}^m)^*K\). Lemma lemma-comparison-via-relative-frobenius, applied to \(j\), shows that the composition isomorphism

\[
R\Gamma_c(\mathbf A^1,K)
\simeq R\Gamma(\mathbf P^1,j_!K)
\]

intertwines the endomorphisms \(T_{c_m}\) induced by \(\psi\) and \(\bar\psi\). The graph is proper on the proper curve \(\mathbf P^1\), so by Theorem \hyperref[bp206]{11.13} the trace on the right is the sum of the local terms at the fixed points. By Proposition \hyperref[bp274]{12.16}, each local term is a trace on a stalk. In particular, the local term at infinity vanishes, because the stalk of \(j_!K\) there is zero. The other fixed points and stalk maps are those of \(\mathbf A^1_{\mathbf F_q}\). This proves the formula for the affine line.

\emph{Closed subschemes of the affine line.} Let \(X_0=\operatorname{Spec}\mathbf F_q[t]/I\). If \(I=0\), the previous paragraph applies. Otherwise \(I\) contains a nonzero polynomial, so \(X_0\) is finite over \(\mathbf F_q\) and the case of finite schemes applies. This also covers nonreduced subschemes.

\emph{Induction on affine coordinates.} We push \(K\) forward along a coordinate projection. This does not change the global side. The local side over a point of the base becomes the global side of a fibre lying in \(\mathbf A^1\). Suppose that the assertion holds for closed subschemes of \(\mathbf A^{d-1}\) over every finite field, and let \(X_0\hookrightarrow\mathbf A^d_{\mathbf F_q}\) be closed. Let \(f_0:X_0\to Y_0=\mathbf A^{d-1}_{\mathbf F_q}\) forget the last coordinate, and let \(f:X\to Y\) be its base change to \(k\). The complex \(L=Rf_!K\) lies in \(D^b_{\mathrm{ctf}}(Y,A)\) by Proposition \bpout{7.1}, and we equip it with the Frobenius correspondence \(\psi'\) of Lemma lemma-comparison-via-relative-frobenius. As Lemma lemma-comparison-via-relative-frobenius identifies the endomorphisms \(T_{c_m}\), the induction hypothesis shows that

\[
\operatorname{tr}_A(T_{c_m}(\psi_m)\mid R\Gamma_c(X,K))
=\sum_{y\in Y_0(\mathbf F_{q^m})}
\operatorname{tr}_A(\psi'_{m,\bar y}\mid L_{\bar y}).
\]

By Lemma lemma-comparison-via-relative-frobenius, each summand is the trace on the cohomology with compact support of the fibre \(X_{0,y}\), a closed subscheme of \(\mathbf A^1_{\mathbf F_{q^m}}\). Here the endomorphism is given by the restriction of \(\psi_m\) to the fibre and the first power of the Frobenius over \(\mathbf F_{q^m}\). The case of closed subschemes of the affine line then turns this summand into the sum of the traces on stalks over \(X_{0,y}(\mathbf F_{q^m})\). As \(y\) varies, these sets partition \(X_0(\mathbf F_{q^m})\), which proves the affine case.

\emph{General schemes.} Both sides are additive for a decomposition into an open subscheme and its closed complement: the global side by Lemma \hyperref[bp283]{12.19}, and the local side because the rational points split. It is therefore enough to cut \(X_0\) into affine pieces. Choose a finite affine open cover \(U_1,\ldots,U_N\) of \(X_0\), let \(W_i=\bigcup_{j\leq i}U_j\) and \(V_i=W_i\setminus W_{i-1}\), with the reduced closed structure. Each \(V_i\) is closed in the affine \(U_i\), hence affine, and it is the closed complement of \(W_{i-1}\) in \(W_i\). Applying Lemma \hyperref[bp283]{12.19} successively, we see that the traces add over this stratification. The affine case applied to each \(V_i\) now gives the equality for all separated \(X_0\) of finite type.
\end{proof}


% object: sec-ell-adic-coefficients-5 (blueprint/chap-The-trace-formula/sec-ell-adic-coefficients-5/_section.md), content hidden
\subsection*{\texorpdfstring{12.5\quad $\ell$-adic coefficients}{12.5 ell-adic coefficients}}
\phantomsection\label{bp289}%
\addcontentsline{toc}{subsection}{\texorpdfstring{12.5\quad $\ell$-adic coefficients}{12.5 ell-adic coefficients}}
% text: prose-12-20
\phantomsection\label{bp290}%
To pass from finite coefficient rings to \(\mathcal O_E\), \(E\) and \(\overline{\mathbf Q}_\ell\) we follow the steps of Proposition \bpout{7.12}. Since the statement is now an equality of traces in \(\Lambda\) rather than an isomorphism, we carry the steps out directly. First, for \(\mathcal O_E\) we reduce modulo \(\pi_E^n\). This works because both sides are traces of perfect complexes, traces commute with reduction, and \(\mathcal O_E\) embeds in \(\varprojlim_n\mathcal O_E/\pi_E^n\). Next, for \(E\) we clear denominators: a Frobenius correspondence \(\psi_E\) with coefficients \(E\) need not come from an integral one; only some multiple \(\pi_E^N\psi_E\) does, and this integral model need not be invertible even when \(\psi_E\) is (Lemma \hyperref[bp294]{12.22}). This is why we prove the formula for correspondences that are not invertible. Finally, for \(\overline{\mathbf Q}_\ell\) we descend to a finite extension of \(\mathbf Q_\ell\).

The Frobenius comparison \(\kappa_m\) of Lemma \hyperref[bp234]{12.3} extends to the coefficient rings \(\mathcal O_E\), \(E\) and \(\overline{\mathbf Q}_\ell\). Let \(K\in\Dbc(X,\mathcal O_E)\) have tower of reductions \((K_n)\) modulo \(\pi_E^n\) (Definition \bpout{2.61}). Both \((\Phi_X^m)^*\) and \(\operatorname{Fr}_X^{m*}\) preserve derived completeness and commute with reduction, so each is the derived limit of its values on the \(K_n\). The finite comparisons \(\kappa_m\) come from a single functorial identification on the classical site. Hence they are compatible with reduction of coefficients, and their derived limit is the integral \(\kappa_m\). The finite comparisons are compatible with iteration and with the units of the pullbacks, and therefore so is the integral one. It follows that the integral \(\kappa_m\) identifies the endomorphisms \(T_{c_m}\) with the geometric Frobenius operators as in Lemma \hyperref[bp234]{12.3}. For \(E\) we localize on integral models. By the localized Hom comparison (9.21) the result does not depend on the choice of model. For \(\overline{\mathbf Q}_\ell\) we descend from finite extensions.

% object: proposition-integral-trace-formula (blueprint/chap-The-trace-formula/sec-ell-adic-coefficients-5/proposition-integral-trace-formula/_section.md)
\phantomsection\label{bp291}%
\begin{bpstatementit}{Proposition~12.21}{ (Integral trace formula)}
Let \(K\in\Dbc(X,\mathcal O_E)\) and let \(\theta:\Phi_X^*K\to K\) be any morphism, that is, a Frobenius correspondence with integral coefficients. Define its iterates \(\theta_m\), the endomorphisms \(T_{c_m}(\theta_m)\) and the stalk maps \(\theta_{m,\bar y}\) by the formulas of Definition \hyperref[bp232]{12.2}. Then for every \(m\geq1\) we have

\[
\operatorname{tr}_{\mathcal O_E}
\bigl(T_{c_m}(\theta_m)\mid R\Gamma_c(X_{\mathrm{pro}\acute{e}t},K)\bigr)
=\sum_{y\in X_0(\mathbf F_{q^m})}
\operatorname{tr}_{\mathcal O_E}
\bigl(\theta_{m,\bar y}\mid K_{\bar y}\bigr),
\tag{12.1}
\]

where both sides are categorical traces of perfect complexes over \(\mathcal O_E\). Their reductions modulo \(\pi_E^n\) are the corresponding traces of Theorem \hyperref[bp287]{12.20} for \(K_n\) and the reduction \(\bar\theta_n:\Phi_X^*K_n\to K_n\) of \(\theta\).
\end{bpstatementit}

% proof: proof/proposition-integral-trace-formula
\begin{proof}[{Proof}]\phantomsection\label{bp292}%
We show that both sides reduce modulo \(\pi_E^n\) to the two sides of the finite formula, by identifying the reductions of the complexes, of the operators and of the traces. This suffices, since an element of \(\mathcal O_E\) that vanishes modulo every \(\pi_E^n\) is zero.

\emph{Complexes and operators.} By Lemma \hyperref[bp105]{10.1}, \(R\Gamma_c(X_{\mathrm{pro}\acute{e}t},K)\) and the geometric stalks are perfect, and Proposition \bpout{7.10} provides canonical isomorphisms

\[
\begin{aligned}
R\Gamma_c(X_{\mathrm{pro}\acute{e}t},K)\otimes^{\mathbf L}_{\mathcal O_E}\mathcal O_E/\pi_E^n
&\simeq R\Gamma_c(X_{\mathrm{pro}\acute{e}t},K_n),\\
K_{\bar y}\otimes^{\mathbf L}_{\mathcal O_E}\mathcal O_E/\pi_E^n&\simeq(K_n)_{\bar y}.
\end{aligned}
\]

Under these isomorphisms the reductions of the integral operators are the operators of \(\bar\theta_n\), because extension of scalars commutes with the pullbacks, with the units defining the endomorphisms \(T_{c_m}\), and with composition, and hence with the iterates \(\theta_m\). Moreover Theorem \hyperref[bp287]{12.20} applies for every \(n\), because \(K_n\in D^b_{\mathrm{ctf}}(X,\mathcal O_E/\pi_E^n)\).

\emph{Traces.} Represent each of the two integral complexes by a bounded complex \(P^\bullet\) of finite projective \(\mathcal O_E\)-modules and its endomorphism by a chain map \(v^\bullet\). By Proposition \hyperref[bp176]{11.2} its trace is

\[
\operatorname{tr}_{\mathcal O_E}(v)
=\sum_i(-1)^i\operatorname{Tr}(v^i\mid P^i).
\]

The reduction of this expression modulo \(\pi_E^n\) is the trace on \(P^\bullet\otimes_{\mathcal O_E}\mathcal O_E/\pi_E^n\), as each \(P^i\) is finite projective.

\emph{Conclusion.} The reduction of each side of (12.1) is therefore the corresponding side of the finite formula for \((K_n,\bar\theta_n)\). Hence the difference of the two sides lies in \(\pi_E^n\mathcal O_E\) for every \(n\). It is zero because \(\mathcal O_E\hookrightarrow\varprojlim_n\mathcal O_E/\pi_E^n\) is injective.
\end{proof}

% text: prose-12-21
\phantomsection\label{bp293}%
To pass from \(\mathcal O_E\) to \(E\) we need an integral Frobenius correspondence whose rationalization is the given one, up to a power of \(\pi_E\). The next lemma provides it and records how the operators of both sides of the trace formula change under this rescaling.

% object: lemma-integral-models-rational-frobenius-correspondences (blueprint/chap-The-trace-formula/sec-ell-adic-coefficients-5/lemma-integral-models-rational-frobenius-correspondences/_section.md)
\phantomsection\label{bp294}%
\begin{bpstatementit}{Lemma~12.22}{ (Integral models for rational Frobenius correspondences)}
Let \(K_E\in\Dbc(X,E)\) and let \(\psi_E:\Phi_X^*K_E\to K_E\) be any morphism. There exist \(K_0\in\Dbc(X,\mathcal O_E)\), an isomorphism \(K_0\otimes^{\mathbf L}_{\mathcal O_E}E\simeq K_E\), an integer \(N\geq0\), and a Frobenius correspondence \(\theta:\Phi_X^*K_0\to K_0\) such that

\[
\psi_E=\pi_E^{-N}(\theta\otimes^{\mathbf L}_{\mathcal O_E}E).
\]

For every \(m\geq1\) the operators of Definition \hyperref[bp232]{12.2} satisfy

\[
\begin{aligned}
T_{c_m}((\psi_E)_m)&=\pi_E^{-Nm}
\bigl(T_{c_m}(\theta_m)\otimes^{\mathbf L}_{\mathcal O_E}E\bigr),\\
(\psi_E)_{m,\bar y}&=\pi_E^{-Nm}
\bigl(\theta_{m,\bar y}\otimes^{\mathbf L}_{\mathcal O_E}E\bigr).
\end{aligned}
\tag{12.2}
\]
\end{bpstatementit}

% proof: proof/lemma-integral-models-rational-frobenius-correspondences
\begin{proof}[{Proof}]\phantomsection\label{bp295}%
Morphisms in the rational category are integral morphisms divided by a power of \(\pi_E\). This gives \(K_0\) and \(\theta\). The formulas for the operators hold because the operators of Definition \hyperref[bp232]{12.2} commute with extension of scalars.

By Proposition \bpout{3.27}, extension of scalars induces an equivalence

\[
\Dbc(X,\mathcal O_E)[1/\ell]\simeq\Dbc(X,E),
\]

where the source has the objects of \(\Dbc(X,\mathcal O_E)\) and Hom groups tensored with \(E\). By essential surjectivity, \(K_E\) has an integral model \(K_0\). The model constructed there is built by finite cones and shifts from the integral generators of Lemma lemma-integral-models, and hence in general it is not a lattice in a lisse sheaf.

Since inverse image preserves integral constructibility, \(\Phi_X^*K_0\) is an integral model of \(\Phi_X^*K_E\). By the description of the Hom groups above, applied to the pair \((\Phi_X^*K_0,K_0)\), we may write \(\psi_E=\pi_E^{-N}(\theta\otimes^{\mathbf L}_{\mathcal O_E}E)\). Since pullback is \(\mathcal O_E\)-linear, each of the \(m\) factors of the iterate contributes \(\pi_E^{-N}\), and so \((\psi_E)_m=\pi_E^{-Nm}(\theta_m\otimes^{\mathbf L} E)\). By Proposition \bpout{7.10}, the change-of-coefficients maps for cohomology with compact support and for stalks respect the units of the pullback adjunctions. Hence the pullback map \((\Phi_X^m)^*\) on \(R\Gamma_c\) for \(E\) is the extension of scalars of the integral one. The same holds for the identification \(((\Phi_X^m)^*K_E)_{\bar y}\simeq (K_E)_{\bar y}\) at a fixed point \(\bar y\) of \(\Phi_X^m\), because the comparison for inverse images is compatible with composition and with extension of scalars. Applying \(T_{c_m}\) and taking stalks at \(\bar y\) in the formula for \((\psi_E)_m\) gives (12.2).
\end{proof}

% object: remark-integral-frobenius-correspondences-not-invertible (blueprint/chap-The-trace-formula/sec-ell-adic-coefficients-5/remark-integral-frobenius-correspondences-not-invertible/_section.md)
\phantomsection\label{bp296}%
\begin{bpstatement}{Remark~12.23}{ (Integral Frobenius correspondences need not be invertible)}
In Lemma \hyperref[bp294]{12.22}, \(K_0\) need not be a lattice in a lisse sheaf, and \(\theta\) need not be invertible even when \(\psi_E\) is. The factor \(\pi_E^{-N}\) is needed because \(\psi_E\) itself need not be integral: on a point, the automorphism \(\ell^{-1}\) of \(E\) preserves no lattice. Even with this factor, \(\theta\) cannot always be chosen invertible, already on a point. On \(E^2\), the two eigenvalues of \(\operatorname{diag}(1,\ell)\) cannot both become units after multiplying the operator by a single scalar. So Lemma \hyperref[bp294]{12.22} only makes \(\pi_E^N\psi_E\) integral, and the reductions of \(\theta\) may fail to be invertible. For this reason Theorem \hyperref[bp287]{12.20} is stated for arbitrary Frobenius correspondences.
\end{bpstatement}

% text: prose-12-22
\phantomsection\label{bp297}%
With the integral formula and integral models of rational Frobenius correspondences, we obtain the trace formula stated at the beginning of this chapter for all coefficient rings.

% object: theorem-grothendieck-lefschetz-trace-formula (blueprint/chap-The-trace-formula/sec-ell-adic-coefficients-5/theorem-grothendieck-lefschetz-trace-formula/_section.md)
\phantomsection\label{bp298}%
\begin{bpstatementit}{Theorem~12.24}{ (Grothendieck–Lefschetz trace formula)}
Let \(X_0\) be separated and of finite type over \(\mathbf F_q\), put \(X=X_0\times_{\mathbf F_q}\overline{\mathbf F}_q\), and suppose that \(\ell\nmid q\). Let

\[
\Lambda\in\{A,\mathbf Z_\ell,\mathcal O_E,\mathbf Q_\ell,E,\overline{\mathbf Q}_\ell\},
\qquad K\in D^b_{\mathrm{ctf}}(X,\Lambda),
\]

with a Weil structure as in Definition \hyperref[bp156]{10.18}. Here \(A\) is any finite commutative \(\ell\)-primary ring, and for integral and rational coefficients, \(D^b_{\mathrm{ctf}}(X,\Lambda)=\Dbc(X,\Lambda)\). Let \(\operatorname{Frob}_{X,K}\) be the geometric Frobenius on \(R\Gamma_c(X_{\mathrm{pro}\acute{e}t},K)\), and let \(\operatorname{Frob}_{y,m}\) be the Frobenius acting on the stalk \(K_{\bar y}\) at \(y\in X_0(\mathbf F_{q^m})\). Then for every \(m\geq1\) we have, in \(\Lambda\),

\[
\operatorname{tr}_{\Lambda}
\bigl(\operatorname{Frob}_{X,K}^m\mid R\Gamma_c(X_{\mathrm{pro}\acute{e}t},K)\bigr)
=\sum_{y\in X_0(\mathbf F_{q^m})}
\operatorname{tr}_{\Lambda}
\bigl(\operatorname{Frob}_{y,m}\mid K_{\bar y}\bigr).
\tag{12.3}
\]

The complexes inside the traces are perfect over \(\Lambda\), and the formula is compatible with the change-of-coefficients maps of Proposition \bpout{7.10}. For \(\Lambda=\mathbf Q_\ell,E,\overline{\mathbf Q}_\ell\), it reads

\[
\sum_i(-1)^i\operatorname{Tr}
\bigl(\operatorname{Frob}_{X,K}^m\mid H_c^i(X_{\mathrm{pro}\acute{e}t},K)\bigr)
=\sum_{y\in X_0(\mathbf F_{q^m})}\sum_j(-1)^j
\operatorname{Tr}\bigl(\operatorname{Frob}_{y,m}\mid H^j(K_{\bar y})\bigr),
\tag{12.4}
\]

where all sums over degrees are finite.
\end{bpstatementit}

% text: prose-12-23
\phantomsection\label{bp299}%
For finite and integral rings, (12.3) is an identity of categorical traces of perfect complexes, which makes sense even when the cohomology modules are not projective.

% proof: proof/theorem-grothendieck-lefschetz-trace-formula
\begin{proof}[{Proof of Theorem 12.24}]\phantomsection\label{bp300}%
In each case we turn the Weil structure into a Frobenius correspondence by composing with \(\kappa_1\), and we use Lemma \hyperref[bp234]{12.3} to see that this does not change the operators. We then treat the coefficient rings in the order finite, integral, rational, algebraic, each case using the previous one.

For finite \(A\), we apply Theorem \hyperref[bp287]{12.20} to \(\psi=\phi\kappa_1\). Then Lemma \hyperref[bp234]{12.3} identifies \(T_{c_m}(\psi_m)\) with \(\operatorname{Frob}_{X,K}^m\) and, at each fixed point, \(\psi_{m,\bar y}\) with \(\operatorname{Frob}_{y,m}\). Similarly, for \(\mathcal O_E\), composing the Weil structure with \(\kappa_1\) we get a Frobenius correspondence \(\Phi_X^*K\to K\). By the integral form of Lemma \hyperref[bp234]{12.3}, the endomorphisms it induces on \(R\Gamma_c\) and on the stalks are those induced by the Weil structure. The identity therefore follows from Proposition \hyperref[bp291]{12.21}. In particular, taking \(E=\mathbf Q_\ell\) gives the case of \(\mathbf Z_\ell\).

For coefficients \(E\), the Weil structure composed with \(\kappa_1\) is a Frobenius correspondence \(\psi_E\), and by the rational form of Lemma \hyperref[bp234]{12.3} the endomorphisms it induces on \(R\Gamma_c\) and on the stalks are those induced by the Weil structure. We apply Lemma \hyperref[bp294]{12.22} to \(\psi_E\), apply (12.1) to \((K_0,\theta)\), extend scalars to \(E\), and multiply by \(\pi_E^{-Nm}\). By linearity of categorical traces and their compatibility with extension of scalars (Proposition \hyperref[bp176]{11.2}), together with (12.2), this gives the identity for every \(m\). The case \(E=\mathbf Q_\ell\) gives the case of \(\mathbf Q_\ell\).

For \(\overline{\mathbf Q}_\ell\), complexes and morphisms are colimits of their versions over finite extensions, so it is enough to find the Weil complex over some finite \(E\) and to check that the answer does not depend on the choice. By Proposition \bpout{7.11} the complex \(K\), its Weil structure \(\phi\) and \(\phi^{-1}\) descend to a common finite extension \(E/\mathbf Q_\ell\). After a further finite extension the two identities saying that these maps are inverse to each other hold as well. Thus we get a Weil complex at a finite stage, and by extension of scalars its formula over \(E\) gives (12.3). If we choose another representative, any isomorphism compatible with the Weil structures descends to a common finite extension, so the resulting operators are conjugate. Since the trace is compatible with extension of scalars and invariant under conjugation, the result does not depend on the stage or on the representative.

The complexes are perfect by Lemma \hyperref[bp105]{10.1}. Compatibility with the change-of-coefficients maps holds because the pullbacks, the endomorphisms \(T_{c_m}\) and the iterated Frobenius correspondences commute with the comparisons of Proposition \bpout{7.10}. Finally, over a field the categorical trace is the alternating trace on cohomology by Proposition \hyperref[bp176]{11.2}, and this proves (12.4).
\end{proof}


% object: sec-The-function-sheaf-dictionary-2 (blueprint/chap-The-trace-formula/sec-The-function-sheaf-dictionary-2/_section.md), content hidden
\subsection*{\texorpdfstring{12.6\quad The function–sheaf dictionary}{12.6 The function–sheaf dictionary}}
\phantomsection\label{bp301}%
\addcontentsline{toc}{subsection}{\texorpdfstring{12.6\quad The function–sheaf dictionary}{12.6 The function–sheaf dictionary}}
% text: prose-12-24
\phantomsection\label{bp302}%
Under the trace formula, \(Rf_!\) becomes summation over fibres and \(\otimes\) becomes multiplication of functions on rational points. This is Grothendieck's function–sheaf dictionary [\hyperref[bp331]{SGA4½}, Sommes trig.] (Section \bpout{1.3}), and in Part III we compute with sheaves through their functions in this way.

% object: definition-trace-function (blueprint/chap-The-trace-formula/sec-The-function-sheaf-dictionary-2/definition-trace-function/_section.md)
\phantomsection\label{bp303}%
\begin{bpstatement}{Definition~12.25}{ (Trace function)}
Let \((K,\phi)\) be a Weil complex on \(X\) as in Theorem \hyperref[bp298]{12.24}, and let \(m\geq1\). The \emph{trace function} of \(K\) is

\[
t_{K,m}:X_0(\mathbf F_{q^m})\longrightarrow\Lambda,\qquad
t_{K,m}(x)=\operatorname{tr}_\Lambda\bigl(\operatorname{Frob}_{x,m}
\mid K_{\bar x}\bigr).
\]
\end{bpstatement}

% text: prose-12-25
\phantomsection\label{bp304}%
For \(\Lambda\) a field, \(t_{K,m}(x)\) is the alternating sum of the traces of Frobenius on the stalks \(\mathcal H^i(K)_{\bar x}\) (Proposition \hyperref[bp176]{11.2}). For example, the constant sheaf \(\Lambda_X\) with its constant Weil structure has trace function \(1\). In terms of trace functions, the theorem says that \(\operatorname{tr}_\Lambda(\operatorname{Frob}_{X,K}^m
\mid R\Gamma_c(X_{\mathrm{pro}\acute{e}t},K))\) is the sum of \(t_{K,m}\) over \(X_0(\mathbf F_{q^m})\). Applied to the fibres of a morphism, it shows that direct image with compact support sums trace functions over the fibres.

% object: corollary-trace-functions-direct-image-compact (blueprint/chap-The-trace-formula/sec-The-function-sheaf-dictionary-2/corollary-trace-functions-direct-image-compact/_section.md)
\phantomsection\label{bp305}%
\begin{bpstatementit}{Corollary~12.26}{ (Trace functions and direct image with compact support)}
Let \(f_0:X_0\to Y_0\) be a morphism of separated \(\mathbf F_q\)-schemes of finite type, with geometric base change \(f:X\to Y\), and let \((K,\phi)\) be a Weil complex on \(X\) with \(K\in D^b_{\mathrm{ctf}}(X,\Lambda)\) for an admissible \(\Lambda\). Give \(Rf_!K\in D^b_{\mathrm{ctf}}(Y,\Lambda)\) the Weil structure induced by \(\phi\) (after Definition \hyperref[bp156]{10.18}). Then for every \(m\geq1\) and \(y\in Y_0(\mathbf F_{q^m})\),

\[
t_{Rf_!K,m}(y)=\sum_{\substack{x\in X_0(\mathbf F_{q^m})\\ f_0(x)=y}}
t_{K,m}(x).
\]
\end{bpstatementit}

% proof: proof/corollary-trace-functions-direct-image-compact
\begin{proof}[{Proof}]\phantomsection\label{bp306}%
The left side is a trace on a stalk of \(Rf_!K\). By base change that stalk is the compactly supported cohomology of a fibre. Hence the corollary is the trace formula for the fibre, once we check that the Frobenius operators correspond.

Fix \(y\) and let \(X_{0,y}\) be the fibre of \(f_0\) over \(y\). It is a separated scheme of finite type over \(\mathbf F_{q^m}\), and its rational points are the \(x\in X_0(\mathbf F_{q^m})\) with \(f_0(x)=y\). Restrict \(K\) to its geometric fibre \(X_{\bar y}\) and give the restriction the Weil structure over \(\mathbf F_{q^m}\) obtained from the \(m\)-th iterate of \(\phi\). By base change for compact supports, \((Rf_!K)_{\bar y}\) is \(R\Gamma_c(X_{\bar y},K|_{X_{\bar y}})\), and \(\operatorname{Frob}_{y,m}\) is the geometric Frobenius operator of this Weil complex. For finite \(A\) this is part (4) of Lemma \hyperref[bp280]{12.18}, applied to the Frobenius correspondence \(\phi\kappa_1\) and translated by Lemma \hyperref[bp234]{12.3}. By the comparisons of Proposition \bpout{7.10} it extends to the other coefficient rings. Since the Frobenius acting on the stalk of the Weil complex on the fibre at a rational point \(x\) of \(X_{0,y}\) is \(\operatorname{Frob}_{x,m}\), the stated identity is Theorem \hyperref[bp298]{12.24} for \(X_{0,y}\) over \(\mathbf F_{q^m}\), with exponent \(1\).
\end{proof}

% object: example-trace-function-squaring (blueprint/chap-The-trace-formula/sec-The-function-sheaf-dictionary-2/example-trace-function-squaring/_section.md)
\phantomsection\label{bp307}%
\begin{bpstatement}{Example~12.27}{ (The squaring map)}
Let \(q\) be odd, let \(f_0:\mathbf A^1_{\mathbf F_q}\to\mathbf A^1_{\mathbf F_q}\) be \(x\mapsto x^2\), and let \(K=\Lambda_X\) with its constant Weil structure, so that \(t_{K,m}=1\). The fibre of \(f_0\) over \(y\in\mathbf F_{q^m}\) has \(1+\chi_m(y)\) rational points, where \(\chi_m\) is the quadratic character of \(\mathbf F_{q^m}^*\), extended by \(\chi_m(0)=0\). Hence

\[
t_{Rf_!\Lambda,m}(y)=1+\chi_m(y).
\]

The trace function of \(Rf_!\Lambda\) is not constant, although the sheaf we pushed forward is: over \(\mathbf G_m\) the map is a finite étale double cover, and \(\chi_m(y)\) records whether the fibre over \(y\) splits into two rational points or is a single point of degree \(2\). Summing over \(y\) recovers the trace formula on the source. Since \(\chi_m\) takes the values \(\pm1\) equally often on \(\mathbf F_{q^m}^*\), the sum is \(q^m=\#\mathbf A^1(\mathbf F_{q^m})\).
\end{bpstatement}

% text: prose-12-26
\phantomsection\label{bp308}%
The other operations on functions also have counterparts on sheaves. These counterparts, together with the corollary, are the entries of the dictionary that we use. In this dictionary, morphisms come from morphisms of separated \(\mathbf F_q\)-schemes of finite type, and the operations carry the Weil structures described after Definition \hyperref[bp156]{10.18}.

% object: proposition-function-sheaf-dictionary (blueprint/chap-The-trace-formula/sec-The-function-sheaf-dictionary-2/proposition-function-sheaf-dictionary/_section.md)
\phantomsection\label{bp309}%
\begin{bpstatementit}{Proposition~12.28}{ (The function–sheaf dictionary)}
Let \(\Lambda\) be as in Theorem \hyperref[bp298]{12.24}, let \(m\geq1\), and let \(K\), \(L\) be Weil complexes in \(D^b_{\mathrm{ctf}}(X,\Lambda)\).

\begin{enumerate}
\item For \(g_0:W_0\to X_0\) and \(w\in W_0(\mathbf F_{q^m})\), \(t_{g^*K,m}(w)=t_{K,m}(g_0(w))\).
\item \(t_{K\otimes^{\mathbf L}_\Lambda L,m}=t_{K,m}\,t_{L,m}\).
\item For \(f_0:X_0\to Y_0\) and \(y\in Y_0(\mathbf F_{q^m})\), \(t_{Rf_!K,m}(y)=\sum_{x\in X_0(\mathbf F_{q^m}),\,f_0(x)=y}t_{K,m}(x)\).
\item For \(U_0\subset X_0\) open with closed complement \(Z_0\), and \(j:U\to X\), \(i:Z\to X\) the inclusions, \(t_{K,m}=t_{j_!j^*K,m}+t_{i_*i^*K,m}\), where \(t_{j_!j^*K,m}\) is \(t_{K,m}\) on \(U_0(\mathbf F_{q^m})\) and \(0\) on \(Z_0(\mathbf F_{q^m})\).
\item \(\sum_{x\in X_0(\mathbf F_{q^m})}t_{K,m}(x)
  =\operatorname{tr}_\Lambda\bigl(\operatorname{Frob}_{X,K}^m\mid
  R\Gamma_c(X_{\mathrm{pro}\acute{e}t},K)\bigr)\).
\end{enumerate}
\end{bpstatementit}

% proof: proof/proposition-function-sheaf-dictionary
\begin{proof}[{Proof}]\phantomsection\label{bp310}%
Part (5) is Theorem \hyperref[bp298]{12.24} and part (3) is Corollary \hyperref[bp305]{12.26}.

For (1), the stalk of \(g^*K\) at \(\bar w\) is \(K_{\bar x}\) with \(x=g_0(w)\), and the Weil structure of \(g^*K\) is \(g^*\phi\) composed with the base change isomorphism. Under this identification of stalks the \(m\)-th iterate of the Weil structure of \(g^*K\) at \(\bar w\) is the \(m\)-th iterate of \(\phi\) at \(\bar x\), so the two Frobenius endomorphisms agree.

For (2), the stalk of \(K\otimes^{\mathbf L}_\Lambda L\) at \(\bar x\) is \(K_{\bar x}\otimes^{\mathbf L}_\Lambda L_{\bar x}\), since pullback to \(\bar x\) commutes with derived tensor products, and its Frobenius is the tensor product of the two Frobenius endomorphisms. These stalks are perfect (Lemma \hyperref[bp105]{10.1}), so we may represent them by bounded complexes \(P\), \(Q\) of finite projective modules and the endomorphisms by chain maps \(v\), \(w\). Now the term of degree \(n\) of \(P\otimes Q\) is \(\bigoplus_{i+j=n}P^i\otimes Q^j\), the trace of \(v^i\otimes w^j\) on \(P^i\otimes Q^j\) is \(\operatorname{Tr}(v^i)\operatorname{Tr}(w^j)\), and \((-1)^n=(-1)^i(-1)^j\). By the chain formula of Proposition \hyperref[bp176]{11.2} the trace on \(P\otimes Q\) is therefore the product of the traces on \(P\) and on \(Q\).

For (4), note that \(j_!=Rj_!\) and \(i_*=Ri_!\), with the same Weil structures. By (3), \(t_{j_!j^*K,m}(x)\) is the sum of \(t_{j^*K,m}\) over the fibre of \(j_0\) at \(x\). This fibre is the point \(x\) if \(x\in U_0\), and it is empty otherwise. By (1), \(t_{j^*K,m}(x)=t_{K,m}(x)\) for \(x\in U_0(\mathbf F_{q^m})\). The same argument applies to \(i_*i^*K\), and the two functions add up to \(t_{K,m}\).
\end{proof}

% text: prose-12-27
\phantomsection\label{bp311}%
The Fourier transform of Chapter \bpout{16} is one more entry, and the proposition shows why it corresponds to the discrete transform. Let \(p_1,p_2\) be the projections of \(\mathbf A^1\times\mathbf A^1\), let \(\mu(x,y)=xy\), let \(\psi:\mathbf F_q\to E^\times\) be a nontrivial additive character, and let \(\mathcal K_\psi=\mu^*\mathcal L_\psi\), where \(\mathcal L_\psi\) is the Artin–Schreier sheaf of Lemma \bpout{16.1} (proved in Chapter \bpout{16} below). Here \(\psi\) is a character, not a Frobenius correspondence. With this kernel the transform is \(\operatorname{FT}_\psi K=Rp_{2!}(p_1^*K\otimes\mathcal K_\psi)[1]\). For \(y\in\mathbf F_{q^m}\) we compute

\[
\begin{aligned}
t_{\operatorname{FT}_\psi K,m}(y)
&=-\,t_{Rp_{2!}(p_1^*K\otimes\mathcal K_\psi),m}(y)
&&\text{(shift)}\\
&=-\sum_{x\in\mathbf F_{q^m}}t_{p_1^*K\otimes\mathcal K_\psi,m}(x,y)
&&\text{by (3)}\\
&=-\sum_{x\in\mathbf F_{q^m}}t_{p_1^*K,m}(x,y)\,t_{\mu^*\mathcal L_\psi,m}(x,y)
&&\text{by (2)}\\
&=-\sum_{x\in\mathbf F_{q^m}}t_{K,m}(x)\,
\psi\bigl(\operatorname{Tr}_{\mathbf F_{q^m}/\mathbf F_q}(xy)\bigr)
&&\text{by (1) and Lemma~16.1.1.}
\end{aligned}
\]

The first equality holds because a shift changes the sign of a trace, by the chain formula of Proposition \hyperref[bp176]{11.2}. Each of the other three steps is one operation in the formula \(\hat t(y)=\sum_xt(x)\psi(xy)\): pull \(t\) back to the plane, multiply by the kernel, and sum over the fibres of \(p_2\).


% object: sec-Rationality-and-the-functional-equation (blueprint/chap-The-trace-formula/sec-Rationality-and-the-functional-equation/_section.md), content hidden
\subsection*{\texorpdfstring{12.7\quad Rationality and the functional equation}{12.7 Rationality and the functional equation}}
\phantomsection\label{bp312}%
\addcontentsline{toc}{subsection}{\texorpdfstring{12.7\quad Rationality and the functional equation}{12.7 Rationality and the functional equation}}
% text: prose-12-28
\phantomsection\label{bp313}%
Over a field of characteristic zero, the traces of all powers of an endomorphism determine its characteristic polynomial. Taking the trace formula for all \(m\geq1\), we therefore obtain an identity of power series: a product over closed points equals an alternating product of characteristic polynomials of Frobenius on cohomology. This step divides by \(m\), so in the determinant formulas below \(\Lambda\) is one of the fields \(\mathbf Q_\ell\), \(E\), \(\overline{\mathbf Q}_\ell\). Characteristic zero is needed: on \(\mathbf F_p^{\oplus p}\) every power of the identity has trace \(p=0\), as for the zero endomorphism, but \(\det(1-T\cdot\operatorname{id})=(1-T)^p=1-T^p\) is not \(1\).

For \(K=\Lambda_X\) the product over cohomology is finite, so \(Z(X_0,T)\) is rational, and for smooth proper \(X_0\) duality pairs the factors in degrees \(i\) and \(2d-i\) (Corollary \hyperref[bp325]{12.33}).

% object: lemma-logarithm-determinant (blueprint/chap-The-trace-formula/sec-Rationality-and-the-functional-equation/lemma-logarithm-determinant/_section.md)
\phantomsection\label{bp314}%
\begin{bpstatementit}{Lemma~12.29}{ (Logarithm of a determinant)}
Let \(L\) be a field of characteristic zero. For an endomorphism \(v\) of a finite-dimensional \(L\)-vector space \(V\) we have

\[
\det(1-Tv\mid V)^{-1}
=\exp\left(\sum_{m\geq1}
\operatorname{Tr}(v^m\mid V)\frac{T^m}{m}\right)
\quad\text{in }1+TL[[T]].
\tag{12.5}
\]

For an endomorphism \(v\) of a perfect complex \(P\) in \(D(L)\), put

\[
\det(1-Tv\mid P)
:=\prod_i\det(1-Tv\mid H^i(P))^{(-1)^i}.
\]

Then (12.5) holds for \(P\), with \(\operatorname{Tr}(v^m\mid V)\) replaced by the categorical trace \(\operatorname{tr}_L(v^m\mid P)\).
\end{bpstatementit}

% proof: proof/lemma-logarithm-determinant
\begin{proof}[{Proof}]\phantomsection\label{bp315}%
For a single eigenvalue \(\alpha\) the identity is \(-\log(1-\alpha T)=\sum_m\alpha^mT^m/m\), and the general case could be reduced to this by triangularizing \(v\) over an algebraic closure. Instead we use logarithmic derivatives, which allows us to stay over \(L\). In \(L[[T]]\), the matrix \(1-Tv\) is invertible, with inverse \(\sum_{r\geq0}T^rv^r\). Moreover, the logarithmic derivative of the determinant of an invertible matrix \(M\) over \(L[[T]]\) is \(\operatorname{Tr}(M^{-1}M')\), and hence

\[
\frac{d}{dT}\log\det(1-Tv)
=-\operatorname{Tr}\bigl(v(1-Tv)^{-1}\bigr)
=-\sum_{m\geq1}\operatorname{Tr}(v^m)T^{m-1}.
\]

Since both logarithms have constant coefficient zero and \(L\) has characteristic zero, we may integrate and exponentiate, and we obtain (12.5). For a perfect complex, the categorical trace over a field is the alternating sum of the traces on the finitely many nonzero cohomology spaces (Proposition \hyperref[bp176]{11.2}), and the assertion follows by multiplying the identities for the \(H^i(P)\).
\end{proof}

% text: prose-12-29
\phantomsection\label{bp316}%
The local factors of the Euler product are indexed by closed points. A closed point \(x\in|X_0|\) of degree \(d\) over \(\mathbf F_q\) corresponds to a Frobenius orbit of \(d\) geometric points, and its Frobenius on the stalk is the Frobenius going once around this orbit.

% object: theorem-trace-formula-euler-product (blueprint/chap-The-trace-formula/sec-Rationality-and-the-functional-equation/theorem-trace-formula-euler-product/_section.md)
\phantomsection\label{bp317}%
\begin{bpstatementit}{Theorem~12.30}{ (The trace formula as an Euler product)}
Let \(X_0\) be separated and of finite type over \(\mathbf F_q\), let \(\Lambda\in\{\mathbf Q_\ell,E,\overline{\mathbf Q}_\ell\}\), and let \(K\in\Dbc(X,\Lambda)\) carry a Weil structure, with stalk Frobenius maps \(\operatorname{Frob}_{y,m}\) as in Theorem \hyperref[bp298]{12.24}. For each closed point \(x\in|X_0|\), let \(\deg(x)=[\kappa(x):\mathbf F_q]\) be its degree, choose one of the \(\deg(x)\) points \(y\in X_0(\mathbf F_{q^{\deg(x)}})\) lying over \(x\), write \(\bar x=\bar y\) for the corresponding geometric point, and let \(\operatorname{Frob}_x=\operatorname{Frob}_{y,\deg(x)}\) be its action on \(K_{\bar x}\). Another choice of \(y\) replaces \(\operatorname{Frob}_x\) by a conjugate. Define

\[
L(X_0,K,T)
:=\prod_{x\in|X_0|}
\det\bigl(1-T^{\deg(x)}\operatorname{Frob}_x\mid K_{\bar x}\bigr)^{-1}.
\tag{12.6}
\]

For every \(n\), only finitely many factors differ from \(1\) modulo \(T^n\), so the product converges \(T\)-adically in \(1+T\Lambda[[T]]\), and

\[
L(X_0,K,T)
=\prod_i\det\bigl(1-T\operatorname{Frob}_{X,K}\mid H_c^i(X_{\mathrm{pro}\acute{e}t},K)\bigr)^{(-1)^{i+1}}.
\tag{12.7}
\]

In particular, \(L(X_0,K,T)\) is the expansion at zero of a rational function in \(\Lambda(T)\) with value \(1\) at zero. If \(X_0\) is proper, we may replace \(H_c^i\) by \(H^i\).
\end{bpstatementit}

% proof: proof/theorem-trace-formula-euler-product
\begin{proof}[{Proof}]\phantomsection\label{bp318}%
Both sides are of the form \(\exp\) of a power series, and we compare the logarithms coefficient by coefficient. The coefficient of \(T^m\) on the cohomological side is the global trace of \(\operatorname{Frob}^m\) divided by \(m\), and on the side of closed points it is the sum of the local traces over \(X_0(\mathbf F_{q^m})\) divided by \(m\). Hence the theorem is the trace formula for all \(m\) at once.

For convergence, note that \(X_0\) has only finitely many closed points of degree at most any given bound \(n\), because each of them lies under a point of the finite set \(X_0(\mathbf F_{q^{n!}})\). The factor at \(x\) in (12.6) is \(1\) modulo \(T^{\deg(x)}\), so the product converges.

Before comparing logarithms we rewrite the sum over points in Theorem \hyperref[bp298]{12.24} in terms of closed points. A closed point of degree \(d\) has exactly \(d\) points over \(\mathbf F_{q^m}\) when \(d\mid m\), and none otherwise. For instance, \(\operatorname{Spec}\mathbf F_{q^2}\) has one closed point, of degree \(2\), with no \(\mathbf F_q\)-points and two \(\mathbf F_{q^2}\)-points. These \(d\) points form one Frobenius orbit, and the Weil structure gives isomorphisms between the stalks at consecutive points of the orbit. The Frobenius on the stalk at a point of the orbit is then the composite of these isomorphisms taken once around the orbit, starting at that point. Moving the starting point cyclically therefore conjugates this composite by one of these isomorphisms. Hence the operators at the \(\deg(x)\) points are conjugate. In the same way, iterating the Weil structure shows that for \(\deg(x)\mid m\) the operator \(\operatorname{Frob}_{y,m}\) at each of the \(\deg(x)\) points \(y\in X_0(\mathbf F_{q^m})\) over \(x\) is conjugate to \(\operatorname{Frob}_x^{m/\deg(x)}\). Since traces are invariant under conjugation, (12.3) reads

\[
\operatorname{tr}_{\Lambda}(\operatorname{Frob}_{X,K}^m\mid R\Gamma_c(X_{\mathrm{pro}\acute{e}t},K))
=\sum_{\substack{x\in|X_0|\\\deg(x)\mid m}}\deg(x)\,
\operatorname{tr}_{\Lambda}
(\operatorname{Frob}_x^{m/\deg(x)}\mid K_{\bar x}).
\tag{12.8}
\]

We multiply (12.8) by \(T^m/m\) and sum over \(m\geq1\). By Lemma \hyperref[bp314]{12.29}, the left side becomes the logarithm of the right side of (12.7). In the sum over closed points on the right, we substitute \(m=r\deg(x)\). Then the multiplicity \(\deg(x)\) cancels the factor \(\deg(x)\) in \(m\), and Lemma \hyperref[bp314]{12.29} applied to \(\operatorname{Frob}_x\) on \(K_{\bar x}\), with \(T^{\deg(x)}\) in place of \(T\), shows that

\[
\sum_{x\in|X_0|}\sum_{r\geq1}
\operatorname{tr}_{\Lambda}(\operatorname{Frob}_x^r\mid K_{\bar x})
\frac{T^{r\deg(x)}}{r}
=\log L(X_0,K,T).
\]

Each coefficient of \(T\) receives only finitely many contributions, so we may regroup the sums. Exponentiating proves (12.7). Rationality holds because the cohomological product has finitely many factors, each a polynomial or the reciprocal of one, with value \(1\) at zero. If \(X_0\) is proper, compact and ordinary cohomology agree.
\end{proof}

% text: prose-12-30
\phantomsection\label{bp319}%
For the constant sheaf every local factor is \((1-T^{\deg(x)})^{-1}\), and the Euler product becomes the zeta function of \(X_0\). This proves rationality of \(Z(X_0,T)\) over \(\mathbf Q_\ell\). A short argument with linear recurrences descends it to \(\mathbf Q\).

% object: corollary-point-counting-zeta-function (blueprint/chap-The-trace-formula/sec-Rationality-and-the-functional-equation/corollary-point-counting-zeta-function/_section.md)
\phantomsection\label{bp320}%
\begin{bpstatementit}{Corollary~12.31}{ (Point counting and the zeta function)}
Give \(\Lambda_X\) the Weil structure of constant sheaves defined after Definition \hyperref[bp156]{10.18}. For every coefficient ring \(\Lambda\) of Theorem \hyperref[bp298]{12.24} and every \(m\geq1\) we have

\[
\operatorname{tr}_{\Lambda}
(\operatorname{Frob}_{X,\Lambda_X}^m\mid R\Gamma_c(X_{\mathrm{pro}\acute{e}t},\Lambda_X))
=\#X_0(\mathbf F_{q^m})\cdot1_{\Lambda}.
\tag{12.9}
\]

For \(\Lambda=\mathbf Q_\ell,E,\overline{\mathbf Q}_\ell\), this gives

\[
\#X_0(\mathbf F_{q^m})
=\sum_i(-1)^i\operatorname{Tr}
(\operatorname{Frob}_{X,\Lambda_X}^m\mid H_c^i(X_{\mathrm{pro}\acute{e}t},\Lambda_X)),
\]

and

\[
\begin{aligned}
Z(X_0,T)
&:=\exp\left(\sum_{m\geq1}
\#X_0(\mathbf F_{q^m})\frac{T^m}{m}\right)\\
&=\prod_{x\in|X_0|}(1-T^{\deg(x)})^{-1}
=\prod_i\det(1-T\operatorname{Frob}_{X,\Lambda_X}\mid H_c^i(X_{\mathrm{pro}\acute{e}t},\Lambda_X))^{(-1)^{i+1}}.
\end{aligned}
\]

The first two series lie in \(1+T\mathbf Z[[T]]\), and the last equality holds in \(\Lambda[[T]]\). Moreover, \(Z(X_0,T)\) belongs to \(\mathbf Q(T)\).
\end{bpstatementit}

% proof: proof/corollary-point-counting-zeta-function
\begin{proof}[{Proof}]\phantomsection\label{bp321}%
The identities over \(\Lambda\) are the trace formula and the Euler product for the constant sheaf. What needs an argument is that \(Z(X_0,T)\) has rational coefficients, which we deduce at the end from the rationality over \(\mathbf Q_\ell\) by linear algebra.

For \(K=\Lambda_X\), every stalk is \(\Lambda\), and every stalk Frobenius \(\operatorname{Frob}_{y,m}\) is its identity, with categorical trace \(1_\Lambda\). Thus (12.9) follows from Theorem \hyperref[bp298]{12.24}, and the formula over the fields follows from its cohomological form (12.4). By the orbit count in the proof of Theorem \hyperref[bp317]{12.30}, we have \(\#X_0(\mathbf F_{q^m})=\sum_{\deg(x)\mid m}\deg(x)\). Expanding each \((1-T^{\deg(x)})^{-1}\) as a geometric series shows that the product over closed points has integer coefficients. By the orbit count, its logarithm is the series defining \(Z(X_0,T)\). The last equality is Theorem \hyperref[bp317]{12.30} for \(K=\Lambda_X\).

Rationality descends from \(\mathbf Q_\ell\) to \(\mathbf Q\) for the following reason. A power series is rational exactly when its coefficients eventually satisfy a linear recurrence with constant coefficients. The existence of such a recurrence is a system of linear equations, and a system with rational coefficients has a rational solution as soon as it has any solution. More precisely, write \(Z(X_0,T)=\sum_{n\geq0}z_nT^n\), with \(z_n\in\mathbf Z\). Rationality over \(\mathbf Q_\ell\) means that there are integers \(d,N\) and \(b_1,\ldots,b_d\in\mathbf Q_\ell\) such that \(z_n+b_1z_{n-1}+\cdots+b_dz_{n-d}=0\) for all \(n>N\). These are linear equations with rational coefficients in the \(d\) unknowns \(b_i\). Their coefficient rows span a finite-dimensional subspace of \(\mathbf Q^{d+1}\), and hence finitely many of the equations imply all the others. This finite rational linear system has a solution over \(\mathbf Q_\ell\), and hence one over \(\mathbf Q\) by Gaussian elimination. With these rational \(b_i\), the product of \(1+b_1T+\cdots+b_dT^d\in\mathbf Q[T]\) and \(Z(X_0,T)\) is a polynomial. Therefore \(Z(X_0,T)\in\mathbf Q(T)\).
\end{proof}

% text: prose-12-31
\phantomsection\label{bp322}%
For example, let \(X_0=\mathbf A^1_{\mathbf F_q}\). Then \(\#X_0(\mathbf F_{q^m})=q^m\), so that \(Z(X_0,T)=(1-qT)^{-1}\). On the cohomological side, \(H^i_c(X,\mathbf Q_\ell)\) vanishes except for \(i=2\), and \(H^2_c(X,\mathbf Q_\ell)\simeq\mathbf Q_\ell(-1)\) is a line on which geometric Frobenius acts by \(q\). The only factor is therefore \((1-qT)^{-1}\), as it should be.

For \(X_0=\mathbf P^1_{\mathbf F_q}\) the cohomology is \(\mathbf Q_\ell\) in degree \(0\) and \(\mathbf Q_\ell(-1)\) in degree \(2\), so \(Z(\mathbf P^1,T)=1/((1-T)(1-qT))\). This matches \(\#\mathbf P^1(\mathbf F_{q^m})=1+q^m\). For a curve of positive genus the middle cohomology contributes as well, and for the running example a single point count determines it.

% object: example-running-example-trace-formula (blueprint/chap-The-trace-formula/sec-Rationality-and-the-functional-equation/example-running-example-trace-formula/_section.md)
\phantomsection\label{bp323}%
\begin{bpstatement}{Example~12.32}{ (The running example: the trace formula)}
For the elliptic curve \(C_0\colon y^2=x^3-x\) over \(\mathbf F_3\), geometric Frobenius acts on \(H^0(C,\mathbf Q_\ell)=\mathbf Q_\ell\) by \(1\) and on \(H^2(C,\mathbf Q_\ell)=\mathbf Q_\ell(-1)\) by \(3\), and \(H^1(C,\mathbf Q_\ell)\) is a plane (Example \hyperref[bp169]{10.23}). The point-counting formula reads

\[
N_m=1-\operatorname{Tr}\bigl(F^m\mid H^1(C,\mathbf Q_\ell)\bigr)+3^m .
\]

From \(N_1=4\) we get \(\operatorname{Tr}(F\mid H^1)=0\). Duality gives \(\det(F\mid H^1)=3\), so \(\det(1-TF\mid H^1)=1+3T^2\). Hence

\[
Z(C_0,T)=\frac{1+3T^2}{(1-T)(1-3T)}.
\]

The formula predicts \(N_2=1+9+6=16\) and \(N_3=1+27-0=28\), and direct counts over \(\mathbf F_9\) and \(\mathbf F_{27}\) confirm this. In general, for a genus-one curve over \(\mathbf F_q\) the single number \(N_1\) determines \(P_1=1-(1+q-N_1)T+qT^2\), and therefore every \(N_m\).
\end{bpstatement}

% text: prose-12-32
\phantomsection\label{bp324}%
For smooth proper \(X_0\), Poincaré duality relates the factors of the determinant formula in complementary degrees, and this gives a functional equation.

% object: corollary-functional-equation (blueprint/chap-The-trace-formula/sec-Rationality-and-the-functional-equation/corollary-functional-equation/_section.md)
\phantomsection\label{bp325}%
\begin{bpstatementit}{Corollary~12.33}{ (Functional equation)}
Let \(X_0/\mathbf F_q\) be smooth and proper of pure dimension \(d\), and let \(\ell\nmid q\). Put \(H^i=H^i(X_{\mathrm{pro}\acute{e}t},\mathbf Q_\ell)\), let \(F_i\) be the Frobenius operator \(F\) on \(H^i\) (Section \hyperref[bp154]{10.5}), and put

\[
b_i=\dim_{\mathbf Q_\ell}H^i,\qquad \chi=\sum_i(-1)^ib_i,\qquad
D=\prod_i\det(F_i)^{(-1)^i}.
\]

Then \(b_{2d-i}=b_i\), \(d\chi\) is even, \(D^2=q^{d\chi}\), and in \(\mathbf Q_\ell(T)\)

\[
Z\!\left(X_0,\frac1{q^dT}\right)=(-1)^{\chi}D\,T^{\chi}\,Z(X_0,T)
=\varepsilon q^{d\chi/2}T^{\chi}Z(X_0,T),
\qquad \varepsilon=(-1)^\chi Dq^{-d\chi/2}\in\{1,-1\}.
\]
\end{bpstatementit}

% proof: proof/corollary-functional-equation
\begin{proof}[{Proof}]\phantomsection\label{bp326}%
Duality gives the relations \(b_{2d-i}=b_i\) and \(D^2=q^{d\chi}\). If we substitute \(1/(q^dT)\) into each factor of the cohomological expression for \(Z(X_0,T)\), duality turns the factor in degree \(i\) into the factor in degree \(2d-i\) times an explicit scalar.

We apply Theorem \hyperref[bp165]{10.21} with \(K=\mathbf Q_\ell\) and \(X\) proper. It gives \(b_{2d-i}=b_i\), and it shows that the eigenvalues of \(F_{2d-i}\) are the \(q^d/\alpha\) for the eigenvalues \(\alpha\) of \(F_i\), with multiplicities. Hence \(\det F_i\cdot\det F_{2d-i}=q^{db_i}\). Since \((-1)^i=(-1)^{2d-i}\), we get

\[
D^2=\prod_i\bigl(\det F_i\cdot\det F_{2d-i}\bigr)^{(-1)^i}
=q^{d\chi}.
\]

Since \(b_{2d-i}=b_i\), the terms of \(\chi\) for \(i\) and \(2d-i\) are equal and cancel modulo \(2\), so \(\chi\equiv b_d\pmod 2\). If \(d\) is even, \(d\chi\) is even anyway. If \(d\) is odd, the pairing on \(H^d\) of each connected component of \(X\) is perfect and alternating (Corollary \hyperref[bp152]{10.17}), so \(b_d\) is even. In both cases \(d\chi\) is even. Then \(\varepsilon^2=1\) because \(D^2=q^{d\chi}\).

Now consider a single factor. Since \(F_i\) is invertible, its eigenvalues \(\alpha\) are nonzero. Pulling \(-\alpha/(q^dT)\) out of each linear factor gives

\[
\det\bigl(1-q^{-d}T^{-1}F_i\bigr)
=\prod_\alpha\Bigl(-\frac{\alpha}{q^dT}\Bigr)\Bigl(1-\frac{q^d}{\alpha}T\Bigr)
=(-1)^{b_i}\frac{\det F_i}{(q^dT)^{b_i}}\det\bigl(1-TF_{2d-i}\bigr).
\]

By Corollary \hyperref[bp320]{12.31}, and since \(H^i_c=H^i\) as \(X_0\) is proper, \(Z(X_0,T)\) is the product of the \(\det(1-TF_i)^{(-1)^{i+1}}\). We substitute \(1/(q^dT)\) for \(T\) and apply the last display to each factor. Since \((-1)^i=(-1)^{2d-i}\), the factors \(\det(1-TF_{2d-i})^{(-1)^{i+1}}\) multiply to \(Z(X_0,T)\), and the scalar factors multiply to

\[
(-1)^{\chi}D^{-1}(q^dT)^{\chi}=(-1)^{\chi}D\,T^{\chi},
\]

by \(D^2=q^{d\chi}\). The second form of the equation is the definition of \(\varepsilon\).
\end{proof}

% text: prose-12-33
\phantomsection\label{bp327}%
For the elliptic curve of Example \hyperref[bp323]{12.32} we have \(d=1\) and \(b_0=b_2=1\), \(b_1=2\). So \(\chi=0\) and \(D=\det(F_1)^{-1}\det(F_2)=3/3=1\). The corollary says \(Z(C_0,1/(3T))=Z(C_0,T)\), and indeed

\[
\frac{1+3(3T)^{-2}}{(1-(3T)^{-1})(1-T^{-1})}
=\frac{(3T^2+1)/(3T^2)}{(3T-1)(T-1)/(3T^2)}
=\frac{1+3T^2}{(1-3T)(1-T)}.
\]

% object: example-functional-equation-projective-line (blueprint/chap-The-trace-formula/sec-Rationality-and-the-functional-equation/example-functional-equation-projective-line/_section.md)
\phantomsection\label{bp328}%
\begin{bpstatement}{Example~12.34}{ (The functional equation of the projective line)}
For \(X_0=\mathbf P^1_{\mathbf F_q}\) we have \(d=1\), \(b_0=b_2=1\), \(b_1=0\), so \(\chi=2\), and \(D=\det(F_0)\det(F_2)=1\cdot q=q\), so that \(\varepsilon=1\). The corollary predicts \(Z(\mathbf P^1,1/(qT))=qT^2Z(\mathbf P^1,T)\). Directly,

\[
\frac1{(1-(qT)^{-1})(1-T^{-1})}
=\frac{qT^2}{(qT-1)(T-1)}
=\frac{qT^2}{(1-T)(1-qT)}.
\]

Here the factor \(T^{\chi}\) is not trivial, unlike for the elliptic curve, and the constant \(q^{d\chi/2}=q\) comes from the twist \((-1)\) on \(H^2\).
\end{bpstatement}

% object: remark-duality-and-weights (blueprint/chap-The-trace-formula/sec-Rationality-and-the-functional-equation/remark-duality-and-weights/_section.md)
\phantomsection\label{bp329}%
\begin{bpstatement}{Remark~12.35}{ (Duality and weights)}
The functional equation is an identity of rational functions, not of power series at \(T=0\), since \(1/(q^dT)\) is not a power series in \(T\).

What duality does not give is the absolute value of the eigenvalues \(\alpha\). In Chapters \bpout{13}–\bpout{Chapter 19} we prove that the eigenvalues of \(F_i\) have absolute value \(q^{i/2}\) under every embedding of \(\overline{\mathbf Q}_\ell\) into \(\mathbf C\). It follows that the factors \(\det(1-TF_i)\) for different \(i\) have no common zeros, and hence we can read them off from \(Z(X_0,T)\) alone. In Chapter \bpout{19} we use this to show that they are polynomials with integer coefficients that do not depend on \(\ell\).
\end{bpstatement}




\clearpage\phantomsection\hypertarget{bp-outside}{}%
\section*{Outside this export}
\addcontentsline{toc}{section}{Outside this export}
References marked \ensuremath{{}^\ast} are to parts of the blueprint this document leaves out:
\begin{description}
\item[Section 1.3] The function–sheaf dictionary, in 1 Introduction (\bpslug{sec-{}\allowbreak{}The-{}\allowbreak{}function-{}\allowbreak{}sheaf-{}\allowbreak{}dictionary})
\item[Corollary 2.29] Continuity of cohomology, in Ⅰ The six operations (\bpslug{corollary-{}\allowbreak{}continuity-{}\allowbreak{}cohomology})
\item[Theorem 2.38] Comparison of direct images, in Ⅰ The six operations (\bpslug{theorem-{}\allowbreak{}comparison-{}\allowbreak{}direct-{}\allowbreak{}images})
\item[Corollary 2.44] Finite Tor amplitude under comparison, in Ⅰ The six operations (\bpslug{corollary-{}\allowbreak{}finite-{}\allowbreak{}tor-{}\allowbreak{}amplitude-{}\allowbreak{}under-{}\allowbreak{}comparison})
\item[Definition 2.45] Cohomological dimension, in Ⅰ The six operations (\bpslug{definition-{}\allowbreak{}cohomological-{}\allowbreak{}dimension})
\item[Definition 2.53] Tate twists, in Ⅰ The six operations (\bpslug{definition-{}\allowbreak{}tate-{}\allowbreak{}twists})
\item[Section 2.13] Continuous cohomology, in Ⅰ The six operations (\bpslug{sec-{}\allowbreak{}Continuous-{}\allowbreak{}cohomology})
\item[Definition 2.61] Derived adic completeness, in Ⅰ The six operations (\bpslug{definition-{}\allowbreak{}derived-{}\allowbreak{}adic-{}\allowbreak{}completeness})
\item[Lemma 2.62] Reduction and derived Nakayama, in Ⅰ The six operations (\bpslug{lemma-{}\allowbreak{}reduction-{}\allowbreak{}derived-{}\allowbreak{}nakayama})
\item[Section 3.2] Constructibility, in Ⅰ The six operations (\bpslug{sec-{}\allowbreak{}Constructibility})
\item[Definition 3.1] Constructible sheaves, in Ⅰ The six operations (\bpslug{definition-{}\allowbreak{}constructible-{}\allowbreak{}sheaves})
\item[Proposition 3.11] The integral constructibility criterion, in Ⅰ The six operations (\bpslug{proposition-{}\allowbreak{}integral-{}\allowbreak{}constructibility-{}\allowbreak{}criterion})
\item[Lemma 3.13] Flat models and generators, in Ⅰ The six operations (\bpslug{lemma-{}\allowbreak{}flat-{}\allowbreak{}models-{}\allowbreak{}generators})
\item[Proposition 3.27] Inverting $\ell$, in Ⅰ The six operations (\bpslug{proposition-{}\allowbreak{}inverting-{}\allowbreak{}ell})
\item[Theorem 5.2] The exceptional inverse image, in Ⅰ The six operations (\bpslug{theorem-{}\allowbreak{}exceptional-{}\allowbreak{}inverse-{}\allowbreak{}image})
\item[Proposition 5.15] Comparison with the étale derived category, in Ⅰ The six operations (\bpslug{proposition-{}\allowbreak{}comparison-{}\allowbreak{}etale-{}\allowbreak{}derived-{}\allowbreak{}category})
\item[Lemma 5.19] Structure of integral constructible complexes, in Ⅰ The six operations (\bpslug{lemma-{}\allowbreak{}structure-{}\allowbreak{}integral-{}\allowbreak{}constructible-{}\allowbreak{}complexes})
\item[Proposition 5.28] Composition and base change for $Rf_!$, in Ⅰ The six operations (\bpslug{proposition-{}\allowbreak{}composition-{}\allowbreak{}base-{}\allowbreak{}change-{}\allowbreak{}rf-{}\allowbreak{}shriek})
\item[Theorem 5.31] Existence of the exceptional inverse image, in Ⅰ The six operations (\bpslug{theorem-{}\allowbreak{}existence-{}\allowbreak{}exceptional-{}\allowbreak{}inverse-{}\allowbreak{}image})
\item[Section 6.1] Conventions, in Ⅰ The six operations (\bpslug{sec-{}\allowbreak{}Conventions-{}\allowbreak{}2})
\item[Lemma 6.3] The cycle class of a Cartier divisor, in Ⅰ The six operations (\bpslug{lemma-{}\allowbreak{}cycle-{}\allowbreak{}class-{}\allowbreak{}cartier-{}\allowbreak{}divisor})
\item[Lemma 6.4] Cohomology of the projective line and the trace, in Ⅰ The six operations (\bpslug{lemma-{}\allowbreak{}cohomology-{}\allowbreak{}projective-{}\allowbreak{}line-{}\allowbreak{}trace})
\item[Proposition 6.7] The adjunction for the projective line, in Ⅰ The six operations (\bpslug{proposition-{}\allowbreak{}adjunction-{}\allowbreak{}projective-{}\allowbreak{}line})
\item[Proposition 6.9] The exceptional inverse image of a smooth morphism, in Ⅰ The six operations (\bpslug{proposition-{}\allowbreak{}exceptional-{}\allowbreak{}inverse-{}\allowbreak{}image-{}\allowbreak{}smooth-{}\allowbreak{}morphism})
\item[Theorem 6.20] $f^!$ with finite coefficients, in Ⅰ The six operations (\bpslug{theorem-{}\allowbreak{}f-{}\allowbreak{}shriek-{}\allowbreak{}finite-{}\allowbreak{}coefficients})
\item[Proposition 6.22] The $\ell$-adic cycle class of the diagonal, in Ⅰ The six operations (\bpslug{proposition-{}\allowbreak{}ell-{}\allowbreak{}adic-{}\allowbreak{}cycle-{}\allowbreak{}class-{}\allowbreak{}diagonal})
\item[Proposition 6.23] Projection formula with rational coefficients, in Ⅰ The six operations (\bpslug{proposition-{}\allowbreak{}projection-{}\allowbreak{}formula-{}\allowbreak{}rational-{}\allowbreak{}coefficients})
\item[Example 6.25] Local cohomology at the origin of the affine line, in Ⅰ The six operations (\bpslug{example-{}\allowbreak{}closed-{}\allowbreak{}support})
\item[Lemma 6.26] Projection formula with integral coefficients, in Ⅰ The six operations (\bpslug{lemma-{}\allowbreak{}projection-{}\allowbreak{}formula-{}\allowbreak{}integral-{}\allowbreak{}coefficients})
\item[Proposition 7.1] The six operations on $\Dctf$, in Ⅰ The six operations (\bpslug{proposition-{}\allowbreak{}six-{}\allowbreak{}operations-{}\allowbreak{}dctf})
\item[Section 7.2] Change of coefficients, in Ⅰ The six operations (\bpslug{sec-{}\allowbreak{}Change-{}\allowbreak{}of-{}\allowbreak{}coefficients})
\item[Definition 7.4] Change-of-coefficients maps, in Ⅰ The six operations (\bpslug{definition-{}\allowbreak{}change-{}\allowbreak{}of-{}\allowbreak{}coefficients-{}\allowbreak{}maps})
\item[Lemma 7.6] Change of coefficients for direct image, in Ⅰ The six operations (\bpslug{lemma-{}\allowbreak{}change-{}\allowbreak{}coefficients-{}\allowbreak{}direct-{}\allowbreak{}image})
\item[Proposition 7.10] The six operations and change of coefficients, in Ⅰ The six operations (\bpslug{proposition-{}\allowbreak{}six-{}\allowbreak{}operations-{}\allowbreak{}change-{}\allowbreak{}coefficients})
\item[Proposition 7.11] $\Dctf(X,\Qbar)$ as a $2$-colimit, in Ⅰ The six operations (\bpslug{proposition-{}\allowbreak{}dctf-{}\allowbreak{}x-{}\allowbreak{}qbar-{}\allowbreak{}2-{}\allowbreak{}colimit})
\item[Section 7.3] Reduction to finite coefficients, in Ⅰ The six operations (\bpslug{sec-{}\allowbreak{}Reduction-{}\allowbreak{}to-{}\allowbreak{}finite-{}\allowbreak{}coefficients})
\item[Proposition 7.12] Isomorphisms reduce to finite coefficients, in Ⅰ The six operations (\bpslug{proposition-{}\allowbreak{}isomorphisms-{}\allowbreak{}reduce-{}\allowbreak{}finite-{}\allowbreak{}coefficients})
\item[Section 8.9] The Künneth formula, in Ⅰ The six operations (\bpslug{sec-{}\allowbreak{}The-{}\allowbreak{}Kunneth-{}\allowbreak{}formula})
\item[Proposition 8.18] Ordinary external Künneth, in Ⅰ The six operations (\bpslug{proposition-{}\allowbreak{}ordinary-{}\allowbreak{}external-{}\allowbreak{}kunneth})
\item[Chapter 13] Weights, in Ⅲ Weights and the Riemann hypothesis (\bpslug{chap-{}\allowbreak{}Weights})
\item[Proposition 13.7] Weights under the usual operations, in Ⅲ Weights and the Riemann hypothesis (\bpslug{proposition-{}\allowbreak{}weights-{}\allowbreak{}under-{}\allowbreak{}usual-{}\allowbreak{}operations})
\item[Proposition 14.2] Kummer, duality and the integral Tate module, in Ⅲ Weights and the Riemann hypothesis (\bpslug{proposition-{}\allowbreak{}kummer-{}\allowbreak{}duality-{}\allowbreak{}integral-{}\allowbreak{}tate-{}\allowbreak{}module})
\item[Chapter 16] The Fourier transform, in Ⅲ Weights and the Riemann hypothesis (\bpslug{chap-{}\allowbreak{}The-{}\allowbreak{}Fourier-{}\allowbreak{}transform})
\item[Lemma 16.1] The Artin–Schreier character sheaf, in Ⅲ Weights and the Riemann hypothesis (\bpslug{lemma-{}\allowbreak{}artin-{}\allowbreak{}schreier-{}\allowbreak{}character-{}\allowbreak{}sheaf})
\item[Chapter 19] Proof of the Weil conjectures, in Ⅲ Weights and the Riemann hypothesis (\bpslug{chap-{}\allowbreak{}Proof-{}\allowbreak{}of-{}\allowbreak{}the-{}\allowbreak{}Weil-{}\allowbreak{}conjectures})
\item[Theorem A.2] Nagata's compactification theorem, in Appendices (\bpslug{theorem-{}\allowbreak{}nagata-{}\allowbreak{}s-{}\allowbreak{}compactification-{}\allowbreak{}theorem})
\item[Corollary C.15] Torsion of the Jacobian and field extension, in Appendices (\bpslug{corollary-{}\allowbreak{}torsion-{}\allowbreak{}jacobian-{}\allowbreak{}field-{}\allowbreak{}extension})
\item[Appendix D] Cohomology of torsion sheaves on curves, in Appendices (\bpslug{chap-{}\allowbreak{}Cohomology-{}\allowbreak{}of-{}\allowbreak{}torsion-{}\allowbreak{}sheaves-{}\allowbreak{}on-{}\allowbreak{}curves})
\item[Proposition D.4] Cohomology of a smooth affine curve, in Appendices (\bpslug{proposition-{}\allowbreak{}cohomology-{}\allowbreak{}smooth-{}\allowbreak{}affine-{}\allowbreak{}curve})
\end{description}
\clearpage\phantomsection%
\section*{References}
\addcontentsline{toc}{section}{References}
\bpbibitem{BBD82}{bp333}{A. Beilinson, J. Bernstein, and P. Deligne. \emph{Faisceaux pervers}, volume 100 of \emph{Astérisque}. 1982.}

\bpbibitem{BS15}{bp334}{B. Bhatt and P. Scholze. The pro-étale topology for schemes. \emph{Astérisque}, 369:99–201, 2015. \href{https://people.mpim-bonn.mpg.de/scholze/proetale.pdf}{Author's manuscript}.}

\bpbibitem{Con07}{bp335}{B. Conrad. Deligne's notes on Nagata compactifications. \emph{Journal of the Ramanujan Mathematical Society}, 22:205–257, 2007. \href{https://math.stanford.edu/\~conrad/papers/nagatafinal.pdf}{Author's manuscript}.}

\bpbibitem{Del74}{bp336}{P. Deligne. La conjecture de Weil. I. \emph{Publ. Math. IHÉS}, 43:273–307, 1974.}

\bpbibitem{Del80}{bp337}{P. Deligne. La conjecture de Weil. II. \emph{Publ. Math. IHÉS}, 52:137–252, 1980.}

\bpbibitem{Den22}{bp338}{C. Deninger. There is no “Weil-”cohomology theory with real coefficients for arithmetic curves. \emph{Ann. Sc. Norm. Super. Pisa Cl. Sci.}, 2022. To appear. arXiv:\href{https://arxiv.org/abs/2204.02714}{2204.02714}.}

\bpbibitem{Deu41}{bp339}{M. Deuring. Die Typen der Multiplikatorenringe elliptischer Funktionenkörper. \emph{Abh. Math. Sem. Univ. Hamburg}, 14:197–272, 1941.}

\bpbibitem{Eke90}{bp340}{T. Ekedahl. On the adic formalism. In \emph{The Grothendieck Festschrift, Vol. II}, volume 87 of \emph{Progress in Mathematics}, pages 197–218. Birkhäuser, 1990.}

\bpbibitem{Has36}{bp341}{H. Hasse. Zur Theorie der abstrakten elliptischen Funktionenkörper. I, II, III. \emph{J. Reine Angew. Math.}, 175:55–62, 69–88, 193–208, 1936.}

\bpbibitem{Hum72}{bp342}{J. E. Humphreys. \emph{Introduction to Lie Algebras and Representation Theory}, volume 9 of \emph{Graduate Texts in Mathematics}. Springer, 1972.}

\bpbibitem{Ill94}{bp343}{L. Illusie. Autour du théorème de monodromie locale. \emph{Astérisque}, 223:9–57, 1994. §4. \href{https://www.numdam.org/item/AST_1994__223__9_0/}{Original article}.}

\bpbibitem{Kat88}{bp344}{N. M. Katz. Travaux de Laumon. \emph{Astérisque}, 161–162:105–132, 1988. Séminaire Bourbaki, exp. 691.}

\bpbibitem{KW01}{bp345}{R. Kiehl and R. Weissauer. \emph{Weil conjectures, perverse sheaves and \(\ell\)-adic Fourier transform}, volume 42 of \emph{Ergebnisse der Mathematik}. Springer, 2001.}

\bpbibitem{Lau87}{bp346}{G. Laumon. Transformation de Fourier, constantes d'équations fonctionnelles et conjecture de Weil. \emph{Publ. Math. IHÉS}, 65:131–210, 1987.}

\bpbibitem{Mil86a}{bp347}{J. S. Milne. \href{https://www.jmilne.org/math/xnotes/AVs.pdf}{Abelian Varieties}. In \emph{Arithmetic Geometry}. 1986. Corrected online version.}

\bpbibitem{Mil86b}{bp348}{J. S. Milne. \href{https://www.jmilne.org/math/xnotes/JVs.pdf}{Jacobian Varieties}. In \emph{Arithmetic Geometry}. 1986. Corrected online version.}

\bpbibitem{Ran39}{bp349}{R. A. Rankin. Contributions to the theory of Ramanujan's function \(\tau(n)\) and similar arithmetical functions. II. \emph{Proc. Cambridge Philos. Soc.}, 35:357–372, 1939.}

\bpbibitem{SGA4}{bp330}{M. Artin, A. Grothendieck, and J.-L. Verdier, editors. \emph{Théorie des topos et cohomologie étale des schémas}, volumes 269, 270, 305 of \emph{Lecture Notes in Mathematics}. Springer, 1972–1973. SGA 4.}

\bpbibitem{SGA4½}{bp331}{P. Deligne. \emph{Cohomologie étale}, volume 569 of \emph{Lecture Notes in Mathematics}. Springer, 1977. SGA 4½; \emph{Théorèmes de finitude}, §3, and \emph{Rapport sur la formule des traces}, §3, and \emph{Sommes trigonométriques}, §1.}

\bpbibitem{SGA5}{bp332}{A. Grothendieck, editor. \emph{Cohomologie \(\ell\)-adique et fonctions \(L\)}, volume 589 of \emph{Lecture Notes in Mathematics}. Springer, 1977. SGA 5.}

\bpbibitem{SGA7II}{bp350}{P. Deligne and N. Katz, editors. \emph{Groupes de monodromie en géométrie algébrique}, volume 340 of \emph{Lecture Notes in Mathematics}. Springer, 1973. SGA 7 II; Exposé XIII, (2.1.7).}

\bpbibitem{Stacks}{bp351}{The Stacks Project Authors. \href{https://stacks.math.columbia.edu}{The Stacks Project}. Individual results are cited by their stable tags in the text.}

\bpbibitem{Var05}{bp352}{Y. Varshavsky. Lefschetz–Verdier trace formula and a generalization of a theorem of Fujiwara. 2005. Version 2. arXiv:\href{https://arxiv.org/abs/math/0505564}{math/0505564}. \href{https://arxiv.org/abs/math/0505564v2}{Author's preprint}.}

\bpbibitem{Wei48a}{bp353}{A. Weil. \emph{Sur les courbes algébriques et les variétés qui s'en déduisent}. Hermann, Paris, 1948.}

\bpbibitem{Wei48b}{bp354}{A. Weil. \emph{Variétés abéliennes et courbes algébriques}. Hermann, Paris, 1948.}

\bpbibitem{Wei49}{bp355}{A. Weil. Numbers of solutions of equations in finite fields. \emph{Bull. Amer. Math. Soc.}, 55:497–508, 1949.}

\end{document}
